% Incognita paper P001, a paper of the owner (AGENTS.md 5.1; CONTRIBUTING.md, section 8). Rules:
% AGENTS.md, section 5 (template preamble; at most 100 pages). Working notes go in notebook/.
% Converted on 2026-11-14 from the author's draft of 5 November 2026 (statements and proofs
% unchanged; form adapted to the library: sections, blueprint, packages).
% Compile: latexmk -pdf -interaction=nonstopmode -halt-on-error -cd papers/<P>/paper.tex (as CI),
% after `python3 scripts/certify.py bib` (papers/library.bib, the library's papers; AGENTS.md 5.6).
% Blueprint conventions: every definition/lemma/proposition/theorem/corollary has a \label;
% statements start with \uses{def:...} when they use definitions; every proof starts with an
% exhaustive \uses{...}; a proof away from its statement has \proves{...}. Never write \lean,
% \leanok or \mathlibok in this file. After acceptance this file is never edited (AGENTS.md 6.6).
\documentclass[11pt]{article}
\usepackage[T1]{fontenc}
\usepackage[utf8]{inputenc}
\usepackage{lmodern}
\usepackage[a4paper,margin=2.5cm]{geometry}
\usepackage{amsmath,amssymb,mathtools}
\usepackage{xcolor}
\usepackage[colorlinks=true,linkcolor=blue!50!black,citecolor=blue!50!black,urlcolor=blue!50!black]{hyperref}
\usepackage{../../templates/paper/blueprint-macros}

% Macros of this paper: \newcommand only.



\newcommand{\dt}[1]{\hypertarget{#1}{}}
\DeclareRobustCommand{\sy}[2]{{\hyperlink{#1}{#2}}}

\newcommand{\R}{\mathbb{R}}
\newcommand{\Nat}{\mathbb{N}}
\newcommand{\E}{\mathbb{E}}
\renewcommand{\P}{\mathbb{P}}
\newcommand{\1}{\mathbf{1}}
\newcommand{\kl}{\mathrm{kl}}
\newcommand{\KL}{\mathrm{KL}}

\newcommand{\RT}{\sy{s:regret}{R_T}}
\newcommand{\Rs}{\sy{s:regret}{R^\star}}
\newcommand{\Cl}{\sy{s:class}{\mathcal C}}
\newcommand{\Cab}{\sy{s:class}{\mathcal C(a,b,\bar\sigma,\sigma_{\mathrm p})}}
\newcommand{\sbar}{\sy{s:class}{\bar\sigma}}
\newcommand{\sigp}{\sy{s:class}{\sigma_{\mathrm p}}}

\newcommand{\Wc}{\sy{s:W}{\mathcal W}}
\newcommand{\Om}{\sy{s:W}{\Omega}}
\newcommand{\Al}{\sy{s:W}{\mathcal A}}
\newcommand{\rk}{\sy{s:risk}{r}}
\newcommand{\lh}{\sy{s:risk}{\ell}}
\newcommand{\VK}{\sy{s:VK}{V}}
\newcommand{\BK}{\sy{s:VK}{B}}
\newcommand{\Vb}{\sy{s:Vb}{V^\flat}}
\newcommand{\Vloc}{\sy{s:Vb}{V^{\mathrm{loc}}}}

\newcommand{\us}{\sy{s:kappa}{u^\star}}
\newcommand{\as}{\sy{s:kappa}{a^\star}}
\newcommand{\kap}{\sy{s:kappa}{\kappa}}
\newcommand{\hb}{\sy{s:kappa}{\mathrm h}}
\newcommand{\lamK}{\sy{s:kappa}{\lambda}}

\newcommand{\SC}{\sy{alg:sc}{\textsc{Screen-Core}}}
\newcommand{\OPT}{\sy{alg:opt}{\textsc{Opt}\textsubscript{2}}}

\newcommand{\Ud}{\sy{s:U}{U}}
\newcommand{\Id}{\sy{s:U}{I}}
\newcommand{\dst}{\sy{s:two}{d^\star}}
\newcommand{\thr}{\sy{s:two}{\vartheta}}
\setlength{\emergencystretch}{3em}

\title{The optimal constant in the minimax regret\\ of stochastic multi-armed bandits}
\author{Adrien Prevost\thanks{Submitted by its author to the Incognita library (roadmap StochasticBandits, direction
  D19) on 2026-11-14. The proofs were checked by AI agents only, not by a human referee, and nothing is formalised in
  Lean yet. The conversion to the format of the library was done with AI assistance (Claude Opus 5.5).}}
\date{2026-11-14}

\begin{document}
\sloppy
\maketitle

\begin{abstract}
The minimax regret of a stochastic $K$-armed bandit with horizon $T$ is of order $\sqrt{KT}$, but its constant is not known: published lower and upper
constants differ by a factor larger than $28$. We show that for every $K\ge2$ and every regular class of reward distributions $\mathcal C$, the minimax regret
satisfies $R^\star_{K,\mathcal C}(T)/\sqrt T\to\bar\sigma\,V_K$. Here $\bar\sigma^2$ is the largest variance in the class, and $V_K$ is the value of a Gaussian game in
continuous time that does not depend on $\mathcal C$. For rewards in $[0,1]$ the constant is $\frac12V_K$. The value has four equivalent descriptions:
minimax over limit policies, supremum of Bayes risks, limit of batched policies, and limit of the local Gaussian minimax regret. It satisfies
$\kappa\lambda_K\sqrt{K-1}\le V_K\le2\sqrt2\sqrt K$ with $\kappa\approx0.9025$ and $\lambda_K\to1$. For Gaussian rewards, the lower bound $\kappa\lambda_K\sqrt{(K-1)T}$ holds
for every $T\ge2.52(K-1)$; this improves the classical constant $\frac1{27}$ by a factor between $6$ and $24$. A two-stage policy, \textsc{Screen-Core}, attains the
constant up to any $\varepsilon>0$: a screening phase removes the arms with non-local gaps, and a near-optimal batched core plays the limit game on a compact set.
That core exists for every $K$ but is computed only for two arms. A computer-assisted computation in double precision, without interval arithmetic, gives
$0.6341\le V_2\le0.7478$; the numerical value is $0.6373$. We also give a practical two-armed policy built on the computed core, without a proved guarantee. In
simulations its worst-case regret is $17\%$ below that of the best of five standard policies and $32\%$ below that of MOSS. On a fixed instance it does worse than
Thompson sampling or IMED, as expected from a policy tuned to the worst gap.
\end{abstract}

\section{Main result}\label{sec:main}\label{sec:intro}

The minimax regret $\Rs_{K,\Cl}(T)$ of a $K$-armed stochastic bandit over a class $\Cl$ of reward distributions (Definitions~\ref{def:regret}, \ref{def:class}) is of order $\sqrt{KT}$; its constant is the subject of this paper. $\VK_K$ is the value of a Gaussian game in continuous time (Definition~\ref{def:risk}).

\begin{theorem}[the constant exists]\label{thm:main}
\uses{def:sc, def:class, def:regret, def:risk}
Let $K\ge2$ and let $\Cl$ be regular with variance $\sbar$. Then
\[
\lim_{T\to\infty}\frac{\Rs_{K,\Cl}(T)}{\sqrt T}=\sbar\,\VK_K .
\]
More precisely: (i) the lower bound $\liminf\ge\sbar\VK_K$ only uses (R2), with an error $O(T^{-1/4})$ uniform over the policies and over the local instances in a compact set;
(ii) the upper bound $\limsup\le\sbar\VK_K$ only uses (R1), and for every $\varepsilon>0$ the policy \SC\ (Definition~\ref{def:sc}), with parameters depending on $\varepsilon$, has
$\sup_{\nu\in\Cl^K}\RT(\SC,\nu)\le(\sbar\VK_K+\varepsilon)\sqrt T$ for $T$ large enough.
\end{theorem}


\paragraph{Constants.} The limit $\sbar\VK_K$ is explicit only through the value $\VK_K$ of the Gaussian limit game
(Definition~\ref{def:risk}), which has no closed form. Theorem~\ref{thm:bounds} gives
$\kap\lamK_K\sqrt{K-1}\le\VK_K\le2\sqrt2\sqrt K$, with the explicit constant $\kap\approx0.9025$, and for two arms
Proposition~\ref{prop:K2} gives $0.6341\le\VK_2\le0.7478$ by a computation in double precision, without interval
arithmetic. In part (i) the error is $O(T^{-1/4})$, uniformly over the policies and over the local instances in a compact
set. In part (ii) the parameters of \SC{} depend on $\varepsilon$, and ``$T$ large enough'' is existential.

\paragraph{Roadmap items.} This paper serves direction D19 of the roadmap StochasticBandits.
(a)~The limit of $\Rs_{K,\Cl}(T)/\sqrt T$ exists and is identified, for every $K\ge2$ and every regular class
(Definition~\ref{def:regular}), which includes the three models of D19: Bernoulli laws, laws on $[0,1]$, Gaussian laws
with bounded means and bounded variance. (b)~The constant is attained up to any $\varepsilon>0$ by the horizon-dependent
policy \SC; its core is a near-optimal batched policy of the limit game, which exists for every $K$ but is computed only
for $K=2$; no anytime policy is given (Section~\ref{sec:open}). (c)~Bounds for every $K$, non-asymptotic for Gaussian
rewards, which improve the published lower constant by a factor between 6 and 24, and the order $\sqrt K$ of $\VK_K$ up to
the factors $0.9025$ and $2\sqrt2$.


\paragraph{Context.}

In a stochastic bandit with $K$ arms and horizon $T$, the minimax regret is of order $\sqrt{KT}$. The lower bound goes back to \cite{ACFS02}. The first policy with a
regret $O(\sqrt{KT})$ without logarithmic factor is MOSS \cite{AB09}: for rewards in $[0,1]$ and a known horizon, its regret is at most $49\sqrt{KT}$. It was later extended to
$1$-sub-Gaussian rewards, with the bound $39\sqrt{KT}+\sum_k\Delta_k$ \cite[Thm~9.1]{LS}, and to anytime \cite{DP16} and asymptotically optimal versions
\cite{MG17,GHMS22}. The best known constants are $\sqrt{2KT}$ for the upper bound with rewards in $[0,1]$ \cite{LS19,ZL19}, and $\frac1{20}\sqrt{KT}$
\cite{ACFS02} or $\frac1{27}\sqrt{(K-1)T}$ \cite[Thm~15.2]{LS} for lower bounds. They differ by a factor larger than $28$. So ``minimax optimal'' has so far
meant optimal up to a constant factor.

This paper asks whether $\Rs/\sqrt{KT}$ has a limit, what it is, and which policy attains it. The worst instances have gaps of order $\sqrt{K/T}$; at this
scale the bandit converges to a Gaussian experiment in continuous time, and the natural candidate for the constant is the value of the corresponding game.
For two arms this programme is old. Vogel \cite{Vog60} proved that the minimax regret is of order $\sqrt N$. Fabius and van Zwet \cite{FvZ70}, then Bather \cite{Bat83},
bounded the constant. Kolnogorov \cite{Kol17,Kol20} derived the limiting Hamilton--Jacobi--Bellman equation and computed $\approx0.637$; he states the identification of the
limit with the minimax value without proof. Adusumilli \cite{Adu25} develops the decision theory of bandit experiments under diffusion asymptotics; for minimax
risk he proves a lower bound, and he lists more than two arms as open.


\paragraph{Contributions.}
\begin{enumerate}
\item \emph{The constant exists, for every $K$} (Theorem~\ref{thm:main}). For every $K\ge2$ and every regular class $\Cl$ (Definition~\ref{def:regular}), $\Rs_{K,\Cl}(T)/\sqrt T\to\sbar\VK_K$, where
$\sbar^2$ is the largest variance in the class. The value $\VK_K$ of the Gaussian limit game has four equivalent descriptions (Proposition~\ref{prop:four}).
For rewards in $[0,1]$, the constant is $\frac12\VK_K$, as for Bernoulli rewards.
\item \emph{Bounds for every $K$} (Theorem~\ref{thm:bounds}): $\kap\lamK_K\sqrt{K-1}\le\VK_K\le2\sqrt2\sqrt K$, with $\kap\approx0.9025$ and $\lamK_K\to1$. For Gaussian
rewards the lower bound is non-asymptotic. The method is that of \cite{GMS19,GMS20}: the fundamental inequality and the joint convexity of $\kl$, applied to a prior
where one arm, at a uniform position, is better than the others. What is new is the optimization of this prior, which yields the constant $\kap$.
\item \emph{An optimal policy} (Section~\ref{sec:algo}). For every $\varepsilon>0$, \SC\ has worst-case regret at most $(\sbar\VK_K+\varepsilon)\sqrt T$ for $T$ large: a screening phase whose
cost does not depend on the gaps, followed by a near-optimal batched policy of the limit game on a compact set. This core exists for every $K$; it is not explicit for $K\ge3$.
For $K=2$ we compute an approximate core and derive a practical policy, \OPT, without a proved guarantee.
\item \emph{Two arms} (Section~\ref{sec:two}): $0.6341\le\VK_2\le0.7478$ by a computer-assisted computation in double precision (no interval arithmetic), and
$\VK_2\approx0.6373$ numerically.
\item \emph{Experiments} (Section~\ref{sec:exp}): worst-case and problem-dependent regret of \OPT, MOSS, KL-UCB++, KL-UCB, Thompson sampling and IMED.
\end{enumerate}


\setcounter{tocdepth}{1}
\tableofcontents


\section{Setting}\label{sec:setting}

The notation of this paper (it does not use the pinned Lean conventions of the roadmap, which formalisation will adopt). Every definition below is a node.

\subsection{Bandits and classes of distributions}
\begin{definition}[regret and minimax regret]\label{def:regret}\dt{s:regret}
\uses{def:class}
Arm $k\in[K]$ gives i.i.d.\ rewards of law $\nu_k$ with mean $\mu_k$; $\Delta_k:=\max_i\mu_i-\mu_k$. A policy $\pi$ plays one arm per round,
possibly with external randomization, and $N_k(T)$ is the number of pulls of arm $k$ in rounds $1,\dots,T$. For a class $\Cl$ of laws on $\R$ let
\[
\RT(\pi,\nu):=\sum_k\Delta_k\,\E[N_k(T)],\qquad \Rs_{K,\Cl}(T):=\inf_\pi\sup_{\nu\in\Cl^K}\RT(\pi,\nu).
\]
\end{definition}
\begin{definition}[the classes $\Cab$]\label{def:class}\dt{s:class}
\uses{def:regret}
For $a<b$ and $0<\sbar\le\sigp$, let $\Cab$ be the set of laws $\nu$ on $\R$ with mean in $[a,b]$, variance at most $\sbar^2$, and
$\E_\nu e^{\lambda(X-\mu)}\le e^{\lambda^2\sigma_{\mathrm p}^2/2}$ for all $\lambda\in\R$. The means must be bounded, otherwise $\Rs=\infty$.
\end{definition}
\begin{definition}[regular class]\label{def:regular}
\uses{def:class}
A class $\Cl$ is \emph{regular with variance $\sbar$} if:
\begin{enumerate}
\item[(R1)] $\Cl\subseteq\Cab$ for some $a<b$ and $\sigp\ge\sbar$;
\item[(R2)] for every $\sigma_0<\sbar$ there is $\nu_0\in\Cl$ with variance at least $\sigma_0^2$ such that either $\Cl$ contains a regular one-parameter exponential family
(a natural exponential family $\nu_\theta\propto e^{\theta x}\nu_0(dx)$ whose natural parameter space contains a neighbourhood of $0$) through
$\nu_0$, on both sides or on one side of $\nu_0$, or $\Cl$ contains the perturbations $(1+tg)\nu_0$ for every bounded centered $g$ and every $|t|$ small enough.
\end{enumerate}
\end{definition}
\noindent Examples: the laws on $[a,b]$, with $\sbar=\frac{b-a}2$ (Bernoulli at the midpoint); the Bernoulli laws, and also the Bernoulli laws of mean at most $\frac12$, with
$\sbar=\frac12$; the Gaussian laws $\mathcal N(\mu,\sigma^2)$ with $\mu\in[a,b]$ and $\sigma\le\sbar$.

\subsection{The Gaussian limit game}\label{sec:game}
\begin{definition}[allocations and the canonical space]\label{def:alloc}\dt{s:W}
Let $\Al$ be the set of allocations: continuous maps $\tau=(\tau_k)_{k\in[K]}$ from $[0,1]$ to $[0,1]^K$, with each $\tau_k$
nondecreasing, $\tau(0)=0$ and $\sum_k\tau_k(t)=t$. Let $\Om:=C_0([0,1],\R^K)\times\Al$, with coordinates $(X,\tau)$ and the filtration $(\mathcal F_t)$ they generate.
\end{definition}
\begin{definition}[limit policies]\label{def:limit}
\uses{def:alloc}
A \emph{limit policy} is a probability $Q$ on $\Om$ under which $X$ is a continuous $(\mathcal F_t)$-martingale with $\langle X_k,X_j\rangle=\1\{k=j\}\tau_k$.
$\Wc$ denotes the set of limit policies.
\end{definition}
\noindent Think of $X_k(t)$ as the cumulative reward of arm $k$, scaled by $\sqrt T$, under the reference instance where all means are equal. The local
instance $h\in\R^K$ means $\mu_k=\mu_0+h_k/\sqrt T$; under $h$ the law of $(X,\tau)$ is $\lh_hQ$, by Girsanov's theorem.

\begin{definition}[risk and values]\label{def:risk}\dt{s:risk}\dt{s:VK}
\uses{def:alloc}
For $h\in\R^K$ let $\lh_h:=\exp\big(\langle h,X(1)\rangle-\frac12\sum_kh_k^2\tau_k(1)\big)$ and
\[
\rk(Q,h):=\E_Q\Big[\lh_h\sum_k\big(\max_ih_i-h_k\big)\tau_k(1)\Big],
\]
\[
\VK_K:=\min_{Q\in\Wc}\ \sup_{h\in\R^K}\rk(Q,h),\qquad \BK_K(\Pi):=\min_{Q\in\Wc}\int\rk(Q,h)\,\Pi(dh).
\]
The minima are attained (Section~\ref{sec:app-game}). A $T$-round policy for Gaussian rewards of variance $1$ is a limit policy, after scaling (Lemma~\ref{lem:embed}).
\end{definition}
\begin{definition}[batched and local values]\label{def:batched}\dt{s:Vb}
\uses{def:risk, def:alloc}
A \emph{batched policy} fixes finitely many decision times $0=t_0<\dots<t_m=1$; at time $t_i$ it splits $[t_i,t_{i+1}]$
between the arms in proportions that are Borel functions of the randomization and of the observed path $X|_{[0,t_i]}$. Reading only the values at the decision times
gives the same risk functions (Lemma~\ref{lem:g-bridge}). For a compact $H\subset\R^K$ let
$\Vb(H):=\inf\{\sup_{h\in H}\rk(\pi,h):\pi\text{ batched}\}$ and $\VK_K(H):=\min_{Q\in\Wc}\sup_{h\in H}\rk(Q,h)$, and let $\Vloc_K:=\sup_M\VK_K([-M,M]^K)$.
\end{definition}

\begin{definition}[the constants $\kap$ and $\lamK_K$]\label{def:kappa}\dt{s:kappa}
Let $\us\approx1.256431$ be the positive root of $e^u=1+2u$, $\as:=\sqrt{2\us}\approx1.585201$,
\[
\kap:=\frac{2\as}{1+(\as)^2}=\sqrt2\,\frac{1-e^{-\us}}{\sqrt{\us}}\approx0.902512,
\]
$\hb(q):=-q\log q-(1-q)\log(1-q)$ the binary entropy, and $\lamK_K:=\exp\big(-K\hb(1/K)/(K-1)\big)=1-O(\frac{\log K}K)$; for instance $\lamK_2=\frac14$, $\lamK_{10}\approx0.70$, $\lamK_{100}\approx0.94$.
\end{definition}

\section{Closest known}\label{sec:closest}

KNOWN.md of the roadmap StochasticBandits gives minimax-optimal algorithms only up to constant factors (entries K20--K24),
and direction D19 lists the published constants: the lower bounds $\frac1{20}\sqrt{KT}$ \cite{ACFS02} and
$\frac1{27}\sqrt{(K-1)T}$ for Gaussian rewards \cite[Theorem~15.2]{LS}, and the upper bounds $49\sqrt{KT}$ for MOSS
\cite{AB09} and $\sqrt{2KT}$ for rewards in $[0,1]$ \cite{ZL19}. The related literature, as described in the author's
manuscript:

\paragraph{Related work.} Diffusion limits of bandit algorithms in the regime of gaps of order $T^{-1/2}$ are studied in \cite{FG21,KW24,KZ21}, and asymptotic
representation theorems for sequential experiments in \cite{HP23,Adu25}. In prediction with expert advice the minimax regret is $\sqrt{(T/2)\log N}$ asymptotically
\cite{CB97}, and for a small number of experts it is the solution of an explicit PDE \cite{Cov66,GPS16,BEZ20}. In the infinitely-many-armed Bayesian bandit with
uniform prior, the optimal constant $\sqrt{2n}$ is known, together with a policy that attains it \cite{BCZ97,BP13}.

\paragraph{Why Theorem~\ref{thm:main} does not follow from these works.}
\begin{itemize}
\item \emph{The published bounds} \cite{ACFS02,LS,AB09,ZL19} bracket $\Rs_K(T)/\sqrt{KT}$ between constants that differ
by a factor larger than $28$; they say nothing about the existence of a limit. Theorem~\ref{thm:bounds} is consistent with
them: in the normalization of \cite[Theorem~15.2]{LS} its lower constant is $\kap\lamK_K\ge\frac1{27}$, and for rewards in
$[0,1]$ its upper bound $\frac12\VK_K\le\sqrt{2K}$ is the constant of \cite{ZL19}.
\item \emph{Two arms} \cite{Vog60,FvZ70,Bat83,Kol17,Kol20}: Vogel proves the order $\sqrt T$, Fabius--van Zwet and Bather
bound the constant, and Kolnogorov derives the limiting Hamilton--Jacobi--Bellman equation and computes $\approx0.637$, but
states the identification of the limit with the minimax value without proof. The attempt to deduce Theorem~\ref{thm:main}
from them fails twice: they treat $K=2$ only, and the identification is precisely the missing step, which needs the
compactness and minimax arguments of Section~\ref{sec:app-game}.
\item \emph{Diffusion asymptotics} \cite{Adu25,HP23,KZ21,FG21,KW24}: Adusumilli develops the decision theory of bandit
experiments in this regime and proves a lower bound for the minimax risk, listing more than two arms as open; the others
study diffusion limits of given algorithms. None of them identifies the minimax constant for $K\ge3$, nor gives a policy
that attains it.
\item \emph{The lower-bound technique} \cite{GMS19,GMS20}: the fundamental inequality and the convexity of $\kl$ applied
to a one-needle prior give $\Omega(\sqrt{KT})$ with a small constant. What is new in Theorem~\ref{thm:bounds} is the
optimization of the prior, which gives the constant $\kap$.
\end{itemize}

\section{Proofs}\label{sec:proofs}

This section states the further results, the policies and the two-armed computations, and Sections~\ref{sec:app-game}--\ref{sec:app-exp} prove them.


\subsection{Further results}\label{sec:results}



\begin{corollary}[examples]\label{cor:examples}
\uses{def:regret, def:risk, def:class}
For every $K\ge2$, $\lim_T\Rs_K(T)/\sqrt T$ equals: $\frac12\VK_K$ for rewards in $[0,1]$, for Bernoulli rewards and for Bernoulli rewards of mean at most $\frac12$;
$\frac{b-a}2\VK_K$ for rewards in $[a,b]$; $\sbar\VK_K$ for Gaussian rewards $\mathcal N(\mu,\sigma^2)$ with $\mu\in[a,b]$ and $\sigma\le\sbar$. In the Gaussian case,
moreover, $\Rs_K(T)\ge\sbar\sqrt T\,\VK_K([-M,M]^K)$ for every $M$ and every $T\ge(2\sbar M/(b-a))^2$.
\end{corollary}

\noindent For the Bernoulli laws of mean at most $\frac12$, the law of largest variance sits at the boundary of the class; the constant does not change because a translated box of
local instances has the same value (Proposition~\ref{prop:four}(d)).

\begin{proposition}[four descriptions of the value]\label{prop:four}
\uses{def:risk, def:batched, def:regret}
For every $K\ge2$:
\begin{enumerate}
\item[(a)] (minimax theorem) $\VK_K=\sup\{\BK_K(\Pi):\Pi$ finitely supported$\}$;
\item[(b)] (batched policies) $\Vb(H)=\VK_K(H)$ for every compact $H$, and $\VK_K=\Vloc_K=\lim_{M\to\infty}\Vb([-M,M]^K)$;
\item[(c)] (local Gaussian minimax regret) for every $M>0$,
\[
\VK_K([-M,M]^K)=\lim_{T\to\infty}\frac1{\sqrt T}\,\inf_\pi\sup_{h\in[-M,M]^K}\RT\big(\pi,(\mathcal N(h_k/\sqrt T,1))_k\big);
\]
\item[(d)] (invariances) $\VK_K(H+c\1)=\VK_K(H)$ for $c\in\R$, there is a minimax limit policy that is symmetric in the arms and invariant by translation, and $\VK_K\le\VK_{K+1}$.
\end{enumerate}
\end{proposition}

\begin{theorem}[bounds on the value]\label{thm:bounds}
\uses{def:kappa, def:risk, def:regret}
For every $K\ge2$,
\[
\kap\,\lamK_K\,\sqrt{K-1}\ \le\ \VK_K\ \le\ 2\sqrt2\,\sqrt K .
\]
Moreover, for unit-variance Gaussian rewards with means in $[0,1]$ and every $T\ge2\us(K-1)$, $\Rs_K(T)\ge\kap\lamK_K\sqrt{(K-1)T}$.
\end{theorem}

\noindent In the normalization of \cite[Thm~15.2]{LS}, the constant $\frac1{27}\approx0.037$ becomes $\kap\lamK_K$: $0.23$ ($K=2$), $0.35$ ($K=3$), $0.63$ ($K=10$),
$0.85$ ($K=100$), and it tends to $\kap$. Hence $0.9025\le\liminf_K\VK_K/\sqrt K\le\limsup_K\VK_K/\sqrt K\le2\sqrt2$.

\begin{remark}[many arms]\label{rem:many}
The prior behind $\kap$ puts one good arm at a uniform position. Its Bayes risk, divided by $\sqrt{KT}$, tends to $\frac2a(1-e^{-a^2/2})$ as $T\to\infty$ then $K\to\infty$; it
is attained by a breadth-first policy that lowers all arms together, and its maximum over $a$ is $\kap$. This prior is probably not least favorable: after finding a
good arm, a policy cannot know whether a better one exists. In a limit search model, a restart policy has worst case about $1.12$ over all configurations of good arms
(numerical). Whether $\lim_K\VK_K/\sqrt K$ exists is open.
\end{remark}

\paragraph{How the proof of Theorem~\ref{thm:main} works.} \emph{Lower bound} (Section~\ref{sec:app-lower}). A $T$-round policy for an exponential family can be simulated inside the
limit game with almost the same risk. Each arm's rewards are embedded in its own Brownian motion by Skorokhod's embedding, which is adapted: each embedded
reward is a stopping time of that arm's path. The likelihood ratios of the family and of the Gaussian model are compared on the full arrays of rewards, so the
comparison does not depend on the policy. \emph{Upper bound} (Section~\ref{sec:app-upper}). The screening of \SC\ costs $O(\sqrt{\delta\log(1/\delta)})\sqrt T$ whatever the gaps. After it,
the surviving arms lie in a box of local instances of size $O(\sqrt{AK/\delta})$, where a functional central limit theorem transfers the risk of a continuous batched core
to its regret. \emph{Identification} (Section~\ref{sec:app-game}). The limit policies form a compact convex set on which the risk is continuous, so Ky Fan's theorem gives the
minimax theorem. Discrete and batched policies are dense in this set, and, given the values of a Brownian motion on a grid, the path in between is a Brownian bridge
that does not depend on the drift. Hence batched, discrete and continuous-time policies have the same value.


\subsection{Policies}\label{sec:algo}

\begin{definition}[algorithm: \textsc{Screen-Core}$(\delta,A,\pi^\circ)$ for $\Cab$]\label{def:sc}\dt{alg:sc}
\uses{def:class}
\leavevmode\par\smallskip\noindent\fbox{\begin{minipage}{0.96\linewidth}\small
\noindent\hspace*{0.0em}\parbox[t]{\dimexpr\linewidth-0.0em\relax}{\raggedright \textbf{Input:} horizon $T$, $K$, $\sbar$, $\sigp$, screening fraction $\delta\in(0,\frac12]$, level $A\ge8$, continuous batched limit policy $\pi^\circ$}\par
\noindent\hspace*{0.0em}\parbox[t]{\dimexpr\linewidth-0.0em\relax}{\raggedright $N\gets\lfloor\delta T/K\rfloor$, $S_0\gets[K]$ \hfill(Stage 1: screening)}\par
\noindent\hspace*{0.0em}\parbox[t]{\dimexpr\linewidth-0.0em\relax}{\raggedright \textbf{for} $n=1,\dots,N$ \textbf{do}}\par
\noindent\hspace*{1.3em}\parbox[t]{\dimexpr\linewidth-1.3em\relax}{\raggedright pull once each arm of $S_{n-1}$; $\hat\mu_k(n)\gets$ mean of the first $n$ rewards of arm $k$}\par
\noindent\hspace*{1.3em}\parbox[t]{\dimexpr\linewidth-1.3em\relax}{\raggedright $S_n\gets\{k\in S_{n-1}:\hat\mu_k(n)+\rho_n\ge\max_{i\in S_{n-1}}\hat\mu_i(n)-\rho_n\}$, with $\rho_n:=\sigp\sqrt{\frac2n\big(A+2\log\frac Nn\big)}$}\par
\noindent\hspace*{0.0em}\parbox[t]{\dimexpr\linewidth-0.0em\relax}{\raggedright $S\gets S_N$; $i_0\gets$ an empirical best arm of $S$; $\hat\mu_{\mathrm{ref}}\gets\hat\mu_{i_0}(N)$}\par
\noindent\hspace*{0.0em}\parbox[t]{\dimexpr\linewidth-0.0em\relax}{\raggedright for $k\in S$: $\hat\sigma_k^2\gets$ empirical variance of arm $k$; $w_k\gets(\sbar^2-\hat\sigma_k^2)_+$}\par
\noindent\hspace*{0.0em}\parbox[t]{\dimexpr\linewidth-0.0em\relax}{\raggedright $c(k)\gets k$ if $k\in S$, $c(k)\gets i_0$ otherwise \hfill(Stage 2: core)}\par
\noindent\hspace*{0.0em}\parbox[t]{\dimexpr\linewidth-0.0em\relax}{\raggedright run the discretization of $\pi^\circ$ on the $T'$ remaining rounds, centered at $\hat\mu_{\mathrm{ref}}$ and scaled by $\sbar\sqrt{T'}$, on $K$ virtual arms: virtual arm $k$ is a pull of $c(k)$, returned with an added independent $\mathcal N(0,w_{c(k)})$ noise}\par
\end{minipage}}
\end{definition}


\noindent Stage 1 makes the regret uniform over the gaps; afterwards the core only faces a compact box of local instances, where a near-optimal batched policy exists by
Proposition~\ref{prop:four}(b). Eliminated arms are replaced by copies of $i_0$, so the core is always a $K$-armed policy. The added noise brings every variance to
$\sbar^2$, the worst case. Theorem~\ref{thm:main}(ii) holds with $\delta\to0$, $A\approx6\log(1/\delta)$ and a core whose risk on the box is at most $\Vb+\varepsilon$
(Section~\ref{sec:app-upper}).

\paragraph{Two arms: a practical policy.}\dt{alg:opt} For $K=2$ we compute (Section~\ref{sec:two}) the Bayes policy of the two-point prior that is least favorable
numerically. It is described by a threshold $\thr(\tau_1,\tau_2)$ on the contrast statistic $\Ud$. Definition~\ref{def:opt2} uses it as a core; it is the version used in the
experiments. \emph{Theorem~\ref{thm:main} does not cover \OPT}, for four reasons:
\begin{itemize}
\item its threshold rule is discontinuous, while the transfer theorem of Section~\ref{sec:app-upper} needs continuous rules;
\item its risk grows with the gap, so it is near-optimal only on a fixed box, whereas \SC\ needs a near-optimal core on boxes whose size grows as $\delta\to0$;
\item the least favorable prior is only supported numerically;
\item \OPT\ keeps the data of its first phase, and its elimination test runs during the whole horizon.
\end{itemize}
\OPT\ is a practical approximation of the optimal policy, and the experiments only measure it.

\begin{definition}[algorithm: \textsc{Opt}$_2$ (two arms, known $\sigma$ and $T$)]\label{def:opt2}
\leavevmode\par\smallskip\noindent\fbox{\begin{minipage}{0.96\linewidth}\small
\noindent\hspace*{0.0em}\parbox[t]{\dimexpr\linewidth-0.0em\relax}{\raggedright \textbf{Input:} horizon $T$, $\sigma$, table $\thr$ on the grid $\varepsilon\Nat^2$ ($\varepsilon=0.0025$), level $A$ ($A=4$ in the experiments)}\par
\noindent\hspace*{0.0em}\parbox[t]{\dimexpr\linewidth-0.0em\relax}{\raggedright \textbf{for} $t=0,\dots,T-1$ \textbf{do}}\par
\noindent\hspace*{1.3em}\parbox[t]{\dimexpr\linewidth-1.3em\relax}{\raggedright $\tau_k\gets N_k/T$; $\Ud\gets(S_1N_2-S_2N_1)/((N_1+N_2)\,\sigma\sqrt T)$ \hfill($S_k$: sum of the rewards of arm $k$)}\par
\noindent\hspace*{1.3em}\parbox[t]{\dimexpr\linewidth-1.3em\relax}{\raggedright $\rho_k\gets\sigma\sqrt{\frac2{N_k}\big(A+2\log^+\frac T{2N_k}\big)}$; if $\hat\mu_i-\rho_i>\hat\mu_j+\rho_j$ once, play $i$ until the end}\par
\noindent\hspace*{1.3em}\parbox[t]{\dimexpr\linewidth-1.3em\relax}{\raggedright \textbf{if} $\min_kN_k<\lceil\varepsilon T\rceil$ \textbf{then}}\par
\noindent\hspace*{2.6em}\parbox[t]{\dimexpr\linewidth-2.6em\relax}{\raggedright pull the arm with fewer pulls}\par
\noindent\hspace*{1.3em}\parbox[t]{\dimexpr\linewidth-1.3em\relax}{\raggedright \textbf{else}}\par
\noindent\hspace*{2.6em}\parbox[t]{\dimexpr\linewidth-2.6em\relax}{\raggedright \ pull arm $1$ if $\Ud\ge\thr(\tau_1,\tau_2)$, arm $2$ otherwise (grid point nearest to $(\tau_1,\tau_2)$)}\par
\end{minipage}}
\end{definition}



\subsection{Two arms}\label{sec:two}

\paragraph{Reduction.}\dt{s:U} Write the local means as $h_1=m+v$, $h_2=m-v$, and use horizon fractions $\tau_k:=N_k/T$, $t:=\tau_1+\tau_2$. With a flat prior on $m$, the
posterior of $v$ is proportional to $\rho(v)\exp(2v\Ud-2v^2\Id)$, where
\[
\Ud:=\frac{X_1\tau_2-X_2\tau_1}t,\qquad \Id:=\frac{\tau_1\tau_2}t .
\]
Under $v$, $\Ud=2v\Id+B(\Id)$ for a Brownian motion $B$: the policy only chooses the speed of the information clock, $d\Id=(\tau_2/t)^2d\tau_1$ or $(\tau_1/t)^2d\tau_2$.
This is the formulation of \cite{Kol17}.

\paragraph{Least favorable prior and policy.}\dt{s:two} For the two-point prior $v=\pm d$, the Bayes risk solves a backward recursion in $(\tau_1,\tau_2,\Ud)$, which we solve with
an exact Gaussian convolution (FFT) for decisions held during a fraction $\varepsilon$ of the horizon, and extrapolate to $\varepsilon\to0$ (Table~\ref{tab:eps}). The worst half-gap
is $\dst\approx1.576$, that is a gap $\approx3.15/\sqrt T$. The batched recursion only uses policies that read $\Ud$, so its value is at least the Bayes risk of
the two-point prior, which is itself at most $\VK_2$; the extrapolation $0.63734$ estimates this Bayes risk, and equals $\VK_2$ if the two-point prior is least favorable.
The Bayes policy pulls arm $1$ iff $\Ud\ge\thr(\tau_1,\tau_2)$, with $\thr(\tau,\tau)=0$ and
$\thr(\tau_2,\tau_1)=-\thr(\tau_1,\tau_2)$. The threshold is negative when arm $1$ is under-sampled, for instance $\thr(0.1,0.3)\approx-0.20$ and $\thr(0.0025,0.5)\approx-0.33$:
the policy keeps pulling an under-sampled arm whose posterior probability of being the best is only about $0.2$. Its frequentist risk is maximal at $v=\dst$
(Table~\ref{tab:saddle}), which is numerical evidence that the two-point prior is least favorable on boxes.

\begin{table}[h]\centering\small
\begin{tabular}{lcccccc}
\hline
$\varepsilon$ & $0.02$ & $0.01$ & $0.005$ & $0.0025$ & $0.00125$ & $0.000625$\\ \hline
$\sup_d$ Bayes risk & $0.651723$ & $0.644212$ & $0.640671$ & $0.638969$ & $0.638141$ & $0.637735$\\
first-order extrapolation & --- & $0.636701$ & $0.637130$ & $0.637267$ & $0.637313$ & $0.637329$\\
\hline
\end{tabular}
\caption{Two arms: Bayes risk of the worst two-point prior with decisions held during $\varepsilon$ (grid $\Delta u=0.0125$; halving it changes the values by less than $10^{-5}$).
The extrapolation gives $0.63734\pm2\cdot10^{-5}$.}\label{tab:eps}
\end{table}

\begin{table}[h]\centering\small
\begin{tabular}{lcccccccccc}
\hline
$v$ & $0.25$ & $0.5$ & $1$ & $1.4$ & $1.55$ & $\mathbf{1.582}$ & $1.62$ & $1.8$ & $2.5$ & $6$\\ \hline
risk & $0.2219$ & $0.3893$ & $0.5820$ & $0.6345$ & $0.6388$ & $\mathbf{0.6389}$ & $0.6388$ & $0.6340$ & $0.5837$ & $0.4614$\\
\hline
\end{tabular}
\caption{Frequentist risk (normalized by $\sqrt T$) of the Bayes policy for $d=1.582$, $\varepsilon=0.0025$, at the half-gap $v$.}\label{tab:saddle}
\end{table}

\begin{proposition}[value for two arms]\label{prop:K2}
\uses{def:risk, def:regret}
$0.6341\le\VK_2\le0.7478$. Hence, for rewards in $[0,1]$, $0.3170\le\lim_T\Rs_2(T)/\sqrt T\le0.3739$.
\end{proposition}

\noindent The lower bound comes from a proper prior, normal on the common mean and two-point on the half-gap, together with a relaxation in which each block of length $\varepsilon$ is revealed
as soon as it is started. That relaxation is computed by a scheme in which every step is a lower bound. The upper bound is the worst-case risk of a balanced sequential test
followed by commitment, written as series with explicit remainders. The numerical value $0.6373$ is not certified: the risk of a threshold policy jumps at the
threshold. Kolnogorov's values $0.637$ and $d\approx1.57$ are close to $2/\pi$ and $\pi/2$, but the differences are more than ten times the extrapolation error.


\subsection{Experiments}\label{sec:exp}
We compare \OPT\ (Definition~\ref{def:opt2}) with two policies that are minimax optimal up to a constant, MOSS \cite{AB09} and KL-UCB++ \cite{MG17}, and with three
policies that are asymptotically optimal on fixed instances: KL-UCB \cite{CGMMS13}, Thompson sampling with a flat prior \cite{AG13} and IMED \cite{HT15}.
Rewards are Gaussian with known variance $1$. In each run the positions of the arms are drawn at random, so that tie-breaking rules play no role.
The precise indices, the parameters and the numbers of runs are in Section~\ref{sec:app-exp}; the whole study takes about three minutes on a laptop.


oindent(The figure of the author's version is not reproduced: the format of the library has no graphics. Its data and the scripts that produce it are in the notebook.)


\begin{table}[h]\centering\footnotesize\setlength{\tabcolsep}{4pt}
\begin{tabular}{lcccccc}
\hline
 & \OPT & IMED & Thompson & KL-UCB++ & MOSS & KL-UCB\\ \hline
(a) $\max\RT/\sqrt T$ & $\mathbf{0.677}\pm.033$ & $0.820\pm.025$ & $0.872\pm.037$ & $0.923\pm.021$ & $0.989\pm.022$ & $1.097\pm.020$\\
(b) $\max\RT/\sqrt{KT}$ & --- & $1.541\pm.027$ & $1.423\pm.037$ & $\mathbf{1.275}\pm.032$ & $1.373\pm.029$ & $2.138\pm.024$\\
(c) $\RT$, $T=10^5$ & $128.8\pm2.5$ & $\mathbf{58.0}\pm2.8$ & $59.8\pm3.5$ & $79.6\pm3.3$ & $88.4\pm3.0$ & $117.4\pm4.1$\\
\hline
\end{tabular}
\caption{Worst case over the grid of gaps (a, b) and regret on a fixed instance (c), with one standard error. The maximum of noisy estimates is biased upwards, by about one
standard error, for every policy alike.}\label{tab:exp}
\end{table}

\paragraph{Worst case, two arms.} The worst-case regret of \OPT\ is $0.677\sqrt T$, against $0.820\sqrt T$ for the best competitor (IMED) and $0.989\sqrt T$ for MOSS: a reduction
of $17\%$ and $32\%$. It is within $7\%$ of the asymptotic value $\VK_2\approx0.637$; the difference comes from the finite horizon ($\lceil\varepsilon T\rceil=25$ initial
pulls per arm) and from the upward bias of the maximum. The curves of the other policies peak at larger gaps, between $4/\sqrt T$ and $6/\sqrt T$: they explore too
little when the gap is a few times $1/\sqrt T$. \OPT\ peaks at $3/\sqrt T$, close to the least favorable gap $2\dst\approx3.15/\sqrt T$.

\paragraph{Worst case, ten arms.} We have no implementable optimal policy for $K=10$, so panel (b) only measures how far the existing policies are from what is achievable on
these instances. The best of them, KL-UCB++, reaches $1.275\sqrt{KT}$, whereas the myopic Bayes policy of this prior stays below $0.771\sqrt{KT}$ for every height. The proved lower
bound for this prior is $0.597\sqrt{KT}$ at most. The minimax constant $\VK_{10}/\sqrt{10}$ lies between $0.597$ (proved) and $1.275$ (up to the finite horizon). If the myopic policy is
Bayes-optimal, the existing policies lose about $40\%$ ($1-0.771/1.275\approx0.395$) on these instances alone. Other instances may be harder, so this says nothing about where
$\VK_{10}$ lies in that interval.

\paragraph{Fixed instances.} Panel (c) shows the price of minimax optimality. At $T=10^3$ the gap $0.2$ is local ($0.2\sqrt T\approx6$) and \OPT\ is the best policy. At
$T=10^5$ it is the worst: it still explores as if the gap could be $3.15/\sqrt T$, until its elimination test fires. The cost is of order $|\thr|\sqrt T$, where
$|\thr|\approx0.3$ is the size of its threshold when one arm is rarely pulled. Thompson sampling and
IMED, which are asymptotically optimal on fixed instances, have the smallest regret there. The Lai--Robbins slope $2\log T/\Delta$ is only reached asymptotically.

\subsection{Discussion}\label{sec:disc}
\paragraph{Worst case versus fixed instances.} \OPT\ is tuned for the worst gap: against a much larger gap it still explores as if the gap could be $3.15/\sqrt T$,
until its elimination test fires. A policy that is minimax optimal with the exact constant and also asymptotically optimal on fixed instances, in the sense of Lai and
Robbins, would combine the core with a switch to KL-UCB once the gap is clearly non-local, as KL-UCB-switch does with MOSS \cite{GHMS22}. We leave this open.

\begin{remark}[questions]
(i) Compute $\VK_3$, even numerically: the limit game has $2K-1=5$ state variables after the invariances. (ii) Does $\lim_K\VK_K/\sqrt K$ exist, and is it larger than
$\kap$? (iii) Is there a closed form for $\VK_2$, and is the two-point prior least favorable on $\R^2$? (iv) Is there an anytime policy with the optimal constant?
\end{remark}


\newcommand{\gSmp}{\sy{g:not}{\Lambda_K}}
\newcommand{\gGp}{\sy{g:not}{\Delta}}
\newcommand{\gsp}[1]{\sy{g:not}{\|#1\|_{\mathrm{sp}}}}
\newcommand{\ggh}{\sy{g:not}{g}}
\newcommand{\gFh}{\sy{g:not}{F}}
\newcommand{\gPf}{\sy{g:not}{\mathcal P_{\mathrm f}}}
\newcommand{\gOs}{\sy{g:pol}{\Omega^\circ}}
\newcommand{\gPh}{\sy{g:pol}{\mathbb P}}
\newcommand{\gEh}{\sy{g:pol}{\mathbb E}}
\newcommand{\gWi}{\sy{g:pol}{\mathbb W}}
\newcommand{\gDc}{\sy{g:pol}{\mathcal D}}
\newcommand{\gBc}{\sy{g:pol}{\mathcal D^\flat}}
\newcommand{\gQp}[1]{\sy{g:pol}{Q^{#1}}}
\newcommand{\gG}{\sy{g:cell}{\mathcal G}}
\newcommand{\ggam}{\sy{g:cell}{\gamma}}
\newcommand{\gMc}{\sy{g:mart}{\mathcal M}}
\newcommand{\gPhi}{\sy{g:act}{\Phi}}
\newcommand{\gTc}{\sy{g:act}{T}}
\newcommand{\gRs}{\sy{g:act}{R}}
\newcommand{\gPs}{\sy{g:act}{P}}
\newcommand{\gJ}{\sy{g:act}{J}}

\section{The limit game}\label{sec:app-game}


\medskip
\noindent This appendix proves the facts on the limit game of Section~\ref{sec:game} that the body and the other appendices use:
\begin{itemize}
\item the class $\Wc$ of limit policies is convex and compact, the risk $(Q,h)\mapsto\rk(Q,h)$ is jointly continuous, and the minima in the definitions of $\VK_K$,
$\VK_K(H)$ and $\BK_K(\Pi)$ are attained (Lemmas~\ref{lem:g-compact}, \ref{lem:g-cont}, \ref{lem:g-min});
\item discrete randomized policies and batched policies are limit policies with the same risk (Lemma~\ref{lem:g-disc}), and both are dense in $\Wc$
(Lemmas~\ref{lem:g-step} and~\ref{lem:g-dense}); a $T$-round policy for Gaussian rewards is a limit policy after scaling (Lemma~\ref{lem:embed});
\item Proposition~\ref{prop:four}(a), (b) and (d) (Section~\ref{sec:app-g-proof}), and the fact that continuous batched policies suffice (Corollary~\ref{cor:g-cont}).
\end{itemize}
Part~(c) of Proposition~\ref{prop:four} is not proved here (Remark~\ref{rem:g-c}).

\paragraph{Skeleton.} \emph{(a)} is Ky Fan's minimax theorem \cite{Fan53}. It needs a compact convex set of policies on which the risk is lower semicontinuous and
affine, and nothing on the side of the priors, so $h$ may range over the non-compact set $\R^K$. Compactness is Lemma~\ref{lem:g-compact}; continuity and affinity are
Lemma~\ref{lem:g-cont}.
\emph{(b)} A batched policy is a limit policy with the same risk function (Lemma~\ref{lem:g-disc}), so $\VK_K(H)\le\Vb(H)$. Conversely, a minimax policy on $H$ is a
weak limit of discrete policies (Lemma~\ref{lem:g-step}), and their risks converge uniformly on the compact $H$ (Lemma~\ref{lem:g-min}). A discrete policy reads the
whole observed path, a batched policy only its values at the decision times; the two have the same risk functions (Lemma~\ref{lem:g-bridge}), because given these
values the path in between is made of Brownian bridges whose law does not depend on the instance, and the batched policy can simulate them. The identity
$\VK_K=\Vloc_K$ then follows from compactness.
\emph{(d)} Translations and relabellings of the arms act on $\Wc$ (Lemma~\ref{lem:g-act}). Averaging a minimax policy over translations (Hunt--Stein,
Lemma~\ref{lem:g-avg}), then over relabellings, gives an invariant minimax policy. A $(K+1)$-armed limit policy plays a $K$-armed game by treating arm $K+1$ as a
second copy of arm~$1$, which gives $\VK_K\le\VK_{K+1}$.

\paragraph{Dependencies.}
\begin{center}\small
\begin{tabular}{l}
Lem.~\ref{lem:g-like} $\to$ Lem.~\ref{lem:g-compact}, \ref{lem:g-cont} $\to$ Lem.~\ref{lem:g-min} $\xrightarrow{\ \text{Fan}\ }$ Prop.~\ref{prop:four}(a);\\[1mm]
Lem.~\ref{lem:g-compact} $\to$ Lem.~\ref{lem:g-disc} $\to$ Lem.~\ref{lem:embed};\quad
Lem.~\ref{lem:g-disc} $\to$ Lem.~\ref{lem:g-bridge};\\[1mm]
Lem.~\ref{lem:g-disc} $\to$ Lem.~\ref{lem:g-real} $\to$ Lem.~\ref{lem:g-step} $\to$ Lem.~\ref{lem:g-dense} $\to$ Cor.~\ref{cor:g-cont};\\[1mm]
Lem.~\ref{lem:g-min}, \ref{lem:g-step}, \ref{lem:g-bridge} $\to$ Prop.~\ref{prop:four}(b);\quad
Lem.~\ref{lem:g-like}, \ref{lem:g-compact} $\to$ Lem.~\ref{lem:g-act} $\to$ Lem.~\ref{lem:g-avg} $\to$ Prop.~\ref{prop:four}(d).
\end{tabular}
\end{center}

\subsection{Setting and policies}\label{sec:app-g-setting}

\paragraph{Notation.}\dt{g:not} Throughout, $K\ge2$ is fixed, except in Lemma~\ref{lem:g-act}(iii) and in the proof of Proposition~\ref{prop:four}(d), where a
superscript $^{(K)}$ indicates the number of arms. We use the objects $\Om$, $\Al$, $(X,\tau)$, $(\mathcal F_t)$, $\Wc$, $\lh_h$, $\rk$, $\VK_K$, $\VK_K(H)$, $\BK_K$,
$\Vb$, $\Vloc_K$ of Section~\ref{sec:game}. In addition:
\begin{itemize}
\item $\gSmp:=\{p\in[0,1]^K:\sum_kp_k=1\}$ is the simplex; $e_k$ is the $k$-th unit vector and $\1=(1,\dots,1)$;
\item $|h|$ is the Euclidean norm of $h\in\R^K$ and $|h|_1$ its $\ell^1$ norm; for a path $x$, $\|x\|_{[0,t]}:=\sup_{s\le t}|x(s)|$;
\item $\gGp_k(h):=\max_ih_i-h_k$ is the gap of arm $k$ and $\gsp h:=\max_kh_k-\min_kh_k$ the spread of $h$;
\item for $(x,a)\in\Om$, $\ggh_h(a):=\sum_k\gGp_k(h)\,a_k(1)$ is the regret and $\gFh_h:=\lh_h\,\ggh_h$, so that $\rk(Q,h)=\E_Q[\gFh_h]$;
\item $\lh_h(t):=\exp\big(\langle h,X(t)\rangle-\frac12\sum_kh_k^2\tau_k(t)\big)$ for $t\in[0,1]$, so that $\lh_h=\lh_h(1)$;
\item $\gPf$ is the set of finitely supported probability measures on $\R^K$.
\end{itemize}
Each $a\in\Al$ satisfies $0\le a_k(t)-a_k(s)\le t-s$ for $s\le t$, so $\Al$ is compact (Arzelà--Ascoli) and $\Om$ is Polish. The set $\mathcal P(\Om)$ of Borel
probability measures carries the weak topology; it is metrizable and separable, so compactness and closures can be checked along sequences.

\paragraph{Two standard facts.}
\begin{itemize}
\item[(UA)] \emph{Usual augmentation.} For $Q\in\mathcal P(\Om)$ let $(\mathcal F^Q_t)$ be the $Q$-completion of $(\bigcap_{s>t}\mathcal F_s)_{t<1}$, $\mathcal F^Q_1$ the
completion of $\mathcal F_1$. An $(\mathcal F_t)$-adapted integrable process $M$ with continuous paths is an $(\mathcal F_t,Q)$-martingale if and only if it is an
$(\mathcal F^Q_t,Q)$-martingale. ``If'' is the tower property. ``Only if'': for $s<t$ and $A\in\bigcap_{v>s}\mathcal F_v$, $\E_Q[(M_t-M_{s+1/n})\1_A]=0$ for $n$ large;
$M_{s+1/n}=\E_Q[M_t\mid\mathcal F_{s+1/n}]$ is uniformly integrable and tends to $M_s$ pointwise, hence in $L^1$; and null sets do not matter. All the stochastic calculus
below (brackets, Burkholder--Davis--Gundy, Novikov, Girsanov, martingale representation) is done on $(\Om,(\mathcal F^Q_t),Q)$.
\item[(V)] \emph{Vitali under weak convergence.} Let $\Theta$ be a metric space, $f:\Theta\times\Om\to\R$ continuous, $\theta_n\to\theta$ in $\Theta$, and $Q_n\to Q$ weakly with
$\sup_n\E_{Q_n}[f(\theta_n,\cdot)^2]<\infty$. Then $\E_Q[f(\theta,\cdot)^2]<\infty$ and $\E_{Q_n}[f(\theta_n,\cdot)]\to\E_Q[f(\theta,\cdot)]$. Indeed, by Skorokhod's
representation theorem there are random elements $\omega_n\sim Q_n$ and $\omega\sim Q$ on one probability space with $\omega_n\to\omega$ a.s.; then
$f(\theta_n,\omega_n)\to f(\theta,\omega)$ a.s.; Fatou's lemma gives the first claim, and Vitali's theorem (the family is bounded in $L^2$) the second.
\end{itemize}

\paragraph{Policies.} Section~\ref{sec:game} defines the limit policies as laws under the reference instance $h=0$. The policies below are defined on a space that
carries the arms themselves, and have one risk function for every instance.

\begin{definition}[strong model; discrete and batched policies]\label{def:g-pol}\dt{g:pol}
\uses{def:risk, def:alloc}
Let $\gOs:=[0,1]\times C([0,\infty),\R)^K$ with coordinates $(U,Y_1,\dots,Y_K)$ and, for $h\in\R^K$, $\gPh_h:=\mathrm{Leb}\otimes\bigotimes_k\gWi_{h_k}$, where $\gWi_c$
is the law of $(cs+W(s))_{s\ge0}$ for a standard Brownian motion $W$; $\gEh_h$ is the expectation under $\gPh_h$. Under $\gPh_h$, arm $k$ observed during an own time
$s$ shows $Y_k(s)$, a Brownian motion with drift $h_k$, and $U$ is an independent uniform randomization.
\begin{enumerate}
\item[(a)] A \emph{discrete policy} $\pi$ is a grid $0=t_0<t_1<\dots<t_m=1$ and Borel maps $F_i:[0,1]\times C([0,t_i],\R^K)\to\gSmp$, $0\le i<m$. It acts recursively:
$\tau(0):=0$ and, for $t\in[t_i,t_{i+1}]$,
\[
p_i:=F_i\big(U,X|_{[0,t_i]}\big),\qquad \tau(t):=\tau(t_i)+(t-t_i)\,p_i,\qquad X_k(t):=Y_k(\tau_k(t)).
\]
Its risk is $\rk(\pi,h):=\gEh_h[\ggh_h(\tau)]=\gEh_h\big[\sum_k\gGp_k(h)\tau_k(1)\big]$, and $\gQp\pi\in\mathcal P(\Om)$ is the law of $(X,\tau)$ under $\gPh_0$.
$\gDc$ denotes the set of discrete policies.
\item[(b)] A \emph{batched policy} is a discrete policy whose maps have the form $F_i(u,x)=f_i\big(u,x(t_1),\dots,x(t_i)\big)$ for Borel maps
$f_i:[0,1]\times(\R^K)^i\to\gSmp$. $\gBc\subset\gDc$ denotes the set of batched policies.
\item[(c)] A discrete (resp.\ batched) policy has \emph{continuous rules} if every $F_i$ (resp.\ every $f_i$) is continuous.
\end{enumerate}
\end{definition}

\noindent The recursion is well defined: $(X,\tau)$ on $[0,t_i]$ only uses $p_0,\dots,p_{i-1}$, so $X|_{[0,t_i]}$ is known when $p_i$ is chosen (for $i=0$ it is the
constant path $0$). The allocation $\tau$ lies in $\Al$, and $X_k(t_i)=Y_k(\tau_k(t_i))$ is the cumulative observation of arm $k$ at the decision time $t_i$.

\begin{remark}[two readings of ``batched'']\label{rem:g-reading}
In Section~\ref{sec:game}, the proportions of a batched policy are functions of ``the randomization and the observations so far''. Definition~\ref{def:g-pol}(b)
reads ``observations so far'' as the values $X(t_1),\dots,X(t_i)$. Section~\ref{sec:game} and the core of \SC\ (Section~\ref{sec:app-upper}) read the whole path $X|_{[0,t_i]}$:
with this reading the batched policies are the discrete policies. Both readings give the same risk functions (Lemma~\ref{lem:g-bridge}), hence the same value $\Vb(H)$.
\end{remark}

\subsection{Why is $\Wc$ compact, and why is the risk continuous?}\label{sec:app-g-compact}

\begin{lemma}[likelihoods]\label{lem:g-like}
\uses{def:alloc, def:risk}
\textit{Role.} Finiteness and uniform integrability of the risk (Lemma~\ref{lem:g-cont}); change of instance (Lemma~\ref{lem:g-act}).

\noindent\textit{Setting.} $Q\in\Wc$, $h\in\R^K$, $p\ge1$.

\noindent Then:
\begin{enumerate}
\item[(a)] $(\lh_h(t))_{t\in[0,1]}$ is a continuous $Q$-martingale, and $\E_Q[\lh_h]=1$;
\item[(b)] $\E_Q[\lh_h^p]\le e^{p(p-1)|h|^2/2}$; in particular $\E_Q[\lh_h^2]\le e^{|h|^2}$;
\item[(c)] under the probability $\lh_h\cdot Q$, the processes $X_k-h_k\tau_k$ are continuous $(\mathcal F^Q_t)$-martingales with
$\langle X_k-h_k\tau_k,X_j-h_j\tau_j\rangle=\1\{k=j\}\tau_k$.
\end{enumerate}
\end{lemma}
\begin{proof}
\uses{def:risk}
(a) $M:=\langle h,X\rangle$ is a continuous $Q$-martingale with $\langle M\rangle_t=\sum_kh_k^2\tau_k(t)\le|h|^2$. Hence $\E_Q[e^{\langle M\rangle_1/2}]\le e^{|h|^2/2}<\infty$,
and by Novikov's criterion \cite[Cor.~3.5.13]{KS91} the stochastic exponential $e^{M_t-\langle M\rangle_t/2}$, which is $\lh_h(t)$, is a martingale with mean $1$.

(b) We bound $\lh_h^p$ by a likelihood of the same form:
\begin{align*}
\lh_h^p&=\lh_{ph}\cdot\exp\Big(\frac{p^2-p}2\sum_kh_k^2\tau_k(1)\Big) &&\text{(expand the exponent)}\\
&\le\lh_{ph}\,e^{p(p-1)|h|^2/2} &&\text{(}p\ge1\text{, }\tau_k(1)\le1\text{)},
\end{align*}
and $\E_Q[\lh_{ph}]=1$ by (a).

(c) By (a), $\lh_h\cdot Q$ is a probability with density process $\lh_h(t)=e^{M_t-\langle M\rangle_t/2}$. By Girsanov's theorem \cite[Thm~VIII.1.4]{RY99}, for each $k$,
$X_k-\langle X_k,M\rangle=X_k-h_k\tau_k$ is a continuous local martingale under $\lh_h\cdot Q$, and brackets do not change under an equivalent change of measure. The
bracket of $X_k-h_k\tau_k$ is bounded by $1$, so it is a true martingale (Burkholder--Davis--Gundy).
\end{proof}

\noindent By (c), $\lh_h\cdot Q$ describes the same limit policy when arm $k$ has drift $h_k$; for discrete policies this is proved exactly in
Lemma~\ref{lem:g-disc}(c). This is the meaning of the formula for $\rk(Q,h)$.

\begin{lemma}[$\Wc$ is convex and compact]\label{lem:g-compact}
\uses{def:alloc}
\textit{Role.} The compact side of Fan's theorem in Proposition~\ref{prop:four}(a); existence of the minima (Lemma~\ref{lem:g-min}); limit points in the proofs of (b) and (d).

\noindent\textit{Setting.}\dt{g:mart} $\gMc:=\{X_k,\ X_k^2-\tau_k,\ X_kX_j:\ k\ne j\in[K]\}$, a set of processes on $\Om$.

\noindent Then:
\begin{enumerate}
\item[(i)] $Q\in\mathcal P(\Om)$ belongs to $\Wc$ if and only if, for every $M\in\gMc$, every $0\le u\le t\le1$ and every bounded $\mathcal F_u$-measurable $Z$,
$\E_Q|M_t|<\infty$ and $\E_Q[(M_t-M_u)Z]=0$;
\item[(ii)] there is a universal constant $C_4$ such that $\E_Q[(X_k(t)-X_k(u))^4]\le C_4\,(t-u)^2$ for every $Q\in\Wc$, $k\in[K]$ and $u\le t$;
\item[(iii)] $\Wc$ is convex and compact.
\end{enumerate}
\end{lemma}
\begin{proof}
\uses{def:alloc}
\textit{Idea.} The martingale conditions are linear in $Q$ and pass to weak limits under uniform moment bounds; tightness comes from the
Burkholder--Davis--Gundy inequality. The limit is a martingale for the filtration of $(X,\tau)$ jointly, which is exactly the definition of $\Wc$.

Step 1: \emph{the characterization (i).} Let $Q\in\Wc$. On $(\Om,(\mathcal F^Q_t),Q)$, the processes $X_k^2-\tau_k$ and $X_kX_j$ ($k\ne j$) are continuous local
martingales, by definition of the brackets. By the Burkholder--Davis--Gundy inequality \cite[Thm~3.3.28]{KS91}, $\E_Q[\sup_t|X_k(t)|^4]\le C_4\E_Q[\tau_k(1)^2]\le C_4$.
So every $M\in\gMc$ is dominated by the integrable variable $\sum_k\sup_tX_k(t)^2+2$, hence is a true martingale; it is $(\mathcal F_t)$-adapted, so it is an
$(\mathcal F_t)$-martingale (fact (UA)). Conversely, if the conditions hold, the processes of $\gMc$ are continuous $(\mathcal F^Q_t)$-martingales (fact (UA)); since
$\tau_k$ is continuous, adapted and nondecreasing, uniqueness of the bracket gives $\langle X_k\rangle=\tau_k$ and $\langle X_k,X_j\rangle=0$.

Step 2: \emph{moments (ii).} Let $Q\in\Wc$ and $u\le t$. The process $N(s):=X_k(u+s)-X_k(u)$, $s\in[0,t-u]$, is a continuous martingale with
$\langle N\rangle_s=\tau_k(u+s)-\tau_k(u)\le s$. The Burkholder--Davis--Gundy inequality gives $\E_Q[N(t-u)^4]\le C_4\E_Q[\langle N\rangle_{t-u}^2]\le C_4(t-u)^2$.

Step 3: \emph{convexity.} The conditions of (i) are preserved by convex combinations, since $\E_{\lambda Q+(1-\lambda)Q'}[\,\cdot\,]=\lambda\E_Q[\,\cdot\,]+(1-\lambda)\E_{Q'}[\,\cdot\,]$.

Step 4: \emph{tightness.} By (ii) and $|y|^4\le K\sum_ky_k^4$ on $\R^K$, $\E_Q|X(t)-X(u)|^4\le K^2C_4(t-u)^2$ for every $Q\in\Wc$, and $X(0)=0$. The Kolmogorov--Chentsov
criterion \cite[Problem~2.4.11]{KS91} makes the laws of $X$ under $Q\in\Wc$ tight in $C_0([0,1],\R^K)$; $\tau$ takes values in the compact $\Al$. So $\Wc$ is tight, hence
relatively compact by Prokhorov's theorem.

Step 5: \emph{closedness.} Let $Q_n\in\Wc$ converge weakly to $Q$. Fix $M\in\gMc$, $u\le t$, and $Z:=\zeta(\omega(\cdot\wedge u))$ with $\zeta$ bounded and continuous on
$\Om$, where $\omega(\cdot\wedge u):=(x(\cdot\wedge u),a(\cdot\wedge u))$. The function $(M_t-M_u)Z$ is continuous on $\Om$, because $M_t$ is a polynomial in $(x(t),a(t))$.
By (ii) with $u=0$ and Cauchy--Schwarz, $\sup_n\E_{Q_n}[(M_t-M_u)^2Z^2]\le\|\zeta\|_\infty^2\sup_n\E_{Q_n}[(M_t-M_u)^2]<\infty$. Fact (V) gives, first, $\E_Q[M_t^2]<\infty$,
and second,
\[
\E_Q[(M_t-M_u)Z]=\lim_n\E_{Q_n}[(M_t-M_u)Z]=0 .
\]
Since $M_t-M_u\in L^1(Q)$, the bounded $\mathcal F_u$-measurable $Z$ with $\E_Q[(M_t-M_u)Z]=0$ form a vector space closed under bounded monotone limits. It contains the
multiplicative class $\{\zeta(\omega(\cdot\wedge u))\}$, which generates $\mathcal F_u=\sigma(\omega(\cdot\wedge u))$; by the monotone class theorem it contains all bounded
$\mathcal F_u$-measurable $Z$. By (i), $Q\in\Wc$.
\end{proof}

\begin{lemma}[the risk is continuous and affine]\label{lem:g-cont}
\uses{def:alloc, def:risk}
\textit{Role.} Lower semicontinuity and affinity needed by Fan's theorem; uniform convergence of the risks on compact sets (Lemma~\ref{lem:g-min}).

\noindent\textit{Setting.} $Q,Q'\in\Wc$, $h\in\R^K$, $\lambda\in[0,1]$.

\noindent Then $0\le\rk(Q,h)\le\gsp h$, the map $(Q,h)\mapsto\rk(Q,h)$ is continuous on $\Wc\times\R^K$, and
$\rk(\lambda Q+(1-\lambda)Q',h)=\lambda\rk(Q,h)+(1-\lambda)\rk(Q',h)$.
\end{lemma}
\begin{proof}
\uses{lem:g-like, def:risk, def:alloc}
\textit{Bound.} Since $0\le\ggh_h(\tau)\le\gsp h\sum_k\tau_k(1)=\gsp h$ and $\E_Q[\lh_h]=1$ (Lemma~\ref{lem:g-like}(a)), $0\le\rk(Q,h)\le\gsp h$.

\textit{Continuity.} The map $(h,\omega)\mapsto\gFh_h(\omega)$ is continuous on $\R^K\times\Om$: it only depends on $x(1)$ and $a(1)$. Let $(Q_n,h_n)\to(Q,h)$. By
Lemma~\ref{lem:g-like}(b),
\[
\sup_n\E_{Q_n}\big[\gFh_{h_n}^2\big]\le\sup_n\gsp{h_n}^2\,e^{|h_n|^2}<\infty,
\]
so fact (V), with $\Theta=\R^K$, gives $\rk(Q_n,h_n)=\E_{Q_n}[\gFh_{h_n}]\to\E_Q[\gFh_h]=\rk(Q,h)$.

\textit{Affinity.} $\rk(\cdot,h)=\E_\cdot[\gFh_h]$ is the integral of a fixed nonnegative function.
\end{proof}

\begin{lemma}[the minima are attained]\label{lem:g-min}
\uses{def:risk, def:alloc}
\textit{Role.} Gives a meaning to the minima in the definitions of $\VK_K$, $\VK_K(H)$ and $\BK_K(\Pi)$ (Section~\ref{sec:game}); continuity in the proofs of
Proposition~\ref{prop:four}(b) and Corollary~\ref{cor:g-cont}.

\noindent\textit{Setting.} $H\subset\R^K$ nonempty; $\Pi$ a Borel probability measure on $\R^K$.

\noindent Then:
\begin{enumerate}
\item[(i)] $Q\mapsto\sup_{h\in H}\rk(Q,h)$ is lower semicontinuous on $\Wc$; if $H$ is compact, then $\sup_{h\in H}|\rk(Q_n,h)-\rk(Q,h)|\to0$ whenever $Q_n\to Q$ in $\Wc$, so
this map is continuous and finite;
\item[(ii)] $Q\mapsto\int\rk(Q,h)\,\Pi(dh)\in[0,\infty]$ is affine and lower semicontinuous on $\Wc$; if $\int\gsp h\,\Pi(dh)<\infty$, it is finite and continuous;
\item[(iii)] hence the minima in $\VK_K(H):=\min_Q\sup_{h\in H}\rk(Q,h)$ (in particular $\VK_K=\VK_K(\R^K)$) and in $\BK_K(\Pi):=\min_Q\int\rk(Q,h)\,\Pi(dh)$ are attained.
\end{enumerate}
\end{lemma}
\begin{proof}
\uses{lem:g-cont, lem:g-compact, def:risk, def:alloc}
(i) A supremum of continuous functions (Lemma~\ref{lem:g-cont}) is lower semicontinuous. Let $H$ be compact and $Q_n\to Q$. If the uniform convergence failed, there would be
$\varepsilon>0$ and $h_n\in H$ with $|\rk(Q_n,h_n)-\rk(Q,h_n)|\ge\varepsilon$ along a subsequence; extracting $h_n\to h\in H$, both $\rk(Q_n,h_n)$ and $\rk(Q,h_n)$ tend to
$\rk(Q,h)$ by joint continuity (Lemma~\ref{lem:g-cont}), a contradiction. Finiteness: $\rk(Q,h)\le\gsp h$ is bounded on $H$.

(ii) $h\mapsto\rk(Q,h)$ is continuous, hence measurable, and nonnegative. Affinity follows from Lemma~\ref{lem:g-cont}. If $Q_n\to Q$, then $\rk(Q_n,h)\to\rk(Q,h)$ for every
$h$, and Fatou's lemma gives $\int\rk(Q,h)\Pi(dh)\le\liminf_n\int\rk(Q_n,h)\Pi(dh)$. If $\int\gsp h\,\Pi(dh)<\infty$, dominated convergence (with $\rk\le\gsp h$) gives
convergence instead.

(iii) A lower semicontinuous function with values in $[0,\infty]$ attains its minimum on the compact set $\Wc$ (Lemma~\ref{lem:g-compact}).
\end{proof}

\subsection{Which policies are limit policies?}\label{sec:app-g-which}

\begin{lemma}[discrete policies are limit policies]\label{lem:g-disc}
\uses{def:g-pol, def:alloc, def:risk}
\textit{Role.} Puts the discrete and batched policies inside the game, with the same risk functions; used in Lemmas~\ref{lem:embed}, \ref{lem:g-real},
\ref{lem:g-dense}, \ref{lem:g-bridge} and in Proposition~\ref{prop:four}(b).

\noindent\textit{Setting.}\dt{g:cell} $\pi\in\gDc$ with grid $(t_i)_{i\le m}$ and maps $(F_i)$; $h\in\R^K$. Let $\gG_i:=\sigma\big(U,X|_{[0,t_i]}\big)$, let
$Y^{(i)}_k(s):=Y_k(\tau_k(t_i)+s)-Y_k(\tau_k(t_i))$ for $s\ge0$, let $b_i:=(t_{i+1}-t_i)\,p_i$ be the own times given to the arms in cell $i$, and let
$c_i:=\big(X(t_i+s(t_{i+1}-t_i))-X(t_i)\big)_{s\in[0,1]}$ be the cell path. For $b\in[0,1]^K$, $\ggam_h(b,\cdot)$ is the law on $C([0,1],\R^K)$ of
$\big(b_kh_ks+\sqrt{b_k}\,W_k(s)\big)_{s\in[0,1],k\in[K]}$, for independent standard Brownian motions $W_k$, and $\ggam:=\ggam_0$.

\noindent Then:
\begin{enumerate}
\item[(a)] under $\gPh_h$, conditionally on $\gG_i$, the processes $Y^{(i)}_1,\dots,Y^{(i)}_K$ are independent with laws $\gWi_{h_1},\dots,\gWi_{h_K}$; hence the conditional law of
$c_i$ given $\gG_i$ is $\ggam_h(b_i,\cdot)$;
\item[(b)] $\gQp\pi\in\Wc$;
\item[(c)] for every bounded measurable $\Psi$ on $[0,1]\times C([0,1],\R^K)$, $\gEh_h[\Psi(U,X)]=\gEh_0\big[\lh_h(X,\tau)\,\Psi(U,X)\big]$. In particular the law of $(X,\tau)$
under $\gPh_h$ is $\lh_h\cdot\gQp\pi$, and $\rk(\pi,h)=\rk(\gQp\pi,h)$.
\end{enumerate}
\end{lemma}
\begin{proof}
\uses{lem:g-compact, def:g-pol, def:alloc, def:risk}
\textit{Idea.} Within a cell, the proportions are fixed given the past, and a Brownian motion with drift restarts afresh after a fixed time; so each cell adds
independent Brownian pieces, whose likelihood ratio is given by the Cameron--Martin formula.

Step 1: \emph{(a), by induction on $i$.} For $i=0$, $\gG_0=\sigma(U)$ and $Y^{(0)}=Y$ is independent of $U$ with law $\bigotimes_k\gWi_{h_k}$. Assume (a) for some $i<m$. The
proportions $p_i$, hence $b_i$, are $\gG_i$-measurable. On $[t_i,t_{i+1}]$,
\[
X_k(t)=X_k(t_i)+Y^{(i)}_k\big((t-t_i)p_{ik}\big),
\]
so $\gG_{i+1}\subset\gG'_{i+1}:=\gG_i\vee\sigma\big(Y^{(i)}_k(\cdot\wedge b_{ik}):k\in[K]\big)$, and $Y^{(i+1)}_k=Y^{(i)}_k(b_{ik}+\cdot)-Y^{(i)}_k(b_{ik})$. For a fixed
$b\in[0,1]^K$ and $y$ of law $\bigotimes_k\gWi_{h_k}$, the stopped paths $(y_k(\cdot\wedge b_k))_k$ are independent of the shifted increments
$\theta_by:=(y_k(b_k+\cdot)-y_k(b_k))_k$, and $\theta_by$ has law $\bigotimes_k\gWi_{h_k}$ (independent increments). Let $\Phi$ be bounded and $\gG_i$-measurable, and
$\psi_1,\psi_2$ bounded measurable. We show that $Y^{(i+1)}$ is independent of $\gG'_{i+1}$:
\begin{align*}
&\gEh_h\Big[\Phi\,\psi_1\big((Y^{(i)}_k(\cdot\wedge b_{ik}))_k\big)\,\psi_2\big(Y^{(i+1)}\big)\Big]\\
&=\gEh_h\Big[\Phi\int\psi_1\big((y_k(\cdot\wedge b_{ik}))_k\big)\,\psi_2(\theta_{b_i}y)\,\textstyle\bigotimes_k\gWi_{h_k}(dy)\Big] &&\text{((a) for }i\text{; }b_i\text{ is }\gG_i\text{-measurable)}\\
&=\gEh_h\Big[\Phi\int\psi_1\big((y_k(\cdot\wedge b_{ik}))_k\big)\,\textstyle\bigotimes_k\gWi_{h_k}(dy)\Big]\displaystyle\int\psi_2\,d\textstyle\bigotimes_k\gWi_{h_k} &&\text{(independence at the fixed time }b_i\text{)}\\
&=\gEh_h\Big[\Phi\,\psi_1\big((Y^{(i)}_k(\cdot\wedge b_{ik}))_k\big)\Big]\displaystyle\int\psi_2\,d\textstyle\bigotimes_k\gWi_{h_k} &&\text{((a) for }i\text{ again).}
\end{align*}
The products $\Phi\,\psi_1(\cdots)$ of indicators form a $\pi$-system generating $\gG'_{i+1}$. Hence $Y^{(i+1)}$ is independent of $\gG'_{i+1}$ with law
$\bigotimes_k\gWi_{h_k}$, and a fortiori conditionally on $\gG_{i+1}\subset\gG'_{i+1}$. This is (a) for $i+1$. For the cell path, $c_{ik}(s)=Y^{(i)}_k(sb_{ik})$, and
$(y(sb))_{s\in[0,1]}$ has law $\ggam_h$ when $y\sim\gWi_{h_k}$ (Brownian scaling).

Step 2: \emph{(b).} Let $h=0$, $t_i\le u\le t\le t_{i+1}$ and $\gG_u:=\sigma(U,X|_{[0,u]})$. Let
\[
D_k:=X_k(t)-X_k(u)=Y^{(i)}_k\big((t-t_i)p_{ik}\big)-Y^{(i)}_k\big((u-t_i)p_{ik}\big).
\]
By Step~1, and the same independence at the $\gG_i$-measurable time $(u-t_i)p_i$, the $D_k$ are independent conditionally on $\gG_u$, with
$D_k\sim\mathcal N(0,(t-u)p_{ik})$. Moreover $(t-u)p_{ik}=\tau_k(t)-\tau_k(u)$ is $\gG_u$-measurable. Hence
\[
\gEh_0[D_k\mid\gG_u]=0,\qquad \gEh_0[D_k^2\mid\gG_u]=\tau_k(t)-\tau_k(u),\qquad \gEh_0[D_kD_j\mid\gG_u]=0\quad(k\ne j).
\]
Chaining over cells, $\gEh_0[X_k(t)^2]=\gEh_0[\tau_k(t)]\le1$, so all products below are integrable. Since
\begin{gather*}
X_k^2(t)-\tau_k(t)-\big(X_k^2(u)-\tau_k(u)\big)=D_k^2+2X_k(u)D_k-\big(\tau_k(t)-\tau_k(u)\big),\\
X_kX_j(t)-X_kX_j(u)=D_kD_j+X_k(u)D_j+X_j(u)D_k,
\end{gather*}
every $M\in\gMc$ satisfies $\gEh_0[M_t-M_u\mid\gG_u]=0$ for $u,t$ in one cell, and, by the tower property, for all $u\le t$. The canonical filtration of $(X,\tau)$ is
contained in $(\gG_u)$, because $\tau|_{[0,u]}$ is a function of $(U,X|_{[0,u]})$. By Lemma~\ref{lem:g-compact}(i), $\gQp\pi\in\Wc$.

Step 3: \emph{(c).} Let $\lh_h(t)$ be computed on $(X,\tau)$ and $H_i:=(U,X|_{[0,t_i]})$. We show by induction on $i$ that
\[
\gEh_h[\Psi(H_i)]=\gEh_0\big[\lh_h(t_i)\,\Psi(H_i)\big]\qquad\text{for every bounded measurable }\Psi .
\]
For $i=0$, $\lh_h(0)=1$ and $U$ has the same law under $\gPh_h$ and $\gPh_0$. By the Cameron--Martin theorem, $\ggam_h(b,\cdot)$ has density
$\exp\big(\langle h,c(1)\rangle-\frac12\sum_kh_k^2b_k\big)$ with respect to $\ggam(b,\cdot)$ at the path $c$. Write a bounded measurable function of $H_{i+1}$ as
$\tilde\Psi(H_i,c_i)$, and let $\Lambda_i(c):=\exp\big(\langle h,c(1)\rangle-\frac12\sum_kh_k^2b_{ik}\big)$. The integral in the second line below is a function of $H_i$
(since $b_i$ is), which allows the induction hypothesis in the third line. Then
\begin{align*}
\gEh_h\big[\tilde\Psi(H_i,c_i)\big]&=\gEh_h\Big[\int\tilde\Psi(H_i,c)\,\ggam_h(b_i,dc)\Big] &&\text{((a) under }\gPh_h\text{)}\\
&=\gEh_h\Big[\int\tilde\Psi(H_i,c)\,\Lambda_i(c)\,\ggam(b_i,dc)\Big] &&\text{(Cameron--Martin)}\\
&=\gEh_0\Big[\lh_h(t_i)\int\tilde\Psi(H_i,c)\,\Lambda_i(c)\,\ggam(b_i,dc)\Big] &&\text{(induction hypothesis)}\\
&=\gEh_0\big[\lh_h(t_i)\,\Lambda_i(c_i)\,\tilde\Psi(H_i,c_i)\big] &&\text{((a) under }\gPh_0\text{)}\\
&=\gEh_0\big[\lh_h(t_{i+1})\,\tilde\Psi(H_i,c_i)\big] &&\text{(definition of }\lh_h(t)\text{)},
\end{align*}
where the last line uses $c_i(1)=X(t_{i+1})-X(t_i)$ and $b_i=\tau(t_{i+1})-\tau(t_i)$. At $i=m$ this is the claim, since $\lh_h(1)=\lh_h(X,\tau)$. As $\tau$ is a function of $(U,X)$, the law of $(X,\tau)$ under $\gPh_h$ is $\lh_h\cdot\gQp\pi$, and
$\rk(\pi,h)=\gEh_h[\ggh_h(\tau)]=\E_{\gQp\pi}[\lh_h\ggh_h]=\rk(\gQp\pi,h)$.
\end{proof}

\noindent Lemma~\ref{lem:g-disc} applies to batched policies, since $\gBc\subset\gDc$: the rule $f_i(U,X(t_1),\dots,X(t_i))$ is a Borel function of $(U,X|_{[0,t_i]})$,
evaluation maps being continuous. So every batched policy $\pi$ has $\gQp\pi\in\Wc$ and $\rk(\pi,\cdot)=\rk(\gQp\pi,\cdot)$.

\begin{lemma}[$T$-round Gaussian policies are limit policies]\label{lem:embed}
\uses{def:g-pol, def:alloc, def:risk, def:regret}
\textit{Role.} The limit game contains every $T$-round Gaussian problem: lower half of Proposition~\ref{prop:four}(c), and the Gaussian lower bound of
Corollary~\ref{cor:examples}.

\noindent\textit{Setting.} $T\ge1$, $\mu_0\in\R$, $\sigma>0$. A \emph{$T$-round policy} $\pi$ is a family of Borel maps $\pi_n:[0,1]\times([K]\times\R)^{n-1}\to[K]$,
$n\le T$; with $U$ uniform on $[0,1]$ it plays $A_n:=\pi_n(U,A_1,R_1,\dots,A_{n-1},R_{n-1})$, where $R_n$ is the reward of round $n$. For $h\in\R^K$ let
$\nu^{T,h}:=\big(\mathcal N(\mu_0+\sigma h_k/\sqrt T,\sigma^2)\big)_{k\in[K]}$, and let $N_k(T)$ be the number of pulls of arm $k$.

\noindent Then there is a batched policy $\pi^\flat$ with grid $t_i:=i/T$ such that $Q_\pi:=\gQp{\pi^\flat}\in\Wc$ satisfies, for every $h\in\R^K$:
\begin{enumerate}
\item[(i)] the law of $N(T)/T$ under $\pi$ and $\nu^{T,h}$ is the law of $\tau(1)$ under $\lh_h\cdot Q_\pi$;
\item[(ii)] $\RT(\pi,\nu^{T,h})=\sigma\sqrt T\,\rk(Q_\pi,h)$.
\end{enumerate}
Consequently, for every $H\subset\R^K$ and every Borel probability $\Pi$ on $\R^K$,
\[
\inf_\pi\sup_{h\in H}\RT(\pi,\nu^{T,h})\ \ge\ \sigma\sqrt T\,\VK_K(H),\qquad \inf_\pi\int\RT(\pi,\nu^{T,h})\,\Pi(dh)\ \ge\ \sigma\sqrt T\,\BK_K(\Pi).
\]
\end{lemma}
\begin{proof}
\uses{lem:g-disc, def:g-pol, def:risk, def:regret, def:alloc}
\textit{Idea.} Give each batch of length $1/T$ to the arm chosen by $\pi$, and feed $\pi$ with the scaled increments of the observed path: they form a table of
independent Gaussian rewards, read exactly as $\pi$ reads them.

Step 1: \emph{the batched policy.} For $u\in[0,1]$ and $z_1,\dots,z_i\in\R^K$, set $z_0:=0$, $a_1:=\pi_1(u)$ and recursively, for $j\le i$,
\[
\rho_j:=\mu_0+\sigma\sqrt T\,\big(z_{j,a_j}-z_{j-1,a_j}\big),\qquad a_{j+1}:=\pi_{j+1}(u,a_1,\rho_1,\dots,a_j,\rho_j).
\]
Let $f_i(u,z_1,\dots,z_i):=e_{a_{i+1}}$: a Borel map. Let $\pi^\flat$ be the batched policy with grid $t_i=i/T$ and rules $f_i$, and $A_{i+1}:=a_{i+1}$ computed from
$(U,X(t_1),\dots,X(t_i))$.

Step 2: \emph{the rewards form a reward table of $\nu^{T,h}$.} Run $\pi^\flat$ under $\gPh_h$. Batch $i$ is given entirely to arm $A_{i+1}$, so
$\tau(t_{i+1})=\tau(t_i)+e_{A_{i+1}}/T$. Hence $T\tau_k(t_i)=N_k(i):=\#\{j\le i:A_j=k\}$ and $T\tau(1)=N(T)$. Let
\[
X_{k,j}:=\mu_0+\sigma\sqrt T\,\big(Y_k(j/T)-Y_k((j-1)/T)\big),\qquad k\in[K],\ j\ge1.
\]
During batch $i$ the own clock of $A:=A_{i+1}$ goes from $N_A(i)/T$ to $(N_A(i)+1)/T$, so $X_A(t_{i+1})-X_A(t_i)=Y_A((N_A(i)+1)/T)-Y_A(N_A(i)/T)$, and the reward
$\rho_{i+1}$ that $\pi^\flat$ feeds to $\pi$ is $X_{A,N_A(i)+1}$: the next unread entry of row $A$. Under $\gPh_h$, the increments $Y_k(j/T)-Y_k((j-1)/T)$ are independent
$\mathcal N(h_k/T,1/T)$, so the $X_{k,j}$ are independent with $X_{k,j}\sim\mathcal N(\mu_0+\sigma h_k/\sqrt T,\sigma^2)$, and independent of $U$. So
$(U,A_1,\rho_1,\dots,A_T,\rho_T)$ under $\gPh_h$ has the law of the interaction of $\pi$ with $\nu^{T,h}$ in the reward-table model, which is the law of the interaction
itself \cite[Sec.~4.6]{LS}.

Step 3: \emph{conclusion.} (i) By Step~2, $N(T)/T$ under $\nu^{T,h}$ has the law of $\tau(1)$ under $\gPh_h$, which is the law of $\tau(1)$ under $\lh_h\cdot Q_\pi$ by
Lemma~\ref{lem:g-disc}(c). (ii) The gaps of $\nu^{T,h}$ are $\sigma\gGp_k(h)/\sqrt T$, so
\begin{align*}
\RT(\pi,\nu^{T,h})&=\sum_k\frac{\sigma\gGp_k(h)}{\sqrt T}\,\E\big[N_k(T)\big] &&\text{(definition of }\RT\text{)}\\
&=\sigma\sqrt T\,\gEh_h\Big[\sum_k\gGp_k(h)\tau_k(1)\Big]=\sigma\sqrt T\,\rk(\pi^\flat,h) &&\text{((i): }N(T)=T\tau(1)\text{ in law)}\\
&=\sigma\sqrt T\,\rk(Q_\pi,h) &&\text{(Lemma~\ref{lem:g-disc}(c)).}
\end{align*}
The consequences follow from $Q_\pi\in\Wc$: $\sup_{h\in H}\rk(Q_\pi,h)\ge\VK_K(H)$ and $\int\rk(Q_\pi,h)\Pi(dh)\ge\BK_K(\Pi)$ (the integrand is continuous in $h$).
\end{proof}

\subsection{Why are batched policies enough?}\label{sec:app-g-dense}

\begin{lemma}[realization of piecewise constant controls]\label{lem:g-real}
\uses{def:g-pol}
\textit{Role.} Turns the step controls of Lemma~\ref{lem:g-step} into discrete policies.

\noindent\textit{Setting.} $(\tilde\Omega,\tilde{\mathcal F},\tilde\P)$ carries a filtration $(\tilde{\mathcal G}_t)_{t\in[0,1]}$, a $K$-dimensional
$(\tilde{\mathcal G}_t)$-Brownian motion $B$, a grid $0=t_0<\dots<t_m=1$, and $\tilde{\mathcal G}_{t_i}$-measurable $p_i\in\gSmp$. Let $\tau(0):=0$, $X(0):=0$ and, for
$t\in[t_i,t_{i+1}]$, $\tau(t):=\tau(t_i)+(t-t_i)p_i$ and $X_k(t):=X_k(t_i)+\sqrt{p_{ik}}\,(B_k(t)-B_k(t_i))$.

\noindent Then there is $\pi\in\gDc$ on the same grid with $\gQp\pi=\mathrm{Law}_{\tilde\P}(X,\tau)$.
\end{lemma}
\begin{proof}
\uses{lem:g-disc, def:g-pol}
\textit{Idea.} Reproduce, cell by cell, the conditional law of the next proportions given the past, by a randomization; the cell paths then have the same conditional
law on both sides.

Step 1: \emph{cell kernel.} Let $c_i:=(X(t_i+s(t_{i+1}-t_i))-X(t_i))_{s\in[0,1]}$ and $b_i:=(t_{i+1}-t_i)p_i$. Since $B(t_i+\cdot)-B(t_i)$ is a Brownian motion independent
of $\tilde{\mathcal G}_{t_i}$ and $p_i$ is $\tilde{\mathcal G}_{t_i}$-measurable, the conditional law of $c_i$ given $\tilde{\mathcal G}_{t_i}$ is $\ggam(b_i,\cdot)$ (Brownian
scaling), the kernel of Lemma~\ref{lem:g-disc}(a).

Step 2: \emph{decision kernels.} Let $H_i:=(X|_{[0,t_i]},p_0,\dots,p_{i-1})$, a random element of a Polish space, and let $\kappa_i(\cdot\mid H_i)$ be a regular conditional
law of $p_i$ given $H_i$. By the randomization lemma \cite[Lemma~3.22]{Kal02} there are Borel maps $G_i$ with $G_i(\vartheta,\eta)\sim\kappa_i(\cdot\mid\eta)$ for $\vartheta$
uniform on $[0,1]$. Let $U\mapsto(U_0,\dots,U_{m-1})$ be a Borel map sending the uniform law to the law of i.i.d.\ uniforms. Let $\pi$ play
$p'_i:=G_i\big(U_i,X'|_{[0,t_i]},p'_0,\dots,p'_{i-1}\big)$, where $X'$ is its observed path; recursively, $p'_i$ is a Borel function of $(U,X'|_{[0,t_i]})$, so
$\pi\in\gDc$.

Step 3: \emph{induction.} Let $H'_i:=(X'|_{[0,t_i]},p'_0,\dots,p'_{i-1})$ under $\gPh_0$. Suppose $H'_i$ and $H_i$ have the same law (true for $i=0$). First,
$(H'_i,p'_i)$ and $(H_i,p_i)$ have the same law: $U_i$ is independent of $H'_i$ (a function of $U_0,\dots,U_{i-1}$ and $Y$), so the conditional law of $p'_i$ given $H'_i$ is
$\kappa_i(\cdot\mid H'_i)$. Second, adding the cell path keeps the equality: given $(H_i,p_i)$, which is $\tilde{\mathcal G}_{t_i}$-measurable, $c_i$ has law $\ggam(b_i,\cdot)$
by Step~1; given $(H'_i,p'_i)$, which is $\sigma(U,X'|_{[0,t_i]})$-measurable, the cell path of $\pi$ has law $\ggam(b'_i,\cdot)$ by Lemma~\ref{lem:g-disc}(a). Since
$H_{i+1}$ and $H'_{i+1}$ are the same function of (previous history, proportions, cell path), they have the same law. At $i=m$, $(X,\tau)$ and $(X',\tau')$ are the same
function of $H_m$ and $H'_m$, so $\gQp\pi=\mathrm{Law}_{\tilde\P}(X,\tau)$.
\end{proof}

\begin{lemma}[discrete policies are dense]\label{lem:g-step}
\uses{def:alloc, def:g-pol}
\textit{Role.} Density of discrete policies, used in Proposition~\ref{prop:four}(b) and in Lemma~\ref{lem:g-dense}.

\noindent\textit{Setting.} $Q\in\Wc$.

\noindent Then there are $\pi_n\in\gDc$ (with grids $(i/n)_{i\le n}$) such that $\gQp{\pi_n}\to Q$ weakly. In particular $\Wc$ is the weak closure of
$\{\gQp\pi:\pi\in\gDc\}$.
\end{lemma}
\begin{proof}
\uses{lem:g-real, lem:g-disc, lem:g-compact, def:alloc, def:g-pol}
\textit{Idea.} Represent $X$ as a stochastic integral $\int\sqrt u\,dB$ against a Brownian motion; replace the rate $u$ by its average over the previous cell of a grid,
which is adapted and piecewise constant; keep the same $B$; realize the result as a discrete policy (Lemma~\ref{lem:g-real}).

Step 1: \emph{rate and representation.} Each $\tau_k$ is nondecreasing and $\sum_k\tau_k(t)=t$, so $0\le\tau_k(t)-\tau_k(s)\le t-s$: $\tau$ is Lipschitz and
$\tau(t)=\int_0^tu(r)\,dr$ with $u(r)\in\gSmp$ for a.e.\ $r$. A progressive version is
\[
u(t):=\iota\Big(\limsup_{n\to\infty}n\big(\tau(t)-\tau((t-1/n)\vee0)\big)\Big),
\]
with the $\limsup$ taken coordinatewise and $\iota:[-\infty,\infty]^K\to\gSmp$ the Borel map equal to the identity on $\gSmp$ and to $\frac1K\1$ elsewhere; at every
point of differentiability of $\tau$, which is a.e., the $\limsup$ is the derivative, which lies in $\gSmp$. So $\langle X_k,X_j\rangle_t=\1\{k=j\}\int_0^tu_k(r)\,dr$. By the
representation theorem for continuous martingales \cite[Thm~3.4.2]{KS91}, there is an extension $(\tilde\Omega,\tilde{\mathcal F},(\tilde{\mathcal G}_t),\tilde\P)$ of
$(\Om,\mathcal F_1^Q,(\mathcal F^Q_t),Q)$ carrying a $K$-dimensional $(\tilde{\mathcal G}_t)$-Brownian motion $B$ with $X_k=\int_0^\cdot\sqrt{u_k}\,dB_k$. On it, $(X,\tau)$
has law $Q$ and $\tau$ is $(\tilde{\mathcal G}_t)$-adapted.

Step 2: \emph{delayed averages.} For $n\ge2$ let $\delta:=1/n$ and
\[
u^n(t):=\tfrac1K\1\ \text{ on }[0,\delta),\qquad u^n(t):=\frac{\tau(i\delta)-\tau((i-1)\delta)}\delta\ \text{ on }[i\delta,(i+1)\delta),\quad1\le i\le n-1 .
\]
Then $u^n$ is $\gSmp$-valued, and its value on the cell $[i\delta,(i+1)\delta)$ is $\tilde{\mathcal G}_{i\delta}$-measurable. We
show
\begin{equation}\label{eq:g-L1}
\tilde\E\int_0^1|u^n(t)-u(t)|_1\,dt\ \longrightarrow\ 0 .
\end{equation}
Fix $\omega$. For $v\in L^1([0,1],\R^K)$ let $A_nv$ be the function equal on $[i\delta,(i+1)\delta)$ to the average of $v$ over $[(i-1)\delta,i\delta)$, $i\ge1$; so $u^n=A_nu$
on $[\delta,1]$. The operator $A_n$ is a contraction from $L^1([0,1])$ to $L^1([\delta,1])$, and $|A_nv(t)-v(t)|\le\varpi_v(2\delta)$ for $t\ge\delta$ when $v$ is continuous with modulus
of continuity $\varpi_v$. Given $\varepsilon>0$, choose a continuous $v$ with $\int_0^1|u-v|_1\le\varepsilon$. Then
\begin{align*}
\int_0^1|u^n-u|_1&\le2\delta+\int_\delta^1|A_nu-A_nv|_1+\int_\delta^1|A_nv-v|_1+\int_\delta^1|v-u|_1 &&\text{(}|\cdot|_1\le2\text{ on }\gSmp\text{)}\\
&\le2\delta+2\varepsilon+\varpi_v(2\delta) &&\text{(contraction, continuity)}.
\end{align*}
So $\int_0^1|u^n-u|_1\to0$ for every $\omega$, and dominated convergence (the integrand is at most $2$) gives \eqref{eq:g-L1}.

Step 3: \emph{convergence.} Let $\tau^n:=\int_0^\cdot u^n$ and $X^n_k:=\int_0^\cdot\sqrt{u^n_k}\,dB_k$. We bound the distance to $(X,\tau)$:
\begin{align*}
\tilde\E\Big[\sup_t|X^n(t)-X(t)|^2\Big]&\le\sum_k\tilde\E\Big[\sup_t|X^n_k(t)-X_k(t)|^2\Big] &&\\
&\le4\sum_k\tilde\E\int_0^1\big(\sqrt{u^n_k}-\sqrt{u_k}\big)^2dt &&\text{(Doob, Itô isometry)}\\
&\le4\,\tilde\E\int_0^1|u^n-u|_1\,dt\ \longrightarrow\ 0 &&\text{(}(\sqrt a-\sqrt b)^2\le|a-b|\text{; \eqref{eq:g-L1})},
\end{align*}
and $\sup_t|\tau^n(t)-\tau(t)|_1\le\int_0^1|u^n-u|_1\,dt\to0$ in $L^1(\tilde\P)$. Hence $(X^n,\tau^n)\to(X,\tau)$ in probability in $\Om$, and
$\mathrm{Law}(X^n,\tau^n)\to Q$ weakly.

Step 4: \emph{realization.} On the cell $[i\delta,(i+1)\delta]$, $\tau^n$ moves linearly with the $\tilde{\mathcal G}_{i\delta}$-measurable proportions $p_i:=u^n(i\delta)$,
and $X^n_k(t)-X^n_k(i\delta)=\sqrt{p_{ik}}\,(B_k(t)-B_k(i\delta))$. Lemma~\ref{lem:g-real} gives $\pi_n\in\gDc$ with $\gQp{\pi_n}=\mathrm{Law}(X^n,\tau^n)$. The last claim
follows with Lemma~\ref{lem:g-disc}(b) and the closedness of $\Wc$ (Lemma~\ref{lem:g-compact}).
\end{proof}

\begin{lemma}[batched policies with continuous rules are dense]\label{lem:g-dense}
\uses{def:alloc, def:sc, def:g-pol}
\textit{Role.} Density of batched policies in $\Wc$; continuous rules for the core of \SC\ (Corollary~\ref{cor:g-cont}).

\noindent\textit{Setting.} $Q\in\Wc$.

\noindent Then there are batched policies $\pi_n$ with continuous rules such that $\gQp{\pi_n}\to Q$ weakly. In particular $\Wc$ is the weak closure of
$\{\gQp\pi:\pi\in\gBc\text{ with continuous rules}\}$.
\end{lemma}
\begin{proof}
\uses{lem:g-step, def:g-pol, def:alloc}
\textit{Idea.} Lusin's theorem makes the decision maps of a discrete policy continuous outside an event of small probability. Then refine the grid and feed the
continuous maps with the linear interpolation of the values of the path on the fine grid: by continuity, the decisions converge.

Step 1: \emph{continuous decision maps.} Let $\pi\in\gDc$ with grid $(t_i)_{i\le m}$ and maps $F_i$, and let $\eta>0$. Let $\mu_i$ be the law under $\gPh_0$ of the input
$(U,X|_{[0,t_i]})$ of $F_i$, a Borel probability on the Polish space $E_i:=[0,1]\times C([0,t_i],\R^K)$. Lusin's theorem \cite[Thm~17.12]{Kec95} gives a closed $C_i\subset E_i$
with $\mu_i(E_i\setminus C_i)\le\eta/m$ on which $F_i$ is continuous. Extend each coordinate of $F_i|_{C_i}$ continuously to $E_i$ (Tietze), and compose with the Euclidean
projection on the closed convex set $\gSmp$; the result $\tilde F_i$ is continuous, $\gSmp$-valued, and equal to $F_i$ on $C_i$. Let $\tilde\pi\in\gDc$ use the maps
$\tilde F_i$. Run $\pi$ and $\tilde\pi$ on the same $(U,Y)$: if they chose the same proportions in cells $0,\dots,i-1$, their paths coincide on $[0,t_i]$, and they choose
the same proportions in cell $i$ if the input of $\pi$ lies in $C_i$. Hence the two pairs $(X,\tau)$ differ only on an event of $\gPh_0$-probability at most
$\sum_i\mu_i(E_i\setminus C_i)\le\eta$, and $|\E_{\gQp\pi}\varphi-\E_{\gQp{\tilde\pi}}\varphi|\le2\|\varphi\|_\infty\eta$ for every bounded measurable $\varphi$.

Step 2: \emph{refinement.} Let $\tilde\pi\in\gDc$ have continuous maps $\tilde F_i$ on the grid $(t_i)$. For $n\ge1$ let $G_n:=\{t_i\}_{i\le m}\cup\{j/n\}_{j\le n}$, a grid of
mesh at most $1/n$. For a path $x$ let $\iota^i_nx\in C([0,t_i],\R^K)$ be the linear interpolation of the values $x(s)$, $s\in G_n\cap[0,t_i]$. Let $\pi_n$ be the batched
policy on the grid $G_n$ which, on every cell of $G_n$ contained in $[t_i,t_{i+1}]$, plays the proportions $\tilde F_i(U,\iota^i_nX)$. These proportions are continuous
functions of $U$ and of the values $X(s)$, $s\in G_n$, $s\le t_i$, which are known at the start of the cell; so $\pi_n\in\gBc$ has continuous rules. Run $\tilde\pi$ and $\pi_n$
on the same $(U,Y)$, with paths $(\tilde X,\tilde\tau)$ and $(X^n,\tau^n)$. We show by induction on $i$ that, for every $\omega$,
$\|X^n-\tilde X\|_{[0,t_i]}+\|\tau^n-\tilde\tau\|_{[0,t_i]}\to0$. This holds for $i=0$. If it holds for $i$, then
\begin{align*}
\|\iota^i_nX^n-\tilde X\|_{[0,t_i]}&\le\|\iota^i_nX^n-\iota^i_n\tilde X\|_{[0,t_i]}+\|\iota^i_n\tilde X-\tilde X\|_{[0,t_i]} &&\text{(triangle inequality)}\\
&\le\|X^n-\tilde X\|_{[0,t_i]}+\varpi_{\tilde X}(1/n)\ \longrightarrow\ 0 &&\text{(}\iota^i_n\text{ is }1\text{-Lipschitz)},
\end{align*}
where $\varpi_{\tilde X}$ is the modulus of continuity of $\tilde X$ on $[0,1]$, and the second term bounds the interpolation error on a grid of mesh at most $1/n$. By continuity of $\tilde F_i$, the proportions $p^n_i:=\tilde F_i(U,\iota^i_nX^n)$ tend to
$\tilde p_i:=\tilde F_i(U,\tilde X|_{[0,t_i]})$. So $\tau^n\to\tilde\tau$ uniformly on $[0,t_{i+1}]$, and $X^n_k=Y_k\circ\tau^n_k\to Y_k\circ\tilde\tau_k=\tilde X_k$ uniformly
on $[0,t_{i+1}]$ ($Y_k$ is uniformly continuous on $[0,1]$). At $i=m$, $(X^n,\tau^n)\to(\tilde X,\tilde\tau)$ in $\Om$ for every $\omega$, so
$\gQp{\pi_n}\to\gQp{\tilde\pi}$ weakly.

Step 3: \emph{diagonal argument.} Let $d$ be the bounded-Lipschitz metric on $\mathcal P(\Om)$, which metrizes the weak topology; by Step~1, $d(\gQp\pi,\gQp{\tilde\pi})\le2\eta$.
By Lemma~\ref{lem:g-step} there are $\pi_n\in\gDc$ with $d(\gQp{\pi_n},Q)\le1/n$. Step~1 with $\eta=1/n$ gives $\tilde\pi_n$ with continuous maps and
$d(\gQp{\tilde\pi_n},\gQp{\pi_n})\le2/n$, and Step~2 gives $\pi^\flat_n\in\gBc$ with continuous rules and $d(\gQp{\pi^\flat_n},\gQp{\tilde\pi_n})\le1/n$. Hence
$\gQp{\pi^\flat_n}\to Q$. The last claim follows as in Lemma~\ref{lem:g-step}.
\end{proof}

\begin{lemma}[reading the whole path does not help]\label{lem:g-bridge}
\uses{def:g-pol, def:risk}
\textit{Role.} In Proposition~\ref{prop:four}(b), replaces the discrete policies of Lemma~\ref{lem:g-step}, which read the whole observed path, by batched policies,
which read only its values at the decision times; shows that the two readings of Remark~\ref{rem:g-reading} give the same risk functions.

\noindent\textit{Setting.} $\pi\in\gDc$ with grid $(t_i)_{i\le m}$ and maps $(F_i)$.

\noindent Then there is $\pi^\flat\in\gBc$ on the same grid such that, for every $h\in\R^K$, the allocation $\tau$ has the same law under $\gPh_h$ for $\pi^\flat$ as for
$\pi$. In particular $\rk(\pi^\flat,h)=\rk(\pi,h)$ for every $h\in\R^K$.
\end{lemma}
\begin{proof}
\uses{lem:g-disc, def:g-pol, def:alloc, def:risk}
\textit{Idea.} Given the values of a Brownian motion on a grid, the path in between consists of independent Brownian bridges, whose law does not depend on the drift.
So $\pi^\flat$ rebuilds a fictitious path from the values at the decision times and simulated bridges, and applies the maps of $\pi$ to it; under $h=0$ the fictitious
path has the law of the true one, and both risks are computed under $h=0$ with the same likelihood, which only depends on the final values.

Step 1: \emph{construction.} Let $\beta=(\beta_k)_{k\in[K]}$ be $K$ independent standard Brownian bridges on $[0,1]$, $\beta_k(s):=W_k(s)-sW_k(1)$. Recall that $\beta$ is
independent of $W(1)$, so that, for $b\in[0,1]^K$ and $\xi\sim\mathcal N(0,\mathrm{diag}(b))$ independent of $\beta$,
\begin{equation}\label{eq:g-bridge}
\big(s\,\xi_k+\sqrt{b_k}\,\beta_k(s)\big)_{s\in[0,1],k\in[K]}\ \text{ has law }\ \ggam(b,\cdot),
\end{equation}
because $\sqrt{b_k}W_k(s)=s\sqrt{b_k}W_k(1)+\sqrt{b_k}\beta_k(s)$. Let $U\mapsto(U',V_0,\dots,V_{m-1})$ be a Borel map sending the uniform law to the law of i.i.d.\ uniforms,
and $\zeta:[0,1]\to C([0,1],\R^K)$ a Borel map with $\zeta(V)\sim\mathrm{Law}(\beta)$ for $V$ uniform \cite[Lemma~3.22]{Kal02}; let $\beta_i:=\zeta(V_i)$. The policy
$\pi^\flat$ observes $X$ and builds a path $\hat X$ recursively: $\hat X(0):=0$; at $t_i$ it plays
\[
p^\flat_i:=F_i\big(U',\hat X|_{[0,t_i]}\big),\qquad b_i:=(t_{i+1}-t_i)\,p^\flat_i,
\]
and at $t_{i+1}$ it sets, with $\Delta_i:=X(t_{i+1})-X(t_i)$,
\[
\hat X\big(t_i+s(t_{i+1}-t_i)\big):=\hat X(t_i)+s\,\Delta_i+\big(\sqrt{b_{ik}}\,\beta_{i,k}(s)\big)_{k\in[K]},\qquad s\in[0,1].
\]
By induction $\hat X(t_i)=X(t_i)$ (bridges vanish at $0$ and $1$), and $p^\flat_i$ is a Borel function of $(U,X(t_1),\dots,X(t_i))$; so $\pi^\flat\in\gBc$.

Step 2: \emph{the fictitious path has the right law under $\gPh_0$.} Let $\hat H_i:=(U',\hat X|_{[0,t_i]})$ for $\pi^\flat$ and $H_i:=(U,X|_{[0,t_i]})$ for $\pi$, both under
$\gPh_0$. We show by induction that they have the same law; this holds for $i=0$. Let $\hat{\mathcal G}_i:=\sigma(U',V_0,\dots,V_{i-1},X(t_1),\dots,X(t_i))$ (for $\pi^\flat$);
$\hat H_i$ and $b_i$ are $\hat{\mathcal G}_i$-measurable. Three facts:
\begin{itemize}
\item \emph{the increment.} By Lemma~\ref{lem:g-disc}(a) applied to $\pi^\flat$ with $h=0$, conditionally on $\sigma(U,X|_{[0,t_i]})\supset\hat{\mathcal G}_i$, $\Delta_i\sim\mathcal N(0,\mathrm{diag}(b_i))$;
by the tower property, the same holds conditionally on $\hat{\mathcal G}_i$.
\item \emph{the bridges.} $\hat{\mathcal G}_i$ and $\Delta_i$ are functions of $(U',V_0,\dots,V_{i-1},Y)$ (the proportions $p^\flat_j$, $j\le i$, only use the bridges
$\beta_0,\dots,\beta_{i-1}$). Since $V_i$ is independent of $(U',(V_j)_{j\ne i},Y)$, $\beta_i$ is independent of $\sigma(\hat{\mathcal G}_i,\Delta_i)$.
\item \emph{the cell path.} By the two previous facts and \eqref{eq:g-bridge}, the cell path $\hat c_i:=(\hat X(t_i+s(t_{i+1}-t_i))-\hat X(t_i))_{s\in[0,1]}$ has
conditional law $\ggam(b_i,\cdot)$ given $\hat{\mathcal G}_i$, with $b_i=(t_{i+1}-t_i)F_i(\hat H_i)$; hence also given $\sigma(\hat H_i)$.
\end{itemize}
For $\pi$, Lemma~\ref{lem:g-disc}(a) gives the same kernel: the cell path $c_i$ has conditional law $\ggam((t_{i+1}-t_i)F_i(H_i),\cdot)$ given $H_i$. Since $\hat H_{i+1}$
and $H_{i+1}$ are the same function of $(\hat H_i,\hat c_i)$ and $(H_i,c_i)$, and $\hat H_i$, $H_i$ have the same law, $\hat H_{i+1}$ and $H_{i+1}$ have the same law.
At $i=m$: $(U',\hat X)$ under $\pi^\flat$ has the law of $(U,X)$ under $\pi$, and the allocation of $\pi^\flat$ is the same function $\Theta(U',\hat X)$ as the
allocation $\Theta(U,X)$ of $\pi$ (piecewise linear, with proportions $F_i$ of the inputs).

Step 3: \emph{all instances.} Let $h\in\R^K$ and $\varphi\ge0$ measurable on $\Al$. Write $\tau^\flat$ for the allocation of $\pi^\flat$. We transport the computation
to $h=0$:
\begin{align*}
\gEh_h\big[\varphi(\tau^\flat)\big]&=\gEh_0\big[\lh_h(X,\tau^\flat)\,\varphi(\tau^\flat)\big] &&\text{(Lemma~\ref{lem:g-disc}(c) for }\pi^\flat\text{)}\\
&=\gEh_0\big[\lh_h(\hat X,\tau^\flat)\,\varphi(\tau^\flat)\big] &&\text{(}\hat X(1)=X(1)\text{; }\lh_h\text{ only uses the values at time }1\text{)}\\
&=\gEh_0\big[\lh_h(X,\tau)\,\varphi(\tau)\big]\ \text{ for }\pi &&\text{(Step 2: same function of a pair with the same law)}\\
&=\gEh_h[\varphi(\tau)]\ \text{ for }\pi &&\text{(Lemma~\ref{lem:g-disc}(c) for }\pi\text{)}.
\end{align*}
So $\tau$ has the same law under $\gPh_h$ for both policies, and $\rk(\pi^\flat,h)=\gEh_h[\ggh_h(\tau^\flat)]=\rk(\pi,h)$.
\end{proof}

\subsection{Which symmetries does the game have?}\label{sec:app-g-sym}

\begin{lemma}[translations, relabellings, duplication]\label{lem:g-act}
\uses{def:alloc, def:risk}
\textit{Role.} The invariances of Proposition~\ref{prop:four}(d).

\noindent\textit{Setting.}\dt{g:act} $Q\in\Wc$, $c\in\R$, $\sigma$ a permutation of $[K]$. Let $\gPhi_c(x,a):=(x+c\,a,\,a)$ on $\Om$, and let $\gTc_cQ$ be the image of
$\lh_{-c\1}\cdot Q$ by $\gPhi_c$. Let $\gRs_\sigma(x,a):=((x_{\sigma(k)})_k,(a_{\sigma(k)})_k)$ and $\gPs_\sigma Q$ the image of $Q$ by $\gRs_\sigma$; write
$(h\circ\sigma^{-1})_k:=h_{\sigma^{-1}(k)}$. For $K+1$ arms let $\gJ:\Om^{(K+1)}\to\Om^{(K)}$,
\[
\gJ(x',a'):=\Big(\big(x'_1+x'_{K+1},x'_2,\dots,x'_K\big),\ \big(a'_1+a'_{K+1},a'_2,\dots,a'_K\big)\Big).
\]
In words: $\gTc_cQ$ runs $Q$ on the observations shifted by $-c$ per unit of time, so that $Q$ sees the instance $h-c\1$; $\gPs_\sigma Q$ relabels the arms;
$\gJ_\#Q'$ runs a $(K+1)$-armed policy $Q'$ on $K$ arms, arm $K+1$ being a second copy of arm $1$.

\noindent Then:
\begin{enumerate}
\item[(i)] $\gTc_cQ\in\Wc$ and $\rk(\gTc_cQ,h)=\rk(Q,h-c\1)$ for every $h$; $\gTc_a\gTc_c=\gTc_{a+c}$ and $\gTc_0=\mathrm{id}$, so $\gTc_c$ is a bijection of $\Wc$; $\gTc_c$ is
linear (it commutes with convex combinations and with the mixtures of (iv)) and weakly continuous on $\Wc$; and $c\mapsto\gTc_cQ(A)$ is Borel for every Borel $A\subset\Om$;
\item[(ii)] $\gPs_\sigma Q\in\Wc$, $\rk(\gPs_\sigma Q,h)=\rk(Q,h\circ\sigma^{-1})$, $\gPs_\rho\gPs_\sigma=\gPs_{\sigma\circ\rho}$, and $\gPs_\sigma\gTc_c=\gTc_c\gPs_\sigma$;
\item[(iii)] for $Q'\in\Wc^{(K+1)}$, $\gJ_\#Q'\in\Wc^{(K)}$ and $\rk^{(K)}(\gJ_\#Q',h)=\rk^{(K+1)}\big(Q',(h_1,\dots,h_K,h_1)\big)$ for every $h\in\R^K$;
\item[(iv)] if $(Q_\theta)_{\theta\in\Theta}$ is a family in $\Wc$ such that $\theta\mapsto Q_\theta(A)$ is measurable for every Borel $A$, and $\lambda$ is a probability on $\Theta$,
then $\bar Q:=\int Q_\theta\,\lambda(d\theta)\in\Wc$ and $\rk(\bar Q,h)=\int\rk(Q_\theta,h)\,\lambda(d\theta)$.
\end{enumerate}
\end{lemma}
\begin{proof}
\uses{lem:g-compact, lem:g-like, def:alloc, def:risk}
\textit{Idea.} Each operation is a change of variables on $\Om$, possibly combined with Girsanov's change of instance; it preserves the martingale conditions of
Lemma~\ref{lem:g-compact}(i), and the likelihoods transform as the instance does.

Step 1: \emph{mixtures, (iv).} For $\varphi\ge0$ measurable, $\E_{\bar Q}\varphi=\int\E_{Q_\theta}\varphi\,\lambda(d\theta)$ (true for indicators, then by monotone convergence).
By Lemma~\ref{lem:g-compact}(ii) with $u=0$ and Cauchy--Schwarz, $\E_{Q'}|M_t|\le1+C_4^{1/4}+C_4^{1/2}$ for every $Q'\in\Wc$ and $M\in\gMc$. So, for $M\in\gMc$ and bounded $\mathcal F_u$-measurable $Z$,
$\E_{\bar Q}|M_t|<\infty$ and, by Fubini, $\E_{\bar Q}[(M_t-M_u)Z]=\int\E_{Q_\theta}[(M_t-M_u)Z]\,\lambda(d\theta)=0$. Lemma~\ref{lem:g-compact}(i) gives $\bar Q\in\Wc$. The risk
formula is the first identity with $\varphi=\gFh_h$.

Step 2: \emph{translations, (i).} By definition of $\gTc_cQ$, for every measurable $\varphi\ge0$ on $\Om$,
\begin{equation}\label{eq:g-Tc}
\E_{\gTc_cQ}\varphi=\E_Q\big[\lh_{-c\1}\,\varphi\circ\gPhi_c\big].
\end{equation}
We check the claims of (i) one by one.
\begin{itemize}
\item \emph{$\gTc_cQ\in\Wc$.} Let $P':=\lh_{-c\1}\cdot Q$, a probability by Lemma~\ref{lem:g-like}(a). By Lemma~\ref{lem:g-like}(c) with $h=-c\1$, the processes
$\tilde X_k:=X_k+c\tau_k$ are continuous $P'$-martingales with $\langle\tilde X_k,\tilde X_j\rangle=\1\{k=j\}\tau_k$. As in Step~1 of Lemma~\ref{lem:g-compact}, the
processes $M\circ\gPhi_c$, $M\in\gMc$, are then $P'$-martingales for $(\mathcal F_t)$. The map $\gPhi_c$ is a homeomorphism of $\Om$ (inverse $\gPhi_{-c}$), and
$\gPhi_c(\omega)(\cdot\wedge t)$ only depends on $\omega(\cdot\wedge t)$, so $\gPhi_c^{-1}(\mathcal F_t)=\mathcal F_t$. Hence, for bounded $\mathcal F_u$-measurable $Z$,
\[
\E_{\gTc_cQ}[(M_t-M_u)Z]=\E_{P'}\big[(M_t\circ\gPhi_c-M_u\circ\gPhi_c)\,Z\circ\gPhi_c\big]=0,
\]
and $\gTc_cQ\in\Wc$ by Lemma~\ref{lem:g-compact}(i).
\item \emph{The risk.} For $(x,a)\in\Om$, using $\sum_ka_k(1)=1$,
\begin{align}
&\log\lh_{-c\1}(x,a)+\log\lh_h(\gPhi_c(x,a)) &&\notag\\
&=-c\sum_kx_k(1)-\frac{c^2}2+\sum_kh_k\big(x_k(1)+c\,a_k(1)\big)-\frac12\sum_kh_k^2a_k(1) &&\text{(definitions)}\notag\\
&=\sum_k(h_k-c)\,x_k(1)-\frac12\sum_k(h_k-c)^2a_k(1)=\log\lh_{h-c\1}(x,a) &&\text{(expand }(h_k-c)^2\text{)}.\label{eq:g-shift}
\end{align}
The gaps do not change under translation, so $\ggh_h\circ\gPhi_c=\ggh_h=\ggh_{h-c\1}$. With $\varphi=\gFh_h$ in \eqref{eq:g-Tc} and \eqref{eq:g-shift},
$\rk(\gTc_cQ,h)=\E_Q[\lh_{h-c\1}\ggh_{h-c\1}]=\rk(Q,h-c\1)$.
\item \emph{The group law.} For measurable $\varphi\ge0$,
\begin{align*}
\E_{\gTc_a\gTc_cQ}\varphi&=\E_{\gTc_cQ}\big[\lh_{-a\1}\,\varphi\circ\gPhi_a\big] &&\text{(\eqref{eq:g-Tc} for }\gTc_a\text{)}\\
&=\E_Q\big[\lh_{-c\1}\,(\lh_{-a\1}\circ\gPhi_c)\,\varphi\circ\gPhi_a\circ\gPhi_c\big] &&\text{(\eqref{eq:g-Tc} for }\gTc_c\text{)}\\
&=\E_Q\big[\lh_{-(a+c)\1}\,\varphi\circ\gPhi_{a+c}\big]=\E_{\gTc_{a+c}Q}\varphi &&\text{(\eqref{eq:g-shift} with }h=-a\1\text{)},
\end{align*}
using also $\gPhi_a\circ\gPhi_c=\gPhi_{a+c}$. Clearly $\gTc_0=\mathrm{id}$.
\item \emph{Linearity} (\eqref{eq:g-Tc} is linear in $Q$, also for the mixtures of (iv), by Fubini).
\item \emph{Continuity in $Q$.} Let $Q_n\to Q$ in $\Wc$ and $\varphi$ be bounded and continuous. The function $\lh_{-c\1}\,\varphi\circ\gPhi_c$ is continuous, and
$\E_{Q_n}[(\lh_{-c\1}\varphi\circ\gPhi_c)^2]\le\|\varphi\|_\infty^2e^{Kc^2}$ by Lemma~\ref{lem:g-like}(b). Fact (V) and \eqref{eq:g-Tc} give $\E_{\gTc_cQ_n}\varphi\to\E_{\gTc_cQ}\varphi$.
\item \emph{Measurability in $c$} (the map $(c,\omega)\mapsto\lh_{-c\1}(\omega)\1_A(\gPhi_c(\omega))$ is Borel on $\R\times\Om$; apply Fubini--Tonelli to \eqref{eq:g-Tc}).
\end{itemize}

Step 3: \emph{relabellings, (ii).} The map $\gRs_\sigma$ is a homeomorphism of $\Om$, preserves $(\mathcal F_t)$, and maps $\gMc$ onto itself ($X_k\circ\gRs_\sigma=X_{\sigma(k)}$,
$\tau_k\circ\gRs_\sigma=\tau_{\sigma(k)}$); so $\gPs_\sigma Q\in\Wc$ by Lemma~\ref{lem:g-compact}(i). Reindexing the sums by $j=\sigma(k)$, and since $\max_ih_i$ is invariant,
$\lh_h\circ\gRs_\sigma=\lh_{h\circ\sigma^{-1}}$ and $\ggh_h\circ\gRs_\sigma=\ggh_{h\circ\sigma^{-1}}$, which gives the risk. Next, $\gRs_\rho\circ\gRs_\sigma=\gRs_{\sigma\circ\rho}$, and
$\gRs_\sigma\circ\gPhi_c=\gPhi_c\circ\gRs_\sigma$ with $\lh_{-c\1}\circ\gRs_\sigma=\lh_{-c\1}$, which gives the two identities.

Step 4: \emph{duplication, (iii).} Let $Q'\in\Wc^{(K+1)}$ with canonical process $(X',\tau')$ and filtration $(\mathcal F'_t)$, and let $(X,\tau):=\gJ(X',\tau')$, that is
\[
X_1=X'_1+X'_{K+1},\quad \tau_1=\tau'_1+\tau'_{K+1},\qquad X_k=X'_k,\quad \tau_k=\tau'_k\quad(2\le k\le K).
\]
\begin{itemize}
\item \emph{$\gJ_\#Q'\in\Wc^{(K)}$.} Under $Q'$ the $X_k$ are continuous martingales, and, since $\langle X'_1,X'_{K+1}\rangle=0$,
\begin{gather*}
\langle X_1\rangle=\tau'_1+\tau'_{K+1}=\tau_1,\qquad \langle X_1,X_k\rangle=\langle X'_1,X'_k\rangle+\langle X'_{K+1},X'_k\rangle=0\quad(k\ge2),\\
\langle X_k,X_j\rangle=\1\{k=j\}\tau_k\quad(k,j\ge2).
\end{gather*}
As in Step~1 of Lemma~\ref{lem:g-compact}, the processes $M\circ\gJ$, $M\in\gMc^{(K)}$, are $(\mathcal F'_t,Q')$-martingales. Since $\gJ(\omega')(\cdot\wedge t)$ only depends on
$\omega'(\cdot\wedge t)$, we have $\gJ^{-1}(\mathcal F_t)\subset\mathcal F'_t$; so, for bounded $\mathcal F_u$-measurable $Z$,
\[
\E_{\gJ_\#Q'}[(M_t-M_u)Z]=\E_{Q'}\big[(M_t\circ\gJ-M_u\circ\gJ)\,Z\circ\gJ\big]=0,
\]
and $\gJ_\#Q'\in\Wc^{(K)}$ by Lemma~\ref{lem:g-compact}(i).
\item \emph{The risk.} Let $h':=(h_1,\dots,h_K,h_1)$. Then $\max_kh'_k=\max_kh_k$ and $\gGp_{K+1}(h')=\gGp_1(h)$. Regrouping the terms of arms $1$ and $K+1$,
\begin{align*}
\sum_{k\le K}h_kX_k(1)-\frac12\sum_{k\le K}h_k^2\tau_k(1)&=\sum_{k\le K+1}h'_kX'_k(1)-\frac12\sum_{k\le K+1}(h'_k)^2\tau'_k(1),\\
\sum_{k\le K}\gGp_k(h)\,\tau_k(1)&=\sum_{k\le K+1}\gGp_k(h')\,\tau'_k(1).
\end{align*}
So $\lh^{(K)}_h\circ\gJ=\lh^{(K+1)}_{h'}$ and $\ggh^{(K)}_h\circ\gJ=\ggh^{(K+1)}_{h'}$, hence $\rk^{(K)}(\gJ_\#Q',h)=\rk^{(K+1)}(Q',h')$.
\end{itemize}
\end{proof}

\begin{lemma}[Hunt--Stein averaging]\label{lem:g-avg}
\uses{def:alloc, def:risk}
\textit{Role.} Gives a translation-invariant minimax policy in Proposition~\ref{prop:four}(d).

\noindent\textit{Setting.} $Q\in\Wc$ and $V:=\sup_h\rk(Q,h)\in[0,\infty]$; for $C>0$, $\bar Q_C:=\frac1{2C}\int_{-C}^C\gTc_cQ\,dc$.

\noindent Then $\bar Q_C\in\Wc$ and $\rk(\bar Q_C,h)=\frac1{2C}\int_{-C}^C\rk(Q,h-c\1)\,dc\le V$. Moreover $(\bar Q_C)_{C>0}$ has weak limit points as $C\to\infty$, and every
such limit point $\bar Q$ satisfies $\gTc_a\bar Q=\bar Q$ for every $a\in\R$ and $\sup_h\rk(\bar Q,h)\le V$.
\end{lemma}
\begin{proof}
\uses{lem:g-act, lem:g-compact, lem:g-cont, def:alloc, def:risk}
\textit{Idea.} An average over a long interval of translations is almost invariant under a fixed translation; the defect is of order $|a|/C$ in total variation and
disappears in the limit.

Step 1: \emph{the average.} By Lemma~\ref{lem:g-act}(i), $c\mapsto\gTc_cQ(A)$ is Borel, so Lemma~\ref{lem:g-act}(iv) with $\lambda$ uniform on $[-C,C]$ gives $\bar Q_C\in\Wc$ and
$\rk(\bar Q_C,h)=\frac1{2C}\int_{-C}^C\rk(\gTc_cQ,h)\,dc=\frac1{2C}\int_{-C}^C\rk(Q,h-c\1)\,dc\le V$.

Step 2: \emph{almost invariance.} Let $a\in\R$. By linearity of $\gTc_a$ and $\gTc_a\gTc_c=\gTc_{a+c}$ (Lemma~\ref{lem:g-act}(i)),
\[
\gTc_a\bar Q_C=\frac1{2C}\int_{-C}^C\gTc_{a+c}Q\,dc=\frac1{2C}\int_{a-C}^{a+C}\gTc_cQ\,dc .
\]
The intervals $[a-C,a+C]$ and $[-C,C]$ differ by a set of length at most $2|a|$, so $|\gTc_a\bar Q_C(A)-\bar Q_C(A)|\le|a|/C$ for every Borel $A$, and
$|\E_{\gTc_a\bar Q_C}\varphi-\E_{\bar Q_C}\varphi|\le2\|\varphi\|_\infty|a|/C$ for bounded measurable $\varphi$.

Step 3: \emph{limit points.} The set $\Wc$ is compact (Lemma~\ref{lem:g-compact}), so there are $C_n\to\infty$ and $\bar Q$ with $\bar Q_{C_n}\to\bar Q$. For bounded continuous
$\varphi$, using the continuity of $\gTc_a$ (Lemma~\ref{lem:g-act}(i)) and Step~2,
\[
\big|\E_{\gTc_a\bar Q}\varphi-\E_{\bar Q}\varphi\big|=\lim_n\big|\E_{\gTc_a\bar Q_{C_n}}\varphi-\E_{\bar Q_{C_n}}\varphi\big|\le\lim_n\frac{2\|\varphi\|_\infty|a|}{C_n}=0,
\]
so $\gTc_a\bar Q=\bar Q$. Finally $\rk(\bar Q,h)=\lim_n\rk(\bar Q_{C_n},h)\le V$ by Lemma~\ref{lem:g-cont} and Step~1.
\end{proof}

\subsection{Proof of Proposition~\ref{prop:four}(a), (b), (d)}\label{sec:app-g-proof}

\begin{proof}[Proof of Proposition~\ref{prop:four}(a)]
\proves{prop:four}
\uses{lem:g-compact, lem:g-cont, def:risk, def:alloc}
\textit{Setting.} $K\ge2$. \textit{Goal.} $\VK_K=\sup\{\BK_K(\Pi):\Pi\in\gPf\}$.

\textit{Idea.} Ky Fan's theorem needs compactness and lower semicontinuity on one side only, and affinity on both; the side of the priors needs no topology.

Recall Fan's theorem \cite[Thm~2]{Fan53}: let $\mathsf X$ be a compact Hausdorff space, $\mathsf Y$ an arbitrary set, and $f:\mathsf X\times\mathsf Y\to\R$ such that
$f(\cdot,y)$ is lower semicontinuous for every $y$; if $f$ is convex on $\mathsf X$ and concave on $\mathsf Y$ (in Fan's sense: for $x_1,x_2\in\mathsf X$ and
$\lambda\in[0,1]$ there is $x_0$ with $f(x_0,\cdot)\le\lambda f(x_1,\cdot)+(1-\lambda)f(x_2,\cdot)$, and symmetrically on $\mathsf Y$), then
$\min_x\sup_yf(x,y)=\sup_y\min_xf(x,y)$. Take $\mathsf X:=\Wc$, $\mathsf Y:=\gPf$, and
\[
f(Q,\Pi):=\int\rk(Q,h)\,\Pi(dh)=\sum_{h\in\mathrm{supp}\,\Pi}\Pi(\{h\})\,\rk(Q,h).
\]
We check the hypotheses:
\begin{itemize}
\item the set $\Wc$ is compact (Lemma~\ref{lem:g-compact}) and Hausdorff ($\mathcal P(\Om)$ is metrizable);
\item the function $f$ is real valued: $0\le f(Q,\Pi)\le\max_{h\in\mathrm{supp}\,\Pi}\gsp h$ (Lemma~\ref{lem:g-cont});
\item the map $f(\cdot,\Pi)$ is continuous, as a finite combination of the continuous maps $\rk(\cdot,h)$ (Lemma~\ref{lem:g-cont});
\item the function $f$ is affine in $Q$ on the convex set $\Wc$ (Lemmas~\ref{lem:g-compact}, \ref{lem:g-cont}) and affine in $\Pi$ on the convex set $\gPf$; an affine function is convex and
concave in Fan's sense (take $x_0:=\lambda x_1+(1-\lambda)x_2$).
\end{itemize}
Hence $\min_Q\sup_\Pi f(Q,\Pi)=\sup_\Pi\min_Qf(Q,\Pi)$. On the left, $\sup_{\Pi\in\gPf}f(Q,\Pi)=\sup_h\rk(Q,h)$: ``$\ge$'' with Dirac masses, ``$\le$'' because $f(Q,\Pi)$ is an
average of values $\rk(Q,h)$. So the left side is $\VK_K$. On the right, $\min_Qf(Q,\Pi)=\BK_K(\Pi)$. Hence $\VK_K=\sup_{\Pi\in\gPf}\BK_K(\Pi)$.
\end{proof}

\noindent The identity holds in $[0,\infty]$; the finiteness of $\VK_K$ comes from Theorem~\ref{thm:bounds} and is not used here. Moreover, if $\Pi$ is any Borel probability
and $Q^\star$ is minimax, $\BK_K(\Pi)\le\int\rk(Q^\star,h)\Pi(dh)\le\VK_K$; so $\VK_K$ is also the supremum of $\BK_K$ over all Borel probabilities on $\R^K$.

\begin{proof}[Proof of Proposition~\ref{prop:four}(b)]
\proves{prop:four}
\uses{lem:g-disc, lem:g-min, lem:g-step, lem:g-bridge, lem:g-compact, lem:g-cont, def:batched, def:risk, def:g-pol, def:alloc}
\textit{Setting.} $H\subset\R^K$ compact. \textit{Goal.} $\Vb(H)=\VK_K(H)$; then $\VK_K=\Vloc_K=\lim_{M\to\infty}\Vb([-M,M]^K)$.

\textit{Idea.} Batched policies are limit policies, which gives one inequality. A minimizer for $\VK_K(H)$ is a limit of discrete policies, whose risk functions are
those of batched policies; this gives the other. Compactness passes from boxes to $\R^K$.

Step 1: \emph{$\VK_K(H)\le\Vb(H)$.} For $\pi\in\gBc$, $\gQp\pi\in\Wc$ and $\rk(\pi,\cdot)=\rk(\gQp\pi,\cdot)$ (Lemma~\ref{lem:g-disc}(b), (c)). So
$\sup_{h\in H}\rk(\pi,h)\ge\VK_K(H)$, and taking the infimum over $\pi$ gives the claim.

Step 2: \emph{$\Vb(H)\le\VK_K(H)$.} Let $Q_H$ attain $\VK_K(H)$ (Lemma~\ref{lem:g-min}(iii)). By Lemma~\ref{lem:g-step} there are $\pi_n\in\gDc$ with $\gQp{\pi_n}\to Q_H$. By
Lemma~\ref{lem:g-bridge} there are $\pi^\flat_n\in\gBc$ with $\rk(\pi^\flat_n,\cdot)=\rk(\pi_n,\cdot)$. Hence
\begin{align*}
\sup_{h\in H}\rk(\pi^\flat_n,h)&=\sup_{h\in H}\rk(\pi_n,h)=\sup_{h\in H}\rk(\gQp{\pi_n},h) &&\text{(Lemmas~\ref{lem:g-bridge} and~\ref{lem:g-disc}(c))}\\
&\xrightarrow[n\to\infty]{}\ \sup_{h\in H}\rk(Q_H,h)=\VK_K(H) &&\text{(Lemma~\ref{lem:g-min}(i), }H\text{ compact)}.
\end{align*}
So $\Vb(H)\le\VK_K(H)$.

Step 3: \emph{from boxes to $\R^K$.} Let $H_M:=[-M,M]^K$. By Steps~1--2, $\Vb(H_M)=\VK_K(H_M)$, which is nondecreasing in $M$ and at most $\VK_K$. So the limit exists and
equals $\Vloc_K=\sup_M\VK_K(H_M)\le\VK_K$. Conversely, let $Q_M$ attain $\VK_K(H_M)$, and let $Q^\ast$ be the limit of $(Q_{M_n})$ along some $M_n\to\infty$
(Lemma~\ref{lem:g-compact}). For a fixed $h\in\R^K$, $h\in H_{M_n}$ for $n$ large, so
\begin{align*}
\rk(Q^\ast,h)&=\lim_n\rk(Q_{M_n},h) &&\text{(Lemma~\ref{lem:g-cont})}\\
&\le\lim_n\VK_K(H_{M_n})=\Vloc_K &&\text{(}h\in H_{M_n}\text{ for }n\text{ large)}.
\end{align*}
Taking the supremum over $h$ gives $\VK_K\le\sup_h\rk(Q^\ast,h)\le\Vloc_K$. Hence $\VK_K=\Vloc_K=\lim_M\Vb(H_M)$.
\end{proof}

\begin{proof}[Proof of Proposition~\ref{prop:four}(d)]
\proves{prop:four}
\uses{lem:g-act, lem:g-min, lem:g-avg, lem:g-cont, def:risk, def:alloc}
\textit{Setting.} $K\ge2$, $c\in\R$, $H\subset\R^K$ nonempty. \textit{Goal.} (1) $\VK_K(H+c\1)=\VK_K(H)$; (2) a minimax $\hat Q\in\Wc$ with $\gTc_a\hat Q=\hat Q$ and
$\gPs_\sigma\hat Q=\hat Q$ for all $a$, $\sigma$; (3) $\VK_K\le\VK_{K+1}$.

Step 1: \emph{translated sets.} The map $\gTc_c$ is a bijection of $\Wc$ (Lemma~\ref{lem:g-act}(i)), and for every $Q\in\Wc$
\[
\sup_{h\in H+c\1}\rk(\gTc_cQ,h)=\sup_{h\in H+c\1}\rk(Q,h-c\1)=\sup_{h'\in H}\rk(Q,h').
\]
Taking the minimum over $Q$ (equivalently over $\gTc_cQ$) gives $\VK_K(H+c\1)=\VK_K(H)$.

Step 2: \emph{an invariant minimax policy.} Let $Q^\star$ be minimax (Lemma~\ref{lem:g-min}(iii)) and $\bar Q$ a limit point given by Lemma~\ref{lem:g-avg}: $\gTc_a\bar Q=\bar Q$
for all $a$ and $\sup_h\rk(\bar Q,h)\le\VK_K$. Let $\hat Q:=\frac1{K!}\sum_\sigma\gPs_\sigma\bar Q$, which is in $\Wc$ by convexity. Then:
\begin{itemize}
\item $\gPs_\rho\hat Q=\frac1{K!}\sum_\sigma\gPs_{\sigma\circ\rho}\bar Q=\hat Q$ for every $\rho$ ($\sigma\mapsto\sigma\circ\rho$ is a bijection; Lemma~\ref{lem:g-act}(ii));
\item $\gTc_a\hat Q=\frac1{K!}\sum_\sigma\gPs_\sigma\gTc_a\bar Q=\hat Q$ (linearity of $\gTc_a$ and $\gTc_a\gPs_\sigma=\gPs_\sigma\gTc_a$, Lemma~\ref{lem:g-act}(i), (ii));
\item $\rk(\hat Q,h)=\frac1{K!}\sum_\sigma\rk(\bar Q,h\circ\sigma^{-1})\le\VK_K$ (Lemma~\ref{lem:g-cont}, Lemma~\ref{lem:g-act}(ii)).
\end{itemize}
So $\hat Q$ is minimax, symmetric in the arms and invariant by translation; in particular $\rk(\hat Q,h\circ\sigma)=\rk(\hat Q,h+c\1)=\rk(\hat Q,h)$ for all $h$, $\sigma$, $c$.

Step 3: \emph{monotonicity in $K$.} Let $Q'$ be minimax for $K+1$ arms and $Q:=\gJ_\#Q'\in\Wc^{(K)}$ (Lemma~\ref{lem:g-act}(iii)). For every $h\in\R^K$,
\[
\rk^{(K)}(Q,h)=\rk^{(K+1)}\big(Q',(h_1,\dots,h_K,h_1)\big)\le\VK_{K+1},
\]
so $\VK_K\le\sup_h\rk^{(K)}(Q,h)\le\VK_{K+1}$.
\end{proof}

\begin{corollary}[continuous batched policies suffice]\label{cor:g-cont}
\uses{def:sc, def:batched, def:risk, def:g-pol}
\textit{Role.} The core of \SC\ must be a batched policy with continuous rules (Section~\ref{sec:app-upper}); this corollary says that this restriction costs nothing.

\noindent\textit{Setting.} $H\subset\R^K$ compact; $\Pi$ a Borel probability on $\R^K$ with $\int\gsp h\,\Pi(dh)<\infty$.

\noindent Then
\begin{gather*}
\Vb(H)=\VK_K(H)=\inf\Big\{\sup_{h\in H}\rk(\pi,h):\pi\in\gBc\text{ with continuous rules}\Big\},\\
\BK_K(\Pi)=\inf\Big\{\int\rk(\pi,h)\,\Pi(dh):\pi\in\gBc\text{ with continuous rules}\Big\}.
\end{gather*}
\end{corollary}
\begin{proof}
\uses{prop:four, lem:g-disc, lem:g-min, lem:g-dense, def:batched, def:risk, def:g-pol, def:alloc}
The identity $\Vb(H)=\VK_K(H)$ is Proposition~\ref{prop:four}(b). The two infima on the right are at least $\VK_K(H)$ and $\BK_K(\Pi)$, because
$\gQp\pi\in\Wc$ and $\rk(\pi,\cdot)=\rk(\gQp\pi,\cdot)$ for $\pi\in\gBc$ (Lemma~\ref{lem:g-disc}). Conversely, let $Q_H$ attain $\VK_K(H)$ and $Q_\Pi$ attain $\BK_K(\Pi)$ (Lemma~\ref{lem:g-min}(iii)). By
Lemma~\ref{lem:g-dense} there are batched policies $\pi_n$, $\pi'_n$ with continuous rules and $\gQp{\pi_n}\to Q_H$, $\gQp{\pi'_n}\to Q_\Pi$. By Lemma~\ref{lem:g-min}(i), (ii),
$\sup_{h\in H}\rk(\gQp{\pi_n},h)\to\VK_K(H)$ and $\int\rk(\gQp{\pi'_n},h)\Pi(dh)\to\BK_K(\Pi)$.
\end{proof}

\begin{remark}[on part (c)]\label{rem:g-c}
The lower half of Proposition~\ref{prop:four}(c) holds for every $T$: by Lemma~\ref{lem:embed} with $\mu_0=0$, $\sigma=1$ and $H=[-M,M]^K$,
\[
T^{-1/2}\inf_\pi\sup_{h\in[-M,M]^K}\RT\big(\pi,(\mathcal N(h_k/\sqrt T,1))_k\big)\ \ge\ \VK_K([-M,M]^K).
\]
The upper half needs a $T$-round discretization of a near-optimal batched policy with continuous rules (Corollary~\ref{cor:g-cont}) and a transfer theorem; it is
proved in Section~\ref{sec:app-upper} (Proposition~\ref{prop:u-local}).
\end{remark}

\section{Proof of the lower bound}\label{sec:app-lower}

\newcommand{\lvp}{\sy{l:ef}{\varpi}}
\newcommand{\lTh}{\sy{l:ef}{\Theta}}
\newcommand{\lpsi}{\sy{l:ef}{\psi}}
\newcommand{\lthz}{\sy{l:ef}{\theta_0}}
\newcommand{\lmuz}{\sy{l:ef}{\mu_0}}
\newcommand{\lsd}{\sy{l:ef}{\varsigma}}
\newcommand{\lnuT}{\sy{l:local}{\nu^{(T,h)}}}
\newcommand{\lM}{\sy{l:space}{M}}
\newcommand{\lFo}{\sy{l:space}{\mathfrak F_1}}
\newcommand{\lG}{\sy{l:space}{\mathfrak G}}
\newcommand{\lphi}{\sy{l:space}{\varphi}}
\newcommand{\lJ}{\sy{l:const}{J}}
\newcommand{\lcth}{\sy{l:const}{c_3}}
\newcommand{\lnub}{\sy{l:const}{\bar\nu}}
\newcommand{\lmq}{\sy{l:const}{m_4}}
\newcommand{\lvb}{\sy{l:const}{\bar v}}
\newcommand{\lTone}{\sy{l:const}{T_1}}
\newcommand{\letT}{\sy{l:const}{\eta_T}}
\newcommand{\lepT}{\sy{l:const}{\varepsilon_T}}
\newcommand{\lGam}{\sy{l:stair}{\Gamma}}
\newcommand{\lsimal}{\sy{l:sim}{simulated allocation}}
\newcommand{\lsimpol}{\sy{l:sim}{simulated limit policy}}
\newcommand{\lL}{\sy{l:lr}{L}}
\newcommand{\lDT}{\sy{l:lr}{D_T}}
\newcommand{\lcg}{\sy{l:tilt}{c_g}}
\newcommand{\lnug}{\sy{l:tilt}{\nu^{(T,h)}_g}}
\newcommand{\lepTp}{\sy{l:tilt}{\varepsilon'_T}}
\providecommand{\gOs}{\sy{g:pol}{\Omega^\circ}}
\providecommand{\gPh}{\sy{g:pol}{\mathbb P}}
\providecommand{\gEh}{\sy{g:pol}{\mathbb E}}
\providecommand{\gWi}{\sy{g:pol}{\mathbb W}}
\providecommand{\gDc}{\sy{g:pol}{\mathcal D}}
\providecommand{\gQp}[1]{\sy{g:pol}{Q^{#1}}}
\providecommand{\gSmp}{\sy{g:not}{\Lambda_K}}
\providecommand{\gGp}{\sy{g:not}{\Delta}}


\medskip
\noindent This appendix proves Theorem~\ref{thm:main}(i) and the lower bounds in Corollary~\ref{cor:examples}. The reader is assumed to know the stack model of
bandits and standard facts on Brownian motion (strong Markov property, optional stopping, Wald's identities, the Cameron--Martin formula), which are used without
proof. Three facts are cited from Section~\ref{sec:app-game}: discrete policies are limit policies (Lemma~\ref{lem:g-disc}), $\Wc$ is compact
(Lemma~\ref{lem:g-compact}), and the risk is continuous (Lemma~\ref{lem:g-cont}); together with Proposition~\ref{prop:four}(b) ($\VK_K=\Vloc_K$) and~(d)
(translation invariance).

\paragraph{Skeleton.} The lower bound reduces to a transfer theorem (Theorem~\ref{thm:lowEF}). Take a law $\nu_0$ of the class with standard deviation
$\varsigma$, embedded in an exponential family, and the local instances whose arm $k$ has mean $\approx\mu_0+\varsigma h_k/\sqrt T$. For every $T$-round policy $\pi$
there is \emph{one} limit policy $\tilde Q$, which does not depend on $h$, such that
\[
\Big|\frac{\RT(\pi,\nu^{(T,h)})}{\varsigma\sqrt T}-\rk(\tilde Q,h)\Big|=O(T^{-1/4})\qquad\text{uniformly over }\pi\text{ and over }h\text{ in a compact set }S.
\]
Taking the supremum over $h\in S$ gives $\sup_{h\in S}\RT(\pi,\nu^{(T,h)})\ge\varsigma\sqrt T\,(\VK_K(S)-O(T^{-1/4}))$ for every $\pi$; boxes $S=[-M,M]^K$ exhaust
$\VK_K$ by Proposition~\ref{prop:four}(b). The policy $\tilde Q$ is $\pi$ itself, run on rewards read off Brownian paths. The proof has three ingredients, with their costs:
\begin{itemize}
\item \emph{Rewards read off Brownian paths} (Lemmas~\ref{lem:skor} and~\ref{lem:timechange}). The $j$-th reward of arm $k$ is the increment of a Brownian motion $Y_k$
between two stopping times of $Y_k$ (Skorokhod's embedding). This embedding is adapted, so the simulated policy does not anticipate (the Komlós--Major--Tusnády coupling
would not do: its dyadic construction uses future values of the path). The time change is off by $O(T^{-1/2})$ in fractions of the horizon.
\item \emph{The simulated policy is a limit policy} (Lemma~\ref{lem:simW}), with no loss: rounding the stopping times up to a grid gives discrete policies, which are
limit policies, and the rounding vanishes in the limit.
\item \emph{Likelihood ratios of full arrays} (Steps~3 and~6 of the proof of Theorem~\ref{thm:lowEF}). Both risks are expectations under the reference instance: the regret
with the likelihood ratio $\lL_h$ of the full array of $T$ rewards per arm, which does not depend on the policy, the limit risk with the Cameron--Martin density
$\lM_h$. For an exponential family, $\log\lL_h$ only depends on the row sums, which are values of $Y_k$ at a time $1+O(T^{-1/2})$; so $\lL_h-\lM_h=O(T^{-1/4})$ in $L^1$.
\end{itemize}
Classes stable under bounded perturbations (the second case of (R2)) are handled by the same proof (Proposition~\ref{prop:lowtilt}).

\paragraph{Dependencies.}
\begin{center}\small
\begin{tabular}{l}
Lem.~\ref{lem:skor} $\to$ Lem.~\ref{lem:timechange} $\to$ Thm~\ref{thm:lowEF} $\leftarrow$ Lem.~\ref{lem:simW};\\[1mm]
Lem.~\ref{lem:staircase} $+$ Lemmas~\ref{lem:g-disc}, \ref{lem:g-compact}, \ref{lem:g-cont} $\to$ Lem.~\ref{lem:simW};\quad
proof of Thm~\ref{thm:lowEF} $\to$ Prop.~\ref{prop:lowtilt};\\[1mm]
Thm~\ref{thm:lowEF} $+$ Prop.~\ref{prop:lowtilt} $+$ Prop.~\ref{prop:four}(b), (d) $\to$ Thm~\ref{thm:main}(i) $\to$ Cor.~\ref{cor:examples} (lower bounds);\\[1mm]
Cor.~\ref{cor:u-gauss} $\to$ Cor.~\ref{cor:examples} (Gaussian bound).
\end{tabular}
\end{center}

\subsection{Setting}\label{sec:app-low-setting}

\paragraph{Policies and the stack model.}\dt{l:stack} A \emph{$T$-round policy} is a family $\pi=(\pi_n)_{n\le T}$ of Borel maps
$\pi_n:[0,1]\times([K]\times\R)^{n-1}\to[K]$. With $U_0$ uniform on $[0,1]$ (the external randomization), it plays $A_n:=\pi_n(U_0,A_1,\zeta_1,\dots,A_{n-1},\zeta_{n-1})$,
where $\zeta_i$ is the reward of round $i$, and $N_k(n):=\sum_{i\le n}\1\{A_i=k\}$. We use the stack model \cite[Chapter~4]{LS}: an array $(Z_{k,j})_{k\in[K],j\ge1}$ of
independent rewards with $Z_{k,j}\sim\nu_k$, independent of $U_0$, and $\zeta_n:=Z_{A_n,N_{A_n}(n)}$. It gives the same law to $(U_0,A_1,\zeta_1,\dots,A_T,\zeta_T)$ as the
interaction protocol, hence the same regret. In this model, $N_k(T)$ is a Borel function of the \emph{full array} $(U_0,(Z_{k,j})_{k\in[K],j\le T})$.

\paragraph{Exponential families.}\dt{l:ef} Let $\lvp$ be a $\sigma$-finite Borel measure on $\R$, not concentrated on one point, such that
$\lTh:=\{\theta\in\R:\int e^{\theta x}\lvp(dx)<\infty\}$ has a nonempty interior. For $\theta\in\operatorname{int}\lTh$ let
\[
\lpsi(\theta):=\log\int e^{\theta x}\,\lvp(dx),\qquad \nu_\theta(dx):=e^{\theta x-\lpsi(\theta)}\,\lvp(dx).
\]
Standard facts \cite{Bro86}: $\lpsi$ is $C^\infty$ on $\operatorname{int}\lTh$; $\nu_\theta$ has mean $\lpsi'(\theta)$ and variance $\lpsi''(\theta)>0$; for
$\theta,\theta'\in\operatorname{int}\lTh$,
\[
\frac{d\nu_\theta}{d\nu_{\theta'}}(x)=\exp\big((\theta-\theta')x-\lpsi(\theta)+\lpsi(\theta')\big);
\]
and $\int e^{\lambda x}\nu_\theta(dx)=e^{\lpsi(\theta+\lambda)-\lpsi(\theta)}$ when $\theta+\lambda\in\lTh$. In particular $\nu_\theta$ has all its moments. Fix
$\lthz\in\operatorname{int}\lTh$ and let
\[
\lmuz:=\lpsi'(\lthz),\qquad \lsd:=\lpsi''(\lthz)^{1/2}.
\]

\paragraph{Reading of (R2).} We use (R2) in the following precise form.

\begin{definition}[the two cases of (R2)]\label{def:lowR2}
\uses{def:class}
Let $\Cl$ be a class of laws on $\R$ and $\nu_0$ a law with a finite mean $\mu_0$.
\begin{enumerate}
\item[(a)] $\Cl$ \emph{contains an exponential family through $\nu_0$ on both sides} if there are a family as above, $\lthz\in\operatorname{int}\lTh$ with $\nu_{\lthz}=\nu_0$, and
$\beta_0>0$ such that $\nu_\theta\in\Cl$ for every $\theta\in(\lthz-\beta_0,\lthz+\beta_0)$; \emph{on one side} if the same holds for every $\theta\in(\lthz-\beta_0,\lthz]$,
or for every $\theta\in[\lthz,\lthz+\beta_0)$.
\item[(b)] $\Cl$ \emph{contains the bounded perturbations of $\nu_0$} if for every bounded Borel $g$ with $\int g\,d\nu_0=0$ there is $t_g>0$ such that $(1+tg)\,\nu_0$ is a
probability measure in $\Cl$ for every $|t|\le t_g$.
\end{enumerate}
\end{definition}

\noindent Condition (R2) of Definition~\ref{def:regular} asks: for every $\sigma_0<\sbar$, a law $\nu_0\in\Cl$ with variance at least $\sigma_0^2$ such that (a) or (b) holds.

\paragraph{Local instances.}\dt{l:local} For $h\in\R^K$ and $T\ge1$ such that $\lthz+h_k/(\lsd\sqrt T)\in\operatorname{int}\lTh$ for all $k$, let
\[
\lnuT:=\big(\nu_{\lthz+h_k/(\lsd\sqrt T)}\big)_{k\in[K]} .
\]
The mean of arm $k$ is $\lpsi'(\lthz+h_k/(\lsd\sqrt T))=\lmuz+\lsd h_k/\sqrt T+O(1/T)$: the local parameter $h_k$ is the mean in units of $\lsd/\sqrt T$.

\paragraph{The probability space.}\dt{l:space} We work on the space $\gOs=[0,1]\times C([0,\infty),\R)^K$ of Definition~\ref{def:g-pol}, with coordinates
$(U,Y_1,\dots,Y_K)$ and the probabilities $\gPh_h=\mathrm{Leb}\otimes\bigotimes_k\gWi_{h_k}$: under $\gPh_h$ the $Y_k$ are independent Brownian motions with drifts
$h_k$, independent of the uniform variable $U$. We fix a Borel map $U\mapsto\big(U_0,(U_{k,j})_{k\in[K],j\ge1},(U'_{k,j})_{k\in[K],j\ge1}\big)$ under which these variables
are i.i.d.\ uniform on $[0,1]$ (a Borel isomorphism between $[0,1]$ and $[0,1]^{\Nat}$ carrying the Lebesgue measure to the product measure). For $k\in[K]$ and $u\ge0$
let
\[
\lG^k_u:=\sigma(U)\vee\sigma\big(Y_k(v):v\le u\big),\qquad \lFo:=\sigma\big(U,\ Y_k(v):v\le1,\ k\in[K]\big):
\]
$\lG^k$ is the filtration of arm $k$'s own path and of the randomization. For $h\in\R^K$ and $x\in\gSmp$ let
\[
\lM_h:=\prod_{k\in[K]}\exp\Big(h_kY_k(1)-\frac{h_k^2}2\Big),\qquad \lphi_h(x):=\sum_{k\in[K]}\gGp_k(h)\,x_k,\qquad \gGp_k(h)=\max_ih_i-h_k .
\]
By the Cameron--Martin formula, applied to each arm ($U$ being independent of the arms), $\gPh_h=\lM_h\cdot\gPh_0$ on $\lFo$. The loss $\lphi_h(x)$ is the regret per unit
of time of the proportions $x$ in the instance $h$, so that $\rk(Q,h)=\E_Q[\lh_h\,\lphi_h(\tau(1))]$ for $Q\in\Wc$.

\paragraph{Two facts on the loss} (used in Step~7 of the proof of Theorem~\ref{thm:lowEF}). Let $|h|_\infty\le H$ and $x,y\in\gSmp$. Then
\begin{equation}\label{eq:low-L}
0\le\lphi_h(x)\le2H,\qquad |\lphi_h(x)-\lphi_h(y)|\le H\,\|x-y\|_1 .
\end{equation}
\begin{itemize}
\item $0\le\lphi_h(x)\le2H$ (each gap $\gGp_k(h)$ lies in $[0,2H]$, and $\sum_kx_k=1$).
\item $|\lphi_h(x)-\lphi_h(y)|\le H\|x-y\|_1$ (since $\sum_k(x_k-y_k)=0$, $\lphi_h(x)-\lphi_h(y)=\sum_k(\gGp_k(h)-c)(x_k-y_k)$ with $c$ the midpoint of the range
$[0,\max_ih_i-\min_ih_i]$ of the gaps, and $|\gGp_k(h)-c|\le\frac12(\max_ih_i-\min_ih_i)\le H$).
\end{itemize}

\paragraph{Discrete policies.} A discrete policy (Definition~\ref{def:g-pol}(a)) is a grid $0=t_0<\dots<t_m=1$ and Borel maps $F_i$ giving the proportions
$p_i=F_i(U,X|_{[0,t_i]})\in\gSmp$ of cell $[t_i,t_{i+1}]$, with $X_k(t):=Y_k(\tau_k(t))$. By Lemma~\ref{lem:g-disc}, the law $\gQp\pi$ of $(X,\tau)$ under $\gPh_0$ is in
$\Wc$, and $\rk(\gQp\pi,h)=\gEh_h[\lphi_h(\tau(1))]$ for every $h$.

\subsection{The transfer statements}\label{sec:app-low-statements}

\begin{theorem}[transfer for an exponential family]\label{thm:lowEF}\dt{l:const}
\uses{def:alloc, def:regret, def:risk}
\textit{Role.} The core of the lower bound: the first case of (R2) in the proof of Theorem~\ref{thm:main}(i).

\noindent\textit{Setting.} $K\ge2$; a family with $\lthz$, $\lmuz$, $\lsd$ as in Section~\ref{sec:app-low-setting}; $\beta>0$ with $\lJ:=[\lthz-\beta,\lthz+\beta]\subset\operatorname{int}\lTh$
and $\lcth:=\max_{\lJ}|\lpsi'''|$; $\lnub$ the law of $(X-\lmuz)/\lsd$ for $X\sim\nu_{\lthz}$, $\lmq:=\int y^4\,\lnub(dy)$ and $\lvb:=(\frac53+2\sqrt{2/3})\,\lmq$; a compact
$S\subset\R^K$ with $H:=\max_{h\in S}|h|_\infty>0$; and
\begin{gather*}
\lTone:=\max\Big\{\Big(\frac{2H}{\lsd\beta}\Big)^2,\Big(\frac{2\lcth H}{\lsd^3}\Big)^2\Big\},\qquad \letT:=\frac{\lcth H}{\lsd^3\sqrt T},\\
\lepT:=4KH^2e^{KH^2}\Big(\frac{\lvb}T\Big)^{1/4}+\frac23\,\frac{K\lcth H^4e^{KH^2}}{\lsd^3\sqrt T}+4KHe^{KH^2/2}\Big(\frac{\lvb}T\Big)^{1/2}.
\end{gather*}
Then for every integer $T\ge\lTone$ and every $T$-round policy $\pi$, the instances $\lnuT$, $h\in S$, are defined, and there is a limit policy $\tilde Q\in\Wc$, which
does not depend on $h$, such that:
\begin{enumerate}
\item[(i)] for every $h\in S$,
\begin{gather*}
\Big|\frac{\RT(\pi,\lnuT)}{\lsd\sqrt T}-\rk(\tilde Q,h)\Big|\ \le\ 2H\letT+\frac32\,\lepT,\\
\RT(\pi,\lnuT)\ \ge\ \lsd\sqrt T\,(1-\letT)\big(\rk(\tilde Q,h)-\lepT\big);
\end{gather*}
\item[(ii)] $\displaystyle\inf_\pi\sup_{h\in S}\frac{\RT(\pi,\lnuT)}{\lsd\sqrt T}\ \ge\ (1-\letT)\big(\VK_K(S)-\lepT\big)$;
\item[(iii)] for every Borel probability $\Pi$ on $S$, $\displaystyle\inf_\pi\int\frac{\RT(\pi,\lnuT)}{\lsd\sqrt T}\,\Pi(dh)\ \ge\ (1-\letT)\big(\BK_K(\Pi)-\lepT\big)$.
\end{enumerate}
\end{theorem}
\noindent(See the proof in Section~\ref{sec:app-low-proofEF}.) For $T\ge\lTone$, $\letT\le\frac12$, and $\lepT=O(T^{-1/4})$. The constants depend only on $(K,H,\lsd,\lcth,\lmq)$, so
the error is uniform over the policies and over $h\in S$. The factor $e^{KH^2}$ makes the bound useless at moderate $T$; only the limit is used. The limit policy $\tilde Q$
is the \lsimpol\ of $\pi$ (Definition~\ref{def:lowsim}) for the marks of Skorokhod's embedding.

\begin{proposition}[transfer for bounded perturbations]\label{prop:lowtilt}\dt{l:tilt}
\uses{def:alloc, def:regret, def:risk}
\textit{Role.} The second case of (R2) in the proof of Theorem~\ref{thm:main}(i).

\noindent\textit{Setting.} $K\ge2$; a law $\nu_0$ with finite mean $\mu_0$; a bounded Borel $g$ with $\int g\,d\nu_0=0$, $I:=\int g^2d\nu_0>0$, $G:=\sup|g|$,
$\dot m:=\int(x-\mu_0)\,g(x)\,\nu_0(dx)>0$, and $\lcg:=\dot m/\sqrt I$. For $|t|\le1/G$, let $\nu^g_t:=(1+tg)\,\nu_0$ (a probability measure) and
$\lnug:=(\nu^g_{h_k/\sqrt{IT}})_{k\in[K]}$. Let $S$ and $H$ be as in Theorem~\ref{thm:lowEF}, $\bar G:=G/\sqrt I$, let $\lnub$ be the law of $g(X)/\sqrt I$ for $X\sim\nu_0$, let
$\lmq$ and $\lvb$ be defined from this $\lnub$ as in Theorem~\ref{thm:lowEF}, and
\begin{gather*}
\lambda'_T:=KH\Big(\frac{\lvb}T\Big)^{1/4}+\frac{KH^2}2\Big(\frac{\lmq}T\Big)^{1/2}+\frac23\,\frac{KH^3\bar G^3}{\sqrt T},\\
\lepTp:=4He^{KH^2/2}\lambda'_T+4KHe^{KH^2/2}\Big(\frac{\lvb}T\Big)^{1/2}.
\end{gather*}
Then for every integer $T\ge4H^2\bar G^2$ and every $T$-round policy $\pi$ there is $\tilde Q\in\Wc$, which does not depend on $h$, such that
\[
\Big|\frac{\RT(\pi,\lnug)}{\lcg\sqrt T}-\rk(\tilde Q,h)\Big|\ \le\ \lepTp\qquad\text{for every }h\in S .
\]
Moreover $\lcg\le\big(\int(x-\mu_0)^2\nu_0(dx)\big)^{1/2}$.
\end{proposition}
\noindent(See the proof in Section~\ref{sec:app-low-prooftilt}.) Here $\lmq\le\bar G^4<\infty$ since $|g|/\sqrt I\le\bar G$, and $\lepTp=O(T^{-1/4})$. The last claim is
Cauchy--Schwarz, $\dot m\le(\int(x-\mu_0)^2d\nu_0)^{1/2}\sqrt I$: the exponential tilt, whose efficiency is $\lsd$, is the least favourable direction.

\subsection{Proof of Theorem~\ref{thm:main}(i) and of the lower bounds of Corollary~\ref{cor:examples}}\label{sec:app-low-main}

\begin{proof}[Proof of Theorem~\ref{thm:main}(i)]
\proves{thm:main}
\uses{thm:lowEF, prop:lowtilt, prop:four, def:lowR2, def:class, def:regret, def:risk, def:alloc, def:batched}
\textit{Setting.} $K\ge2$ and a class $\Cl$ satisfying (R2) with the constant $\sbar$.

\textit{Goal.} $\liminf_T\Rs_{K,\Cl}(T)/\sqrt T\ge\sbar\,\VK_K$. The error $O(T^{-1/4})$, uniform over the policies and over the local instances in a compact set, is
the one of Theorem~\ref{thm:lowEF}(i) and Proposition~\ref{prop:lowtilt}.

\textit{Idea.} For one law $\nu_0$ given by (R2) and one box of local instances around it, Theorem~\ref{thm:lowEF} or Proposition~\ref{prop:lowtilt} bounds the
worst-case regret over the box from below by the value of the box. Boxes exhaust $\VK_K$ (Proposition~\ref{prop:four}(b)), and $\sigma_0\uparrow\sbar$.

\textit{Plan.} Step~1 treats one law $\nu_0$ and one box, in the three cases of Definition~\ref{def:lowR2}; Step~2 takes the suprema.

Step 1: \emph{one law and one box.} Fix $\sigma_0<\sbar$, a law $\nu_0\in\Cl$ with variance at least $\sigma_0^2$ as in (R2), and $M>0$. We show that for every $s<\sigma_0$,
\begin{equation}\label{eq:low-onebox}
\liminf_{T\to\infty}\frac{\Rs_{K,\Cl}(T)}{\sqrt T}\ \ge\ s\,\VK_K\big([-M,M]^K\big).
\end{equation}
\emph{Case 1: an exponential family on both sides.} Let $\lthz$, $\beta_0$ be as in Definition~\ref{def:lowR2}(a); then $\lsd^2=\lpsi''(\lthz)$ is the variance of $\nu_0$, so
$\lsd\ge\sigma_0$. Choose $\beta\in(0,\beta_0]$ with $[\lthz-\beta,\lthz+\beta]\subset\operatorname{int}\lTh$, let $S:=[-M,M]^K$ (so $H=M$), and let $\lTone$, $\letT$, $\lepT$ be
as in Theorem~\ref{thm:lowEF}. For $T\ge\lTone$ and $h\in S$, $|h_k|/(\lsd\sqrt T)\le\beta/2<\beta_0$, so $\lnuT\in\Cl^K$. Hence
\begin{align*}
\Rs_{K,\Cl}(T)&\ge\inf_\pi\sup_{h\in S}\RT(\pi,\lnuT) &&\text{(}\lnuT\in\Cl^K\text{ for }h\in S\text{)}\\
&\ge\lsd\sqrt T\,(1-\letT)\big(\VK_K(S)-\lepT\big) &&\text{(Theorem~\ref{thm:lowEF}(ii))}.
\end{align*}
Divide by $\sqrt T$ and let $T\to\infty$: $\letT,\lepT\to0$, so the liminf is at least $\lsd\VK_K(S)\ge\sigma_0\VK_K(S)$ (values are nonnegative). This is
\eqref{eq:low-onebox}.

\emph{Case 2: an exponential family on one side.} Assume $\nu_\theta\in\Cl$ for $\theta\in(\lthz-\beta_0,\lthz]$, and take $S:=[-2M,0]^K$, so $H=2M$. For $T\ge\lTone$ (computed
with $H=2M$) and $h\in S$, $\lthz+h_k/(\lsd\sqrt T)\in[\lthz-\beta/2,\lthz]$, so $\lnuT\in\Cl^K$. The chain of Case~1 gives a liminf at least $\lsd\VK_K([-2M,0]^K)$, and, by Proposition~\ref{prop:four}(d) (translation by $c=-M$),
\[
\VK_K\big([-2M,0]^K\big)=\VK_K\big([-M,M]^K-M\1\big)=\VK_K\big([-M,M]^K\big).
\]
If instead $\nu_\theta\in\Cl$ for $\theta\in[\lthz,\lthz+\beta_0)$, take $S:=[0,2M]^K$ and $c=M$.

\emph{Case 3: bounded perturbations.} Fix $s<\sigma_0$. For $L>0$ let
\[
g_L(x):=(x-\mu_0)\,\1\{|x-\mu_0|\le L\}-c_L,\qquad c_L:=\int(y-\mu_0)\,\1\{|y-\mu_0|\le L\}\,\nu_0(dy).
\]
We check the hypotheses of Proposition~\ref{prop:lowtilt} for $g_L$, for $L$ large:
\begin{itemize}
\item $g_L$ is bounded and $\int g_L\,d\nu_0=0$ (definition of $c_L$);
\item $\dot m_L:=\int(x-\mu_0)g_L(x)\nu_0(dx)=\int(x-\mu_0)^2\1\{|x-\mu_0|\le L\}\nu_0(dx)$ (since $\int(x-\mu_0)\nu_0(dx)=0$), and $\dot m_L\uparrow\int(x-\mu_0)^2d\nu_0\ge\sigma_0^2$
as $L\to\infty$ (monotone convergence); fix $L$ with $\dot m_L>s^2$;
\item $I_L:=\int g_L^2d\nu_0=\dot m_L-c_L^2\le\dot m_L$ (expand the square), and $I_L>0$ (otherwise $g_L=0$ $\nu_0$-a.s., and then $\dot m_L=0$);
\item hence $\lcg=\dot m_L/\sqrt{I_L}\ge\sqrt{\dot m_L}>s$ for $g=g_L$.
\end{itemize}
Let $t_{g_L}$ be given by Definition~\ref{def:lowR2}(b) and $S:=[-M,M]^K$. For $T\ge\max\{4M^2\bar G^2,\ M^2/(I_Lt_{g_L}^2)\}$ and $h\in S$, $|h_k|/\sqrt{I_LT}\le t_{g_L}$, so
$\lnug\in\Cl^K$. By Proposition~\ref{prop:lowtilt}, for every policy $\pi$, with its $\tilde Q$,
\begin{align*}
\sup_{h\in S}\RT(\pi,\lnug)&\ge\lcg\sqrt T\,\sup_{h\in S}\big(\rk(\tilde Q,h)-\lepTp\big) &&\text{(Proposition~\ref{prop:lowtilt})}\\
&\ge\lcg\sqrt T\,\big(\VK_K(S)-\lepTp\big) &&\text{(}\tilde Q\in\Wc\text{, definition of }\VK_K(S)\text{)}.
\end{align*}
As in Case~1, the liminf of $\Rs_{K,\Cl}(T)/\sqrt T$ is at least $\lcg\VK_K(S)\ge s\VK_K(S)$.

Step 2: \emph{conclusion.} In all cases \eqref{eq:low-onebox} holds for every $s<\sigma_0$ and every $M>0$. Then:
\begin{align*}
\liminf_{T\to\infty}\frac{\Rs_{K,\Cl}(T)}{\sqrt T}&\ge\sigma_0\sup_{M>0}\VK_K\big([-M,M]^K\big) &&\text{(\eqref{eq:low-onebox}, }s\uparrow\sigma_0\text{)}\\
&=\sigma_0\,\Vloc_K=\sigma_0\,\VK_K &&\text{(definition of }\Vloc_K\text{; Proposition~\ref{prop:four}(b))}.
\end{align*}
Letting $\sigma_0\uparrow\sbar$ gives $\liminf_T\Rs_{K,\Cl}(T)/\sqrt T\ge\sbar\,\VK_K$.
\end{proof}

\noindent The proof never uses the minimax theorem (Proposition~\ref{prop:four}(a)); it uses (R1) nowhere. If $\VK_K$ is replaced by $\Vloc_K$, it does not use
Section~\ref{sec:app-game} beyond Lemmas~\ref{lem:g-disc}, \ref{lem:g-compact}, \ref{lem:g-cont} and Proposition~\ref{prop:four}(d).

\begin{proof}[Proof of the lower bounds of Corollary~\ref{cor:examples}]
\proves{cor:examples}
\uses{thm:main, def:lowR2, cor:u-gauss, lem:embed, def:regret, def:risk, def:class}
\textit{Goal.} For each class of the corollary, $\liminf_T\Rs_K(T)/\sqrt T$ is at least the announced constant times $\VK_K$; and the Gaussian non-asymptotic bound.
By Theorem~\ref{thm:main}(i), it suffices to check (R2) with the announced $\sbar$.
\begin{itemize}
\item \emph{Laws on $[a,b]$, $\sbar=\frac{b-a}2$.} Take $\lvp:=\delta_a+\delta_b$, so $\lTh=\R$, $\lpsi(\theta)=\log(e^{\theta a}+e^{\theta b})$, and $\nu_\theta$ gives the mass
$p(\theta):=e^{\theta b}/(e^{\theta a}+e^{\theta b})$ to $b$. Its variance is $\lpsi''(\theta)=(b-a)^2p(\theta)(1-p(\theta))$, so at $\lthz:=0$, $\nu_0=\frac12(\delta_a+\delta_b)$ and
$\lsd=\frac{b-a}2=\sbar$. Every $\nu_\theta$ is a law on $[a,b]$: Definition~\ref{def:lowR2}(a) holds on both sides, for every $\sigma_0<\sbar$, with the same $\nu_0$.
\item \emph{Laws on $[0,1]$, and Bernoulli laws, $\sbar=\frac12$.} The case $a=0$, $b=1$: the family above is the Bernoulli family, contained in both classes.
\item \emph{Bernoulli laws of mean at most $\frac12$, $\sbar=\frac12$.} With $a=0$, $b=1$, $p(\theta)\le\frac12$ iff $\theta\le0$: the class contains $\nu_\theta$ for
$\theta\in(-1,0]$, so Definition~\ref{def:lowR2}(a) holds on one side at $\lthz=0$, with $\lsd=\frac12$.
\item \emph{Gaussian laws $\mathcal N(\mu,\sigma^2)$, $\mu\in[a,b]$, $\sigma\le\sbar$.} Take $\lvp:=\mathcal N(0,\sbar^2)$, so $\lTh=\R$, $\lpsi(\theta)=\sbar^2\theta^2/2$ and
$\nu_\theta=\mathcal N(\sbar^2\theta,\sbar^2)$. At $\lthz:=(a+b)/(2\sbar^2)$, $\nu_0=\mathcal N(\frac{a+b}2,\sbar^2)$, $\lsd=\sbar$, and $\nu_\theta$ is in the class for
$|\theta-\lthz|\le(b-a)/(2\sbar^2)$: Definition~\ref{def:lowR2}(a) holds on both sides.
\end{itemize}
The non-asymptotic Gaussian bound $\Rs_K(T)\ge\sbar\sqrt T\,\VK_K([-M,M]^K)$ for $T\ge(2\sbar M/(b-a))^2$ is Corollary~\ref{cor:u-gauss} (Section~\ref{sec:app-upper}),
which rests on Lemma~\ref{lem:embed}.
\end{proof}

\noindent\emph{The boundary case.} For the Bernoulli laws of mean at most $\frac12$, the law of largest variance is at the boundary of the class, and only one side of the
family is available. This does not lower the bound: the box $[-2M,0]^K$ of allowed local instances is a translate of $[-M,M]^K$, with the same value
(Proposition~\ref{prop:four}(d)). More generally, Definition~\ref{def:lowR2} gives a sufficient condition for the constant $\sbar$; we do not claim that it is necessary.

\subsection{How are rewards read off a Brownian path?}\label{sec:app-low-skor}

\begin{lemma}[Skorokhod embedding of a random walk]\label{lem:skor}
\textit{Role.} Builds the rewards of each arm from its own Brownian path, without anticipation; used in Step~1 of the proof of Theorem~\ref{thm:lowEF} and in the
proof of Proposition~\ref{prop:lowtilt}.

\noindent\textit{Setting.} A law $\lnub$ on $\R$ with $\int y\,\lnub(dy)=0$, $\int y^2\,\lnub(dy)=1$ and $\lmq:=\int y^4\,\lnub(dy)<\infty$; $\lvb:=(\frac53+2\sqrt{2/3})\,\lmq$; a standard
Brownian motion $B$ and i.i.d.\ uniform variables $(V_j)_{j\ge1}$ independent of $B$; $\mathfrak B_u:=\sigma(V_j:j\ge1)\vee\sigma(B(v):v\le u)$.

\noindent Then there are a.s.\ finite $\mathfrak B$-stopping times $0=:\rho_0\le\rho_1\le\rho_2\le\cdots$ such that:
\begin{enumerate}
\item[(i)] the pairs $\big(\rho_j-\rho_{j-1},\,B(\rho_j)-B(\rho_{j-1})\big)$, $j\ge1$, are i.i.d., and $B(\rho_1)\sim\lnub$;
\item[(ii)] $\E\rho_1=1$ and $\E\rho_1^2\le\lvb$.
\end{enumerate}
\end{lemma}

\begin{proof}
\uses{thm:lowEF}
\textit{Construction (cited).} By Skorokhod's randomized embedding \cite[Section~37]{Bil95}, \cite{Obl04}, there is a Borel map $v\mapsto(a(v),b(v))$ from $[0,1]$ to
$\{(a,b)\in\R^2:a\le0\le b\}$ such that, with
\[
\rho(w,v):=\inf\big\{u\ge0:\ w(u)\notin(a(v),b(v))\big\}
\]
the exit time of the interval $(a(v),b(v))$ by a path $w$ (it is $0$ when this interval does not contain $0$), and $V$ uniform and independent of $B$, we have
$B(\rho(B,V))\sim\lnub$ and $\E\rho(B,V)=\int y^2\,\lnub(dy)=1$. (The cited construction draws the pair $(a,b)$ from a law that depends on $\lnub$; any law on $\R^2$ is
the law of a Borel function of a uniform variable.) Let $\theta_sw:=w(s+\cdot)-w(s)$ and, recursively,
\[
\rho_j:=\rho_{j-1}+\rho\big(\theta_{\rho_{j-1}}B,\,V_j\big)\qquad(j\ge1).
\]
\textit{Stopping times and (i).} Assume that $\rho_{j-1}$ is an a.s.\ finite $\mathfrak B$-stopping time. Since $V_j$ is $\mathfrak B_0$-measurable, $\rho_j$ is the hitting
time, after $\rho_{j-1}$, of the closed set $\R\setminus(a(V_j),b(V_j))$ by the continuous adapted process $B-B(\rho_{j-1})$, hence a $\mathfrak B$-stopping time. The variables
$(V_j)$ are independent of $B$, so $B$ is a $(\mathfrak B_u)$-Brownian motion, and by the strong Markov property \cite{RY99} $\theta_{\rho_{j-1}}B$ is a Brownian motion
independent of $\mathfrak B_{\rho_{j-1}}$. The previous pairs and $V_j$ are $\mathfrak B_{\rho_{j-1}}$-measurable, and $V_j$ is independent of the previous pairs (which are
functions of $B$ and $V_1,\dots,V_{j-1}$). So the previous pairs, $V_j$ and $\theta_{\rho_{j-1}}B$ are independent, the last two with the same laws as $V_1$ and $B$. The
$j$-th pair is the same function of $(\theta_{\rho_{j-1}}B,V_j)$ as the first pair is of $(B,V_1)$; hence it is independent of the previous pairs and has the law of the
first pair. In particular $\rho_j-\rho_{j-1}$ has the law of $\rho_1$, which is integrable, so $\rho_j$ is a.s.\ finite, and the induction proceeds. This proves (i).

\textit{Bound on $\E\rho_1^2$.} Write $\rho:=\rho_1$ and $\rho^{(n)}:=\rho\wedge n$. The process $B_t^4-6tB_t^2+3t^2$ is a martingale, and $\rho^{(n)}$ is bounded. With
$x_n:=\E[(\rho^{(n)})^2]^{1/2}$ and $y_n:=\E[B_{\rho^{(n)}}^4]^{1/2}$, both finite,
\begin{align*}
3\,x_n^2&=6\,\E\big[\rho^{(n)}B_{\rho^{(n)}}^2\big]-\E\big[B_{\rho^{(n)}}^4\big] &&\text{(optional stopping at the bounded time }\rho^{(n)}\text{)}\\
&\le6\,x_ny_n-y_n^2 &&\text{(Cauchy--Schwarz).}
\end{align*}
So $3x_n^2-6x_ny_n+y_n^2\le0$, and $x_n$ is at most the larger root of $3x^2-6xy_n+y_n^2$:
\[
x_n\ \le\ \Big(1+\sqrt{2/3}\Big)\,y_n .
\]
Next, $y_n^2\le\lmq$. Indeed $\E\rho=1$, so the stopped martingale $(B_{t\wedge\rho})_t$ is bounded in $L^2$ ($\E B_{t\wedge\rho}^2=\E[t\wedge\rho]\le1$), hence uniformly
integrable, and optional stopping gives $B_{\rho^{(n)}}=\E[B_\rho\mid\mathfrak B_{\rho^{(n)}}]$; by Jensen's inequality,
$\E[B_{\rho^{(n)}}^4]\le\E[B_\rho^4]=\int y^4\,\lnub(dy)=\lmq$. Finally, by monotone convergence,
\[
\E\rho^2=\lim_nx_n^2\ \le\ \Big(1+\sqrt{2/3}\Big)^2\lmq=\Big(\frac53+2\sqrt{2/3}\Big)\lmq=\lvb .\qedhere
\]
\end{proof}

\begin{lemma}[time-change errors]\label{lem:timechange}
\textit{Role.} Bounds the two errors of the proof of Theorem~\ref{thm:lowEF}: (i) the gap between the numbers of pulls and the simulated allocation (Step~5), (ii) the
gap between the row sums and the values of the paths at time $1$ (Step~6).

\noindent\textit{Setting.} As in Lemma~\ref{lem:skor}; $v:=\operatorname{Var}\rho_1$, so $v\le\E\rho_1^2\le\lvb$.

\noindent Then for every integer $T\ge1$:
\begin{enumerate}
\item[(i)] $\E\big[\max_{n\le T}(\rho_n-n)^2\big]\le4Tv$;
\item[(ii)] $\E\big[(B(\rho_T)-B(T))^2\big]=\E|\rho_T-T|\le\sqrt{Tv}$.
\end{enumerate}
\end{lemma}
\begin{proof}
\uses{lem:skor}
(i) By Lemma~\ref{lem:skor}(i)--(ii), $(\rho_n-n)_{n\ge0}$ is a sum of i.i.d.\ centred increments of variance $v$, hence a martingale with $\E[(\rho_T-T)^2]=Tv$. Doob's
$L^2$ inequality gives $\E[\max_{n\le T}(\rho_n-n)^2]\le4\,\E[(\rho_T-T)^2]=4Tv$.

(ii) Let $\sigma_1:=\rho_T\wedge T$ and $\sigma_2:=\rho_T\vee T$. They are $\mathfrak B$-stopping times with $\sigma_1\le\sigma_2$, $\E\sigma_2\le\E\rho_T+T=2T<\infty$, and
$|B(\rho_T)-B(T)|=|B(\sigma_2)-B(\sigma_1)|$. Since $\E\sigma_2<\infty$, the stopped martingale $(B_{t\wedge\sigma_2})_t$ is bounded in $L^2$, so
$\E[B(\sigma_2)\mid\mathfrak B_{\sigma_1}]=B(\sigma_1)$ (optional stopping). Hence
\begin{align*}
\E\big[(B(\sigma_2)-B(\sigma_1))^2\big]&=\E[B(\sigma_2)^2]-\E[B(\sigma_1)^2] &&\text{(}\E[B(\sigma_2)B(\sigma_1)]=\E[B(\sigma_1)^2]\text{)}\\
&=\E\sigma_2-\E\sigma_1=\E|\rho_T-T| &&\text{(Wald's second identity, }\E\sigma_i<\infty\text{)}\\
&\le\E\big[(\rho_T-T)^2\big]^{1/2}=\sqrt{Tv} &&\text{(Cauchy--Schwarz).}\qedhere
\end{align*}
\end{proof}

\subsection{Why is a simulated policy a limit policy?}\label{sec:app-low-sim}

The rewards of Lemma~\ref{lem:skor} are read at stopping times of each arm's own path. We run a $T$-round policy on such rewards, describe the resulting
allocation as a path in $\R^K$, and show that its law under $\gPh_0$ is a limit policy whose risk is computed with the density $\lM_h$.

\begin{lemma}[staircase paths]\label{lem:staircase}\dt{l:stair}
\uses{def:alloc}
\textit{Role.} In Lemma~\ref{lem:simW}, rounding the stopping times changes the allocation uniformly little; in Step~5 of the proof of Theorem~\ref{thm:lowEF}, the
form of the simulated allocation.

\noindent\textit{Setting.} $T\ge1$; a sequence of arms $\alpha=(\alpha_1,\dots,\alpha_{T+1})\in[K]^{T+1}$; round lengths $\ell=(\ell_0,\dots,\ell_{T-1})\in[0,\infty)^T$;
$s_n:=\sum_{r<n}\ell_r$ and $w_n:=\sum_{r<n}\ell_r\,e_{\alpha_{r+1}}$ for $0\le n\le T$, and $s_{T+1}:=\infty$. The staircase path $\lGam(\alpha,\ell):[0,\infty)\to\R^K$ is
\[
\lGam(\alpha,\ell)(t):=w_n+(t-s_n)\,e_{\alpha_{n+1}}\qquad\text{for }t\in[s_n,s_{n+1}),\ 0\le n\le T .
\]
\noindent Then $\lGam(\alpha,\ell)$ is continuous, each coordinate is nondecreasing, $\sum_k\lGam_k(\alpha,\ell)(t)=t$ (so $\lGam(\alpha,\ell)|_{[0,1]}\in\Al$), and for every
$\ell'\in[0,\infty)^T$,
\[
\sup_{t\ge0}\big\|\lGam(\alpha,\ell)(t)-\lGam(\alpha,\ell')(t)\big\|_1\ \le\ 2\,\|\ell-\ell'\|_1 .
\]
\end{lemma}
\begin{proof}
\textit{Properties.} The intervals $[s_n,s_{n+1})$ partition $[0,\infty)$ (some may be empty). On each of them one coordinate increases at unit speed; at $s_{n+1}$ the left
limit is $w_n+\ell_ne_{\alpha_{n+1}}=w_{n+1}$, so the path is continuous; and $\sum_k\lGam_k(\alpha,\ell)(t)=s_n+(t-s_n)=t$. In particular $\|\lGam(\alpha,\ell)(t)-\lGam(\alpha,\ell)(u)\|_1=t-u$
for $u\le t$.

\textit{One round length.} Let $\ell'$ differ from $\ell$ only in coordinate $r$, with $\ell'_r=\ell_r+d$, $d\ge0$ (the case $d<0$ is symmetric). Write $\gamma:=\lGam(\alpha,\ell)$,
$\gamma':=\lGam(\alpha,\ell')$ and $u:=s_{r+1}$, the end of round $r$ under $\ell$. We check $\|\gamma'(t)-\gamma(t)\|_1\le2d$ for all $t$:
\begin{itemize}
\item $t<u$: the two paths have the same rounds before $u$, so $\gamma'(t)=\gamma(t)$;
\item $u\le t<u+d$: $\gamma'(t)=\gamma(u)+(t-u)e_{\alpha_{r+1}}$ (round $r$ lasts longer under $\ell'$), so
$\|\gamma'(t)-\gamma(t)\|_1\le(t-u)+\|\gamma(u)-\gamma(t)\|_1=2(t-u)\le2d$;
\item $t\ge u+d$: under $\ell'$ the later rounds are those of $\ell$ shifted by $d$ in time, and $w'_n=w_n+de_{\alpha_{r+1}}$ for $n\ge r+1$, so $\gamma'(t)=\gamma(t-d)+de_{\alpha_{r+1}}$
and $\|\gamma'(t)-\gamma(t)\|_1\le\|\gamma(t-d)-\gamma(t)\|_1+d=2d$.
\end{itemize}
\textit{General case.} Change the coordinates of $\ell$ into those of $\ell'$ one at a time and use the triangle inequality: the bound is $2\sum_r|\ell_r-\ell'_r|$.
\end{proof}

\begin{definition}[simulated allocation]\label{def:lowsim}\dt{l:sim}
\uses{def:g-pol}
\textit{Setting.} $T\ge1$ and a $T$-round policy $\pi$. The \emph{simulation data} are, for each arm $k$, nondecreasing $\lG^k$-stopping times
$0=\varrho^k_0\le\varrho^k_1\le\dots\le\varrho^k_T$, finite $\gPh_0$-a.s.\ (the \emph{marks}), and a Borel map $z:\R\times[0,1]\to\R$ (the \emph{reward map}).

\noindent The \emph{simulated rewards} are $Z_{k,j}:=z\big(\sqrt T\,(Y_k(\varrho^k_j)-Y_k(\varrho^k_{j-1})),\,U'_{k,j}\big)$ on $\{\varrho^k_j<\infty\}$, and $Z_{k,j}:=0$
otherwise ($k\in[K]$, $1\le j\le T$). Run $\pi$ on the stack $(Z_{k,j})$ with randomization $U_0$: this defines $A_n$, $\zeta_n$, $N_k(n)$ for $n\le T$ as in
Section~\ref{sec:app-low-setting}; let $A_{T+1}:=1$. Round $n+1$ runs arm $A_{n+1}$ from its current mark to its next mark; on the event $F$ that all the marks are finite,
its length is
\[
\ell_n:=\varrho^{k}_{N_k(n)+1}-\varrho^{k}_{N_k(n)},\qquad k:=A_{n+1},\qquad 0\le n<T .
\]
The \emph{\lsimal} is $\tau:=\lGam(A,\ell)|_{[0,1]}$ on $F$, and $\tau(t):=te_1$ off $F$. The \emph{simulated observation process} is $X_k(t):=Y_k(\tau_k(t))$, and the
\emph{\lsimpol} is the law of $(X,\tau)$ under $\gPh_0$.
\end{definition}

\noindent In words: the policy reads its $j$-th reward from arm $k$ when arm $k$'s own clock reaches $\varrho^k_j$, which it can detect, since $\varrho^k_j$ is a stopping
time of arm $k$'s path; it never looks ahead. After $T$ rounds, arm $1$ is run until the end. For the $T$-round Gaussian policies of Lemma~\ref{lem:embed}, the marks are
$\varrho^k_j=j/T$; here they are random.

\begin{lemma}[simulated policies are limit policies]\label{lem:simW}
\uses{def:lowsim, def:alloc, def:risk, def:g-pol}
\textit{Role.} Puts the simulated policy of Theorem~\ref{thm:lowEF} inside the limit game, and computes its risk under the reference instance (Steps~2--3 of the proof).

\noindent\textit{Setting.} $T$, $\pi$, marks and reward map as in Definition~\ref{def:lowsim}; $\tau$, $X$ and $Q$ the simulated allocation, observation process and limit
policy.

\noindent Then $Q\in\Wc$ and, for every $h\in\R^K$,
\[
\rk(Q,h)=\gEh_0\big[\lM_h\,\lphi_h(\tau(1))\big].
\]
\end{lemma}
\begin{proof}
\uses{lem:g-disc, lem:staircase, lem:g-compact, lem:g-cont, def:g-pol, def:alloc, def:risk}
\textit{Idea.} Round every mark up to the grid $2^{-m}\Nat$. The policy that waits for the rounded marks decides only at grid times, so it is a discrete policy, which is
a limit policy with a known risk (Lemma~\ref{lem:g-disc}). Rounding moves the allocation by at most $4T2^{-m}$ (Lemma~\ref{lem:staircase}); compactness and continuity
(Lemmas~\ref{lem:g-compact} and~\ref{lem:g-cont}) pass to the limit $m\to\infty$.

\textit{Plan.} Step~1 defines the rounded allocation; Step~2 shows that it is the allocation of a discrete policy; Step~3 proves the convergence; Step~4 concludes.

Step 1: \emph{the rounded allocation.} Fix $m\ge1$ and let $\varrho^{k,m}_j:=\lceil2^m\varrho^k_j\rceil2^{-m}$ (and $\infty$ if $\varrho^k_j=\infty$). These are
nondecreasing in $j$, with $\varrho^{k,m}_0=0$ and
\[
\varrho^k_j\le\varrho^{k,m}_j\le\varrho^k_j+2^{-m}.
\]
On $F$, let $\ell^m_n$ be defined as $\ell_n$ with the rounded marks, and $\tau^m:=\lGam(A,\ell^m)|_{[0,1]}$, with the \emph{same} arms $A$: the rewards, hence the
decisions, are still read at the true marks. The rounds start and end at the times $s^m_n:=\sum_{r<n}\ell^m_r=\sum_k\varrho^{k,m}_{N_k(n)}$ (telescoping over the rounds of
each arm), which belong to $2^{-m}\Nat$.

Step 2: \emph{$\tau^m$ is the allocation of a discrete policy.} Let $t_i:=i2^{-m}$, $0\le i\le2^m$. Consider the discrete policy $\pi^m$ with grid $(t_i)$ that gives each
cell $[t_i,t_{i+1}]$ entirely to one arm $\kappa_i$, computed at time $t_i$ as follows: from $n:=0$, while $n<T$ and the arm $k:=A_{n+1}$ satisfies
$\varrho^k_{N_k(n)+1}\le\tau^m_k(t_i)$ (round $n+1$ is complete), read $\zeta_{n+1}=Z_{k,N_k(n)+1}$, compute $A_{n+2}$, and increase $n$; then $\kappa_i:=A_{n+1}$. Here
$\tau^m(t_i)$ denotes the allocation of $\pi^m$ itself, and $A_{T+1}=1$.
\begin{itemize}
\item \emph{The rule is Borel in $(U,X^m|_{[0,t_i]})$}, where $X^m_k(t):=Y_k(\tau^m_k(t))$. The allocation on $[0,t_i]$ is determined by $\kappa_0,\dots,\kappa_{i-1}$, and
arm $k$'s path $Y_k$ on $[0,\tau^m_k(t_i)]$ is read off $X^m_k$ on $[0,t_i]$. Since $\varrho^k_j$ is a $\lG^k$-stopping time and $Z_{k,j}$ is
$\lG^k_{\varrho^k_j}$-measurable (it is a function of $U$ and of the path stopped at $\varrho^k_j$), the variables $\1\{\varrho^k_j\le s\}$ and $Z_{k,j}\1\{\varrho^k_j\le s\}$
are $\lG^k_s$-measurable, hence Borel functions of $(U,Y_k|_{[0,s]})$ (Doob--Dynkin). The own times $s=\tau^m_k(t_i)$ range over the countable set $2^{-m}\Nat$, and the
loop has at most $T$ iterations. By induction on $i$, $\kappa_i$ is a Borel function of $(U,X^m|_{[0,t_i]})$, so $\pi^m\in\gDc$.
\item \emph{On $F$, its allocation is $\lGam(A,\ell^m)$.} By induction on $i$: assume the two allocations agree on $[0,t_i]$. Round $n$ of $\lGam(A,\ell^m)$, with arm $k:=A_n$
and $j:=N_k(n)$, moves arm $k$'s own time from $\varrho^{k,m}_{j-1}$ to $\varrho^{k,m}_j$; its own time at a grid time is in $2^{-m}\Nat$, and for $s\in2^{-m}\Nat$,
$s\ge\varrho^{k,m}_j$ iff $s\ge\varrho^k_j$. Hence the test of the loop holds for a started round iff this round is complete at $t_i$, and the loop ends with
$n=n_i$, the number of rounds complete at $t_i$. Each round starts and ends on the grid (Step~1), so round $n_i+1$ (or the final run of arm $A_{T+1}=1$ if
$n_i=T$) covers the whole cell $[t_i,t_{i+1}]$, and $\kappa_i=A_{n_i+1}$ is the arm that $\lGam(A,\ell^m)$ moves on this cell.
\end{itemize}
From now on, $\tau^m$ denotes the allocation of $\pi^m$ (which equals $\lGam(A,\ell^m)|_{[0,1]}$ on $F$), and $Q_m:=\gQp{\pi^m}$ the law of $(X^m,\tau^m)$ under $\gPh_0$.

Step 3: \emph{convergence.} On $F$, Lemma~\ref{lem:staircase} and Step~1 give, with $k:=A_{n+1}$ and $j:=N_k(n)$ in each term,
\begin{align*}
\sup_{t\le1}\|\tau^m(t)-\tau(t)\|_1&\le2\sum_{n<T}|\ell^m_n-\ell_n| &&\text{(Lemma~\ref{lem:staircase})}\\
&\le2\sum_{n<T}\Big(\big|\varrho^{k,m}_{j+1}-\varrho^k_{j+1}\big|+\big|\varrho^{k,m}_j-\varrho^k_j\big|\Big)\ \le\ 4T2^{-m} &&\text{(Step~1)}.
\end{align*}
Since $\tau^m_k(t)$ and $\tau_k(t)$ lie in $[0,1]$ and $Y_k$ is uniformly continuous on $[0,1]$, $\sup_{t\le1}|X^m_k(t)-X_k(t)|\to0$. So $(X^m,\tau^m)\to(X,\tau)$ in
$\Om$ on $F$, that is $\gPh_0$-a.s., and $Q_m\to Q$ weakly.

Step 4: \emph{conclusion.} The own times of $\pi^m$ stay in $[0,1]$, so, by induction over the cells, $X^m|_{[0,1]}$ and $\tau^m(1)$ are functions of $U$ and of the paths
on $[0,1]$: $\lphi_h(\tau^m(1))$ is $\lFo$-measurable. Hence $Q\in\Wc$, because $Q_m\in\Wc$ (Lemma~\ref{lem:g-disc}(b)) and $\Wc$ is closed (Lemma~\ref{lem:g-compact}), and
\begin{align*}
\rk(Q,h)&=\lim_{m\to\infty}\rk(Q_m,h) &&\text{(Lemma~\ref{lem:g-cont}, }Q_m\to Q\text{ in }\Wc\text{)}\\
&=\lim_{m\to\infty}\gEh_h\big[\lphi_h(\tau^m(1))\big] &&\text{(Lemma~\ref{lem:g-disc}(c))}\\
&=\lim_{m\to\infty}\gEh_0\big[\lM_h\,\lphi_h(\tau^m(1))\big] &&\text{(}\gPh_h=\lM_h\cdot\gPh_0\text{ on }\lFo\text{)}\\
&=\gEh_0\big[\lM_h\,\lphi_h(\tau(1))\big] &&\text{(Step~3, dominated convergence)}.
\end{align*}
In the last line, $\lphi_h$ is continuous on $\gSmp$ with $0\le\lphi_h\le\max_ih_i-\min_ih_i$, and $\gEh_0\lM_h=1$.
\end{proof}

\subsection{Proof of Theorem~\ref{thm:lowEF}}\label{sec:app-low-proofEF}

\begin{proof}[Proof of Theorem~\ref{thm:lowEF}]
\proves{thm:lowEF}
\uses{lem:simW, lem:skor, def:lowsim, lem:staircase, lem:timechange, lem:g-cont, def:g-pol, def:risk, def:alloc, def:regret}
\textit{Setting.} The family, $\lthz$, $\beta$, $\lJ$, $\lcth$, $\lnub$, $\lmq$, $\lvb$, $S$, $H$ as in the theorem; an integer $T\ge\lTone$; a $T$-round policy $\pi$. For $h\in S$ let
\[
\delta_k:=\frac{h_k}{\lsd\sqrt T},\qquad \theta_k:=\lthz+\delta_k\qquad(k\in[K]).
\]
Since $T\ge(2H/(\lsd\beta))^2$, $2|\delta_k|\le2H/(\lsd\sqrt T)\le\beta$, so $\lthz+2\delta_k\in\lJ$ and $\lnuT$ is defined; since $T\ge(2\lcth H/\lsd^3)^2$, $\letT\le\frac12$.

\textit{Goal.} A limit policy $\tilde Q$, independent of $h$, with $|\mathcal E_h|\le\lepT$ below, where
\[
\mathcal E_h:=\gEh_0\big[\lL_h\,\lphi_h(N(T)/T)\big]-\rk(\tilde Q,h),
\]
and the gap bounds \eqref{eq:low-gap}; (i)--(iii) follow at the end.

\textit{Idea.} Read the rewards off Brownian paths by Skorokhod's embedding, so that $\pi$ becomes a limit policy (Lemma~\ref{lem:simW}); then compare the two risks under
$\gPh_0$, where the regret carries the likelihood ratio $\lL_h$ of the full array and the limit risk the Cameron--Martin density $\lM_h$.

\textit{Plan.} Steps~1--2 build the simulated policy; Step~3 writes both risks under $\gPh_0$; Step~4 compares the gaps, Step~5 the allocations, Step~6 the likelihood ratios;
Step~7 concludes.

Step 1: \emph{an embedded stack of rewards.} Under $\gPh_0$, $B_k(u):=\sqrt T\,Y_k(u/T)$ is a standard Brownian motion, and $(U_{k,j})_{j\ge1}$ are i.i.d.\ uniform, independent of
$B_k$. Let $(\rho^k_j)_{j\ge0}$ be the stopping times of Lemma~\ref{lem:skor} for $(B_k,(U_{k,j})_j,\lnub)$, and
\[
\varrho^k_j:=\rho^k_j/T,\qquad Z_{k,j}:=\lmuz+\lsd\,\big(B_k(\rho^k_j)-B_k(\rho^k_{j-1})\big)=\lmuz+\lsd\sqrt T\,\big(Y_k(\varrho^k_j)-Y_k(\varrho^k_{j-1})\big).
\]
The marks $\varrho^k_j$ are $\lG^k$-stopping times, finite $\gPh_0$-a.s.: $\{\rho^k_j\le uT\}\in\sigma(U_{k,\cdot})\vee\sigma(B_k(v):v\le uT)$, and
$\sigma(B_k(v):v\le uT)=\sigma(Y_k(v):v\le u)$. If $Y\sim\lnub$ then $\lmuz+\lsd Y\sim\nu_{\lthz}$. The arms are independent under $\gPh_0$ and use disjoint coordinates of the
randomization, so Lemma~\ref{lem:skor}(i) gives:
\begin{equation}\label{eq:low-law}
\begin{minipage}{0.8\linewidth}\centering
under $\gPh_0$, $\big(U_0,(Z_{k,j})_{k\in[K],j\ge1}\big)$ has the law of the randomization\\ and the stack of the instance $(\nu_{\lthz},\dots,\nu_{\lthz})$.
\end{minipage}
\end{equation}
By telescoping, for every $n\ge0$,
\begin{equation}\label{eq:low-sum}
\sum_{j\le n}(Z_{k,j}-\lmuz)=\lsd\,B_k(\rho^k_n)=\lsd\sqrt T\,Y_k(\varrho^k_n).
\end{equation}

Step 2: \emph{the simulated policy.} Let $\tau$, $X$ and $\tilde Q$ be the simulated allocation, observation process and limit policy (Definition~\ref{def:lowsim}) of $\pi$, with
the marks $\varrho^k_j$ and the reward map $z(y,u):=\lmuz+\lsd y$, so that the simulated rewards are the $Z_{k,j}$ of Step~1. By Lemma~\ref{lem:simW},
\begin{equation}\label{eq:low-limrisk}
\tilde Q\in\Wc,\qquad \rk(\tilde Q,h)=\gEh_0\big[\lM_h\,\lphi_h(\tau(1))\big]\quad\text{for every }h\in\R^K .
\end{equation}
The construction does not involve $h$, so $\tilde Q$ does not depend on $h$. The variables $A_n$, $N_k(n)$ are those of $\pi$ run on the stack $(Z_{k,j})$ with randomization $U_0$.

Step 3: \emph{the regret as an expectation under $\gPh_0$.}\dt{l:lr} Let
\[
\lL_h:=\prod_{k\in[K]}\prod_{j\le T}\frac{d\nu_{\theta_k}}{d\nu_{\lthz}}(Z_{k,j})=\prod_{k\in[K]}\exp\Big(\delta_k\sum_{j\le T}(Z_{k,j}-\lmuz)-T\big[\lpsi(\theta_k)-\lpsi(\lthz)-\delta_k\lmuz\big]\Big).
\]
In the stack model, this function of the full array is the density of the law of the full array under $\lnuT$ with respect to its law under $(\nu_{\lthz},\dots,\nu_{\lthz})$,
and $N_k(T)$ is a function of the full array. With \eqref{eq:low-law},
\begin{equation}\label{eq:low-regret}
\RT(\pi,\lnuT)=\sum_{k\in[K]}\Delta_k\,\gEh_0\big[\lL_h\,N_k(T)\big],\qquad \Delta_k:=\lpsi'(\theta_\star)-\lpsi'(\theta_k),\quad \theta_\star:=\lthz+\max_j\delta_j .
\end{equation}
Here $\Delta_k$ is the gap of arm $k$ in $\lnuT$: $\lpsi'$ is increasing ($\lpsi''>0$), so the best arm is one with the largest $h_k$.

Step 4: \emph{the gaps.} We show
\begin{equation}\label{eq:low-gap}
\lsd\sqrt T\,(1+\letT)\,\gEh_0\big[\lL_h\,\lphi_h(N(T)/T)\big]\ \ge\ \RT(\pi,\lnuT)\ \ge\ \lsd\sqrt T\,(1-\letT)\,\gEh_0\big[\lL_h\,\lphi_h(N(T)/T)\big].
\end{equation}
On $\lJ$, $|\lpsi''(\theta)-\lsd^2|\le\lcth|\theta-\lthz|$ (mean value theorem), and every $\theta$ between $\theta_k$ and $\theta_\star$ has $|\theta-\lthz|\le H/(\lsd\sqrt T)$. So
\begin{align*}
\Delta_k&=\int_{\theta_k}^{\theta_\star}\lpsi''(\theta)\,d\theta &&\\
&\ge(\theta_\star-\theta_k)\Big(\lsd^2-\frac{\lcth H}{\lsd\sqrt T}\Big) &&\text{(bound on }\lpsi''\text{ on }\lJ\text{)}\\
&=\frac{\lsd\,(1-\letT)}{\sqrt T}\,\gGp_k(h) &&\text{(}\theta_\star-\theta_k=\gGp_k(h)/(\lsd\sqrt T)\text{; definition of }\letT\text{)},
\end{align*}
and in the same way $\Delta_k\le\lsd(1+\letT)\gGp_k(h)/\sqrt T$. Since $\lL_h\ge0$, $N_k(T)\ge0$ and $\sum_k\gGp_k(h)N_k(T)=T\,\lphi_h(N(T)/T)$, \eqref{eq:low-regret} gives
\eqref{eq:low-gap}.

Step 5: \emph{the allocations are close.} Let $\lDT:=\frac1T\sum_{k\in[K]}\max_{n\le T}|\rho^k_n-n|$. We show
\begin{equation}\label{eq:low-alloc}
\Big\|\frac{N(T)}T-\tau(1)\Big\|_1\le2\lDT\quad\gPh_0\text{-a.s.},\qquad \|\lDT\|_{L^2(\gPh_0)}\le2K\sqrt{\lvb/T}.
\end{equation}
On $F$ (all marks finite, a $\gPh_0$-a.s.\ event), $\tau=\lGam(A,\ell)$ passes through the points $\sigma^{(n)}:=(\varrho^k_{N_k(n)})_k$ at the times
$s_n=\sum_k\varrho^k_{N_k(n)}$ (Lemma~\ref{lem:staircase}). Two cases:
\begin{itemize}
\item if $s_T\ge1$, then $\tau(1)\le\sigma^{(T)}$ componentwise (the path is nondecreasing), and $\|\sigma^{(T)}-\tau(1)\|_1=s_T-1$;
\item if $s_T<1$, then $\tau(1)=\sigma^{(T)}+(1-s_T)e_1$, and $\|\sigma^{(T)}-\tau(1)\|_1=1-s_T$.
\end{itemize}
In both cases, since $\sum_kN_k(T)=T$ and $N_k(T)\le T$,
\begin{gather*}
\|\sigma^{(T)}-\tau(1)\|_1=|s_T-1|=\Big|\sum_k\frac{\rho^k_{N_k(T)}-N_k(T)}T\Big|\le\lDT,\\
\Big\|\sigma^{(T)}-\frac{N(T)}T\Big\|_1=\sum_k\frac{|\rho^k_{N_k(T)}-N_k(T)|}T\le\lDT .
\end{gather*}
The triangle inequality gives the first claim. For the second, by Lemma~\ref{lem:timechange}(i) applied to $B_k$ (with $v\le\lvb$), each of the $K$ terms of $\lDT$ has
$L^2(\gPh_0)$ norm at most $2\sqrt{T\lvb}/T$.

Step 6: \emph{the likelihood ratios are close.} We show
\begin{equation}\label{eq:low-lr}
\gEh_0|\lL_h-\lM_h|\le2e^{KH^2}\lambda_T,\qquad \lambda_T:=KH\Big(\frac{\lvb}T\Big)^{1/4}+\frac{K\lcth H^3}{6\lsd^3\sqrt T}.
\end{equation}
Let $q(\delta):=\lpsi(\lthz+\delta)-\lpsi(\lthz)-\delta\lpsi'(\lthz)-\frac{\lsd^2}2\delta^2$; by Taylor's formula, $|q(\delta)|\le\lcth|\delta|^3/6$ for $|\delta|\le\beta$. Since
$T\lsd^2\delta_k^2/2=h_k^2/2$ and, by \eqref{eq:low-sum}, $\delta_k\sum_{j\le T}(Z_{k,j}-\lmuz)=h_kY_k(\varrho^k_T)$,
\[
\log\lL_h-\log\lM_h=\sum_{k\in[K]}\Big[h_k\big(Y_k(\varrho^k_T)-Y_k(1)\big)-T\,q(\delta_k)\Big].
\]
We bound the two terms:
\begin{itemize}
\item $\|Y_k(\varrho^k_T)-Y_k(1)\|_{L^2(\gPh_0)}=\|B_k(\rho^k_T)-B_k(T)\|_{L^2(\gPh_0)}/\sqrt T\le(T\lvb)^{1/4}/\sqrt T=(\lvb/T)^{1/4}$ (Lemma~\ref{lem:timechange}(ii));
\item $T|q(\delta_k)|\le T\,\lcth H^3/(6\lsd^3T^{3/2})=\lcth H^3/(6\lsd^3\sqrt T)$ (Taylor, $|\delta_k|\le H/(\lsd\sqrt T)$).
\end{itemize}
Hence $\|\log\lL_h-\log\lM_h\|_{L^2(\gPh_0)}\le\lambda_T$. Next, $\gEh_0\lM_h^2=e^{|h|^2}\le e^{KH^2}$, and
\begin{align*}
\gEh_0\lL_h^2&=\prod_k\exp\Big(T\big[\lpsi(\lthz+2\delta_k)-2\lpsi(\lthz+\delta_k)+\lpsi(\lthz)\big]\Big) &&\text{(\eqref{eq:low-law}; Laplace transform)}\\
&\le\prod_k\exp\Big(T\delta_k^2\big(\lsd^2+2\lcth|\delta_k|\big)\Big) &&\text{(second difference)}\\
&\le\prod_k\exp\big(h_k^2(1+2\letT)\big)\ \le\ e^{2KH^2} &&\text{(}T\delta_k^2=h_k^2/\lsd^2\text{, }\letT\le\tfrac12\text{)}.
\end{align*}
In the first line, the Laplace transform of $\nu_{\lthz}$ is finite at $2\delta_k$ because $\lthz+2\delta_k\in\lJ$. In the second, the second difference equals
$\delta_k^2\lpsi''(\vartheta_k)$ for some $\vartheta_k$ between $\lthz$ and $\lthz+2\delta_k$, and $\lpsi''(\vartheta_k)\le\lsd^2+2\lcth|\delta_k|$ on $\lJ$. In the third,
$2\lcth|\delta_k|/\lsd^2\le2\lcth H/(\lsd^3\sqrt T)=2\letT$.
Since $|e^a-e^b|\le|a-b|(e^a+e^b)$, Cauchy--Schwarz gives
$\gEh_0|\lL_h-\lM_h|\le\lambda_T\big(\|\lL_h\|_{L^2}+\|\lM_h\|_{L^2}\big)\le2e^{KH^2}\lambda_T$.

Step 7: \emph{conclusion.} We bound $\mathcal E_h$. Adding and subtracting,
\begin{align*}
\mathcal E_h&=\gEh_0\big[(\lL_h-\lM_h)\,\lphi_h(N(T)/T)\big] &&\\
&\qquad+\gEh_0\big[\lM_h\big(\lphi_h(N(T)/T)-\lphi_h(\tau(1))\big)\big] &&\text{(\eqref{eq:low-limrisk})}\\
|\mathcal E_h|&\le2H\,\gEh_0|\lL_h-\lM_h|+H\,\gEh_0\big[\lM_h\,\|N(T)/T-\tau(1)\|_1\big] &&\text{(facts \eqref{eq:low-L})}\\
&\le4He^{KH^2}\lambda_T+2H\,\|\lM_h\|_{L^2}\|\lDT\|_{L^2} &&\text{(\eqref{eq:low-lr}; \eqref{eq:low-alloc}; Cauchy--Schwarz)}\\
&\le4He^{KH^2}\lambda_T+4KHe^{KH^2/2}\sqrt{\lvb/T} &&\text{(}\|\lM_h\|_{L^2}\le e^{KH^2/2}\text{; \eqref{eq:low-alloc})}\\
&=\lepT &&\text{(definitions of }\lambda_T\text{, }\lepT\text{)}.
\end{align*}
Hence $|\gEh_0[\lL_h\lphi_h(N(T)/T)]-\rk(\tilde Q,h)|\le\lepT$ for every $h\in S$. Also $0\le\rk(\tilde Q,h)\le2H$, by \eqref{eq:low-limrisk}, \eqref{eq:low-L} and
$\gEh_0\lM_h=1$. With \eqref{eq:low-gap} and $\letT\le\frac12$:
\begin{gather*}
\frac{\RT(\pi,\lnuT)}{\lsd\sqrt T}\le(1+\letT)\big(\rk(\tilde Q,h)+\lepT\big)\le\rk(\tilde Q,h)+2H\letT+\frac32\lepT,\\
\frac{\RT(\pi,\lnuT)}{\lsd\sqrt T}\ge(1-\letT)\big(\rk(\tilde Q,h)-\lepT\big)\ge\rk(\tilde Q,h)-2H\letT-\lepT .
\end{gather*}
This is (i), in both forms.

(ii) For every $\pi$, by (i) and $\tilde Q\in\Wc$,
$\sup_{h\in S}\RT(\pi,\lnuT)/(\lsd\sqrt T)\ge(1-\letT)\big(\sup_{h\in S}\rk(\tilde Q,h)-\lepT\big)\ge(1-\letT)\big(\VK_K(S)-\lepT\big)$, since
$\VK_K(S)=\min_{Q\in\Wc}\sup_{h\in S}\rk(Q,h)$. Take the infimum over $\pi$.

(iii) Both sides of (i) are continuous in $h\in S$, hence measurable: $h\mapsto\rk(\tilde Q,h)$ by Lemma~\ref{lem:g-cont}; $h\mapsto\RT(\pi,\lnuT)$ by \eqref{eq:low-regret},
since $\Delta_k$ is continuous in $h$ and $\gEh_0[\lL_hN_k(T)]$ is continuous by dominated convergence ($N_k(T)\le T$, $\lL_h$ is continuous in $h$, and for $h\in S$,
$\lL_h\le C\prod_k\exp\big(\frac\beta2\sum_{j\le T}|Z_{k,j}-\lmuz|\big)$ with $C$ independent of $h$, which is integrable because $\lthz\pm\frac\beta2\in\operatorname{int}\lTh$).
Integrating the lower bound of (i) over $\Pi$,
$\int\RT(\pi,\lnuT)\Pi(dh)/(\lsd\sqrt T)\ge(1-\letT)\big(\int\rk(\tilde Q,h)\Pi(dh)-\lepT\big)\ge(1-\letT)\big(\BK_K(\Pi)-\lepT\big)$. Take the infimum over $\pi$.
\end{proof}

\subsection{Proof of Proposition~\ref{prop:lowtilt}}\label{sec:app-low-prooftilt}

\begin{proof}[Proof of Proposition~\ref{prop:lowtilt}]
\proves{prop:lowtilt}
\uses{thm:lowEF, lem:skor, def:lowsim, lem:timechange, def:alloc, def:regret, def:risk, def:g-pol}
\textit{Setting.} $\nu_0$, $\mu_0$, $g$, $I$, $G$, $\dot m$, $\lcg$, $S$, $H$, $\bar G$ as in the proposition; $\bar g:=g/\sqrt I$, which is centred, of variance $1$ under $\nu_0$,
with $|\bar g|\le\bar G$; $\lnub$ the law of $\bar g(X)$, $X\sim\nu_0$; an integer $T\ge4H^2\bar G^2$; a $T$-round policy $\pi$. For $h\in S$ let $\delta_k:=h_k/\sqrt{IT}$;
then $|\delta_kg|\le|h_k|\bar G/\sqrt T\le\frac12$.

\textit{Goal.} $\tilde Q\in\Wc$, independent of $h$, with $|\RT(\pi,\lnug)/(\lcg\sqrt T)-\rk(\tilde Q,h)|\le\lepTp$ for $h\in S$.

\textit{Idea.} The proof of Theorem~\ref{thm:lowEF}, with the standardized score $\bar g(X)$ embedded in place of the standardized reward, the reward itself being drawn from its
conditional law given the score. We list the changes; Steps~2 and~5 are unchanged.

Step 1: \emph{embedding.} Let $(y,A)\mapsto\kappa(y,A)$ be a regular conditional law of $X$ given $\bar g(X)$ under $X\sim\nu_0$, and $z(y,u):=\inf\{x:\kappa(y,(-\infty,x])\ge u\}$
its quantile function, a Borel map with $z(y,U')\sim\kappa(y,\cdot)$ for $U'$ uniform. Since the joint law of $(\bar g(X),X)$ is $\lnub(dy)\kappa(y,dx)$ and $\bar g(X)=\bar g(X)$,
$\kappa(y,\{\bar g=y\})=1$ for $\lnub$-a.e.\ $y$. Let $\rho^k_j$, $\varrho^k_j:=\rho^k_j/T$ be as in Step~1 of the proof of Theorem~\ref{thm:lowEF}, for this $\lnub$, and
\[
Z_{k,j}:=z\big(B_k(\rho^k_j)-B_k(\rho^k_{j-1}),\,U'_{k,j}\big).
\]
By Lemma~\ref{lem:skor}(i), the pairs (increment, $U'_{k,j}$) are i.i.d.\ in $j$, independent across arms and of $U_0$, with first coordinate of law $\lnub$. Hence the $Z_{k,j}$
are i.i.d.\ with law $\int\lnub(dy)\kappa(y,\cdot)=\nu_0$, so \eqref{eq:low-law} holds with $\nu_0$ in place of $\nu_{\lthz}$; and $\bar g(Z_{k,j})=B_k(\rho^k_j)-B_k(\rho^k_{j-1})$
$\gPh_0$-a.s., so
\[
\sum_{j\le n}\bar g(Z_{k,j})=B_k(\rho^k_n)=\sqrt T\,Y_k(\varrho^k_n)\qquad\gPh_0\text{-a.s.}
\]
The rewards are of the form of Definition~\ref{def:lowsim}, with this reward map $z$; Step~2 is unchanged, and \eqref{eq:low-limrisk} holds.

Step 3: \emph{regret under $\gPh_0$.} Let $\lL_h:=\prod_k\prod_{j\le T}(1+\delta_kg(Z_{k,j}))$, the density of the full array under $\lnug$ with respect to its law under
$\nu_0^{\otimes}$ (since $d\nu^g_t/d\nu_0=1+tg$). Then \eqref{eq:low-regret} holds, with $\Delta_k$ the gaps of $\lnug$.

Step 4: \emph{gaps, exactly.} The mean of $\nu^g_t$ is $\int x(1+tg(x))\nu_0(dx)=\mu_0+t\dot m$ (because $\int g\,d\nu_0=0$), increasing in $t$. Hence
$\Delta_k=\dot m\,\gGp_k(h)/\sqrt{IT}=\lcg\,\gGp_k(h)/\sqrt T$, and \eqref{eq:low-regret} gives the identity
\[
\RT(\pi,\lnug)=\lcg\sqrt T\,\gEh_0\big[\lL_h\,\lphi_h(N(T)/T)\big].
\]

Step 6: \emph{likelihood ratios.} Write $\log(1+u)=u-\frac{u^2}2+q(u)$, with $|q(u)|\le\sum_{n\ge3}\frac{|u|^n}3=\frac{|u|^3}{3(1-|u|)}\le\frac23|u|^3$ for $|u|\le\frac12$, and let
$Q_k:=\frac1T\sum_{j\le T}\bar g(Z_{k,j})^2$. Since $\delta_kg=h_k\bar g/\sqrt T$, Step~1 gives $\sum_{j\le T}\delta_kg(Z_{k,j})=h_kY_k(\varrho^k_T)$ and
$\frac12\sum_{j\le T}\delta_k^2g(Z_{k,j})^2=\frac{h_k^2}2Q_k$; so
\[
\log\lL_h-\log\lM_h=\sum_{k\in[K]}\Big[h_k\big(Y_k(\varrho^k_T)-Y_k(1)\big)-\frac{h_k^2}2\,(Q_k-1)+\sum_{j\le T}q\big(\delta_kg(Z_{k,j})\big)\Big].
\]
We bound the three terms:
\begin{itemize}
\item $\|Y_k(\varrho^k_T)-Y_k(1)\|_{L^2}\le(\lvb/T)^{1/4}$, as in Step~6 of the proof of Theorem~\ref{thm:lowEF} (Lemma~\ref{lem:timechange}(ii));
\item $\|Q_k-1\|_{L^2}=(\operatorname{Var}\bar g(Z_{k,1})^2/T)^{1/2}\le(\lmq/T)^{1/2}$ (the $\bar g(Z_{k,j})^2$ are i.i.d.\ with mean $1$ and second moment $\lmq$);
\item $\sum_{j\le T}|q(\delta_kg(Z_{k,j}))|\le T\cdot\frac23\big(H\bar G/\sqrt T\big)^3=\frac23H^3\bar G^3/\sqrt T$.
\end{itemize}
Hence $\|\log\lL_h-\log\lM_h\|_{L^2(\gPh_0)}\le\lambda'_T$. Next,
\[
\gEh_0\lL_h^2=\prod_k\Big(\int(1+\delta_kg)^2d\nu_0\Big)^T=\prod_k\big(1+\delta_k^2I\big)^T\le\exp\Big(\sum_kT\delta_k^2I\Big)=e^{|h|^2}\le e^{KH^2},
\]
using $\int g\,d\nu_0=0$, $\int g^2d\nu_0=I$ and $1+x\le e^x$. With $\gEh_0\lM_h^2\le e^{KH^2}$ and $|e^a-e^b|\le|a-b|(e^a+e^b)$, Cauchy--Schwarz gives
$\gEh_0|\lL_h-\lM_h|\le2e^{KH^2/2}\lambda'_T$.

Step 7: \emph{conclusion.} As in Step~7 of the proof of Theorem~\ref{thm:lowEF}, with the bound of Step~6,
\[
\Big|\gEh_0\big[\lL_h\lphi_h(N(T)/T)\big]-\rk(\tilde Q,h)\Big|\le2H\cdot2e^{KH^2/2}\lambda'_T+2H\,e^{KH^2/2}\cdot2K\sqrt{\lvb/T}=\lepTp .
\]
By the identity of Step~4, this is the claim.

\textit{The bound $\lcg\le(\int(x-\mu_0)^2d\nu_0)^{1/2}$.} By Cauchy--Schwarz, $\dot m=\int(x-\mu_0)g(x)\nu_0(dx)\le(\int(x-\mu_0)^2d\nu_0)^{1/2}\sqrt I$.
\end{proof}

\subsection{Notation of this section}\label{sec:app-low-notation}
Every symbol is a link to its definition.
\noindent $\lvp$, $\lTh$, $\lpsi$, $\nu_\theta$\quad base measure, natural parameters, log-partition, exponential family\hfill{\small Sec.~\ref{sec:app-low-setting}}\par
\noindent $\lthz$, $\lmuz$, $\lsd$\quad reference parameter, its mean and standard deviation\hfill{\small Sec.~\ref{sec:app-low-setting}}\par
\noindent $\lnuT$\quad local instance of the family\hfill{\small Sec.~\ref{sec:app-low-setting}}\par
\noindent $\gOs$, $\gPh_h$, $\gWi_c$, $\gDc$, $\gQp\pi$\quad space of the arms, its laws, discrete policies and their laws\hfill{\small Def.~\ref{def:g-pol}}\par
\noindent $\lG^k$, $\lFo$, $\lM_h$, $\lphi_h$\quad filtration of arm $k$, paths up to time $1$, Cameron--Martin density, loss\hfill{\small Sec.~\ref{sec:app-low-setting}}\par
\noindent $\lJ$, $\lcth$, $\lnub$, $\lmq$, $\lvb$\quad neighbourhood of $\lthz$, bound on $\lpsi'''$, standardized law, its fourth moment, bound on $\E\rho_1^2$\hfill{\small Thm~\ref{thm:lowEF}}\par
\noindent $\lTone$, $\letT$, $\lepT$\quad threshold, relative error on the gaps, additive error\hfill{\small Thm~\ref{thm:lowEF}}\par
\noindent $\lGam(\alpha,\ell)$\quad staircase path\hfill{\small Lem.~\ref{lem:staircase}}\par
\noindent $\varrho^k_j$, $Z_{k,j}$, $\tau$, \lsimpol\quad marks, simulated rewards, \lsimal, its law\hfill{\small Def.~\ref{def:lowsim}}\par
\noindent $\lL_h$, $\lDT$\quad likelihood ratio of the full array, time-change error\hfill{\small Proof of Thm~\ref{thm:lowEF}}\par
\noindent $\lcg$, $\lnug$, $\lepTp$\quad efficiency of a bounded perturbation, perturbed local instance, error\hfill{\small Prop.~\ref{prop:lowtilt}}\par

\section{Proof of the upper bound}\label{sec:app-upper}

\newcommand{\uSimp}{\sy{u:batched}{\mathcal S_K}}
\newcommand{\urk}{\sy{u:risk}{r}}
\newcommand{\uDisc}[3]{\sy{u:disc}{#1_{#2}[#3]}}
\newcommand{\uRG}{\sy{u:RG}{R^{\mathrm G}}}
\newcommand{\umus}{\sy{u:setting}{\mu^\star}}
\newcommand{\uN}{\sy{alg:sc}{N}}
\newcommand{\urho}{\sy{alg:sc}{\rho}}
\newcommand{\uS}{\sy{alg:sc}{S}}
\newcommand{\uio}{\sy{alg:sc}{i_0}}
\newcommand{\umuh}{\sy{alg:sc}{\hat\mu}}
\newcommand{\umuref}{\sy{alg:sc}{\hat\mu_{\mathrm{ref}}}}
\newcommand{\usigh}{\sy{alg:sc}{\hat\sigma}}
\newcommand{\uw}{\sy{alg:sc}{w}}
\newcommand{\uc}{\sy{alg:sc}{c}}
\newcommand{\upulls}{\sy{u:virtual}{n}}
\newcommand{\uTp}{\sy{u:virtual}{T'}}
\newcommand{\unu}{\sy{u:virtual}{\tilde\nu}}
\newcommand{\umusS}{\sy{u:virtual}{\mu^\star_{S}}}
\newcommand{\uD}{\sy{u:dev}{D}}
\newcommand{\uCD}{\sy{u:CD}{C_D}}
\newcommand{\upv}{\sy{u:pvar}{p_{\mathrm{var}}}}
\newcommand{\uF}{\sy{u:decomp}{\mathcal F^{\mathrm{scr}}}}
\newcommand{\uRt}{\sy{u:decomp}{\tilde R}}
\newcommand{\uM}{\sy{u:param}{M}}
\newcommand{\uCt}{\sy{u:param}{C_3}}
\newcommand{\uG}{\sy{u:good}{\mathcal G}}


\medskip
This appendix proves the upper bound of Theorem~\ref{thm:main}(ii), in the quantitative form of Theorem~\ref{thm:u-upper}, and
Proposition~\ref{prop:four}(c), as Proposition~\ref{prop:u-local}. The policy is \SC\ (Definition~\ref{def:sc}). The reader is a specialist of bandits;
standard facts on the weak convergence of processes are cited. We use three facts from Section~\ref{sec:app-game}.
\begin{enumerate}
\item[(G1)] \emph{Batched and limit values agree on compact sets} (Proposition~\ref{prop:four}(b)): $\Vb(H)=\VK_K(H)$ for every compact $H\subset\R^K$. In particular
$\Vb(H)\le\VK_K$, because for every $Q\in\Wc$ the supremum of $\rk(Q,\cdot)$ over $H$ is at most its supremum over $\R^K$.
\item[(G2)] \emph{Near-optimal continuous batched policies exist} (Corollary~\ref{cor:g-cont}): for every compact $H\subset\R^K$ and every $\varepsilon>0$ there is a
continuous batched policy $\pi$ (Definition~\ref{def:u-batched}) with $\sup_{h\in H}\urk(\pi,h)\le\Vb(H)+\varepsilon$. Corollary~\ref{cor:g-cont} gives rules that
are continuous jointly in the randomization and the observed path, which implies the continuity required in Definition~\ref{def:u-batched} (in the observed path,
for the uniform norm, at fixed randomization).
\item[(G3)] \emph{The risk of a batched policy is its risk as a limit policy} (Lemma~\ref{lem:g-disc}): for a batched policy $\pi$, the law of its observations and allocations
$(X,\tau)$ under $h$ is $\lh_h$ times their law $Q_\pi$ under $h=0$, so that $\urk(\pi,h)$ (Definition~\ref{def:u-batched}) equals $\rk(Q_\pi,h)$
(Section~\ref{sec:game}). Hence, with (G1), $\Vb(H)$ is the infimum of $\sup_{h\in H}\urk(\pi,h)$ over the batched policies of Definition~\ref{def:u-batched}.
\end{enumerate}

\paragraph{Dependencies.}
\[
\begin{array}{l}
\text{Lemma~\ref{lem:u-dev}}\ \to\ \text{Lemma~\ref{lem:u-survive}}\ \to\ \text{Lemma~\ref{lem:u-cost}};\qquad\text{Lemma~\ref{lem:u-var}}
\qquad\text{(the screening, Section~\ref{sec:app-u-screen})};\\[1mm]
\text{(G2)}+\text{Theorem~\ref{thm:u-transfer}}+\text{(G1)}\ \to\ \text{Proposition~\ref{prop:u-local}}\ \to\ \text{Corollary~\ref{cor:u-gauss}};\\[1mm]
\text{Lemmas~\ref{lem:u-dev}--\ref{lem:u-var}}+\text{Theorem~\ref{thm:u-transfer}}+\text{(G1)--(G3)}\ \to\ \text{Theorem~\ref{thm:u-upper}}\ \to\
\text{Theorem~\ref{thm:main}(ii)}.
\end{array}
\]

\paragraph{Plan.} Section~\ref{sec:app-u-setting} fixes the setting and states the results. Section~\ref{sec:app-u-skeleton} gives the skeleton of the proof of
Theorem~\ref{thm:u-upper}: a decomposition of the regret into three terms, and what each term costs. Sections~\ref{sec:app-u-screen} and~\ref{sec:app-u-transfer}
state the lemmas, grouped by the question they answer. The proofs follow, one subsection each.

\subsection{Setting and statements}\label{sec:app-u-setting}

\paragraph{Conventions.}\dt{u:setting} Rewards follow the reward-table model: independent entries $X_{k,n}\sim\nu_k$ ($k\in[K]$, $n\ge1$), and the $n$-th pull
of arm $k$ returns $X_{k,n}$; the law of the interaction is the same as in the sequential model \cite[Sec.~4.6]{LS}. Let $\umus:=\max_k\mu_k$ and
$\sigma_k^2:=\mathrm{Var}(\nu_k)$. The external randomization of a policy is one variable $U\sim\mathrm{Unif}[0,1]$, independent of the table: the arm played
at round $t$ is $\varphi_t(U,x_1,\dots,x_{t-1})$ for Borel maps $\varphi_t$, where $x_1,x_2,\dots$ are the rewards received. Every policy has this form, by the
randomization lemma \cite[Lemma~3.22]{Kal02} and a Borel isomorphism between $[0,1]$ and $[0,1]^T$. For $\nu\in\Cab$ with mean $\mu$, Chernoff's bound gives
$\P(|X-\mu|>x)\le2e^{-x^2/(2\sigp^2)}$; integrating this tail \cite[Lemma~1.4]{RH17},
\begin{equation}\label{eq:u-moments}
\begin{gathered}
\E_\nu|X-\mu|^p\le p\,(2\sigp^2)^{p/2}\,\Gamma(p/2)\qquad(p>0),\\
\text{hence}\qquad \E_\nu|X-\mu|^3\le3\sqrt{2\pi}\,\sigp^3\le8\sigp^3\qquad\text{and}\qquad\E_\nu(X-\mu)^4\le16\sigp^4 .
\end{gathered}
\end{equation}

\begin{definition}[batched policy, risk with variances]\label{def:u-batched}\dt{u:batched}\dt{u:risk}
Let $\uSimp:=\{\alpha\in[0,1]^K:\sum_k\alpha_k=1\}$. A \emph{batched policy} $\pi=(m,(t_j)_{j\le m},(F_j)_{j\le m})$ consists of $m\ge1$, decision times
$0=t_0<t_1<\dots<t_m=1$, and Borel maps $F_j:[0,1]\times C([0,t_{j-1}],\R^K)\to\uSimp$. It is \emph{continuous} if, for every $j\in[m]$ and every
$u\in[0,1]$, the map $x\mapsto F_j(u,x)$ is continuous for the uniform norm.

Given $(h,v)\in\R^K\times[0,\infty)^K$, let $W_1,\dots,W_K$ be independent standard Brownian motions, $U\sim\mathrm{Unif}[0,1]$ independent of them, and
$Y_k(s):=h_ks+\sqrt{v_k}\,W_k(s)$ for $s\in[0,1]$. Let $\tau(0):=0$ and, for $j=1,\dots,m$ and $t\in[t_{j-1},t_j]$,
\[
\alpha^{(j)}:=F_j\big(U,X|_{[0,t_{j-1}]}\big),\qquad \tau(t):=\tau(t_{j-1})+(t-t_{j-1})\,\alpha^{(j)},\qquad X_k(t):=Y_k\big(\tau_k(t)\big).
\]
The \emph{risk} is
\[
\urk(\pi,h;v):=\E\Big[\sum_k\big(\max_ih_i-h_k\big)\,\tau_k(1)\Big],\qquad \urk(\pi,h):=\urk(\pi,h;\mathbf 1),\quad \mathbf 1:=(1,\dots,1).
\]
\end{definition}

\noindent Since $\alpha^{(j)}\in\uSimp$, each $\tau_k$ is nondecreasing, $\sum_k\tau_k(t)=t$ and $\tau_k(1)\le1$, so $Y_k$ is only read on $[0,1]$. By
induction on $j$, $X|_{[0,t_{j-1}]}$ is a Borel function of $(W,U)$ (the evaluation $(y,s)\mapsto y(s)$ is continuous on $C([0,1])\times[0,1]$), so the risk is
well defined, and $0\le\urk(\pi,h;v)\le\max_ih_i-\min_ih_i$. This is the batched policy of Section~\ref{sec:game} in coordinates: at time $t_{j-1}$ the
proportions $\alpha^{(j)}$ are a Borel function of the randomization and of the observations so far. A rule that reads only the values $x(t_1),\dots,x(t_{j-1})$
at the ends of the batches is continuous as soon as it is continuous in these values. Variances $v\ne\mathbf 1$ only appear in the proofs: they describe arms whose
variance is not exactly $\sbar^2$.

\begin{definition}[discretization]\label{def:u-disc}\dt{u:disc}
Let $\pi=(m,(t_j),(F_j))$ be a batched policy, $T\ge1$, $m_0\in\R$ and $\sigma>0$, and let $R_j:=\lfloor t_jT\rfloor-\lfloor t_{j-1}T\rfloor$. The $T$-round
policy $\uDisc{\pi}{T}{m_0,\sigma}$ draws $U\sim\mathrm{Unif}[0,1]$; its batch $j$ is made of the rounds $\lfloor t_{j-1}T\rfloor+1,\dots,\lfloor t_jT\rfloor$. Let
$N^{(j)}_k$ be the number of pulls of arm $k$ in batches $1,\dots,j$ ($N^{(0)}:=0$), and $x_{k,1},x_{k,2},\dots$ the successive rewards of arm $k$. At the start
of batch $j$ the policy forms:
\begin{itemize}
\item the normalized partial sums of each arm $k$, linearly interpolated: for $u\in[0,N^{(j-1)}_k/T]$,
\[
\hat z_k(u):=\frac1{\sigma\sqrt T}\Big(\sum_{i\le\lfloor uT\rfloor}(x_{k,i}-m_0)+\big(uT-\lfloor uT\rfloor\big)\big(x_{k,\lfloor uT\rfloor+1}-m_0\big)\Big);
\]
\item the clocks $\hat\tau(t_i):=N^{(i)}/T$ for $i\le j-1$, linearly interpolated on each $[t_{i-1},t_i]$;
\item the observed path $\hat X_k(t):=\hat z_k(\hat\tau_k(t))$ for $t\in[0,t_{j-1}]$;
\end{itemize}
and it computes $\hat\alpha^{(j)}:=F_j\big(U,\hat X|_{[0,t_{j-1}]}\big)$. During batch $j$ it pulls each arm $k$ exactly $\lfloor\hat\alpha^{(j)}_kR_j\rfloor$ times,
in increasing order of $k$, and arm $1$ during the remaining $R_j-\sum_k\lfloor\hat\alpha^{(j)}_kR_j\rfloor$ rounds (fewer than $K$).
\end{definition}

\noindent The path $\hat X$ only uses observed rewards: the coefficient of $x_{k,\lfloor uT\rfloor+1}$ vanishes at $u=N^{(j-1)}_k/T$, and
$\hat\tau_k(t)\le N^{(j-1)}_k/T$ for $t\le t_{j-1}$. It is continuous, starts at $0$, and its restriction to $[0,t_{j-1}]$ does not change when later data
arrive. At the ends of the batches, $\hat X_k(t_i)=(\sigma\sqrt T)^{-1}\sum(x-m_0)$, the sum running over the pulls of $k$ in batches $1,\dots,i$: a rule that
reads only the ends of the batches uses only these normalized sums. The last batch ends at round $\lfloor t_mT\rfloor=T$.

\paragraph{Local Gaussian minimax regret.}\dt{u:RG} For a compact $H\subset\R^K$ and $T\ge1$ let
\[
\uRG_H(T):=\inf_\pi\ \sup_{h\in H}\ \RT\big(\pi,(\mathcal N(h_k/\sqrt T,1))_{k\in[K]}\big),
\]
the infimum being over all $T$-round policies. Proposition~\ref{prop:four}(c) is about $\uRG_{[-M,M]^K}(T)/\sqrt T$.

\paragraph{The policy \SC.}\dt{u:virtual} We use the notation of Definition~\ref{def:sc}: $\uN=\lfloor\delta T/K\rfloor$, the radii $\urho_n$, the active sets
$\uS_n$, the survivors $\uS=\uS_\uN$, the reference arm $\uio$ (ties broken by the smallest index), $\umuref=\umuh_{\uio}(\uN)$, the variance estimates
$\usigh_k^2:=\frac1\uN\sum_{n\le\uN}(X_{k,n}-\umuh_k(\uN))^2$, the garbling variances $\uw_k:=(\sbar^2-\usigh_k^2)_+$ and the map $\uc$. Let $\upulls_k$ be the
number of pulls of arm $k$ in Stage~1 and $\uTp:=T-\sum_k\upulls_k$ the number of remaining rounds. Stage~2 runs $\uDisc{\pi^\circ}{\uTp}{\umuref,\sbar}$
(Definition~\ref{def:u-disc}) on $K$ virtual arms: when it asks for virtual arm $k$, it pulls the real arm $\uc(k)$, observes $X$ and receives $X+\xi$ with a fresh
$\xi\sim\mathcal N(0,\uw_{\uc(k)})$. The \emph{virtual instance} is $\unu=(\unu_k)_{k\in[K]}$, where $\unu_k$ is the law of $X+\xi$ with $X\sim\nu_{\uc(k)}$ and
$\xi\sim\mathcal N(0,\uw_{\uc(k)})$ independent: its mean is $\mu_{\uc(k)}$ and its variance $\sigma^2_{\uc(k)}+\uw_{\uc(k)}$. Let
$\umusS:=\max_{i\in\uS}\mu_i$; since $\uc([K])=\uS$, it is also $\max_k\mu_{\uc(k)}$. Stage~1 uses $\sum_k\upulls_k\le K\uN\le\delta T$ rounds, so
$\uTp\ge(1-\delta)T$.

\noindent Any measurable tie-breaking rule can be used for $i_0$; we take the smallest index.

\paragraph{Statements.}

\begin{theorem}[quantitative upper bound]\label{thm:u-upper}
\uses{def:class, def:regret, def:sc, def:risk}
Let $K\ge2$, $a<b$ and $0<\sbar\le\sigp$. Then, for every $\varepsilon\in(0,1]$, there are $\delta\in(0,\frac12]$, $A\ge8$, a continuous batched policy $\pi^\circ$
and $T_1\in\Nat$ such that, for every $T\ge T_1$,
\[
\sup_{\nu\in\Cab^K}\RT\big(\SC(\delta,A,\pi^\circ),\nu\big)\ \le\ \big(\sbar\,\VK_K+(3\sbar+4)\,\varepsilon\big)\sqrt T .
\]
In particular Theorem~\ref{thm:main}(ii) holds: for every class $\Cl$ satisfying (R1) with variance $\sbar$, $\limsup_T\Rs_{K,\Cl}(T)/\sqrt T\le\sbar\VK_K$.
\end{theorem}


\begin{proposition}[local Gaussian minimax regret]\label{prop:u-local}
\uses{def:batched, def:risk}
Let $K\ge2$ and let $H\subset\R^K$ be compact. Then:
\begin{enumerate}
\item[(i)] $\uRG_H(T)\ge\sqrt T\,\Vb(H)=\sqrt T\,\VK_K(H)$ for every $T\ge1$;
\item[(ii)] $\uRG_H(T)/\sqrt T\to\VK_K(H)$ as $T\to\infty$.
\end{enumerate}
For $H=[-M,M]^K$, (ii) is Proposition~\ref{prop:four}(c).
\end{proposition}

\begin{corollary}[Gaussian rewards: a non-asymptotic lower bound]\label{cor:u-gauss}
\uses{def:class, def:regret, def:risk}
Let $K\ge2$, $a<b$, $\sbar>0$, and let $\mathcal C_{\mathrm G}$ be the class of the Gaussian laws $\mathcal N(\mu,\sigma^2)$ with $\mu\in[a,b]$ and
$\sigma\le\sbar$, or its subclass $\sigma=\sbar$. Then, for every $M>0$ and every $T\ge(2\sbar M/(b-a))^2$,
\[
\Rs_{K,\mathcal C_{\mathrm G}}(T)\ \ge\ \sbar\sqrt T\;\VK_K\big([-M,M]^K\big).
\]
\end{corollary}

\subsection{Skeleton of the proof of Theorem~\ref{thm:u-upper}}\label{sec:app-u-skeleton}

Fix $\nu\in\Cab^K$ and $T$. The regret of \SC\ splits into three terms (Step~1 of Section~\ref{sec:app-u-proof-main}):
\begin{equation}\label{eq:u-decomp}
\RT(\SC,\nu)=\underbrace{\E\Big[\sum_k\Delta_k\upulls_k\Big]}_{\text{screening}}+\underbrace{\E\big[\uTp\,(\umus-\umusS)\big]}_{\text{loss of the best arm}}
+\underbrace{\E\big[\uRt\big]}_{\text{core}},\qquad
\uRt:=\sy{s:regret}{R}_{\uTp}\big(\uDisc{\pi^\circ}{\uTp}{\umuref,\sbar},\unu\big),
\end{equation}
where $\uRt$, the regret of the core on the virtual instance, is evaluated at the parameters produced by Stage~1. Let $\uD$ be the total excess of the
empirical means over the radii (Section~\ref{sec:app-u-screen}); by Lemma~\ref{lem:u-dev}, $\E[\uD]\le K\uCD\sigp e^{-A/2}/\sqrt\uN$.
\begin{itemize}
\item \emph{Screening}: at most $K(b-a)+16\sigp\sqrt{AK\delta}\,\sqrt T+8K\uN\,\E[\uD]$, \emph{whatever the gaps} (Lemma~\ref{lem:u-cost}).
\item \emph{Loss of the best arm}: $\umus-\umusS\le2\uD$ (Lemma~\ref{lem:u-survive}(ii)), so this term is at most $2T\,\E[\uD]=O(e^{-A/2}\delta^{-1/2})\sqrt T$.
\item \emph{Core}: on the good event $\uG$ ($\uD\le\sbar/\sqrt T$ and all variance estimates within $\eta\sbar^2$), the virtual instance has local means in the box
$[-\uM,\uM]^K$, with $\uM=O(\sqrt{AK/\delta})$, and variances in $\sbar^2[1-\eta,1+\eta]$; the transfer theorem (Theorem~\ref{thm:u-transfer}) and (G1)--(G2) give
$\uRt\le\sbar\sqrt T(\VK_K+3\varepsilon)$. Off $\uG$, $\uRt\le T(4\urho_\uN+2\uD)$ (Lemma~\ref{lem:u-survive}(iv)), and
\[
\P(\uG^c)\ \le\ \underbrace{\frac{\sqrt T}{\sbar}\,\E[\uD]}_{O(e^{-A/2}\delta^{-1/2})}\ +\ \underbrace{K\,\upv(\uN,\eta)}_{O(K^2/(\delta T\eta^2))}
\qquad\text{(Markov's inequality and Lemma~\ref{lem:u-var})}.
\]
The first term is small when $\delta$ is small (with $A=6\log(1/\delta)$). The second one, the failure of the variance estimates, is small only when $T$ is
large, at fixed $\delta$ and $\eta$; it is absorbed by the choice of $T_1$.
\end{itemize}
The parameters are chosen in this order: $\delta$ and $A$ (screening terms), the box size $\uM$, the core $\pi^\circ$ (by (G1)--(G2)), the tolerance $\eta$ on the
variances (continuity of the risk), and $T_1$ (transfer theorem and variance estimates).

\paragraph{Why the intermediate gaps are harmless.} Write the gap of an arm as $\Delta=g\sbar/\sqrt T$. The screening costs at most
$(b-a)+16\sigp\sqrt{A\uN}+8\uN\uD$ per arm whatever $g$; the arms with $g\gg\sqrt{A/\delta}$ are eliminated, and the survivors stay in a box of local means of
size $\uM=O(\sqrt{AK/\delta})$. So the core only has to be good on this box, and no uniformity of the diffusion approximation as $|h|\to\infty$ is needed.
This matters, because no batched policy is good on all of $\R^K$. The first allocation $\alpha^{(1)}=F_1(U)$ does not depend on $h$; choose $k$ with
$\E[\alpha^{(1)}_k]\ge1/K$ (it exists since $\sum_k\alpha^{(1)}_k=1$). The instance with $h_k<0$ and the other coordinates $0$ has
\[
\urk(\pi,h)\ \ge\ |h_k|\,\E[\tau_k(t_1)]=|h_k|\,t_1\,\E[\alpha^{(1)}_k]\ \ge\ t_1|h_k|/K ,
\]
which tends to $\infty$ with $|h_k|$.

\subsection{The screening}\label{sec:app-u-screen}

\paragraph{Setting of the subsection.}\dt{u:dev} Fix $\nu\in\Cab^K$, $T$ with $\uN\ge1$, and $A>0$, and run Stage~1 of Definition~\ref{def:sc} on the reward table.
For every arm $k$ (also for the arms eliminated early: this is an analysis device) let $\umuh_k(n)$ be the mean of $X_{k,1},\dots,X_{k,n}$; for
$k\in\uS_{n-1}$ it is the empirical mean computed by the algorithm. Let
\[
\uD_k:=\max_{1\le n\le\uN}\big(|\umuh_k(n)-\mu_k|-\urho_n\big)_+,\qquad \uD:=\sum_{k\in[K]}\uD_k .
\]
By definition, for every $k\in[K]$ and $n\in[\uN]$,
\begin{equation}\label{eq:u-basic}
|\umuh_k(n)-\mu_k|\ \le\ \urho_n+\uD_k .
\end{equation}
The radius $\urho_n=\sigp\sqrt{(2/n)(A+2\log(\uN/n))}$ is decreasing in $n\in[\uN]$, since both factors decrease.

\subsubsection*{Why are the radii rarely violated?}

\begin{lemma}[maximal deviation]\label{lem:u-dev}\dt{u:CD}
\uses{def:sc, def:class}
\textit{Role.} Controls $\uD$ in the three terms of~\eqref{eq:u-decomp}: $\E[\uD]=O(e^{-A/2}/\sqrt\uN)$, without a $\log\uN$ factor.

\noindent Let $X_1,X_2,\dots$ be i.i.d.\ with mean $\mu$ and $\E e^{\lambda(X_1-\mu)}\le e^{\lambda^2\sigp^2/2}$ for all $\lambda\in\R$. Let $N\ge1$, $A>0$,
$\rho_n:=\sigp\sqrt{(2/n)(A+2\log(N/n))}$ for $n\in[N]$, $\hat\mu(n)$ the mean of $X_1,\dots,X_n$, and $D':=\max_{1\le n\le N}(|\hat\mu(n)-\mu|-\rho_n)_+$. Then,
for every $x\ge0$,
\[
\P(D'>x)\ \le\ 2e^{-A/2}\sum_{\ell\ge0}2^{-\ell}\exp\Big(-\frac{N2^{-\ell}x^2}{8\sigp^2}\Big).
\]
In particular $\P(D'>0)\le4e^{-A/2}$ and $\E[D']\le\uCD\,\sigp\,e^{-A/2}/\sqrt N$, where $\uCD:=2\sqrt{2\pi}/(1-2^{-1/2})<17.2$.
\end{lemma}

\noindent Applied to arm $k$ (with $N=\uN$, so that $\rho_n=\urho_n$), it bounds $\uD_k$; hence $\E[\uD]\le K\uCD\sigp e^{-A/2}/\sqrt\uN$.

\subsubsection*{Which arms survive the screening, and how good are they?}

\begin{lemma}[survivors]\label{lem:u-survive}
\uses{def:sc}
\textit{Role.} Turns the deviation $\uD$ into deterministic bounds on the eliminations: used in Lemma~\ref{lem:u-cost} (screening term), and in Steps~3--5 of
Section~\ref{sec:app-u-proof-main} (loss of the best arm, box of the virtual instance, core off the good event).

\noindent Fix any realization of the reward table. For $n\in[\uN]$ let $i_n$ be the smallest index in $\arg\max_{i\in\uS_{n-1}}\umuh_i(n)$. Then:
\begin{enumerate}
\item[(i)] $i_n\in\uS_n$, and every $k\in\uS_{n-1}\setminus\uS_n$ satisfies $\mu_{i_n}>\mu_k-\uD_k-\uD_{i_n}$;
\item[(ii)] for every $n\in\{0,\dots,\uN\}$, $\max_{i\in\uS_n}\mu_i\ge\umus-2\uD$; in particular $\umus-\umusS\le2\uD$;
\item[(iii)] for every $n\in[\uN]$ and every $k\in\uS_n$, $\Delta_k\le4\urho_n+4\uD$;
\item[(iv)] for all $i,j\in\uS$, $\mu_i-\mu_j\le4\urho_\uN+2\uD$; and $|\umuref-\mu_{\uio}|\le\urho_\uN+\uD$.
\end{enumerate}
\end{lemma}

\subsubsection*{What does the screening cost, whatever the gaps?}

\begin{lemma}[cost of the screening]\label{lem:u-cost}
\uses{def:class, def:sc}
\textit{Role.} Bounds the first term of~\eqref{eq:u-decomp} uniformly in the gaps; this is what makes the intermediate gaps harmless.

\noindent Assume $A\ge8$ and fix any realization of the reward table. Let $\upulls_k$ be the number of pulls of arm $k$ in Stage~1. Then, for every $k\in[K]$,
\[
\Delta_k\,\upulls_k\ \le\ (b-a)+16\,\sigp\sqrt{A\uN}+8\uN\uD .
\]
\end{lemma}

\subsubsection*{Are the virtual variances close to $\bar\sigma^2$?}

\begin{lemma}[variance estimates]\label{lem:u-var}\dt{u:pvar}
\uses{def:class}
\textit{Role.} Bounds the probability that the garbling of Stage~2 misses the target variance $\sbar^2$ by more than $\eta\sbar^2$; used in Steps~4--5 of
Section~\ref{sec:app-u-proof-main}.

\noindent Let $N\ge1$, $\nu\in\Cab$ with variance $\sigma^2$, $X_1,\dots,X_N$ i.i.d.\ of law $\nu$, $\bar X$ their mean, $\hat\sigma^2:=\frac1N\sum_i(X_i-\bar X)^2$, and
$\eta\in(0,1]$. Then
\[
\P\big(|\hat\sigma^2-\sigma^2|>\eta\sbar^2\big)\ \le\ \upv(N,\eta):=\frac{64\,\sigp^4}{N\eta^2\sbar^4}+\frac2{N\eta}.
\]
\end{lemma}

\subsection{Why does the core behave as in the limit game?}\label{sec:app-u-transfer}

\begin{theorem}[transfer for continuous batched policies]\label{thm:u-transfer}
\uses{def:u-batched, def:regret, def:u-disc}
\textit{Role.} Turns the risk of a continuous batched policy into the regret of its discretization, uniformly over a compact set of local parameters and over the
reward laws with bounded third moments, whatever the reference $m_0$ and the scale $\sigma$. This uniformity lets the core face a random virtual instance
with a random reference (Step~4 of Section~\ref{sec:app-u-proof-main}); also used in Proposition~\ref{prop:u-local}(ii).

\noindent Let $\pi$ be a continuous batched policy (Definition~\ref{def:u-batched}), $\Theta\subset\R^K\times[0,\infty)^K$ compact and $C_3\ge0$. For a
$K$-tuple $\nu$ of laws with finite variances, $m_0\in\R$, $\sigma>0$ and $T\ge1$, let $h_k:=(\mu_k-m_0)\sqrt T/\sigma$ and $v_k:=\mathrm{Var}(\nu_k)/\sigma^2$.
Then:
\begin{enumerate}
\item[(i)] $(h,v)\mapsto\urk(\pi,h;v)$ is continuous on $\R^K\times[0,\infty)^K$;
\item[(ii)] let $\mathcal I_T$ be the set of the triples $(\nu,m_0,\sigma)$ with $(h,v)\in\Theta$ and $\E_{\nu_k}|X-\mu_k|^3\le C_3\sigma^3$ for every $k$; then
\[
\sup_{(\nu,m_0,\sigma)\in\mathcal I_T}\ \bigg|\frac{\RT\big(\uDisc{\pi}{T}{m_0,\sigma},\nu\big)}{\sigma\sqrt T}-\urk(\pi,h;v)\bigg|\ \longrightarrow\ 0\qquad(T\to\infty).
\]
\end{enumerate}
\end{theorem}

\subsection{Proof of Lemma~\ref{lem:u-dev}}
\begin{proof}
\proves{lem:u-dev}
\uses{def:class}
\textit{Setting.} The notation of the lemma; $S_n:=\sum_{i\le n}X_i$ and $x\ge0$. \textit{Goal.} The tail bound, then its two consequences.

\textit{Idea.} Peeling: on the block of indices $n\approx N2^{-\ell}$ the radius contains an extra $2\ell\log2$, which pays for the union bound over the blocks
with a factor $2^{-\ell}$.

Step~1: \emph{we reduce $\{D'>x\}$ to maximal deviations on blocks.} For $\ell\ge0$ let
\[
N_\ell:=N2^{-\ell},\qquad I_\ell:=\{n\in\Nat:N_\ell/2<n\le N_\ell\},\qquad A_\ell:=A+2\ell\log2 .
\]
These blocks cover $[N]$: every $n\in[N]$ lies in $I_\ell$ for $\ell:=\lfloor\log_2(N/n)\rfloor$. Let $n\in I_\ell$. Then:
\begin{itemize}
\item $\log(N/n)\ge\ell\log2$ (since $n\le N_\ell=N2^{-\ell}$);
\item $n>N_\ell/2$ (definition of $I_\ell$).
\end{itemize}
Hence
\begin{align*}
n\rho_n&=\sigp\sqrt{2n\big(A+2\log(N/n)\big)}\ \ge\ \sigp\sqrt{N_\ell A_\ell}, &&\text{(the two facts above)}\\
nx&\ge N_\ell x/2 . &&\text{(since $n>N_\ell/2$)}
\end{align*}
By definition of $D'$, and since $n\rho_n+nx\ge y_\ell$ for $n\in I_\ell$, where $y_\ell:=\sigp\sqrt{N_\ell A_\ell}+N_\ell x/2$,
\[
\{D'>x\}=\bigcup_{n\le N}\big\{|S_n-n\mu|>n\rho_n+nx\big\}\ \subset\ \bigcup_{\ell\ge0}\Big\{\max_{n\le N_\ell}|S_n-n\mu|>y_\ell\Big\}.
\]

Step~2: \emph{we bound each block by Doob's inequality.} For $\lambda>0$, $(e^{\lambda(S_n-n\mu)})_n$ is a nonnegative submartingale. Let $y>0$ and
$m:=\lfloor N_\ell\rfloor$ (the event below is empty if $m=0$). Then
\begin{align*}
\P\Big(\max_{n\le N_\ell}(S_n-n\mu)\ge y\Big)&\le e^{-\lambda y}\,\E e^{\lambda(S_m-m\mu)} &&\text{(Doob's maximal inequality)}\\
&\le e^{-\lambda y+\lambda^2\sigp^2N_\ell/2} &&\text{(sub-Gaussian bound, $m\le N_\ell$)}.
\end{align*}
We take $\lambda:=y/(\sigp^2N_\ell)$, and use the same bound for $-(S_n-n\mu)$:
\begin{align*}
\P\Big(\max_{n\le N_\ell}|S_n-n\mu|>y_\ell\Big)&\le2\exp\Big(-\frac{y_\ell^2}{2\sigp^2N_\ell}\Big)\\
&\le2\exp\Big(-\frac{A_\ell}2-\frac{N_\ell x^2}{8\sigp^2}\Big) &&\text{(since $y_\ell^2\ge\sigp^2N_\ell A_\ell+N_\ell^2x^2/4$)}\\
&=2e^{-A/2}\,2^{-\ell}\exp\Big(-\frac{N2^{-\ell}x^2}{8\sigp^2}\Big) &&\text{(since $e^{-\ell\log2}=2^{-\ell}$)}.
\end{align*}
A union bound over $\ell$, with Step~1, gives the tail bound. At $x=0$ it reads
\[
\P(D'>0)\ \le\ 2e^{-A/2}\sum_{\ell\ge0}2^{-\ell}\ =\ 4e^{-A/2}.
\]

Step~3: \emph{we integrate the tail.} Using $\int_0^\infty e^{-cx^2}dx=\frac12\sqrt{\pi/c}$ with $c=N2^{-\ell}/(8\sigp^2)$, and Tonelli's theorem,
\begin{align*}
\E[D']=\int_0^\infty\P(D'>x)\,dx&\le2e^{-A/2}\sum_{\ell\ge0}2^{-\ell}\,\sigp\sqrt{\frac{2\pi\,2^{\ell}}N} &&\text{(tail bound)}\\
&=\frac{2\sqrt{2\pi}\,\sigp e^{-A/2}}{\sqrt N}\sum_{\ell\ge0}2^{-\ell/2}=\frac{\uCD\,\sigp\,e^{-A/2}}{\sqrt N}. &&\text{(geometric series)}\qedhere
\end{align*}
\end{proof}

\subsection{Proof of Lemma~\ref{lem:u-survive}}
\begin{proof}
\proves{lem:u-survive}
\uses{def:sc}
\textit{Setting.} A realization of the table; \eqref{eq:u-basic} holds for all $k\in[K]$ and $n\in[\uN]$. \textit{Goal.} (i)--(iv).

\textit{Idea.} An arm is eliminated only by an arm that is at most $\uD_k+\uD_i$ worse; following the successive eliminators of an optimal arm loses at most $2\uD$.

Step~1: \emph{(i).} The arm $i_n$ survives cycle $n$, because $\umuh_{i_n}(n)+\urho_n\ge\max_{i\in\uS_{n-1}}\umuh_i(n)-\urho_n$. If $k\in\uS_{n-1}\setminus\uS_n$,
then $\umuh_k(n)+\urho_n<\umuh_{i_n}(n)-\urho_n$, and
\begin{align*}
\mu_{i_n}&\ge\umuh_{i_n}(n)-\urho_n-\uD_{i_n} &&\text{(by \eqref{eq:u-basic})}\\
&>\umuh_k(n)+\urho_n-\uD_{i_n} &&\text{(elimination of $k$)}\\
&\ge\mu_k-\uD_k-\uD_{i_n} &&\text{(by \eqref{eq:u-basic})}.
\end{align*}

Step~2: \emph{(ii).} Let $b_0$ be an optimal arm and, for $n\in[\uN]$, $b_n:=b_{n-1}$ if $b_{n-1}\in\uS_n$ and $b_n:=i_n$ otherwise. By induction $b_n\in\uS_n$, and
by Step~1, $\mu_{b_n}>\mu_{b_{n-1}}-\uD_{b_{n-1}}-\uD_{b_n}$ whenever $b_n\ne b_{n-1}$. Let $c_0,\dots,c_J$ be the successive distinct values of
$(b_0,\dots,b_n)$. They are distinct arms: when the sequence leaves $c_j$ at cycle $n'$, the arm $c_j$ is eliminated, so $c_j\notin\uS_{n''}$ for $n''\ge n'$,
while every later value lies in such an $\uS_{n''}$. Each arm appears at most twice in the following sum, so
\[
\mu_{b_n}\ \ge\ \umus-\sum_{j<J}\big(\uD_{c_j}+\uD_{c_{j+1}}\big)\ \ge\ \umus-2\uD ,
\]
and $\max_{i\in\uS_n}\mu_i\ge\mu_{b_n}$. For $n=\uN$ this is $\umus-\umusS\le2\uD$.

Step~3: \emph{(iii).} Let $n\in[\uN]$ and $k\in\uS_n$. The survival rule, with $b_{n-1}\in\uS_{n-1}$ (Step~2), gives $\umuh_k(n)+\urho_n\ge\umuh_{b_{n-1}}(n)-\urho_n$, so
\begin{align*}
\mu_k&\ge\umuh_k(n)-\urho_n-\uD_k &&\text{(by \eqref{eq:u-basic})}\\
&\ge\umuh_{b_{n-1}}(n)-3\urho_n-\uD_k &&\text{(survival of $k$)}\\
&\ge\mu_{b_{n-1}}-4\urho_n-\uD_k-\uD_{b_{n-1}} &&\text{(by \eqref{eq:u-basic})}\\
&\ge\umus-4\urho_n-4\uD &&\text{(Step~2: $\mu_{b_{n-1}}\ge\umus-2\uD$)}.
\end{align*}

Step~4: \emph{(iv).} Let $i,j\in\uS=\uS_\uN\subset\uS_{\uN-1}$. The survival of $j$ at cycle $\uN$ gives $\umuh_j(\uN)+\urho_\uN\ge\umuh_i(\uN)-\urho_\uN$, so
\begin{align*}
\mu_i&\le\umuh_i(\uN)+\urho_\uN+\uD_i &&\text{(by \eqref{eq:u-basic})}\\
&\le\umuh_j(\uN)+3\urho_\uN+\uD_i &&\text{(survival of $j$)}\\
&\le\mu_j+4\urho_\uN+\uD_i+\uD_j\ \le\ \mu_j+4\urho_\uN+2\uD &&\text{(by \eqref{eq:u-basic})}.
\end{align*}
The second claim is \eqref{eq:u-basic} for $k=\uio$ and $n=\uN$, since $\umuref=\umuh_{\uio}(\uN)$.
\end{proof}

\subsection{Proof of Lemma~\ref{lem:u-cost}}
\begin{proof}
\proves{lem:u-cost}
\uses{lem:u-survive, def:sc, def:class}
\textit{Setting.} $A\ge8$; a realization of the table. For $y>0$ let $\bar n(y):=\max\{n\in[\uN]:\urho_n\ge y\}$, with $\bar n(y):=0$ if the set is empty.
\textit{Goal.} $\Delta_k\upulls_k\le(b-a)+16\sigp\sqrt{A\uN}+8\uN\uD$.

\textit{Idea.} An arm with gap $\Delta>8\uD$ survives a cycle only while $\urho_n>\Delta/8$; this lasts about $\sigp^2(A+\log)/\Delta^2$ cycles, so it costs about
$\sigp^2(A+\log)/\Delta$, and the minimum of this cost and of the full cost $\Delta\uN$ is at most $16\sigp\sqrt{A\uN}$ for every $\Delta$.

Step~1: \emph{we bound the pulls of an arm by those of the cycles it survives.} Arm $k$ is pulled in cycle $n$ if and only if $k\in\uS_{n-1}$, and $\uS_0=[K]$, so
\[
\upulls_k=1+\#\{n\in[\uN-1]:k\in\uS_n\}.
\]
\begin{itemize}
\item \emph{Case $\Delta_k>8\uD$.} Let $n\in[\uN-1]$ with $k\in\uS_n$. Lemma~\ref{lem:u-survive}(iii), which bounds the gap of a survivor of cycle $n$, gives
\[
4\urho_n\ \ge\ \Delta_k-4\uD\ >\ \Delta_k/2 ,
\]
so $\urho_n>\Delta_k/8$, and $n\le\bar n(\Delta_k/8)$ because $\urho$ is decreasing. Hence $\upulls_k\le1+\min(\uN,\bar n(\Delta_k/8))$.
\item \emph{Case $\Delta_k\le8\uD$.} Then $\Delta_k\upulls_k\le8\uN\uD$ (since $\upulls_k\le\uN$).
\end{itemize}
In both cases
\begin{equation}\label{eq:u-cost1}
\Delta_k\upulls_k\ \le\ \Delta_k+\Delta_k\min\big(\uN,\bar n(\Delta_k/8)\big)+8\uN\uD .
\end{equation}

Step~2: \emph{we bound $\bar n$.} Let $y>0$ and $c:=2\sigp^2/y^2$, so that $\urho_n\ge y$ if and only if $n\le c\,(A+2\log(\uN/n))$. Let $n:=\bar n(y)\ge1$.
\begin{itemize}
\item If $n\le c$, then $n\le c\,(A+2\log_+(\uN/c))$ (since $A\ge1$).
\item If $n>c$, then $n\le c\,(A+2\log(\uN/n))\le c\,(A+2\log_+(\uN/c))$ (since $\log(\uN/n)<\log(\uN/c)$).
\end{itemize}
Hence, in all cases,
\[
\bar n(y)\ \le\ c\,\big(A+2\log_+(\uN/c)\big).
\]

Step~3: \emph{we bound the cost uniformly in the gap.} Let $\Delta>0$, $y:=\Delta/8$, so that $c=128\sigp^2/\Delta^2$, and let $x:=\Delta\sqrt\uN/(\sqrt{128}\,\sigp)$; then
$\uN/c=x^2$. Therefore
\begin{align*}
\Delta\min\big(\uN,\bar n(\Delta/8)\big)&\le\min\Big(\Delta\uN,\ \frac{128\sigp^2}\Delta\Big(A+2\log_+\frac{\uN\Delta^2}{128\sigp^2}\Big)\Big) &&\text{(Step~2)}\\
&=\sqrt{128}\,\sigp\sqrt\uN\,\min\Big(x,\ \frac{A+4\log_+x}x\Big) &&\text{(definition of $x$)}\\
&\le\sqrt{128}\,\sigp\sqrt\uN\,\sqrt{2A}=16\,\sigp\sqrt{A\uN} &&\text{(fact below)}.
\end{align*}
The fact: for $A\ge8$ and $x>0$, $\min(x,(A+4\log_+x)/x)\le\sqrt{2A}$. This is clear if $x\le\sqrt{2A}$. If $x>\sqrt{2A}$, then $x>4$ and we set
$\varphi(x):=\sqrt{2A}\,x-A-4\log x$; we have:
\begin{itemize}
\item $\varphi(\sqrt{2A})=A-2\log(2A)\ge0$ (the map $A\mapsto A-2\log(2A)$ is increasing on $[2,\infty)$ and equals $8-2\log16>0$ at $A=8$);
\item $\varphi'(x)=\sqrt{2A}-4/x>0$ for $x\ge\sqrt{2A}$ (since $\sqrt{2A}\ge4$, so $4/x\le1$).
\end{itemize}
So $\varphi\ge0$ on $[\sqrt{2A},\infty)$, that is $(A+4\log x)/x\le\sqrt{2A}$.

Step~4: \emph{conclusion.} Insert Step~3 in~\eqref{eq:u-cost1} and use $\Delta_k\le b-a$, the means lying in $[a,b]$.
\end{proof}

\subsection{Proof of Lemma~\ref{lem:u-var}}
\begin{proof}
\proves{lem:u-var}
\uses{def:class}
\textit{Setting.} $\mu$ is the mean of $\nu$. \textit{Goal.} The tail bound.

Since $\hat\sigma^2-\sigma^2=\big(\frac1N\sum_i(X_i-\mu)^2-\sigma^2\big)-(\bar X-\mu)^2$, the triangle inequality gives
$\{|\hat\sigma^2-\sigma^2|>\eta\sbar^2\}\subset E_1\cup E_2$, where
\[
E_1:=\Big\{\Big|\frac1N\sum_i(X_i-\mu)^2-\sigma^2\Big|>\frac{\eta\sbar^2}2\Big\},\qquad E_2:=\Big\{(\bar X-\mu)^2>\frac{\eta\sbar^2}2\Big\}.
\]
\begin{itemize}
\item $\P(E_1)\le\dfrac{64\,\sigp^4}{N\eta^2\sbar^4}$ (Chebyshev's inequality, with $\mathrm{Var}((X-\mu)^2)\le\E(X-\mu)^4\le16\sigp^4$ by~\eqref{eq:u-moments}).
\item $\P(E_2)\le\dfrac{2\sigma^2}{N\eta\sbar^2}\le\dfrac2{N\eta}$ (Markov's inequality, with $\E(\bar X-\mu)^2=\sigma^2/N$ and $\sigma^2\le\sbar^2$).
\end{itemize}
A union bound concludes.
\end{proof}

\subsection{Proof of Theorem~\ref{thm:u-transfer}}
\begin{proof}
\proves{thm:u-transfer}
\uses{def:u-disc, def:u-batched, def:regret}
\textit{Setting.} $\pi=(m,(t_j),(F_j))$ is a continuous batched policy, $\ell_j:=t_j-t_{j-1}$; $\Theta$ is compact and $C_3\ge0$. \textit{Goal.} (ii), then (i).

\textit{Idea.} Argue along sequences: compactness fixes the limit $(h,v)$; the martingale functional central limit theorem and Skorokhod's representation make
the normalized partial sums converge uniformly almost surely; since the rules are continuous in the observed path, the allocations and the observed paths
converge batch after batch.

\textit{Plan.} Step~1 reduces (ii) to sequences. Step~2 builds a space where the partial sums converge uniformly a.s. Step~3 shows the convergence of the
allocations and of the observed paths, by induction on the batches. Step~4 concludes for the regret. Step~5 proves (i) by the same coupling.

Step~1: \emph{reduction to sequences.} Suppose that (ii) fails. Then there are $\gamma>0$, $T_n\to\infty$ and $(\nu^n,m^n_0,\sigma^n)\in\mathcal I_{T_n}$, with local
parameters $(h^n,v^n)\in\Theta$, such that
\[
\bigg|\frac{\sy{s:regret}{R}_{T_n}\big(\uDisc{\pi}{T_n}{m^n_0,\sigma^n},\nu^n\big)}{\sigma^n\sqrt{T_n}}-\urk(\pi,h^n;v^n)\bigg|\ \ge\ \gamma\qquad\text{for every }n.
\]
Since $\Theta$ is compact, we may assume $(h^n,v^n)\to(h,v)\in\Theta$. Steps~2--4 show that the first term tends to $\urk(\pi,h;v)$, and Step~5 shows that
$\urk(\pi,h^n;v^n)\to\urk(\pi,h;v)$: a contradiction.

Step~2: \emph{we build a space on which the partial sums converge uniformly a.s.} Let $(X^n_{k,i})_{k\in[K],i\ge1}$ be a reward table of $\nu^n$,
$\zeta^n_{k,i}:=(X^n_{k,i}-\mu^n_k)/(\sigma^n\sqrt{T_n})$, and
\[
\xi^n(u):=\Big(\sum_{i\le\lfloor uT_n\rfloor}\zeta^n_{k,i}\Big)_{k\in[K]},\qquad u\ge0 .
\]
The process $\xi^n$ is a martingale for the filtration $\mathcal F^n_u:=\sigma(X^n_{k,i}:k\in[K],\,i\le\lfloor uT_n\rfloor)$, with independent centered
increments; since $\mathrm{Var}(\zeta^n_{k,i})=v^n_k/T_n$ and the coordinates are independent, $\xi^n(\xi^n)^{\top}-A_n$ is a martingale, where
$A_n(u):=(\lfloor uT_n\rfloor/T_n)\,\mathrm{diag}(v^n)$. We check the three conditions of the martingale FCLT \cite[Chap.~7, Thm~1.4(b)]{EK86} on every interval
$[0,u_0]$:
\begin{itemize}
\item the jumps of $A_n$ tend to $0$ (they equal $\mathrm{diag}(v^n)/T_n$, and $v^n$ is bounded on the compact $\Theta$);
\item $A_n(u)\to u\,\mathrm{diag}(v)$ for every $u$ (since $\lfloor uT_n\rfloor/T_n\to u$ and $v^n\to v$);
\item the jumps of $\xi^n$ satisfy $\E[\max_{i\le u_0T_n}(\zeta^n_{k,i})^2]\to0$ for every $k$ (this is the jump condition; it follows from the Lindeberg condition:
for every $c>0$,
\begin{align*}
\E\Big[\max_{i\le u_0T_n}(\zeta^n_{k,i})^2\Big]&\le c^2+\sum_{i\le u_0T_n}\E\big[(\zeta^n_{k,i})^2\1(|\zeta^n_{k,i}|>c)\big] &&\text{(split at $|\zeta^n_{k,i}|=c$)}\\
&\le c^2+\frac1c\sum_{i\le u_0T_n}\E|\zeta^n_{k,i}|^3 &&\text{(on the event, $\zeta^2\le|\zeta|^3/c$)}\\
&\le c^2+\frac{u_0C_3}{c\sqrt{T_n}} &&\text{($\E|X-\mu^n_k|^3\le C_3(\sigma^n)^3$ on $\mathcal I_{T_n}$)},
\end{align*}
so the $\limsup$ is at most $c^2$ for every $c>0$).
\end{itemize}
Hence $\xi^n\Rightarrow\xi:=(\sqrt{v_k}\,W_k)_{k\in[K]}$ in $D_{\R^K}[0,\infty)$ with the Skorokhod topology, where $W_1,\dots,W_K$ are independent standard Brownian
motions: the limit is the continuous Gaussian process with independent increments and covariance $u\,\mathrm{diag}(v)$.

The space $D_{\R^K}[0,\infty)$ is Polish \cite[Chap.~3, Thm~5.6]{EK86}. We apply Skorokhod's representation theorem \cite[Thm~6.7]{Bil99} to the laws of the
$\xi^n$ \emph{alone}: on some probability space $(\Omega_1,\mathcal F_1,\P_1)$ there are random elements $\bar\xi^n$ with the law of $\xi^n$ and $\bar\xi$ with
the law of $\xi$ such that $\bar\xi^n\to\bar\xi$ a.s. Since $\bar\xi$ is continuous, this convergence is uniform on $[0,1]$ \cite[Chap.~3, Sec.~10]{EK86}. We then
take the product $\Omega:=\Omega_1\times[0,1]\times\Omega_2$, where $[0,1]$ carries the Lebesgue measure and its coordinate $U$, and $\Omega_2$ carries $K$
independent standard Brownian motions $B_k$. We set $W_k:=\bar\xi_k/\sqrt{v_k}$ if $v_k>0$ and $W_k:=B_k$ if $v_k=0$; then $W_1,\dots,W_K,U$ are independent
and each $W_k$ is a standard Brownian motion. On $\Omega$,
\[
\bar X^n_{k,i}:=\mu^n_k+\sigma^n\sqrt{T_n}\,\big(\bar\xi^n_k(i/T_n)-\bar\xi^n_k((i-1)/T_n)\big)
\]
is a Borel function of $\bar\xi^n$ (evaluations at fixed times are Borel on $D$), so $(\bar X^n_{k,i})_{k,i}$ is a reward table of $\nu^n$, independent of $U$.
The interaction of $\uDisc{\pi}{T_n}{m^n_0,\sigma^n}$ is a Borel function of its table and of $U$, so its law, hence its regret, is the same when it is run on
$(\bar X^n,U)$. We drop the bars from now on.

Finally, let $\tilde\xi^n_k$ be the linear interpolation of $\xi^n_k$ between the points $i/T_n$. On $[(i-1)/T_n,i/T_n]$ we have
$|\tilde\xi^n_k-\xi^n_k|\le|\zeta^n_{k,i}|$, and, since $\zeta^n_{k,i}=\xi^n_k(i/T_n)-\xi^n_k((i-1)/T_n)$,
\[
\max_{i\le T_n}|\zeta^n_{k,i}|\ \le\ 2\sup_{u\le1}\big|\xi^n_k(u)-\sqrt{v_k}W_k(u)\big|+\omega_k(1/T_n),
\]
where $\omega_k$ is the modulus of continuity of $\sqrt{v_k}W_k$ on $[0,1]$; the right-hand side tends to $0$ a.s.\ (uniform convergence above, and continuity
of $W_k$). Since $\sup_{u\le1}|\tilde\xi^n_k-\sqrt{v_k}W_k|\le\max_{i\le T_n}|\zeta^n_{k,i}|+\sup_{u\le1}|\xi^n_k-\sqrt{v_k}W_k|$, we get, almost surely,
\begin{equation}\label{eq:u-unif}
\sup_{u\le1}\big|\tilde\xi^n_k(u)-\sqrt{v_k}\,W_k(u)\big|\ \longrightarrow\ 0\qquad\text{for every }k\in[K].
\end{equation}

Step~3: \emph{we show that the allocations and the observed paths converge, batch after batch.} Let $N^{n,(j)}$, $\hat\tau^n$, $\hat X^n$ and $\hat\alpha^{n,(j)}$ be the
pull counts, clocks, observed path and allocations of $\uDisc{\pi}{T_n}{m^n_0,\sigma^n}$ (Definition~\ref{def:u-disc}), and let $(\alpha^{(j)},\tau,X)$ be the
quantities of Definition~\ref{def:u-batched} built from $(h,v)$, $W$ and $U$. Since $x-m^n_0=(x-\mu^n_k)+(\mu^n_k-m^n_0)$ and
$(\mu^n_k-m^n_0)/(\sigma^n\sqrt{T_n})=h^n_k/T_n$, the interpolated partial sums of arm $k$ are $\hat z^n_k(u)=\tilde\xi^n_k(u)+h^n_ku$, so that
\begin{equation}\label{eq:u-Xhat}
\hat X^n_k(t)=\tilde\xi^n_k\big(\hat\tau^n_k(t)\big)+h^n_k\,\hat\tau^n_k(t),\qquad X_k(t)=\sqrt{v_k}\,W_k\big(\tau_k(t)\big)+h_k\,\tau_k(t).
\end{equation}
We show by induction on $j\in\{0,\dots,m\}$ that, almost surely,
\[
N^{n,(j)}/T_n\to\tau(t_j)\qquad\text{and}\qquad\sup_{t\le t_j}\big|\hat X^n(t)-X(t)\big|\to0,
\]
where $\hat X^n$ on $[0,t_j]$ is formed with the data of batches $1,\dots,j$. For $j=0$ both sides are $0$. Assume the claim for $j-1$.
\begin{itemize}
\item \emph{Allocations}: $\hat\alpha^{n,(j)}=F_j(U,\hat X^n|_{[0,t_{j-1}]})\to F_j(U,X|_{[0,t_{j-1}]})=\alpha^{(j)}$ (continuity of $F_j(U,\cdot)$ for the uniform norm, and
the induction hypothesis).
\item \emph{Clocks}: $N^{n,(j)}/T_n\to\tau(t_{j-1})+\ell_j\alpha^{(j)}=\tau(t_j)$ (in batch $j$, arm $k$ is pulled between $\hat\alpha^{n,(j)}_kR^n_j-1$ and
$\hat\alpha^{n,(j)}_kR^n_j+K$ times, and $R^n_j/T_n\to\ell_j$). Since $\hat\tau^n$ and $\tau$ are linear on each $[t_{i-1},t_i]$ and converge at the nodes,
$\sup_{t\le t_j}|\hat\tau^n(t)-\tau(t)|\to0$.
\item \emph{Paths}: by~\eqref{eq:u-Xhat}, for $t\le t_j$,
\[
\big|\hat X^n_k(t)-X_k(t)\big|\le\sup_{u\le1}\big|\tilde\xi^n_k(u)-\sqrt{v_k}W_k(u)\big|+\sqrt{v_k}\,\big|W_k(\hat\tau^n_k(t))-W_k(\tau_k(t))\big|+\big|h^n_k\hat\tau^n_k(t)-h_k\tau_k(t)\big|,
\]
and the three terms tend to $0$ uniformly in $t\le t_j$ (by~\eqref{eq:u-unif}; by the uniform continuity of $W_k$ on $[0,1]$ and the uniform convergence of
$\hat\tau^n_k$; since $h^n\to h$ and $\hat\tau^n_k,\tau_k\in[0,1]$).
\end{itemize}

Step~4: \emph{we conclude for the regret.} The last batch ends at round $T_n$, so $N_k(T_n)=N^{n,(m)}_k$, and the gaps of $\nu^n$ are
$\Delta^n_k=\sigma^n(\max_ih^n_i-h^n_k)/\sqrt{T_n}$. Hence
\[
\frac{\sy{s:regret}{R}_{T_n}\big(\uDisc{\pi}{T_n}{m^n_0,\sigma^n},\nu^n\big)}{\sigma^n\sqrt{T_n}}=\E\Big[\sum_k\big(\max_ih^n_i-h^n_k\big)\frac{N^{n,(m)}_k}{T_n}\Big]
\ \longrightarrow\ \E\Big[\sum_k\big(\max_ih_i-h_k\big)\tau_k(1)\Big]=\urk(\pi,h;v),
\]
by dominated convergence: the integrand converges a.s.\ by Step~3, and it is bounded by $\sup\{\max_ih'_i-\min_ih'_i:(h',v')\in\Theta\}<\infty$, since
$\sum_kN^{n,(m)}_k=T_n$.

Step~5: \emph{continuity, (i).} Let $(h^n,v^n)\to(h,v)$ in $\R^K\times[0,\infty)^K$, and build the quantities of Definition~\ref{def:u-batched} for $(h^n,v^n)$ and
for $(h,v)$ from the same $W$ and $U$. Then $Y^n_k(s)=h^n_ks+\sqrt{v^n_k}W_k(s)\to Y_k(s)$ uniformly on $[0,1]$. The induction of Step~3, with
$\tau^n(t_j)=\tau^n(t_{j-1})+\ell_j\alpha^{n,(j)}$ in place of the pull counts and the bound
$|Y^n_k(\tau^n_k(t))-Y_k(\tau_k(t))|\le\sup_{s\le1}|Y^n_k(s)-Y_k(s)|+|Y_k(\tau^n_k(t))-Y_k(\tau_k(t))|$ in place of~\eqref{eq:u-Xhat}, gives
$\sup_{t\le1}|\tau^n(t)-\tau(t)|\to0$ a.s. Dominated convergence, with the bound $\sup_n(\max_ih^n_i-\min_ih^n_i)<\infty$, gives
$\urk(\pi,h^n;v^n)\to\urk(\pi,h;v)$.
\end{proof}


\subsection{Proof of Theorem~\ref{thm:u-upper}}\label{sec:app-u-proof-main}
\begin{proof}
\proves{thm:u-upper}
\uses{lem:u-dev, thm:main, thm:u-transfer, lem:u-cost, lem:u-survive, def:sc, lem:u-var, def:class, def:regret, def:risk, def:u-batched, def:batched, def:u-disc}
\textit{Setting.}\dt{u:param} $K\ge2$, $a<b$, $0<\sbar\le\sigp$; $\uCD$ is the constant of Lemma~\ref{lem:u-dev} and $\uCt:=4\big(8(\sigp/\sbar)^3+2\big)$. Fix
$\varepsilon\in(0,1]$.
\textit{Goal.} Choose $(\delta,A,\pi^\circ)$ and $T_1$ such that $\RT(\SC,\nu)\le\big(\sbar(\VK_K+3\varepsilon)+4\varepsilon\big)\sqrt T$ for all $\nu\in\Cab^K$ and
$T\ge T_1$; then deduce Theorem~\ref{thm:main}(ii).

\textit{Idea.} The three terms of~\eqref{eq:u-decomp} (Section~\ref{sec:app-u-skeleton}): the screening and the loss of the best arm are small once $\delta$ is small
and $A=6\log(1/\delta)$; the core is the discretization of a near-optimal policy for a fixed box, and the transfer theorem applies on the good event.

\textit{Plan.} Step~0 chooses the parameters. Step~1 proves the decomposition~\eqref{eq:u-decomp}. Steps~2--5 bound its terms: screening, loss of the best arm,
core on the good event, core off it. Step~6 adds them up, and Step~7 deduces Theorem~\ref{thm:main}(ii).

Step~0: \emph{we choose the parameters}, in six moves.
\begin{itemize}
\item \emph{Screening.} Choose $\delta\in(0,\frac12]$ and $A\ge8$ such that
\begin{align}
&16\,\sigp\sqrt{AK\delta}\ \le\ \varepsilon, \label{eq:u-P1}\\
&\uCD\,\sigp K^{3/2}e^{-A/2}\Big(8\sqrt\delta+\frac6{\sqrt\delta}+\frac{12\sqrt K(\sigp/\sbar)\sqrt A}{\delta}\Big)\ \le\ \varepsilon. \label{eq:u-P2}
\end{align}
This is possible with $A:=6\log(1/\delta)$ and $\delta$ small: then $e^{-A/2}=\delta^3$, $A\ge8$ as soon as $\delta\le e^{-4/3}$, and both left-hand sides tend to $0$
as $\delta\to0$.
\item \emph{Box.} Let
\[
\uM:=10\,\frac{\sigp}{\sbar}\sqrt{\frac{AK}\delta}+3,\qquad H_\uM:=[-\uM,\uM]^K .
\]
\item \emph{Core.} By (G2), choose a continuous batched policy $\pi^\circ$ with $\sup_{h\in H_\uM}\urk(\pi^\circ,h)\le\Vb(H_\uM)+\varepsilon$. By (G1),
$\Vb(H_\uM)\le\VK_K$, so
\begin{equation}\label{eq:u-core}
\sup_{h\in H_\uM}\urk(\pi^\circ,h)\ \le\ \VK_K+\varepsilon .
\end{equation}
\item \emph{Tolerance on the variances.} By Theorem~\ref{thm:u-transfer}(i), $\urk(\pi^\circ,\cdot\,;\cdot)$ is uniformly continuous on the compact
$H_\uM\times[\frac12,\frac32]^K$. Choose $\eta\in(0,\frac12]$ such that
\[
\urk(\pi^\circ,h;v)\le\urk(\pi^\circ,h)+\varepsilon\qquad\text{for all }h\in H_\uM\text{ and }v\in[1-\eta,1+\eta]^K .
\]
\item \emph{Transfer.} Let $\Theta:=H_\uM\times[1-\eta,1+\eta]^K$, and let $T_0$ be such that the supremum of Theorem~\ref{thm:u-transfer}(ii), with $\pi^\circ$,
$\Theta$ and $\uCt$, is at most $\varepsilon$ for every $T\ge T_0$.
\item \emph{Horizon.} Let $T_1\ge\max(2K/\delta,2T_0)$ be such that, for every $T\ge T_1$,
\[
K(b-a)\le\varepsilon\sqrt T\qquad\text{and}\qquad 8\sigp K\sqrt{AK/\delta}\;\upv\big(\lceil\delta T/(2K)\rceil,\eta\big)\le\varepsilon ;
\]
this is possible since $\upv(n,\eta)\to0$ as $n\to\infty$.
\end{itemize}

Fix $T\ge T_1$ and $\nu\in\Cab^K$. We record three consequences of $T\ge T_1$:
\begin{itemize}
\item $\uN\in[\delta T/(2K),\,\delta T/K]$, hence $\uN\ge\lceil\delta T/(2K)\rceil$ (since $\delta T/K\ge2$ and $\uN=\lfloor\delta T/K\rfloor$ is an integer);
\item $\uTp\ge(1-\delta)T\ge T/2\ge T_0$ (Section~\ref{sec:app-u-setting} and $\delta\le\frac12$);
\item Lemma~\ref{lem:u-dev}, applied to each arm, gives
\begin{equation}\label{eq:u-ED}
\E[\uD]\ \le\ \frac{K\uCD\sigp e^{-A/2}}{\sqrt\uN}\ \le\ \frac{\sqrt2\,K^{3/2}\uCD\sigp e^{-A/2}}{\sqrt{\delta T}}\qquad\text{(since $\uN\ge\delta T/(2K)$)}.
\end{equation}
\end{itemize}

Step~1: \emph{we prove the decomposition~\eqref{eq:u-decomp}.}\dt{u:decomp}
\begin{itemize}
\item \emph{What Stage~1 determines.} Let $\uF:=\sigma(X_{k,n}:k\in[K],\,n\le\uN)$. Stage~1 only reads entries with $n\le\uN$, so $\uS$, $\uio$, $\umuref$,
$(\usigh^2_k)_k$, $(\uw_k)_k$, $(\upulls_k)_k$, $\uTp$ and $\uD$ are $\uF$-measurable. Let $P:=(\uS,\uio,\umuref,(\uw_k)_{k\in\uS},\uTp)$; it takes values in a
finite union of Euclidean spaces, and it determines the map $\uc$.
\item \emph{What Stage~2 reads.} Let $G_1,\dots,G_T$ be i.i.d.\ $\mathcal N(0,1)$, the noise added at round $t$ being $\sqrt{\uw_{\uc(k)}}\,G_t$ when virtual arm $k$
is asked, and let $U^\circ$ be the randomization of the core. The arms of $\uS$ were pulled exactly $\uN$ times in Stage~1, so the core reads the entries
$X_{j,\uN+1},X_{j,\uN+2},\dots$ of each real arm $j\in\uS$, in this order. Hence the regret of Stage~2, $\sum_{t>T-\uTp}\Delta_{A_t}$, equals $\Psi(P,\mathcal Y)$
for a Borel map $\Psi\ge0$, where $\mathcal Y:=\big((X_{k,\uN+n})_{k\in[K],n\ge1},(G_t)_{t\le T},U^\circ\big)$ is independent of $\uF$.
\item \emph{Freezing.} Let $\psi(p):=\E[\Psi(p,\mathcal Y)]$, a Borel function by Tonelli's theorem. Then
\[
\E\big[\Psi(P,\mathcal Y)\,\big|\,\uF\big]=\psi(P).
\]
This holds for $\Psi(p,y)=\1_B(p)\1_C(y)$ by independence, and extends to every Borel $\Psi\ge0$ by the monotone class theorem.
\item \emph{The frozen core.} Fix a value $p$ of $P$. Each reward received by the core for virtual arm $k$ is $X_{\uc(k),\uN+n}+\sqrt{\uw_{\uc(k)}}\,G_t$, where
the index $n$ and the round $t$ are determined by the past: a table entry and a noise variable that have not been used before, and that are independent of the
variables used before. So, conditionally on the past of the core, this reward has law $\unu_k$, and this determines the law of the interaction
\cite[Sec.~4.6]{LS}: the core is $\uDisc{\pi^\circ}{\uTp}{\umuref,\sbar}$ facing the $K$-armed instance $\unu$, which may have repeated arms. Asking for
virtual arm $k$ costs $\Delta_{\uc(k)}=(\umus-\umusS)+(\umusS-\mu_{\uc(k)})$, and the second term is the gap of virtual arm $k$ in $\unu$, since
$\max_k\mu_{\uc(k)}=\umusS$. Summing over the $\uTp$ rounds of Stage~2,
\[
\psi(p)=\uTp\,(\umus-\umusS)+\sy{s:regret}{R}_{\uTp}\big(\uDisc{\pi^\circ}{\uTp}{\umuref,\sbar},\unu\big)\qquad\text{at }p .
\]
\end{itemize}
Taking expectations and adding the regret of Stage~1, $\sum_k\Delta_k\upulls_k$, gives~\eqref{eq:u-decomp}, where $\uRt$ is the Borel function of $P$ obtained by
freezing.

Step~2: \emph{we bound the screening term.} Summing Lemma~\ref{lem:u-cost}, which bounds the cost of each arm whatever its gap, over $k$:
\begin{align*}
\E\Big[\sum_k\Delta_k\upulls_k\Big]&\le K(b-a)+16K\sigp\sqrt{A\uN}+8K\uN\,\E[\uD] &&\text{(Lemma~\ref{lem:u-cost})}\\
&\le K(b-a)+16\sigp\sqrt{AK\delta}\,\sqrt T+8K^2\uCD\sigp e^{-A/2}\sqrt\uN &&\text{($\uN\le\delta T/K$; \eqref{eq:u-ED})}\\
&\le\varepsilon\sqrt T+16\sigp\sqrt{AK\delta}\,\sqrt T+8K^{3/2}\uCD\sigp e^{-A/2}\sqrt\delta\,\sqrt T &&\text{($T\ge T_1$; $\uN\le\delta T/K$)}.
\end{align*}

Step~3: \emph{we bound the loss of the best arm.} By Lemma~\ref{lem:u-survive}(ii), which says that the best survivor is at most $2\uD$ below $\umus$, and since
$\uTp\le T$,
\[
\E\big[\uTp\,(\umus-\umusS)\big]\ \le\ 2T\,\E[\uD]\ \le\ 2\sqrt2\,K^{3/2}\uCD\sigp e^{-A/2}\delta^{-1/2}\sqrt T\qquad\text{(by \eqref{eq:u-ED})}.
\]

Step~4: \emph{we bound the core on the good event.}\dt{u:good} Let $\usigh_k^2$ be the empirical variance of $X_{k,1},\dots,X_{k,\uN}$ for every $k$ (for $k\in\uS$
it is the estimate of Definition~\ref{def:sc}), and
\[
\uG:=\{\uD\le\sbar/\sqrt T\}\cap\{|\usigh_k^2-\sigma_k^2|\le\eta\sbar^2\ \ \forall k\in[K]\}\ \in\ \uF .
\]
On $\uG$ we check that $(\unu,\umuref,\sbar)\in\mathcal I_{\uTp}$ (Theorem~\ref{thm:u-transfer}(ii), with $\Theta$ and $\uCt$ of Step~0 and horizon $\uTp$). Let
$\tilde h_k:=(\mu_{\uc(k)}-\umuref)\sqrt{\uTp}/\sbar$ and $\tilde v_k:=(\sigma^2_{\uc(k)}+\uw_{\uc(k)})/\sbar^2$ be the local parameters of $\unu_k$.
\begin{itemize}
\item \emph{Local means}: $|\tilde h_k|\le\uM$. Indeed $\uc(k)$ and $\uio$ lie in $\uS$, and $\urho_\uN=\sigp\sqrt{2A/\uN}$, so
\begin{align*}
|\tilde h_k|&\le\big(|\mu_{\uc(k)}-\mu_{\uio}|+|\mu_{\uio}-\umuref|\big)\sqrt T/\sbar &&\text{(triangle inequality; $\uTp\le T$)}\\
&\le\big(5\urho_\uN+3\uD\big)\sqrt T/\sbar &&\text{(Lemma~\ref{lem:u-survive}(iv): spread of $\uS$, error of $\umuref$)}\\
&\le\Big(10\,\sigp\sqrt{\frac{AK}{\delta T}}+\frac{3\sbar}{\sqrt T}\Big)\frac{\sqrt T}{\sbar}=\uM &&\text{($\uN\ge\delta T/(2K)$; $\uD\le\sbar/\sqrt T$ on $\uG$)}.
\end{align*}
\item \emph{Variances}: $\tilde v_k\in[1-\eta,1+\eta]$. Let $j:=\uc(k)$; on $\uG$, $|\usigh_j^2-\sigma_j^2|\le\eta\sbar^2$.
\begin{itemize}
\item If $\usigh_j^2\le\sbar^2$: $\tilde v_k\sbar^2=\sbar^2+(\sigma_j^2-\usigh_j^2)\in\sbar^2[1-\eta,1+\eta]$ (since $\uw_j=\sbar^2-\usigh_j^2$).
\item If $\usigh_j^2>\sbar^2$: $\tilde v_k\sbar^2=\sigma_j^2\in[\usigh_j^2-\eta\sbar^2,\sbar^2]\subset\sbar^2[1-\eta,1]$ (since $\uw_j=0$, and $\sigma_j^2\le\sbar^2$ in $\Cab$).
\end{itemize}
\item \emph{Third moments}: for $\tilde X=X+\xi\sim\unu_k$,
\begin{align*}
\E|\tilde X-\mu_j|^3&\le4\big(\E|X-\mu_j|^3+\E|\xi|^3\big) &&\text{(since $(s+t)^3\le4(s^3+t^3)$ for $s,t\ge0$)}\\
&\le4\big(8\sigp^3+2\sbar^3\big)=\uCt\sbar^3 &&\text{(by \eqref{eq:u-moments}; $\E|\xi|^3=2\sqrt{2/\pi}\,\uw_j^{3/2}$, $\uw_j\le\sbar^2$)}.
\end{align*}
\end{itemize}
So $(\tilde h,\tilde v)\in\Theta$, and since $\uTp\ge T_0$, on $\uG$:
\begin{align*}
\uRt&\le\sbar\sqrt{\uTp}\,\big(\urk(\pi^\circ,\tilde h;\tilde v)+\varepsilon\big) &&\text{(Theorem~\ref{thm:u-transfer}(ii), choice of $T_0$)}\\
&\le\sbar\sqrt{\uTp}\,\big(\urk(\pi^\circ,\tilde h)+2\varepsilon\big) &&\text{(choice of $\eta$)}\\
&\le\sbar\sqrt T\,\big(\VK_K+3\varepsilon\big) &&\text{(by \eqref{eq:u-core}; $\uTp\le T$)}.
\end{align*}
Theorem~\ref{thm:u-transfer}(ii) is applied at the random horizon $\uTp$ and to the random triple $(\unu,\umuref,\sbar)$, realization by realization: this is
legitimate because its bound is uniform over $T\ge T_0$ and over $\mathcal I_T$. Hence $\E[\uRt\1_{\uG}]\le\sbar\sqrt T(\VK_K+3\varepsilon)$.

Step~5: \emph{we bound the core off the good event.} The regret of the core is at most $\uTp$ times the largest virtual gap, so by
Lemma~\ref{lem:u-survive}(iv), which bounds the spread of the survivors,
\[
\uRt\ \le\ \uTp\max_{i,j\in\uS}(\mu_i-\mu_j)\ \le\ T\,(4\urho_\uN+2\uD),\qquad\text{hence}\qquad
\E\big[\uRt\1_{\uG^c}\big]\le4T\urho_\uN\,\P(\uG^c)+2T\,\E[\uD].
\]
The probability of the bad event has two parts, the second one being the failure of the variance estimates:
\begin{align*}
\P(\uG^c)&\le\P\big(\uD>\sbar/\sqrt T\big)+\sum_{k\in[K]}\P\big(|\usigh_k^2-\sigma_k^2|>\eta\sbar^2\big)\\
&\le\frac{\sqrt T}{\sbar}\,\E[\uD]+K\,\upv(\uN,\eta) &&\text{(Markov; Lemma~\ref{lem:u-var} for $X_{k,1},\dots,X_{k,\uN}$)}\\
&\le\sqrt2\,K^{3/2}\uCD\frac{\sigp}{\sbar}\,e^{-A/2}\delta^{-1/2}+K\,\upv(\uN,\eta) &&\text{(by \eqref{eq:u-ED})}.
\end{align*}
With $4T\urho_\uN\le8\sigp\sqrt{AK/\delta}\,\sqrt T$ and~\eqref{eq:u-ED} for the last term, we get
\[
\frac{\E[\uRt\1_{\uG^c}]}{\sqrt T}\ \le\ 8\sqrt2\,\uCD K^2\frac{\sigp^2}{\sbar}\,\frac{\sqrt A\,e^{-A/2}}\delta+8\sigp K\sqrt{\frac{AK}\delta}\;\upv(\uN,\eta)
+2\sqrt2\,K^{3/2}\uCD\sigp e^{-A/2}\delta^{-1/2}.
\]

Step~6: \emph{we add the four bounds.} Insert Steps~2--5 in~\eqref{eq:u-decomp}.
\begin{itemize}
\item The term $16\sigp\sqrt{AK\delta}$ is at most $\varepsilon$ (by~\eqref{eq:u-P1}).
\item The terms in $e^{-A/2}$ add up to $\uCD\sigp K^{3/2}e^{-A/2}\big(8\sqrt\delta+4\sqrt2\,\delta^{-1/2}+8\sqrt2\sqrt K(\sigp/\sbar)\sqrt A\,\delta^{-1}\big)$, which is at
most $\varepsilon$ (by~\eqref{eq:u-P2}, since $4\sqrt2\le6$ and $8\sqrt2\le12$).
\item The term $K(b-a)/\sqrt T$ is at most $\varepsilon$, and so is the term in $\upv$ (choice of $T_1$, since $\upv(\cdot,\eta)$ is decreasing and
$\uN\ge\lceil\delta T/(2K)\rceil$).
\end{itemize}
Hence, for all $\nu\in\Cab^K$ and all $T\ge T_1$,
\[
\frac{\RT(\SC,\nu)}{\sqrt T}\ \le\ \sbar\big(\VK_K+3\varepsilon\big)+4\varepsilon=\sbar\VK_K+(3\sbar+4)\varepsilon .
\]

Step~7: \emph{we deduce Theorem~\ref{thm:main}(ii).} Let $\Cl$ satisfy (R1): $\Cl\subseteq\Cab$ for some $a<b$ and $\sigp\ge\sbar$. Given $\varepsilon>0$, apply
Steps~0--6 with $\varepsilon':=\min(1,\varepsilon/(3\sbar+4))$ in place of $\varepsilon$; they give $(\delta,A,\pi^\circ)$ and $T_1$ such that, for $T\ge T_1$,
\begin{align*}
\Rs_{K,\Cl}(T)&\le\sup_{\nu\in\Cl^K}\RT(\SC,\nu) &&\text{(definition of $\Rs$)}\\
&\le\sup_{\nu\in\Cab^K}\RT(\SC,\nu) &&\text{(since $\Cl\subseteq\Cab$)}\\
&\le(\sbar\VK_K+\varepsilon)\sqrt T &&\text{(Step~6 with $\varepsilon'$)}.
\end{align*}
This is the quantitative statement of Theorem~\ref{thm:main}(ii), and letting $\varepsilon\to0$ gives $\limsup_T\Rs_{K,\Cl}(T)/\sqrt T\le\sbar\VK_K$.
\end{proof}

\subsection{Proof of Proposition~\ref{prop:u-local}}
\begin{proof}
\proves{prop:u-local}
\uses{def:u-batched, thm:u-transfer, def:risk, def:regret, def:batched, def:u-disc}
\textit{Setting.} $H\subset\R^K$ is compact. \textit{Goal.} (i) $\uRG_H(T)\ge\sqrt T\,\VK_K(H)$ for every $T$; (ii) $\uRG_H(T)/\sqrt T\to\VK_K(H)$.

\textit{Idea.} A discrete Gaussian policy is a batched policy with $T$ batches, read from Brownian increments, which gives (i) for every $T$. Conversely, the
discretization of a near-optimal continuous batched policy is a $T$-round policy, and the transfer theorem gives (ii).

Step~1: \emph{we show (i): a discrete Gaussian policy is a batched policy.} Let $\pi_d$ be a $T$-round policy, given by Borel maps $\varphi_t$
(Section~\ref{sec:app-u-setting}). We define a batched policy $\pi^\flat$ with $T$ batches of length $1/T$: its decision times are $t_j:=j/T$, and its rules are
$F_j(u,y):=e_{a_j}$ (the $a_j$-th vertex of $\uSimp$), where, recursively,
\[
a_1:=\varphi_1(u),\qquad x_i:=\sqrt T\,\big(y_{a_i}(i/T)-y_{a_i}((i-1)/T)\big),\qquad a_{i+1}:=\varphi_{i+1}(u,x_1,\dots,x_i)\qquad(i<j).
\]
Each $F_j$ is Borel, as a composition of Borel maps and of evaluations of the path $y$. Now run $\pi^\flat$ as in Definition~\ref{def:u-batched}, with $(h,\mathbf 1)$.
\begin{itemize}
\item \emph{Clocks}: only arm $a_i$ runs during batch $i$, at speed $1$, so $T\tau_k(i/T)$ is the number of $i'\le i$ with $a_{i'}=k$.
\item \emph{Rewards}: with $n:=T\tau_{a_i}(i/T)$, the $i$-th reward is
\[
x_i=\sqrt T\,\big(Y_{a_i}(n/T)-Y_{a_i}((n-1)/T)\big)=X_{a_i,n},\qquad\text{where}\quad X_{k,n}:=\sqrt T\,\big(Y_k(n/T)-Y_k((n-1)/T)\big).
\]
\item \emph{Table}: the $X_{k,n}$ are independent, of law $\mathcal N(h_k/\sqrt T,1)$, and independent of $U$ (Brownian increments over intervals of length $1/T$, with
drift $h_k$).
\end{itemize}
So $\pi^\flat$ reads a reward table of the instance $(\mathcal N(h_k/\sqrt T,1))_k$ exactly as $\pi_d$ does, and $T\tau(1)$ has the law of $N(T)$ under $\pi_d$ and
this instance. Since the gaps of this instance are $(\max_ih_i-h_k)/\sqrt T$,
\[
\urk(\pi^\flat,h)=\E\Big[\sum_k\big(\max_ih_i-h_k\big)\frac{N_k(T)}T\Big]=\frac{\RT\big(\pi_d,(\mathcal N(h_k/\sqrt T,1))_k\big)}{\sqrt T}\qquad\text{for every }h\in\R^K .
\]
Taking the supremum over $h\in H$, then the infimum over $\pi_d$, gives $\uRG_H(T)/\sqrt T\ge\Vb(H)$ by (G3), and $\Vb(H)=\VK_K(H)$ by (G1). This bound holds for every
$T\ge1$.

Step~2: \emph{we show (ii).} Let $\varepsilon>0$. By (G2), let $\pi$ be a continuous batched policy with $\sup_{h\in H}\urk(\pi,h)\le\Vb(H)+\varepsilon$. We apply
Theorem~\ref{thm:u-transfer}(ii) with $\Theta:=H\times\{\mathbf 1\}$, $C_3:=2$, $m_0:=0$ and $\sigma:=1$. Its hypotheses hold for every instance
$(\mathcal N(h_k/\sqrt T,1))_k$ with $h\in H$:
\begin{itemize}
\item its local parameters are $(h,\mathbf 1)\in\Theta$ (since $(h_k/\sqrt T-0)\sqrt T/1=h_k$ and the variances are $1$);
\item its third moments are at most $C_3$ (since $\E|G|^3=2\sqrt{2/\pi}<2$ for $G\sim\mathcal N(0,1)$).
\end{itemize}
Let $\gamma_T$ be the supremum of Theorem~\ref{thm:u-transfer}(ii); then $\gamma_T\to0$. Since $\uDisc{\pi}{T}{0,1}$ is a $T$-round policy,
\begin{align*}
\frac{\uRG_H(T)}{\sqrt T}&\le\sup_{h\in H}\frac{\RT\big(\uDisc{\pi}{T}{0,1},(\mathcal N(h_k/\sqrt T,1))_k\big)}{\sqrt T} &&\text{(definition of $\uRG_H$)}\\
&\le\sup_{h\in H}\urk(\pi,h)+\gamma_T &&\text{(Theorem~\ref{thm:u-transfer}(ii))}\\
&\le\Vb(H)+\varepsilon+\gamma_T &&\text{(choice of $\pi$)}.
\end{align*}
Hence $\limsup_T\uRG_H(T)/\sqrt T\le\Vb(H)+\varepsilon$ for every $\varepsilon>0$. With Step~1, $\uRG_H(T)/\sqrt T\to\Vb(H)=\VK_K(H)$, by (G1).
\end{proof}

\subsection{Proof of Corollary~\ref{cor:u-gauss}}
\begin{proof}
\proves{cor:u-gauss}
\uses{prop:u-local, def:class, def:regret, def:risk}
\textit{Setting.} $M>0$, $T\ge(2\sbar M/(b-a))^2$, $\mu_0:=(a+b)/2$ and $H:=[-M,M]^K$. \textit{Goal.} $\Rs_{K,\mathcal C_{\mathrm G}}(T)\ge\sbar\sqrt T\,\VK_K(H)$.

For $h\in H$ we have $\sbar|h_k|/\sqrt T\le\sbar M/\sqrt T\le(b-a)/2$, so the means $\mu_0+\sbar h_k/\sqrt T$ lie in $[a,b]$ and the instance
$\nu^h:=(\mathcal N(\mu_0+\sbar h_k/\sqrt T,\sbar^2))_k$ belongs to $\mathcal C_{\mathrm G}^K$ (also to the subclass $\sigma=\sbar$). The map
$x\mapsto(x-\mu_0)/\sbar$ sends the rewards of $\nu^h$ to those of $(\mathcal N(h_k/\sqrt T,1))_k$ and divides every gap by $\sbar$; composing a policy with this map
or with its inverse is a bijection between $T$-round policies, since $\mu_0$ and $\sbar$ are known. Hence
\begin{align*}
\Rs_{K,\mathcal C_{\mathrm G}}(T)&\ge\inf_\pi\sup_{h\in H}\RT(\pi,\nu^h) &&\text{(sub-family of instances)}\\
&=\sbar\,\uRG_H(T) &&\text{(rescaling)}\\
&\ge\sbar\sqrt T\,\VK_K(H) &&\text{(Proposition~\ref{prop:u-local}(i))}.\qedhere
\end{align*}
\end{proof}

\subsection{Notation of this section}
\noindent $\umus$, $\sigma_k^2$\quad best mean, variance of arm $k$\hfill{\small Sec.~\ref{sec:app-u-setting}}\par
\noindent $\uSimp$, $\pi=(m,(t_j),(F_j))$\quad simplex, batched policy\hfill{\small Def.~\ref{def:u-batched}}\par
\noindent $\urk(\pi,h;v)$, $\urk(\pi,h)$\quad risk of a batched policy with variances $v$; $v=\mathbf 1$\hfill{\small Def.~\ref{def:u-batched}}\par
\noindent $\uDisc{\pi}{T}{m_0,\sigma}$, $\hat X$, $\hat\tau$\quad discretization, its observed path and clocks\hfill{\small Def.~\ref{def:u-disc}}\par
\noindent $\uRG_H(T)$\quad minimax regret over $(\mathcal N(h_k/\sqrt T,1))_k$, $h\in H$\hfill{\small Sec.~\ref{sec:app-u-setting}}\par
\noindent $\uN$, $\urho_n$, $\uS_n$, $\uS$, $\uio$, $\umuref$\quad screening length, radius, active sets, survivors, reference arm and mean\hfill{\small Alg.~\ref{def:sc}}\par
\noindent $\usigh_k^2$, $\uw_k$, $\uc(k)$\quad variance estimate, garbling variance, real arm behind virtual arm $k$\hfill{\small Alg.~\ref{def:sc}}\par
\noindent $\upulls_k$, $\uTp$, $\unu$, $\umusS$\quad screening pulls, remaining rounds, virtual instance, best surviving mean\hfill{\small Sec.~\ref{sec:app-u-setting}}\par
\noindent $\uD_k$, $\uD$\quad excess deviation of arm $k$ over the radii, its sum\hfill{\small Sec.~\ref{sec:app-u-screen}}\par
\noindent $\uCD$\quad $2\sqrt{2\pi}/(1-2^{-1/2})<17.2$\hfill{\small Lemma~\ref{lem:u-dev}}\par
\noindent $\upv(N,\eta)$\quad tail bound of a variance estimate\hfill{\small Lemma~\ref{lem:u-var}}\par
\noindent $\uF$, $\uRt$\quad $\sigma$-field of Stage~1, regret of the core on $\unu$\hfill{\small Sec.~\ref{sec:app-u-proof-main}}\par
\noindent $\uM$, $\uCt$, $\uG$\quad size of the box, third-moment constant, good event\hfill{\small Sec.~\ref{sec:app-u-proof-main}}\par

\section{Proof of Theorem~\ref{thm:bounds}}\label{sec:app-bounds}
\begin{proof}
\proves{thm:bounds}
\uses{cor:examples, def:kappa, def:regret, def:class, def:risk}
\paragraph{Lower bound.} \emph{Idea.} The method is that of \cite{GMS19,GMS20}; what we add is the choice of the prior and of its height, which gives the constant $\kap$.
Put one good arm at a uniform position. By the fundamental inequality of \cite{GMS19}, the probability of pulling the good arm is
controlled by the information gathered on each arm; averaging over the position and using the joint convexity of $\kl$ shows that no policy finds the good arm faster
than an exponential clock of rate $D/K$, up to the factor $\lamK_K$. Summing over the rounds gives the bound.

\noindent\emph{Setting.} Fix $K\ge2$, $T\ge1$ and $a>0$ with $\delta:=a\sqrt{K/T}\le1$. Let $J$ be uniform on $[K]$; arm $J$ has law $F:=\mathcal N(\delta,1)$ and the other arms have law
$G:=\mathcal N(0,1)$. Let $D:=\KL(G,F)=\delta^2/2$. Fix any policy, possibly randomized. \emph{Goal.} Show that, for every $t\in\{0,\dots,T-1\}$,
\begin{equation}\label{eq:search}
\P(A_{t+1}\ne J)\ \ge\ \exp\Big(-\frac{Dt+K\hb(1/K)}{K-1}\Big),
\end{equation}
then sum over $t$.

Step 1: \emph{one arm.} Let $P_0$ be the law of the run (rewards, randomization and actions up to round $t+1$) when all arms have law $G$, and $P_j$ its law when arm $j$ has law $F$.
Let $q_j:=P_0(A_{t+1}=j)$ and $p_j:=P_j(A_{t+1}=j)$. Under $P_0$ and $P_j$ only the rewards of arm $j$ have different laws, so by the chain rule
$\KL(P_0,P_j)=D\,\E_0[N_j(t)]$. By the data-processing inequality applied to the indicator of $\{A_{t+1}=j\}$ (the ``fundamental inequality'' of \cite{GMS19}),
\[
\kl(q_j,p_j)\ \le\ \KL(P_0,P_j)\ =\ D\,\E_0[N_j(t)] .
\]
Step 2: \emph{average over the position.} Let $\bar p:=\frac1K\sum_jp_j=\P(A_{t+1}=J)$. Since $\sum_jq_j=1$, the average of the $q_j$ is $\frac1K$; since $\sum_j\E_0[N_j(t)]=t$,
\begin{align*}
\kl\Big(\frac1K,\bar p\Big)&\le\frac1K\sum_j\kl(q_j,p_j) &&\text{(joint convexity of }\kl\text{)}\\
&\le\frac DK\sum_j\E_0[N_j(t)]=\frac{Dt}K &&\text{(Step 1)}.
\end{align*}
Step 3: \emph{invert.} For $q,p\in[0,1]$,
\[
\kl(q,p)=-\hb(q)+q\log\frac1p+(1-q)\log\frac1{1-p}\ \ge\ (1-q)\log\frac1{1-p}-\hb(q).
\]
With $q=\frac1K$ and Step 2, $(1-\frac1K)\log\frac1{1-\bar p}\le\frac{Dt}K+\hb(\frac1K)$, which is \eqref{eq:search}.

Step 4: \emph{sum.} The Bayes regret of the policy is $\delta\sum_{t=0}^{T-1}\P(A_{t+1}\ne J)$, and the minimax regret is at least the Bayes regret. Let $\beta:=\frac{a^2K}{2(K-1)}$, so that
$\frac D{K-1}t=\beta\frac tT$. A sum of nonincreasing terms is at least the corresponding integral, hence
\begin{align*}
\Rs_K(T)&\ge\delta\,\lamK_K\sum_{t=0}^{T-1}e^{-\beta t/T} &&\text{(\eqref{eq:search}, definition of }\lamK_K\text{)}\\
&\ge a\sqrt{\frac KT}\,\lamK_K\,T\,\frac{1-e^{-\beta}}\beta &&\text{(}\textstyle\sum_{t=0}^{T-1}e^{-\beta t/T}\ge\int_0^Te^{-\beta s/T}ds\text{)}\\
&=\lamK_K\sqrt{(K-1)T}\cdot\sqrt2\,\frac{1-e^{-\beta}}{\sqrt\beta} &&\text{(}a=\sqrt{2\beta(K-1)/K}\text{)}.
\end{align*}
Step 5: \emph{choice of $\beta$.} The function $\beta\mapsto(1-e^{-\beta})/\sqrt\beta$ has derivative of the sign of $2\beta e^{-\beta}-(1-e^{-\beta})$, that is of $1+2\beta-e^\beta$, which is positive on
$(0,\us)$ and negative after; its maximum is at $\beta=\us$, where $\sqrt2(1-e^{-\us})/\sqrt{\us}=\kap$ (Section~\ref{sec:results}). The constraint $\delta\le1$ reads
$T\ge a^2K=2\us(K-1)$. Hence $\Rs_K(T)\ge\kap\lamK_K\sqrt{(K-1)T}$ for every $T\ge2\us(K-1)$.

Step 6: \emph{the value.} Divide by $\sqrt T$ and let $T\to\infty$: by Corollary~\ref{cor:examples} (Gaussian rewards, $\sbar=1$), $\VK_K=\lim_T\Rs_K(T)/\sqrt T\ge\kap\lamK_K\sqrt{K-1}$.

\paragraph{Upper bound.} For rewards in $[0,1]$, $\Rs_K(T)\le\sqrt{2KT}$ for every $T$ \cite{LS19,ZL19}, and $\lim_T\Rs_K(T)/\sqrt T=\frac12\VK_K$ by Corollary~\ref{cor:examples}.
Hence $\frac12\VK_K\le\sqrt{2K}$.
\end{proof}

\begin{remark}
The same prior gives the bound $\P(A_{t+1}=J)\le1-e^{-Dt/K}+\eta_K$ with $\eta_K:=\min_{c>0}(\frac cK+\frac1{e^c-1})$, by a change of measure to the instance without good arm
and a Lagrangian relaxation arm by arm \cite{PD}; it has the same limit and worse constants at finite $K$. Pinsker's inequality in the limit game gives
$\VK_K\ge\frac12(1-\frac1K)^2\sqrt K$, also weaker.
\end{remark}

\section{Two arms: computations and certified bounds}\label{sec:app-K2}

\newcommand{\kU}{\sy{k:UI}{U^{\lambda}}}
\newcommand{\kI}{\sy{k:UI}{I^{\lambda}}}
\newcommand{\kF}{\sy{k:UI}{\mathcal G}}
\newcommand{\kg}{\sy{k:g}{g}}
\newcommand{\kR}{\sy{k:rec}{\mathcal R}}
\newcommand{\kB}{\sy{k:rec}{\mathcal B}}
\newcommand{\kPi}{\sy{k:rec}{\Pi}}
\newcommand{\kPri}{\sy{k:prior}{\Pi_{\lambda,d}}}
\newcommand{\kQ}{\sy{k:prior}{\mathbb Q}}
\newcommand{\kP}[1]{\sy{k:prior}{\mathbb P}_{#1}}
\newcommand{\kBb}{\sy{k:prior}{B^\flat}}
\newcommand{\kW}{\sy{k:W}{\Psi}}
\newcommand{\kL}{\sy{k:W}{L_\varepsilon}}
\newcommand{\kT}{\sy{k:W}{T}^{\mathrm{rem}}}
\newcommand{\kchi}{\sy{k:ker}{\chi}}
\newcommand{\kGam}{\sy{k:ker}{\Gamma}}
\newcommand{\kS}{\sy{k:step}{\mathcal S}}
\newcommand{\kdfp}{\sy{k:alg}{\delta_{\mathrm{fp}}}}
\newcommand{\kTC}{\sy{k:tc}{\mathrm{TC}}}
\newcommand{\kH}{\sy{k:tc}{H}}
\newcommand{\kr}{\sy{k:rho}{r_{\mathrm{TC}}}}
\newcommand{\kEm}{\sy{k:rho}{E_-}}
\newcommand{\kPo}{\sy{k:rho}{P_0}}


\medskip
\noindent This appendix has two parts.
\begin{itemize}
\item Sections~\ref{sec:K2-red} and~\ref{sec:K2-num} give the two-arm reduction and the numerical solution of the batched Bayes recursion. This is what produces
Tables~\ref{tab:eps} and~\ref{tab:saddle} and the threshold table $\thr$ of Definition~\ref{def:opt2}. These computations are numerical and come without
error bounds.
\item Sections~\ref{sec:K2-proof}--\ref{sec:K2-upper} give the computer-assisted proof of Proposition~\ref{prop:K2}, and Section~\ref{sec:K2-open} lists what is
not certified.
\end{itemize}
Throughout, $K=2$, rewards have unit variance, $Z$ denotes a standard Gaussian variable, and $\Phi$, $\varphi$ are the standard Gaussian distribution function and density.

\subsection{The two-arm reduction}\label{sec:K2-red}

\paragraph{Data and observation points.} We use the representation of the limit game in which arm $k$, observed during a time $s$, returns
\[
Y_k(s):=h_ks+W_k(s),
\]
where $W_1,W_2$ are independent standard Brownian motions. At horizon fraction $t$ the data are $X_k(t)=Y_k(\tau_k(t))$; under the local instance $h$ this is the
law $\lh_hQ$ of Section~\ref{sec:game}, by Girsanov's theorem. A batched policy also uses an independent uniform variable $\zeta$, its randomization.
An \emph{observation point} is a pair $o=(o_1,o_2)\in[0,\infty)^2$, read as ``arm $k$ has been observed during $o_k$''. We write $o\le o'$ componentwise and
$t(o):=o_1+o_2$. Let
\[
\kF_o:=\sigma\big(Y_1(s),\,s\le o_1;\ Y_2(s),\,s\le o_2\big)\vee\sigma(\zeta),\qquad
\lh_h(o):=\exp\Big(\sum_{k=1}^2\big(h_kY_k(o_k)-\tfrac12h_k^2o_k\big)\Big).
\]
By the Cameron--Martin formula, $\lh_h(o)$ is the density on $\kF_o$ of the data under $h$ with respect to the data under $h=0$; at $o=\tau(1)$ it is the
likelihood $\lh_h$ of Section~\ref{sec:game}. We write the local means as $h=(m+v,m-v)$: $m$ is the common mean and $v$ the half-gap. The gap is $2|v|$, so for
every policy $\pi$
\[
\rk(\pi,h)=2|v|\,\E_h[\text{time spent on the worse arm}].
\]

\paragraph{Statistics.}\dt{k:UI} For $\lambda\ge0$ and an observation point $o$ with $t(o)+\lambda>0$, let
\[
\kU(o):=\frac12\Big[\big(Y_1(o_1)-Y_2(o_2)\big)-\big(Y_1(o_1)+Y_2(o_2)\big)\frac{o_1-o_2}{t(o)+\lambda}\Big],\qquad
\kI(o):=\frac{o_1o_2+\lambda\,t(o)/4}{t(o)+\lambda}.
\]
The parameter $\lambda$ will be the precision of a normal prior on the common mean, and $\lambda=0$ is the flat prior. For $\lambda=0$ and $o=\tau$ these are the
statistics of Section~\ref{sec:two}:
\[
U^0(\tau)=\frac{X_1\tau_2-X_2\tau_1}{t}=\Ud,\qquad I^0(\tau)=\frac{\tau_1\tau_2}{t}=\Id .
\]

\begin{lemma}[two-arm reduction]\label{lem:K2-gauss}
\uses{def:risk}
Let $\lambda\ge0$, $v\in\R$, and let $o\le o'$ be deterministic observation points with $t(o)+\lambda>0$. Let $\pi_\lambda:=\mathcal N(0,1/\lambda)$ if $\lambda>0$,
and let $\pi_0$ be the Lebesgue measure. Then:
\begin{enumerate}
\item[(i)] (integrating the common mean) there is a random variable $C_\lambda(o)\in(0,\infty)$, which does not depend on $v$, such that
\[
\int\lh_{(m+v,m-v)}(o)\,\pi_\lambda(dm)=C_\lambda(o)\exp\big(2v\,\kU(o)-2v^2\kI(o)\big).
\]
In particular, for a prior $\rho$ on $v$ and the prior $\pi_\lambda$ on $m$, the posterior of $v$ given $\kF_o$ is proportional to $\rho(dv)\exp\big(2v\kU(o)-2v^2\kI(o)\big)$;
\item[(ii)] (proper prior, $\lambda>0$)\dt{k:prior} let $\kP{v}$ be the law of the data and of $\zeta$ when $m\sim\mathcal N(0,1/\lambda)$ and the half-gap is $v$, and let
$\kQ:=\kP{0}$. Then
\[
\frac{d\kP{v}}{d\kQ}\Big|_{\kF_o}=\exp\big(2v\,\kU(o)-2v^2\kI(o)\big),
\]
and under $\kQ$ the increment $\kU(o')-\kU(o)$ is independent of $\kF_o$, with law $\mathcal N\big(0,\kI(o')-\kI(o)\big)$; moreover $\kU(0)=\kI(0)=0$;
\item[(iii)] (fixed instance, $\lambda=0$) for every $h=(m+v,m-v)$, $\Ud(o)=2v\,\Id(o)+N(o)$, where $N(o):=\big(W_1(o_1)o_2-W_2(o_2)o_1\big)/t(o)$ does not depend on $h$.
Moreover, $N(o')-N(o)$ is independent of $\sigma\big(N(p):p\le o,\ t(p)>0\big)$, with law $\mathcal N\big(0,\Id(o')-\Id(o)\big)$.
\end{enumerate}
\end{lemma}
\noindent\textit{Role.} Item~(i) is the reduction of Section~\ref{sec:two}: the posterior of the half-gap depends on the data only through $o$ and the
statistic. Item~(iii) says that a policy that only uses $\Ud$ sees a Brownian motion with drift $2v$ in the clock $\Id$, whatever the common mean; it drives
the numerical recursion of Section~\ref{sec:K2-num} and the test of Section~\ref{sec:K2-upper}. Item~(ii) is the version with a proper prior, used for the
lower bound.

\begin{proof}
\uses{def:risk}
\textit{Step 1: we prove (i).} Let $S:=Y_1(o_1)+Y_2(o_2)$, $D:=Y_1(o_1)-Y_2(o_2)$, $t:=t(o)$ and $P:=t+\lambda$. With $h=(m+v,m-v)$, the exponent of $\lh_h(o)$ is
\[
mS+vD-\tfrac12m^2t-mv(o_1-o_2)-\tfrac12v^2t .
\]
The density of $\pi_\lambda$ is proportional to $e^{-\lambda m^2/2}$ (a constant if $\lambda=0$). Integrating in $m$ a Gaussian function,
\[
\int e^{\,m(S-v(o_1-o_2))-m^2P/2}\,dm=\sqrt{2\pi/P}\,\exp\Big(\frac{(S-v(o_1-o_2))^2}{2P}\Big),
\]
which is finite since $P>0$. The terms of the resulting exponent that depend on $v$ are
\begin{align*}
&vD-\frac{v^2t}{2}+\frac{-2vS(o_1-o_2)+v^2(o_1-o_2)^2}{2P}\\
&\qquad=v\Big(D-S\,\frac{o_1-o_2}{P}\Big)-\frac{v^2}{2P}\big(tP-(o_1-o_2)^2\big) && \text{(collect the powers of $v$)}\\
&\qquad=2v\,\kU(o)-\frac{v^2}{2P}\big(4o_1o_2+\lambda t\big) && \text{(definition of $\kU$; $t^2-(o_1-o_2)^2=4o_1o_2$)}\\
&\qquad=2v\,\kU(o)-2v^2\,\kI(o) && \text{(definition of $\kI$)}.
\end{align*}
The other factors form $C_\lambda(o)$. The posterior statement is Bayes' formula.

\textit{Step 2: we prove (ii).} For $\lambda>0$, $\pi_\lambda$ is a probability, so by (i) the density of $\kP{v}$ on $\kF_o$ with respect to the law under $h=0$ is
$C_\lambda(o)e^{2v\kU(o)-2v^2\kI(o)}$. Dividing by its value at $v=0$ gives the likelihood ratio; $\zeta$ has the same law under every $\kP{v}$. Under $\kQ$ the data
and $\kU(o')$ are jointly Gaussian, so the conditional law of $\kU(o')$ given $\kF_o$ is $\mathcal N(\mu_c,s^2)$, with $\mu_c$ $\kF_o$-measurable and $s^2$ deterministic.
Likelihood ratios on the nested $\sigma$-fields $\kF_o\subset\kF_{o'}$ form a $\kQ$-martingale, so for every $v$, almost surely,
\[
e^{2v\mu_c+2v^2s^2-2v^2\kI(o')}=\E_{\kQ}\big[e^{2v\kU(o')-2v^2\kI(o')}\,\big|\,\kF_o\big]=e^{2v\kU(o)-2v^2\kI(o)} .
\]
This holds almost surely simultaneously for all rational $v$. Both exponents are polynomials in $v$, so their coefficients agree:
\[
\mu_c=\kU(o),\qquad s^2=\kI(o')-\kI(o).
\]
The conditional law of $\kU(o')-\kU(o)$ given $\kF_o$ is therefore $\mathcal N(0,\kI(o')-\kI(o))$ whatever the conditioning, which is the independence claimed.
At $o=0$ the definitions give $\kU(0)=\kI(0)=0$.

\textit{Step 3: we prove (iii).} Substituting $Y_k(o_k)=h_ko_k+W_k(o_k)$ in $\Ud(o)=(Y_1(o_1)o_2-Y_2(o_2)o_1)/t(o)$ gives
$\Ud(o)=(h_1-h_2)\,o_1o_2/t(o)+N(o)=2v\Id(o)+N(o)$. The process $N$ is centered Gaussian. For $p\le o'$ with $t(p)>0$, using
$\mathrm{Cov}(W_k(o'_k),W_k(p_k))=p_k$,
\[
\mathrm{Cov}\big(N(o'),N(p)\big)=\frac{o'_2p_2\,p_1+o'_1p_1\,p_2}{t(o')\,t(p)}=\frac{p_1p_2}{t(p)}=\Id(p).
\]
Hence for $p\le o\le o'$ we get $\mathrm{Cov}(N(o')-N(o),N(p))=\Id(p)-\Id(p)=0$ and $\mathrm{Var}(N(o')-N(o))=\Id(o')-2\Id(o)+\Id(o)=\Id(o')-\Id(o)$.
Jointly Gaussian and uncorrelated variables are independent.
\end{proof}

\paragraph{The information clock.} Differentiating $\kI$,
\[
\frac{\partial\kI}{\partial o_1}=\frac{(o_2+\lambda/2)^2}{(t(o)+\lambda)^2},\qquad \frac{\partial\kI}{\partial o_2}=\frac{(o_1+\lambda/2)^2}{(t(o)+\lambda)^2}.
\]
For $\lambda>0$ both derivatives lie in $(0,1]$, so a block of length $\varepsilon$ on one arm increases $\kI$ by an amount in $(0,\varepsilon]$. For $\lambda=0$ they are
the speeds $(\tau_2/t)^2$ and $(\tau_1/t)^2$ of Section~\ref{sec:two}: the policy only chooses at which speed the clock $\Id$ advances. In particular,
\begin{equation}\label{eq:K2-clock}
\Id(\tau+\varepsilon e_1)-\Id(\tau)=\frac{\varepsilon\,\tau_2^2}{t(t+\varepsilon)},\qquad \Id(\tau+\varepsilon e_2)-\Id(\tau)=\frac{\varepsilon\,\tau_1^2}{t(t+\varepsilon)},
\end{equation}
with $e_1:=(1,0)$, $e_2:=(0,1)$ and $t:=\tau_1+\tau_2$.

\paragraph{Bayes cost rates.}\dt{k:g} For $d>0$ let
\[
\kg_1(u,I):=d\,e^{-2du-2d^2I},\qquad \kg_2(u,I):=d\,e^{2du-2d^2I}.
\]
For the prior ``$v=\pm d$ with probability $\frac12$ each'', arm~$1$ is the worse arm when $v=-d$. By Lemma~\ref{lem:K2-gauss}(ii), for $\lambda>0$ and every event
$E\in\kF_o$,
\begin{equation}\label{eq:K2-cost}
2d\cdot\tfrac12\,\kP{-d}(E)=\E_{\kQ}\big[\1_E\,\kg_1\big(\kU(o),\kI(o)\big)\big],\qquad 2d\cdot\tfrac12\,\kP{+d}(E)=\E_{\kQ}\big[\1_E\,\kg_2\big(\kU(o),\kI(o)\big)\big].
\end{equation}
So the Bayes cost of a block of length $\varepsilon$ on arm $l$ is $\varepsilon\,\kg_l$ evaluated at the current statistic, under the reference law. The same holds for
$\lambda=0$ under $h=(m,m)$, for events determined by the values of $\Ud$ (Lemma~\ref{lem:K2-dp}).

\subsection{The batched Bayes recursion and its numerical solution}\label{sec:K2-num}

\medskip
\noindent This section explains how Tables~\ref{tab:eps} and~\ref{tab:saddle} and the threshold table of Definition~\ref{def:opt2} were computed. The recursion
is exact; its numerical solution has no error bound.

\begin{definition}[Bayes recursion of a two-point prior]\label{def:K2-rec}\dt{k:rec}
Let $d>0$, let $n\ge3$ be an integer and $\varepsilon:=1/n$. For integers $i,j\ge1$ with $i+j\le n$, let $I_{ij}:=\Id(i\varepsilon,j\varepsilon)=\varepsilon\,ij/(i+j)$ and,
by~\eqref{eq:K2-clock},
\[
\big(s^{(1)}_{ij}\big)^2:=\frac{\varepsilon\,j^2}{(i+j)(i+j+1)},\qquad \big(s^{(2)}_{ij}\big)^2:=\frac{\varepsilon\,i^2}{(i+j)(i+j+1)} .
\]
Define $\kR_{ij}:\R\to[0,\infty)$ by $\kR_{ij}:=0$ if $i+j=n$, and otherwise
\[
\kR_{ij}(u):=\min_{l\in\{1,2\}}\Big\{\varepsilon\,\kg_l(u,I_{ij})+\E\big[\kR_{(i,j)+e_l}\big(u+s^{(l)}_{ij}Z\big)\big]\Big\}.
\]
Let $\kB_\varepsilon(d):=2d\varepsilon+\E\big[\kR_{11}\big(\sqrt{\varepsilon/2}\,Z\big)\big]$. Let $\kPi_\varepsilon$ be the set of batched policies that give $\varepsilon$ to each arm on
$[0,2\varepsilon]$, and then, on each $[k\varepsilon,(k+1)\varepsilon]$ with $2\le k\le n-1$, pull one arm chosen as a Borel function of $\zeta$ and of the values of $\Ud$ at the
previous decision times.
\end{definition}

\begin{lemma}[the recursion computes a Bayes risk]\label{lem:K2-dp}
\uses{def:K2-rec, def:risk}
Let $d>0$ and $n\ge3$. Then for every $m\in\R$,
\[
\min_{\pi\in\kPi_\varepsilon}\tfrac12\Big[\rk\big(\pi,(m+d,m-d)\big)+\rk\big(\pi,(m-d,m+d)\big)\Big]=\kB_\varepsilon(d),
\]
and the minimum is attained by the policy that, at the grid point $(i\varepsilon,j\varepsilon)$ and the value $u$ of $\Ud$, pulls arm~$1$ if and only if the first term of the
minimum in Definition~\ref{def:K2-rec} is at most the second one.
\end{lemma}
\noindent\textit{Role.} It says what Table~\ref{tab:eps} reports: the Bayes risk of the two-point prior on the half-gap, flat in the common mean, when decisions are
held during $\varepsilon$.
\begin{proof}
\uses{lem:K2-gauss, def:K2-rec}
\textit{Step 1: the regret does not depend on $m$.} Fix $\pi\in\kPi_\varepsilon$. It uses only $\zeta$ and values of $\Ud$, whose joint law under $(m+v,m-v)$ does not depend
on $m$ (Lemma~\ref{lem:K2-gauss}(iii)). Neither does the regret, which is $2|v|$ times the time spent on the worse arm.

\textit{Step 2: change of measure to $v=0$.} After $k\ge2$ decision times the policy is at a lattice point $(i\varepsilon,j\varepsilon)$ with $i+j=k$, and it has observed
$\Ud$ at the points visited so far. Along every lattice path, these values have independent Gaussian increments, of mean $2v\Delta$ and variance $\Delta$, where
$\Delta$ is the increment of $\Id$ (Lemma~\ref{lem:K2-gauss}(iii)). Their joint density under half-gap $v$ with respect to half-gap $0$ is therefore the telescoping
product $\exp(2v\Ud(o)-2v^2\Id(o))$, with $o$ the current point, on each branch of the path; $\zeta$ has the same law in both cases. Write $\P_h$, $\E_h$ for the law
under the instance $h$. So for a lattice point $o$ and an event $E\subset\{\text{the policy is at }o\}$ determined by $\zeta$ and the values observed so far,
\[
2d\cdot\tfrac12\,\P_{(m-d,m+d)}(E)=\E_{(m,m)}\big[\1_E\,\kg_1(\Ud(o),\Id(o))\big],
\]
and similarly with $\kg_2$ and $(m+d,m-d)$.

\textit{Step 3: backward induction.} The first batch costs $2d\varepsilon$ under both hypotheses, since the worse arm receives $\varepsilon$. By Step~2, each later block on arm $l$
costs $\varepsilon\,\kg_l(\Ud,\Id)$ under $\E_{(m,m)}$. Under $(m,m)$ the increments of $\Ud$ are independent of its past values, so the minimal expected cost-to-go at
$(i\varepsilon,j\varepsilon)$ is a function of the current value of $\Ud$; by backward induction on $i+j$ it is $\kR_{ij}$, and the minimizing choices are those of the statement.
Randomization does not lower the minimum of an expectation. At $(\varepsilon,\varepsilon)$, $\Ud\sim\mathcal N(0,\varepsilon/2)$ under $(m,m)$, which gives $\kB_\varepsilon(d)$.
\end{proof}

\paragraph{Numerical solution.} The recursion is solved diagonal by diagonal, $t=k\varepsilon$ for $k=n-1,\dots,2$.
\begin{itemize}
\item \emph{Grid in $u$.} A periodic grid of $[-5,5)$ with step $\Delta u=0.0125$ (800 points). The Gaussian expectation is applied exactly to the trigonometric
interpolant: the Fourier coefficients are multiplied by $e^{-s^2\omega^2/2}$ (FFT).
\item \emph{Periodization.} Always pulling the same arm $l$ shows, as in Lemma~\ref{lem:K2-upper} below, that $\kR_{ij}(u)\le(1-(i+j)\varepsilon)\,d\,e^{-2d|u|-2d^2I_{ij}}$.
At $|u|=5$ this bound is $e^{-10d}\approx10^{-7}$ times its value at $u=0$, so the periodization is harmless.
\item \emph{What is not exact.} The minimum of the two options has a kink at the threshold, on which the spectral step is not exact; the root expectation is a
quadrature on the grid. No error bound is derived.
\item \emph{Maximization in $d$.} For each $\varepsilon$, $\kB_\varepsilon$ is maximized by the vertex of the parabola through $d_0-0.03$, $d_0$, $d_0+0.03$, where $d_0$ is the
maximizer found for the previous $\varepsilon$ (starting from $1.60$), and then re-evaluated at the vertex.
\item \emph{Effect of the grid.} With the spectral scheme, halving the step from $\Delta u=0.0125$ to $0.00625$ changes the recursion value by $2\cdot10^{-6}$
($\varepsilon=0.02$), $1\cdot10^{-6}$ ($\varepsilon=0.005$), $4\cdot10^{-6}$ ($\varepsilon=0.00125$, $d=1.578$) and $6\cdot10^{-6}$ ($\varepsilon=0.000625$, $d=1.5773$:
$0.637735$ against $0.637729$); outputs in \texttt{scripts/convergence\_dU.txt} and \texttt{scripts/convergence\_dU\_petit\_eps\_sortie.txt}. The file
\texttt{scripts/hjb2\_sup\_dU0025.txt} comes from an earlier explicit finite-difference scheme and is not comparable. The grid error at $\Delta u=0.0125$ is therefore
below $10^{-5}$, small against the extrapolation error in $\varepsilon$.
\end{itemize}

\paragraph{The least favorable two-point prior (Table~\ref{tab:eps}).} The maximizers are $d^\star(\varepsilon)=1.6329$, $1.6019$, $1.5886$, $1.5820$, $1.5788$, $1.5773$
for $\varepsilon=0.02,\dots,0.000625$. The error in $\varepsilon$ is of first order, so Table~\ref{tab:eps} also gives the first-order (Richardson) extrapolations
$2\sup_d\kB_\varepsilon-\sup_d\kB_{2\varepsilon}$. Their successive increments are $4.3$, $1.4$, $0.46$ and $0.16$ (unit $10^{-4}$), roughly geometric with ratio
$\frac13$. Summing this geometric tail from the last value $0.637329$ gives $0.63734$; the error bar $\pm2\cdot10^{-5}$ is a heuristic for the extrapolation in $\varepsilon$
only. The same extrapolation of the maximizers gives $\dst\approx1.576$. At $\varepsilon=0.02$ and $d=1.6$ the recursion gives $0.6516$, which matches the value $0.65$
of \cite{Kol17}.

\paragraph{Saddle-point check (Table~\ref{tab:saddle}).} For $d=1.582$ (the maximizer at $\varepsilon=0.0025$) the decisions of the Bayes policy of
Lemma~\ref{lem:K2-dp} are stored, and its frequentist risk at a half-gap $v$ is computed by the same recursion with these decisions:
\begin{itemize}
\item a block on arm $l$ moves $\Ud$ by $\mathcal N\big(2v s^2,s^2\big)$ with $s=s^{(l)}_{ij}$ (Lemma~\ref{lem:K2-gauss}(iii)); the drift is applied exactly in Fourier, as a phase;
\item a block on the worse arm costs $2|v|\varepsilon$, the first batch costs $2|v|\varepsilon$, and $\Ud(\varepsilon,\varepsilon)\sim\mathcal N(v\varepsilon,\varepsilon/2)$.
\end{itemize}
On the 17 values $v\in[0.25,6]$ of the script, the risk is maximal at $v=1.582$, where it is $0.63894$; the Bayes risk is $0.63897$. Within $3\cdot10^{-5}$, the
two-point prior and its Bayes policy thus form a saddle point of the $\varepsilon$-batched game restricted to $|v|\le6$. This is numerical evidence, not a proof: the
grid in $v$ comes without derivative bounds, and for large $|v|$ the risk tends to infinity (the first batch alone costs $2|v|\varepsilon$). The statement is about boxes,
which is consistent with Proposition~\ref{prop:four}(b).

\paragraph{The threshold table of Definition~\ref{def:opt2}.} The script \texttt{table\_seuils.py} solves the recursion with $d=1.576\approx\dst$, $\varepsilon=0.0025$
and $\Delta u=0.0125$. At each grid point $(i\varepsilon,j\varepsilon)$ with $i,j\ge1$ and $i+j\le n-1$ it stores the smallest grid value of $u$ at which arm~$1$ is chosen: this
is $\thr(i\varepsilon,j\varepsilon)$, set to $+\infty$ if arm~$1$ is never chosen.
\begin{itemize}
\item The rule ``arm $1$ if and only if $\Ud\ge\thr$'' is the Bayes rule of Lemma~\ref{lem:K2-dp} on the rows where the set of grid values $u$ at which arm~$1$ is chosen
is of the form $[\thr,\infty)$; the script counts the other rows.
\item The symmetries $\thr(\tau,\tau)=0$ and $\thr(\tau_2,\tau_1)=-\thr(\tau_1,\tau_2)$ hold up to one grid step. For instance, at $d=1.582$ the decision switches
between $u=-0.0125$ and $u=0$ at $(0.2,0.2)$, between $-0.2125$ and $-0.2$ at $(0.1,0.3)$, and between $0.2$ and $0.2125$ at $(0.3,0.1)$.
\item The table does not depend on $T$ or $\sigma$, because $\Ud$ and the clock are normalized. Definition~\ref{def:opt2} reads it at the grid point nearest to
$(\tau_1,\tau_2)$, in the coordinates $(t/\varepsilon,\tau_1/\varepsilon)$.
\item The posterior log-odds of ``arm $1$ is the better arm'' are $4d\,\Ud$ (Lemma~\ref{lem:K2-gauss}(i)). At $\thr(0.1,0.3)\approx-0.20$ the posterior probability is
$1/(1+e^{4\cdot1.576\cdot0.20})\approx0.22$: the policy keeps pulling the under-sampled arm although it is probably the worse one.
\end{itemize}
\noindent The script finds $24$ rows that are not of this form, out of about $8\cdot10^4$ grid points (output saved in \texttt{sim/table\_seuils\_sortie.txt}); they lie
at the ends of the rows, where the two decisions have nearly equal risks.

\subsection{Proof of Proposition~\ref{prop:K2}: overview}\label{sec:K2-proof}
\begin{proof}
\proves{prop:K2}
\uses{prop:four, lem:K2-box, lem:K2-conv, prop:K2-relax, lem:K2-step, def:K2-inf, prop:K2-Lnum, prop:K2-Unum, lem:K2-series, lem:K2-gauss, lem:K2-upper, lem:K2-semi, lem:K2-lip, lem:K2-tail, def:risk, def:batched, def:K2-W}
\paragraph{Statement.} Recall Proposition~\ref{prop:K2}:
\[
0.6341\ \le\ \VK_2\ \le\ 0.7478 .
\]

\paragraph{Link with $V_2$.} Both bounds go through batched policies on boxes. By Proposition~\ref{prop:four}(b),
\[
\VK_2=\lim_{M\to\infty}\Vb\big([-M,M]^2\big),
\]
the value of the limit game over all limit policies and all $h\in\R^2$ is the limit of the batched values on growing boxes.
\begin{itemize}
\item \emph{Lower bound: weak duality with a proper prior.} A normal prior on the common mean is carried by no box, but it is almost carried by large boxes. Its Bayes
risk over batched policies is therefore a lower bound on $\Vb([-M,M]^2)$ up to a vanishing term, hence on $\VK_2$ (Lemma~\ref{lem:K2-box}).
\item \emph{Upper bound: one policy on every box.} A batched policy whose risk on every box is at most $\sup_v\kr(v)$ bounds $\Vb([-M,M]^2)$ for every $M$, hence $\VK_2$
(Lemma~\ref{lem:K2-conv}).
\end{itemize}
Neither bound uses fully sequential policies directly: Proposition~\ref{prop:four}(b) does that.

\paragraph{In what sense this is a proof.}
\begin{itemize}
\item \emph{Proved, in exact real arithmetic:} $\VK_2\ge\kL(d,\lambda)$ for all $d,\lambda>0$ and $n\ge2$ (Lemma~\ref{lem:K2-box}, Proposition~\ref{prop:K2-relax});
$\kL(d,\lambda)$ is at least every certificate in the sense of Lemma~\ref{lem:K2-step}, in particular at least the output of Definition~\ref{def:K2-inf} for every margin
$\kdfp\ge0$ (Proposition~\ref{prop:K2-Lnum}(i)); $\VK_2\le\sup_v\kr(v)$, and $\sup_v\kr$ is at most the grid bound of Proposition~\ref{prop:K2-Unum}, built from series
with exact remainders, a Lipschitz bound and a tail bound.
\item \emph{Not proved in the strict sense:} the two finite computations, Definition~\ref{def:K2-inf} at $(d,n)=(1.565,800)$ and the series of
Lemma~\ref{lem:K2-series} on a grid of $v$, were run in IEEE double precision (NumPy and SciPy), without interval arithmetic or directed rounding. Their errors are
covered by margins justified by an \emph{a priori} error analysis (Remark~\ref{rem:K2-fp}), which relies on standard, but not formally guaranteed, accuracy of the
FFT and of the elementary functions. Proposition~\ref{prop:K2} is therefore a computer-assisted proof up to the floating-point evaluation of two explicit finite
computations. A machine-checkable certificate would rerun both in interval arithmetic.
\item \emph{Not certified at all:} the value $0.6373$ and the frequentist risk of the Bayes policy (Sections~\ref{sec:K2-num} and~\ref{sec:K2-open}).
\end{itemize}

\paragraph{Skeleton.} Fix $\lambda:=10^{-6}$, $d:=1.565$, $n:=800$ and $\varepsilon:=1/n$ for the lower bound, and $(b,S):=(0.30,0.21)$ for the upper bound. Then
\begin{align*}
\VK_2&\ge\kBb(\kPri) && \text{(Prop.~\ref{prop:four}(b) and Lemma~\ref{lem:K2-box}: a proper prior almost carried by boxes)}\\
&\ge\kL(d,\lambda) && \text{(Prop.~\ref{prop:K2-relax}: blocks revealed as soon as they are started)}\\
&\ge0.634171 && \text{(Prop.~\ref{prop:K2-Lnum}: minorant scheme, margin $\kdfp=10^{-8}$)},\\[3pt]
\VK_2&\le\sup_{v\in\R}\kr(v) && \text{(Prop.~\ref{prop:four}(b) and Lemma~\ref{lem:K2-conv}: test then commit)}\\
&\le0.747721 && \text{(Prop.~\ref{prop:K2-Unum}: series, Lipschitz and tail bounds)}.
\end{align*}
Rounding gives Proposition~\ref{prop:K2}. The roundings $0.634171\to0.6341$ and $0.747721\to0.7478$ leave margins of $7\cdot10^{-5}$ and $8\cdot10^{-5}$.

\paragraph{Dependencies.}
\begin{center}\small
\begin{tabular}{@{}l@{}}
\hline
\textbf{Lower bound} (Section~\ref{sec:K2-lower})\\
Lemma~\ref{lem:K2-gauss}(ii) $\Rightarrow$ cost identity~\eqref{eq:K2-cost} $\Rightarrow$ Proposition~\ref{prop:K2-relax}: $\kBb(\kPri)\ge\kL(d,\lambda)$\\
Proposition~\ref{prop:four}(b) $+$ Lemma~\ref{lem:K2-box} $\Rightarrow$ $\VK_2\ge\kBb(\kPri)$\\
Lemmas~\ref{lem:K2-upper} and~\ref{lem:K2-semi} $\Rightarrow$ Lemma~\ref{lem:K2-step} $\Rightarrow$ Proposition~\ref{prop:K2-Lnum}: $\kL(1.565,10^{-6})\ge0.634171$ at $\varepsilon=1/800$\\
\hline
\textbf{Upper bound} (Section~\ref{sec:K2-upper})\\
Lemma~\ref{lem:K2-gauss}(iii) $\Rightarrow$ risk formula~\eqref{eq:K2-rho}; Lemma~\ref{lem:K2-lip} $\Rightarrow$ Lemma~\ref{lem:K2-conv}: $\VK_2\le\sup_v\kr(v)$\\
Lemmas~\ref{lem:K2-lip}, \ref{lem:K2-series} and~\ref{lem:K2-tail} $\Rightarrow$ Proposition~\ref{prop:K2-Unum}: $\sup_v\kr(v)\le0.747721$\\
\hline
\end{tabular}
\end{center}

\begin{table}[ht]\centering\small
\begin{tabular}{llll}
\hline
computation & parameters & output of the run & certified bound\\ \hline
Definition~\ref{def:K2-inf}, $n=400$ & $d=1.56$ & $0.6317025746$ & $\kL\ge0.631698$\\
Definition~\ref{def:K2-inf}, $n=800$ & $d=1.565$ & $0.6341794644$ & $\kL\ge0.634171$\\
grid bound of Prop.~\ref{prop:K2-Unum} & $b=0.30$, $S=0.21$ & $0.7477208772$ & $\sup_v\kr\le0.747721$\\
\hline
\end{tabular}
\caption{The three computer-assisted computations. Lower bound: $\lambda=10^{-6}$, $h=10^{-3}$, $u_{\max}=3.5$, runs made with the margin $10^{-10}$; the certified
bound uses the margin $\kdfp=10^{-8}$ (Proposition~\ref{prop:K2-Lnum}). CPU times: $63$~s, $240$~s and $30$~s. The first and the third outputs were reproduced to the
last digit by an independent rerun, the second one by an instrumented copy of the code.}\label{tab:K2-runs}
\end{table}
\end{proof}

\subsection{The lower bound}\label{sec:K2-lower}

\subsubsection{Why does a proper prior bound the value from below?}

\begin{definition}[proper prior and Bayes risk over batched policies]\label{def:K2-prior}
\uses{def:risk}
Let $\lambda,d>0$. Let $\kPri$ be the law of $h=(m+v,m-v)$ when $m\sim\mathcal N(0,1/\lambda)$ and $v=\pm d$ with probability $\frac12$ each, independently. The laws
$\kP{v}$ and $\kQ$ are those of Lemma~\ref{lem:K2-gauss}(ii). Let
\[
\kBb(\kPri):=\inf_{\pi\text{ batched}}\int\rk(\pi,h)\,\kPri(dh).
\]
\end{definition}

\begin{lemma}[a proper prior almost carried by boxes]\label{lem:K2-box}
\uses{def:risk}
Let $\lambda,d>0$. Then $\VK_2\ge\kBb(\kPri)$.
\end{lemma}
\noindent\textit{Role.} It replaces the improper flat prior on the common mean by a proper one, at the only price of $\lambda>0$ in the formulas, and it reduces the lower
bound to a Bayes risk over batched policies, which Proposition~\ref{prop:K2-relax} bounds.
\begin{proof}
\uses{prop:four, def:risk, lem:K2-gauss, def:batched}
\textit{Setting.} Fix $M>d$, let $H_M:=[-M,M]^2$, and fix a batched policy $\pi$. \textit{Goal.} Show that
$\sup_{h\in H_M}\rk(\pi,h)\ge\kBb(\kPri)-2d\,\P(|m|>M-d)$.

On the support of $\kPri$ the gap is $2d$, so $0\le\rk(\pi,h)\le2d$. If $|m|\le M-d$ then $h\in H_M$. Hence
\begin{align*}
\kBb(\kPri)\le\int\rk(\pi,h)\,\kPri(dh)
&\le\sup_{h\in H_M}\rk(\pi,h)\,\kPri(H_M)+2d\,\kPri(H_M^c) && \text{(split; $\rk\le2d$)}\\
&\le\sup_{h\in H_M}\rk(\pi,h)+2d\,\P(|m|>M-d) && \text{($H_M^c\subset\{|m|>M-d\}$)}.
\end{align*}
Taking the infimum over $\pi$ gives $\Vb(H_M)\ge\kBb(\kPri)-2d\,\P(|m|>M-d)$. Letting $M\to\infty$ and using $\VK_2=\lim_M\Vb(H_M)$ (Proposition~\ref{prop:four}(b)) gives
the claim.
\end{proof}

\subsubsection{Why does revealing blocks in advance give a lower bound?}

\begin{definition}[block-revealing recursion]\label{def:K2-W}\dt{k:W}
Let $\lambda,d>0$, let $n\ge2$ be an integer and $\varepsilon:=1/n$. For integers $a,b\ge1$ let $o_{ab}:=(a\varepsilon,b\varepsilon)$, $I_{ab}:=\kI(o_{ab})$, and
\[
\big(s^{(1)}_{ab}\big)^2:=I_{a+1,b}-I_{ab},\qquad \big(s^{(2)}_{ab}\big)^2:=I_{a,b+1}-I_{ab},
\]
which lie in $(0,\varepsilon]$ by the information clock. Define $\kW_{ab}:\R\to[0,\infty)$ for $a+b\le n+1$ by $\kW_{ab}:=0$ if $a+b=n+1$, and otherwise
\[
\kW_{ab}(u):=\min_{l\in\{1,2\}}\Big\{\varepsilon\,\kg_l(u,I_{ab})+\E\big[\kW_{(a,b)+e_l}\big(u+s^{(l)}_{ab}Z\big)\big]\Big\}.
\]
Let $\kL(d,\lambda):=\E\big[\kW_{11}\big(\sqrt{I_{11}}\,Z\big)\big]$, and let $\kT_{ab}:=(n+1-a-b)\,\varepsilon$ be the time left in the recursion.
\end{definition}

\noindent This is the recursion of Definition~\ref{def:K2-rec} with $\lambda>0$ and two changes: the observation point is one block \emph{ahead} on both arms, and
there is no initial cost $2d\varepsilon$. Equivalently, it is the recursion of Definition~\ref{def:K2-rec} run up to the diagonal $t=1+\varepsilon$ and started at
$(\varepsilon,\varepsilon)$ for free.

\begin{proposition}[block-revealing relaxation]\label{prop:K2-relax}
\uses{def:risk, def:K2-W}
Let $\lambda,d>0$ and $n\ge2$. Then every batched policy $\pi$ satisfies $\int\rk(\pi,h)\,\kPri(dh)\ge\kL(d,\lambda)$. Hence $\kBb(\kPri)\ge\kL(d,\lambda)$ and, by
Lemma~\ref{lem:K2-box}, $\VK_2\ge\kL(d,\lambda)$.
\end{proposition}
\noindent\textit{Role.} It turns the lower bound into a recursion on a lattice, which a computer can bound from below (Section~\ref{sec:K2-scheme}).

\noindent\textit{Idea.} A batched policy moves in the plane of observation times $(\tau_1,\tau_2)$ along a monotone path. Cut each axis into blocks of length
$\varepsilon$. When the path completes a block of arm~$1$, its other coordinate lies in some block $j$. Everything the policy has used so far is contained in the
information at the point $o:=$ (end of this block, end of block $j$). The expected regret of the completed block can therefore be computed with the
likelihood ratio at $o$, by~\eqref{eq:K2-cost}. The completed blocks form a lattice path driven by an information at least as large as the policy's; the lattice
game with this larger information is exactly the recursion $\kW$.

\begin{proof}
\uses{lem:K2-gauss, def:risk, def:K2-W}
\textit{Setting.} Fix a batched policy $\pi$ with decision times $0=t_0<\dots<t_q=1$. Let $\alpha^{(i)}\in[0,1]$ be the proportion of $[t_i,t_{i+1}]$ given to arm~$1$,
and $\tau^{(i)}$ the observation point at time $t_i$, so that $\tau^{(0)}=0$ and $\tau^{(i+1)}=\tau^{(i)}+(t_{i+1}-t_i)(\alpha^{(i)},1-\alpha^{(i)})$. The order of the pulls inside a batch
changes neither the observations at the decision times nor the regret, so we may assume that each batch samples the arms at constant rates. Let $\gamma$ be the
polygonal path through $\tau^{(0)},\dots,\tau^{(q)}$: a continuous monotone path in $[0,1]^2$ from $0$ to the segment $\{o_1+o_2=1\}$. The data are drawn under
$\kP{\pm d}$. \textit{Goal.} Show that $\int\rk(\pi,h)\,\kPri(dh)\ge\kL(d,\lambda)$.

\textit{Step 1: the lattice path of completed blocks.} For $i\ge0$, block $i$ of arm~$1$ is $[i\varepsilon,(i+1)\varepsilon]$; it is \emph{completed} if $\gamma$ reaches
$\{\tau_1=(i+1)\varepsilon\}$. Let $R_{ij}$ (``right at $(i,j)$'') be the event that it is completed and that the second coordinate of the first such point lies in
$[j\varepsilon,(j+1)\varepsilon)$. Symmetrically, $A_{ij}$ (``up at $(i,j)$'') is the event that block $j$ of arm~$2$ is completed while the first coordinate lies in
$[i\varepsilon,(i+1)\varepsilon)$. If $\gamma$ passes exactly through a corner $((i+1)\varepsilon,(j+1)\varepsilon)$ we count $R_{ij}$, then $A_{i+1,j}$. The completed blocks, in
the order in which they are completed, form a lattice path $(0,0)=p_0,p_1,\dots$ with unit steps: from $p=(i,j)$ the step is to the right on $R_{ij}$ and upward on
$A_{ij}$. At the end $\tau_1+\tau_2=1=n\varepsilon$, so at least $\lfloor\tau_1/\varepsilon\rfloor+\lfloor\tau_2/\varepsilon\rfloor\ge n-1$ blocks are completed: the path has at least
$n-1$ steps.

\textit{Step 2: the steps are measurable with respect to the revealed information.} Let $o:=o_{i+1,j+1}$ and $B_o:=[0,o_1]\times[0,o_2]$. We show that $R_{ij}$, $A_{ij}$
and the event $V_{ij}$ that the lattice path visits $(i,j)$ belong to $\kF_o$. We argue by induction on $l$. On the event that $\tau^{(0)},\dots,\tau^{(l)}\in B_o$, the
proportion $\alpha^{(l)}$ is a Borel function of $\zeta$ and of the paths $Y_k$ on $[0,\tau^{(l')}_k]$, $l'\le l$. These are $\kF_o$-measurable because
$\tau^{(l')}\in B_o$ and the paths are continuous, so that the restriction of a path to a random interval is jointly measurable. Hence $\tau^{(l+1)}$ and the event
$\{\tau^{(l+1)}\in B_o\}$ are $\kF_o$-measurable on this event. So the part of $\gamma$ inside $B_o$, up to its first exit, is $\kF_o$-measurable. Since $\gamma$ is
monotone, the events $R_{ij}$, $A_{ij}$ and $V_{ij}$ are determined by this part of $\gamma$.

\textit{Step 3: cost of the completed blocks.} The time spent on arm~$1$ is at least $\varepsilon$ times the number of completed blocks of arm~$1$, and each such block
corresponds to exactly one right step; the same holds for arm~$2$ and up steps. Let $\iota_k\in\{1,2\}$ be the direction of the step leaving $p_k$ ($1$ for right),
and $o_{p+(1,1)}:=o_{i+1,j+1}$ for $p=(i,j)$. Then
\begin{align*}
\int\rk(\pi,h)\,\kPri(dh)
&=\tfrac12\,\E_{\kP{-d}}\big[2d\,\tau_1^{(q)}\big]+\tfrac12\,\E_{\kP{+d}}\big[2d\,\tau_2^{(q)}\big] && \text{(the two hypotheses)}\\
&\ge\sum_{i,j}\varepsilon\Big(d\,\kP{-d}(R_{ij})+d\,\kP{+d}(A_{ij})\Big) && \text{(completed blocks only)}\\
&=\E_{\kQ}\Big[\sum_{k\ge0}\varepsilon\,\kg_{\iota_k}\big(\kU(o_{p_k+(1,1)}),\kI(o_{p_k+(1,1)})\big)\Big] && \text{(by~\eqref{eq:K2-cost} and Step~2)}.
\end{align*}
The last line uses $R_{ij},A_{ij}\in\kF_{o_{i+1,j+1}}$, proved in Step~2.
Dropping the steps after the $(n-1)$-th, whose terms are non-negative, we keep exactly $n-1$ steps.

\textit{Step 4: comparison with the recursion.} Write $(a,b):=p+(1,1)$. The path visits states $(a,b)$ on the diagonals $a+b=2,\dots,n$, and its step from $(a,b)$ is
$\kF_{o_{ab}}$-measurable (Step~2). We show by backward induction on $a+b$ that, on the event that the path visits $(a,b)$,
\[
\E_{\kQ}\Big[\text{cost of the steps from $(a,b)$ up to the diagonal $n+1$}\ \Big|\ \kF_{o_{ab}}\Big]\ \ge\ \kW_{ab}\big(\kU(o_{ab})\big).
\]
On the diagonal $a+b=n+1$ both sides vanish. Assume the claim on the diagonal $a+b+1$, and let $R$ be the event that the step from $(a,b)$ goes right. On $R$ and the visit of
$(a,b)$, which are $\kF_{o_{ab}}$-measurable events contained in the visit of $(a+1,b)$, the induction hypothesis and the tower property
($\kF_{o_{ab}}\subset\kF_{o_{a+1,b}}$) show that the conditional cost is at least
\[
\varepsilon\,\kg_1\big(\kU(o_{ab}),I_{ab}\big)+\E_{\kQ}\big[\kW_{a+1,b}\big(\kU(o_{a+1,b})\big)\,\big|\,\kF_{o_{ab}}\big].
\]
By Lemma~\ref{lem:K2-gauss}(ii), the conditional expectation equals $\E[\kW_{a+1,b}(u+s^{(1)}_{ab}Z)]$ at $u=\kU(o_{ab})$. The same holds upward. The conditional cost is
therefore at least the minimum of the two options, which is $\kW_{ab}(\kU(o_{ab}))$. The state $(1,1)$ is visited surely, and $\kU(o_{11})\sim\mathcal N(0,I_{11})$ under
$\kQ$ (Lemma~\ref{lem:K2-gauss}(ii) with $o=0$). So the total cost is at least $\E[\kW_{11}(\sqrt{I_{11}}Z)]=\kL(d,\lambda)$.
\end{proof}

\begin{remark}[known common mean]\label{rem:K2-known}
Revealing the common mean $m$ to the policy can only lower its Bayes risk. Once $m$ is known, the information on $v$ at time $t$ is $\tau_1+\tau_2=t$ whatever the
allocation, so the probability of pulling the worse arm at time $t$ is at least the Bayes error $\Phi(-d\sqrt t)$. Hence
\[
\kBb(\kPri)\ \ge\ \int_0^12d\,\Phi(-d\sqrt t)\,dt,
\]
whose maximum over $d$ is $0.5298$, at $d\approx1.25$. This is the elementary lower bound. Proposition~\ref{prop:K2-relax} improves it because it keeps the cost of learning the common mean.
\end{remark}

\subsubsection{How is the recursion bounded from below by a finite computation?}\label{sec:K2-scheme}

\begin{lemma}[an explicit majorant]\label{lem:K2-upper}
\uses{def:K2-W}
For all $(a,b)$ with $a+b\le n+1$ and all $u\in\R$,
\[
\kW_{ab}(u)\le\kT_{ab}\min_{l}\kg_l(u,I_{ab})=\kT_{ab}\,d\,e^{-2d|u|-2d^2I_{ab}} .
\]
\end{lemma}
\noindent\textit{Role.} It bounds the supremum of $\kW_{ab}$ on a grid cell in the interpolation error of Lemma~\ref{lem:K2-semi}.
\begin{proof}
\uses{def:K2-W}
Always choosing the same option $l$ is admissible in the recursion. Its cost is a sum of $\kT_{ab}/\varepsilon$ terms $\varepsilon\,\E[\kg_l(u+\text{Gaussian increments},I)]$. Each
of them equals $\varepsilon\,\kg_l(u,I_{ab})$, because $\E[\kg_l(u+sZ,I+s^2)]=\kg_l(u,I)$ (exponential martingale: $\E e^{\mp2dsZ}=e^{2d^2s^2}$).
\end{proof}

\begin{lemma}[semiconcavity]\label{lem:K2-semi}
\uses{def:K2-W}
For every $(a,b)$, the function $\kW_{ab}$ is continuous and, at every $u\in\R$,
\[
\limsup_{\eta\to0}\frac{\kW_{ab}(u+\eta)+\kW_{ab}(u-\eta)-2\kW_{ab}(u)}{\eta^2}\ \le\ 4d^2\,\kW_{ab}(u).
\]
Consequently, if $x_0<x_1$ and $x_1-x_0=h$, then for every $x\in[x_0,x_1]$
\[
\kW_{ab}(x)\ \ge\ \big(\text{chord of }\kW_{ab}\text{ between }x_0\text{ and }x_1\big)(x)-\frac{d^2h^2}{2}\sup_{[x_0,x_1]}\kW_{ab}.
\]
\end{lemma}
\noindent\textit{Role.} The error of linear interpolation is one-sided and proportional to $\kW$ itself: this is what makes a rigorous minorant cheap.

\noindent\textit{Idea.} The operator $\partial_u^2-4d^2$ kills the cost rates $\kg_l\propto e^{\mp2du}$ and commutes with the Gaussian convolution, so each option is a
supersolution, and a minimum of supersolutions is a supersolution.

\begin{proof}
\uses{lem:K2-upper, def:K2-W}
\textit{Step 1: each option satisfies $A''\le4d^2A$.} Let $\mathcal D:=\partial_u^2-4d^2$. Assume, by backward induction (true on the last diagonal, where $\kW=0$), that the
two next functions $\kW'$ are continuous, bounded by multiples of $e^{-2d|u|}$ (Lemma~\ref{lem:K2-upper}), and satisfy $\mathcal D\kW'\le0$ in the sense of distributions.
Each option $A_l:=\varepsilon\kg_l+\E[\kW'(\cdot+sZ)]$ is $C^\infty$, as a Gaussian convolution with $s>0$. Since $\mathcal D\kg_l=0$ and $\mathcal D$ commutes with the
convolution, $\mathcal DA_l=\E[(\mathcal D\kW')(\cdot+sZ)]\le0$, that is, $A_l''\le4d^2A_l$ pointwise.

\textit{Step 2: the minimum.} Fix $u$ and an option $l$ with $\kW_{ab}(u)=A_l(u)$. Since $\kW_{ab}\le A_l$ everywhere, for some $\xi\in(u-\eta,u+\eta)$,
\begin{align*}
\kW_{ab}(u+\eta)+\kW_{ab}(u-\eta)-2\kW_{ab}(u)
&\le A_l(u+\eta)+A_l(u-\eta)-2A_l(u) && \text{($\kW_{ab}\le A_l$, equality at $u$)}\\
&=A_l''(\xi)\,\eta^2\ \le\ 4d^2A_l(\xi)\,\eta^2 && \text{(Taylor; Step~1)},
\end{align*}
and $A_l(\xi)\to\kW_{ab}(u)$ as $\eta\to0$. This is the claimed bound. The same display shows that the second difference quotients of $\kW_{ab}$ are bounded above, locally
uniformly, by $4d^2\max_l\sup A_l$ over a neighbourhood. For a test function $\phi\ge0$, moving the second difference onto $\phi$ and applying the reverse Fatou lemma
gives $\int\kW_{ab}\phi''\le\int4d^2\kW_{ab}\phi$, that is, $\mathcal D\kW_{ab}\le0$ in the sense of distributions, which closes the induction.

\textit{Step 3: interpolation.} Let $c:=4d^2\sup_{[x_0,x_1]}\kW_{ab}$. By Step~2 the continuous function $f:=\kW_{ab}-cx^2/2$ has an upper second symmetric derivative
$\le0$ on $(x_0,x_1)$, hence also a lower one. By Schwarz's theorem (a continuous function whose lower second symmetric derivative is $\le0$ on an open interval is
concave there; Zygmund, \emph{Trigonometric Series}, Vol.~I, Ch.~IX), $f$ is concave on $[x_0,x_1]$ and lies above its chord. Hence
\[
\kW_{ab}(x)\ \ge\ \text{chord}(x)-\frac c2(x-x_0)(x_1-x)\ \ge\ \text{chord}(x)-\frac{c\,h^2}{8},
\]
and $c h^2/8=\frac{d^2h^2}2\sup_{[x_0,x_1]}\kW_{ab}$.
\end{proof}

\paragraph{Hat functions and their Gaussian expectations.}\dt{k:ker} Fix a grid step $h>0$, a half-width $u_{\max}$ with $N:=2u_{\max}/h\in\Nat$, and the nodes
$u_j:=-u_{\max}+jh$, $j=0,\dots,N$. Let $\kchi(x):=\max(0,1-|x|/h)$ and $\kchi_j(x):=\kchi(x-u_j)$, the hat function of node $j$. For $s>0$ let
$G_s(y):=\E\big[(y+sZ)^+\big]=y\,\Phi(y/s)+s\,\varphi(y/s)$ and
\[
\kGam_s(x):=\E\big[\kchi(x+sZ)\big]=\frac{G_s(x+h)-2G_s(x)+G_s(x-h)}{h}.
\]
The closed form holds because $h\,\kchi(y)=(y+h)^+-2y^++(y-h)^+$. So the Gaussian expectation of a piecewise linear function with nodes $u_j$ is explicit:
$\E[\kchi_j(u_i+sZ)]=\kGam_s(u_i-u_j)$. Moreover $\kGam_s\ge0$ and $\sum_{j=0}^N\kGam_s(u_i-u_j)=\E\big[\sum_j\kchi_j(u_i+sZ)\big]\le1$, since $\sum_j\kchi_j\le1$.

\begin{definition}[step map of the scheme]\label{def:K2-step}\dt{k:step}
\uses{def:K2-W}
Fix $(a,b)$ with $a+b\le n$, and let
\[
e^{ab}_i:=\frac{d^3h^2}{2}\,\kT_{ab}\,e^{-2d^2I_{ab}}\max_{C\ni u_i}e^{-2d\,\mathrm{dist}(0,C)},\qquad i=0,\dots,N,
\]
where $C$ ranges over the grid cells $[u_j,u_{j+1}]$ that contain $u_i$. For $c',c''\in\R^{N+1}$ let
\[
y_i(c',c''):=\min\Big\{\varepsilon\,\kg_1(u_i,I_{ab})+\sum_{j=0}^Nc'_j\,\kGam_{s^{(1)}_{ab}}(u_i-u_j),\ \ \varepsilon\,\kg_2(u_i,I_{ab})+\sum_{j=0}^Nc''_j\,\kGam_{s^{(2)}_{ab}}(u_i-u_j)\Big\},
\]
and define $\kS_{ab}(c',c'')\in\R^{N+1}$ by
\[
\kS_{ab}(c',c'')_i:=y_i(c',c'')-e^{ab}_i\quad(0<i<N),\qquad \kS_{ab}(c',c'')_i:=\min\big(0,\,y_i(c',c'')-e^{ab}_i\big)\quad(i\in\{0,N\}).
\]
\end{definition}

\noindent By Lemmas~\ref{lem:K2-upper} and~\ref{lem:K2-semi}, $e^{ab}_i$ bounds the interpolation error of $\kW_{ab}$ on each cell adjacent to $u_i$. The minimum with $0$ at the
two boundary nodes is what makes the piecewise linear function vanish, or be negative, outside $[-u_{\max},u_{\max}]$, where $\kW_{ab}\ge0$ is all we know.

\begin{definition}[algorithm: Certified minorant of $\kL(d,\lambda)$]\label{def:K2-inf}\dt{k:alg}
\leavevmode\par\smallskip\noindent\fbox{\begin{minipage}{0.96\linewidth}\small
\noindent\hspace*{0.0em}\parbox[t]{\dimexpr\linewidth-0.0em\relax}{\raggedright \textbf{Input:} $d,\lambda>0$, $n\ge2$ ($\varepsilon:=1/n$), grid $u_j=-u_{\max}+jh$, $j=0,\dots,N$, margin $\kdfp\ge0$}\par
\noindent\hspace*{0.0em}\parbox[t]{\dimexpr\linewidth-0.0em\relax}{\raggedright $c^{(ab)}\gets0$ for $a+b=n+1$}\par
\noindent\hspace*{0.0em}\parbox[t]{\dimexpr\linewidth-0.0em\relax}{\raggedright \textbf{for} $k=n,n-1,\dots,2$ \textbf{and} each $(a,b)$ with $a+b=k$ \textbf{do}}\par
\noindent\hspace*{1.3em}\parbox[t]{\dimexpr\linewidth-1.3em\relax}{\raggedright $y_i\gets y_i\big(c^{(a+1,b)},c^{(a,b+1)}\big)-\kdfp$ for $i=0,\dots,N$ \hfill(exact Gaussian expectations)}\par
\noindent\hspace*{1.3em}\parbox[t]{\dimexpr\linewidth-1.3em\relax}{\raggedright $c^{(ab)}_i\gets y_i-e^{ab}_i$ for $0<i<N$; \quad $c^{(ab)}_i\gets\min\big(0,y_i-e^{ab}_i\big)$ for $i\in\{0,N\}$}\par
\noindent\hspace*{0.0em}\parbox[t]{\dimexpr\linewidth-0.0em\relax}{\raggedright }\par
\noindent\hspace*{0.0em}\parbox[t]{\dimexpr\linewidth-0.0em\relax}{\raggedright \textbf{return} $\sum_{j=0}^Nc^{(11)}_j\,\kGam_{\sqrt{I_{11}}}(u_j)-\kdfp$}\par
\end{minipage}}
\end{definition}


\noindent\textit{Implementation.} The sums over $j$ are convolutions, computed by FFT with zero padding to a length at least $N+1+m_{\max}$, so that there is no
wrap-around; the kernel $\kGam_s$ is truncated to $|x|\le h+8.5s$, and $m_{\max}h$ is the largest truncation radius. The truncation and the rounding errors are
counted in Remark~\ref{rem:K2-fp}. The code sets $c^{(ab)}_0=c^{(ab)}_N:=0$ instead of taking the minimum with $0$; Proposition~\ref{prop:K2-Lnum} explains why this is
the same in the runs.

\begin{lemma}[what a certificate is; errors do not accumulate]\label{lem:K2-step}
\uses{def:K2-step, def:K2-W}
\begin{enumerate}
\item[(a)] (monotonicity) If $c'\le\tilde c'$ and $c''\le\tilde c''$ componentwise, then $\kS_{ab}(c',c'')\le\kS_{ab}(\tilde c',\tilde c'')$. If $D\ge0$, then
$\kS_{ab}(c'-D\1,c''-D\1)\ge\kS_{ab}(c',c'')-D\1$.
\item[(b)] (certificate) Let vectors $\hat c^{(ab)}\in\R^{N+1}$, $a,b\ge1$, $a+b\le n+1$, satisfy $\hat c^{(ab)}=0$ if $a+b=n+1$ and
\[
\hat c^{(ab)}\le\kS_{ab}\big(\hat c^{(a+1,b)},\hat c^{(a,b+1)}\big)\qquad\text{componentwise, if }a+b\le n.
\]
Then $\sum_j\hat c^{(ab)}_j\kchi_j\le\kW_{ab}$ on $\R$ for every $(a,b)$, and $\kL(d,\lambda)\ge\sum_j\hat c^{(11)}_j\kGam_{\sqrt{I_{11}}}(u_j)$.
\end{enumerate}
\end{lemma}
\noindent\textit{Role.} Item~(b) says what a computation must deliver: each computed vector is at most the exact image of the previously computed ones. Item~(a) shows
that this is a condition on each step separately, so that errors do not accumulate; it also gives the cost of a larger margin.
\begin{proof}
\uses{lem:K2-upper, lem:K2-semi, def:K2-W, def:K2-step}
\textit{Step 1: we prove (a).} The coefficients $\kGam_s(u_i-u_j)$ are non-negative, and the minimum and $x\mapsto\min(0,x)$ are non-decreasing: this gives monotonicity.
Lowering $c'$ and $c''$ by $D$ lowers each sum by $D\sum_j\kGam_s(u_i-u_j)\le D$, hence lowers $y_i$ by at most $D$; the map $x\mapsto\min(0,x)$ is $1$-Lipschitz.

\textit{Step 2: nodal minorants in (b).} Assume $\sum_j\hat c'_j\kchi_j\le\kW_{a+1,b}$ and $\sum_j\hat c''_j\kchi_j\le\kW_{a,b+1}$ on $\R$, for the two next vectors. The
Gaussian expectation is monotone, so
\[
\E\big[\kW_{a+1,b}(u_i+s^{(1)}_{ab}Z)\big]\ \ge\ \sum_j\hat c'_j\,\E\big[\kchi_j(u_i+s^{(1)}_{ab}Z)\big]=\sum_j\hat c'_j\,\kGam_{s^{(1)}_{ab}}(u_i-u_j),
\]
and similarly for arm~$2$. Taking the minimum of the two options gives $y_i(\hat c',\hat c'')\le\kW_{ab}(u_i)$ for every node.

\textit{Step 3: from the nodes to the line.} Let $C=[u_j,u_{j+1}]$ be a cell and $e_C:=\frac{d^3h^2}2\kT_{ab}e^{-2d^2I_{ab}}e^{-2d\,\mathrm{dist}(0,C)}$. By
Lemma~\ref{lem:K2-upper}, $\sup_C\kW_{ab}\le\kT_{ab}\,d\,e^{-2d^2I_{ab}}e^{-2d\,\mathrm{dist}(0,C)}$, so Lemma~\ref{lem:K2-semi} gives $\kW_{ab}\ge\text{chord}_C(\kW_{ab})-e_C$ on
$C$. By definition $e^{ab}_j,e^{ab}_{j+1}\ge e_C$. Hence, on $C$,
\begin{align*}
\kW_{ab}&\ge\text{chord}_C\big(\kW_{ab}(u_j),\kW_{ab}(u_{j+1})\big)-e_C && \text{(Lemmas~\ref{lem:K2-upper} and~\ref{lem:K2-semi})}\\
&\ge\text{chord}_C\big(y_j-e^{ab}_j,\ y_{j+1}-e^{ab}_{j+1}\big) && \text{(Step~2; $e^{ab}_j,e^{ab}_{j+1}\ge e_C$)}\\
&\ge\text{chord}_C\big(\hat c^{(ab)}_j,\hat c^{(ab)}_{j+1}\big)=\sum_i\hat c^{(ab)}_i\kchi_i && \text{($\hat c^{(ab)}\le\kS_{ab}\le y-e^{ab}$ at every node)}.
\end{align*}
On $[u_0-h,u_0]$ the function $\sum_i\hat c^{(ab)}_i\kchi_i$ equals $\hat c^{(ab)}_0\kchi_0\le0\le\kW_{ab}$, because $\hat c^{(ab)}_0\le\kS_{ab}(\cdot)_0\le0$; the same holds on
$[u_N,u_N+h]$, and the function vanishes outside $[u_0-h,u_N+h]$. So $\sum_i\hat c^{(ab)}_i\kchi_i\le\kW_{ab}$ on $\R$. Backward induction from the last diagonal, where both
sides vanish, proves the first claim of~(b).

\textit{Step 4: the root.} By Definition~\ref{def:K2-W} and the first claim,
\[
\kL(d,\lambda)=\E\big[\kW_{11}(\sqrt{I_{11}}Z)\big]\ \ge\ \sum_j\hat c^{(11)}_j\,\E\big[\kchi_j(\sqrt{I_{11}}Z)\big]=\sum_j\hat c^{(11)}_j\,\kGam_{\sqrt{I_{11}}}(u_j),
\]
using that $\kGam_s$ is even.
\end{proof}

\begin{proposition}[certified value of the lower bound]\label{prop:K2-Lnum}
\uses{def:K2-inf, def:K2-W, def:risk}
\begin{enumerate}
\item[(i)] (exact arithmetic) For every margin $\kdfp\ge0$, the output of Definition~\ref{def:K2-inf}, computed in exact real arithmetic, is at most $\kL(d,\lambda)$.
\item[(ii)] (the runs) Let $\lambda=10^{-6}$, $h=10^{-3}$, $u_{\max}=3.5$ and $\kdfp:=10^{-8}$. The double-precision runs of Definition~\ref{def:K2-inf} with the margin $10^{-10}$
returned $\hat L=0.6317025746$ for $(d,n)=(1.56,400)$ and $\hat L=0.6341794644$ for $(d,n)=(1.565,800)$. Assume that every computed nodal value, and the final sum,
differ from the exact value of the same expression, evaluated on the previously computed vectors, by at most $\kdfp$ (Remark~\ref{rem:K2-fp}). Then
\[
\kL(d,\lambda)\ \ge\ \hat L-n\big(\kdfp-10^{-10}\big),
\]
that is, $\kL(1.56,10^{-6})\ge0.631698$ at $\varepsilon=1/400$ and $\kL(1.565,10^{-6})\ge0.634171$ at $\varepsilon=1/800$. Hence, by Proposition~\ref{prop:K2-relax},
$\VK_2\ge0.6341$.
\end{enumerate}
\end{proposition}
\noindent\textit{Role.} It is the numerical half of the lower bound of Proposition~\ref{prop:K2}.

\noindent\textit{Idea.} By Lemma~\ref{lem:K2-step}(b), it suffices that each computed vector be at most the exact image of the previous ones. In exact arithmetic the
margin only lowers the vectors. In floating point, lowering the computed vectors by a constant per diagonal absorbs the rounding errors of each step, and
Lemma~\ref{lem:K2-step}(a) shows that the constants add up linearly: $n$ steps of $\kdfp-10^{-10}$.

\begin{proof}
\uses{def:K2-inf, lem:K2-step, prop:K2-relax, def:K2-step, def:K2-W, def:risk}
\textit{Step 1: we prove (i).} In exact arithmetic, the vectors of Definition~\ref{def:K2-inf} satisfy, for $0<i<N$,
$c^{(ab)}_i=\kS_{ab}(c^{(a+1,b)},c^{(a,b+1)})_i-\kdfp$, and, for $i\in\{0,N\}$, $c^{(ab)}_i=\min(0,y_i-\kdfp-e^{ab}_i)\le\kS_{ab}(\cdot)_i$. So they satisfy the hypothesis of
Lemma~\ref{lem:K2-step}(b), and the output is $\sum_jc^{(11)}_j\kGam_{\sqrt{I_{11}}}(u_j)-\kdfp\le\kL(d,\lambda)$.

\textit{Step 2: the boundary nodes in the runs.} The computed values $y_i-e^{ab}_i$ at the two boundary nodes, margin $10^{-10}$ included, were checked to be positive on
every diagonal: their minimum $\beta$ is $2.08\cdot10^{-8}$ at $n=400$ and $9.94\cdot10^{-9}$ at $n=800$. So the minimum with $0$ is not active, and the code's
assignment $c_0=c_N=0$ coincides with Definition~\ref{def:K2-inf}.

\textit{Step 3: shifted vectors are a certificate.} Let $\hat c^{(ab)}$ be the vectors computed by the run with the margin $10^{-10}$ (with $\hat c^{(ab)}_0=\hat c^{(ab)}_N=0$),
and let $\phi\le\kdfp$ bound the errors of the assumption. Let $\eta:=\kdfp-10^{-10}$, $p_{ab}:=n+1-a-b$ (the number of diagonals already processed when
$(a,b)$ is), and
\[
\tilde c^{(ab)}:=\hat c^{(ab)}-p_{ab}\,\eta\,\1 .
\]
We check $\tilde c^{(ab)}\le\kS_{ab}(\tilde c^{(a+1,b)},\tilde c^{(a,b+1)})$; the next vectors are shifted by $(p_{ab}-1)\eta$. Let
$x_i:=\kS_{ab}(\hat c^{(a+1,b)},\hat c^{(a,b+1)})_i$ and $\tilde x_i:=\kS_{ab}(\tilde c^{(a+1,b)},\tilde c^{(a,b+1)})_i$. For an interior node $i$,
\begin{align*}
\tilde c^{(ab)}_i
&\le x_i-10^{-10}+\phi-p_{ab}\eta && \text{(rounding error $\le\phi$; run margin $10^{-10}$)}\\
&\le\tilde x_i+(p_{ab}-1)\eta-10^{-10}+\phi-p_{ab}\eta && \text{(Lemma~\ref{lem:K2-step}(a), $D=(p_{ab}-1)\eta$)}\\
&=\tilde x_i+\phi-\kdfp\ \le\ \tilde x_i && \text{($\eta+10^{-10}=\kdfp$, $\phi\le\kdfp$)}.
\end{align*}
For a boundary node, $\tilde c^{(ab)}_i=-p_{ab}\eta\le0$. By Step~2, the exact value of $y_i-e^{ab}_i$ computed from the vectors $\hat c$ is at least $\beta+10^{-10}-\phi$. Hence
\begin{align*}
\tilde x_i
&\ge\min\big(0,\ \beta+10^{-10}-\phi-(p_{ab}-1)\eta\big) && \text{(Lemma~\ref{lem:K2-step}(a))}\\
&\ge-p_{ab}\eta=\tilde c^{(ab)}_i && \text{($\beta+\kdfp\ge\phi$)}.
\end{align*}
On the last diagonal $\tilde c=\hat c=0$. So the vectors $\tilde c$ satisfy the hypothesis of Lemma~\ref{lem:K2-step}(b).

\textit{Step 4: the root.} The run returned $\hat L=\hat\Sigma-10^{-10}$, where $\hat\Sigma$ is the computed value of $\sum_j\hat c^{(11)}_j\kGam_{\sqrt{I_{11}}}(u_j)$.
Hence
\begin{align*}
\kL(d,\lambda)
&\ge\sum_j\tilde c^{(11)}_j\,\kGam_{\sqrt{I_{11}}}(u_j) && \text{(Lemma~\ref{lem:K2-step}(b))}\\
&\ge\sum_j\hat c^{(11)}_j\,\kGam_{\sqrt{I_{11}}}(u_j)-(n-1)\eta && \text{($p_{11}=n-1$; $\textstyle\sum_j\kGam(u_j)\le1$)}\\
&\ge\hat\Sigma-\phi-(n-1)\eta=\hat L+10^{-10}-\phi-(n-1)\eta && \text{(rounding of the final sum)}\\
&\ge\hat L-n\eta && \text{($\eta+10^{-10}=\kdfp\ge\phi$)}.
\end{align*}
With $n\eta=800\cdot9.9\cdot10^{-9}=7.92\cdot10^{-6}$ this gives $\kL\ge0.6341794644-0.0000079200\ge0.634171$ at $n=800$, and $\kL\ge0.6317025746-0.0000039600\ge0.631698$ at
$n=400$. With Proposition~\ref{prop:K2-relax}, $\VK_2\ge0.634171\ge0.6341$.
\end{proof}

\begin{remark}[floating-point errors: an a priori analysis]\label{rem:K2-fp}
By Lemma~\ref{lem:K2-step}, the run must satisfy a condition on each step separately: every computed entry within $\kdfp=10^{-8}$ of its exact value given the previous
computed vectors. The sources of error, per entry, are the following.
\begin{itemize}
\item \emph{FFT convolution.} Zero padding to the length $8192\ge N+1+m_{\max}$ ($N+1=7001$; $m_{\max}=427$ at $n=400$ and $302$ at $n=800$) excludes wrap-around. For an
FFT of length $2^{13}$ the relative error in $\ell^2$ norm is of order $\log_2(2^{13})\,u\approx2\cdot10^{-15}$, with $u=2^{-53}$ (Higham, \emph{Accuracy and Stability of
Numerical Algorithms}, 2nd ed., 2002, Thm~24.2). That theorem is stated for a radix-$2$ FFT with precomputed twiddle factors, not for the \texttt{pocketfft} code used by
SciPy, so we use it up to a constant. Here $\|c\|_2\le\sqrt{N+1}\max|c|$ with $\max|c|\le d\le1.6$ (Lemma~\ref{lem:K2-upper}), and $\|\kGam\|_2\le\|\kGam\|_1\le1$. Including
the cross term (error on the transform of the kernel times $\|c\|_1$), each entry is computed within about $4\cdot10^{-12}$. A direct convolution of one row of length
$7001$ differs from the FFT by $4\cdot10^{-16}$.
\item \emph{Kernel entries.} The closed form of $\kGam_s$ is a second difference of $G_s$ divided by $h$, with cancellation. Against quadrature, the error per entry is at
most $1.4\cdot10^{-13}$, and the sum of the errors over the kernel is $2.3\cdot10^{-11}$ ($s=0.05$, $n=400$) and $1.2\cdot10^{-11}$ ($s=0.035$, $n=800$). Through the
convolution this gives at most $2.3\cdot10^{-11}\max|c|\le4\cdot10^{-11}$ per entry. Negative rounding values of the kernel are set to $0$, a change of the same size.
\item \emph{Clock increments.} The difference $(s^{(l)}_{ab})^2=I_{a+1,b}-I_{ab}$ suffers cancellation, with relative error up to about $10^{-11}$ for $b=1$; the effect on the
convolution is below $10^{-12}$.
\item \emph{Elementary functions and grid.} \texttt{exp}, \texttt{scipy.special.ndtr} and the floating-point nodes $u_j$ are accurate to a few units in the last place; the
system library gives no formal guarantee. The effect is below $10^{-14}$.
\item \emph{Truncation of the kernel} to $|x|\le h+8.5s$. The omitted terms with $c_j\ge0$ only lower the result. The omitted terms with $c_j<0$ have total weight at most
$\P(|Z|>8.5)<2\cdot10^{-17}$, and a priori $c_j\ge-n(\max e+\kdfp)\ge-2\cdot10^{-3}$, using $\max e\le d^3h^2/2\le1.92\cdot10^{-6}$. In the runs no nodal value was negative.
\end{itemize}
The total is below $5\cdot10^{-11}$ per entry, $200$ times smaller than $\kdfp=10^{-8}$; the independent verification estimated the worst case at $2$ to $4\cdot10^{-11}$ per
step. The runs used the margin $10^{-10}$, only $2$ to $5$ times this estimate, which is why Proposition~\ref{prop:K2-Lnum} raises it to $10^{-8}$; the price is
$n(10^{-8}-10^{-10})<8\cdot10^{-6}$, below the rounding margin $7.9\cdot10^{-5}$ between $0.634179$ and $0.6341$. This is an a priori analysis, not an interval
computation. A machine-checkable certificate would rerun Definition~\ref{def:K2-inf} in interval arithmetic (for instance \texttt{arb}) or with directed rounding,
computing $\kGam_s$ by an interval evaluation of $G_s$ or by quadrature, without cancellation.
\end{remark}

\begin{remark}[size of the corrections]\label{rem:K2-size}
At $n=100$ and $d=1.55$, the spectral solver of Section~\ref{sec:K2-num} applied to the recursion $\kW$ gives $0.61610$, and the certified scheme gives $0.61530$ ($h=0.004$)
and $0.61590$ ($h=0.002$): the interpolation correction scales like $nh^2$. At $n=400$ the certified value is within about $2\cdot10^{-4}$ of the spectral one
($0.6319$). The remaining distance to $0.6373$ is the cost of the relaxation, about $2.1\varepsilon$: the spectral values of $\sup_d\kL$ are $0.5962$, $0.6161$, $0.6266$,
$0.6319$ for $\varepsilon=0.02$, $0.01$, $0.005$, $0.0025$. These spectral values are not certified.
\end{remark}

\subsection{The upper bound}\label{sec:K2-upper}

\subsubsection{A sequential test followed by commitment}

\begin{definition}[algorithm: $\kTC(b,S)$: test, then commit (continuous time); $\kTC_\delta(b,S)$: its batched version]\label{def:K2-tc}\dt{k:tc}
\leavevmode\par\smallskip\noindent\fbox{\begin{minipage}{0.96\linewidth}\small
\noindent\hspace*{0.0em}\parbox[t]{\dimexpr\linewidth-0.0em\relax}{\raggedright \textbf{Input:} level $b>0$, clock limit $S\in(0,\frac14]$; for $\kTC_\delta$, a batch length $\delta$ with $1/\delta\in\Nat$ and $4S/\delta\in\Nat$}\par
\noindent\hspace*{0.0em}\parbox[t]{\dimexpr\linewidth-0.0em\relax}{\raggedright pull both arms at equal rates; at time $t$ the clock is $s:=\Id=t/4$ and the statistic is $\Ud_s=\big(Y_1(t/2)-Y_2(t/2)\big)/2$}\par
\noindent\hspace*{0.0em}\parbox[t]{\dimexpr\linewidth-0.0em\relax}{\raggedright $\kH:=\inf\{s\le S:|\Ud_s|\ge b\}$ \quad (for $\kTC_\delta$: the infimum over the batch ends $s\in\frac\delta4\Nat$)}\par
\noindent\hspace*{0.0em}\parbox[t]{\dimexpr\linewidth-0.0em\relax}{\raggedright at clock $\kH\wedge S$ (time $4(\kH\wedge S)$), commit to arm $1$ if $\Ud\ge0$ and to arm $2$ otherwise, until $t=1$}\par
\end{minipage}}
\end{definition}


\noindent With balanced sampling, $\tau_1=\tau_2=t/2$, so $\Id=t/4$ and $\Ud=(Y_1(t/2)-Y_2(t/2))/2$. By Lemma~\ref{lem:K2-gauss}(iii), $\Ud_s=2vs+B_s$ for a standard
Brownian motion $B$, whatever the common mean. The policy $\kTC_\delta(b,S)$ is batched: batches of length $\delta$, proportions $\frac12$ during the test and $0$ or $1$ after it.

\paragraph{Its risk.}\dt{k:rho} Let $v\ge0$ and $\mu:=2v$. During the test the worse arm receives $2(\kH\wedge S)$, half of the time $4(\kH\wedge S)$. After a wrong
commitment it receives the rest of the horizon. Let
\[
\kEm:=\E\big[(1-4\kH)\,\1\{\kH\le S,\ \Ud_{\kH}=-b\}\big],\qquad \kPo:=\P\big(\kH>S,\ \Ud_S<0\big).
\]
Then, for every $h$ with half-gap $v$,
\begin{equation}\label{eq:K2-rho}
\kr(v):=\rk\big(\kTC(b,S),h\big)=2v\Big[2\,\E[\kH\wedge S]+\kEm+(1-4S)\,\kPo\Big],
\end{equation}
and $\kr(-v)=\kr(v)$ by symmetry. The time spent on the worse arm is at most $1$.

\begin{lemma}[Lipschitz in $v$]\label{lem:K2-lip}
\uses{def:risk, def:K2-tc}
For $v\ge0$, $|\kr'(v)|\le2+4v\sqrt S$. The same bound holds for $v\mapsto\rk(\kTC_\delta(b,S),h)$, for every admissible $\delta$.
\end{lemma}
\noindent\textit{Role.} It turns a maximum over a grid of $v$ into a bound on $\sup_v\kr$ (Proposition~\ref{prop:K2-Unum}), and it gives the equicontinuity used in
Lemma~\ref{lem:K2-conv}.
\begin{proof}
\uses{def:K2-tc}
Let $T_w\in[0,1]$ be the time spent on the worse arm. It is a function of the path of $\Ud$ up to the commitment clock $\kH\wedge S\le S$, a bounded stopping time
(for $\kTC_\delta$, replace $\kH$ by its batched version). By Girsanov's theorem, $\E_v[T_w]=\E_0\big[T_w\,e^{\mu\Ud_{\kH\wedge S}-\mu^2(\kH\wedge S)/2}\big]$. Under $\E_0$,
$|\Ud_{\kH\wedge S}|\le b$ for $\kTC$, and $|\Ud_{\kH\wedge S}|\le\sup_{s\le S}|B_s|$, which has Gaussian tails, for $\kTC_\delta$; so we may differentiate under $\E_0$.
Since $\frac d{dv}=2\frac d{d\mu}$,
\begin{align*}
\Big|\frac{d}{dv}\E_v[T_w]\Big|
&=2\,\big|\E_v\big[T_w\,\big(\Ud_{\kH\wedge S}-\mu(\kH\wedge S)\big)\big]\big| && \text{(derivative of the exponential)}\\
&\le2\sqrt{\E_v[T_w^2]}\,\sqrt{\E_v\big[B_{\kH\wedge S}^2\big]} && \text{(Cauchy--Schwarz; $\Ud_s-\mu s=B_s$)}\\
&\le2\sqrt{\E_v[\kH\wedge S]}\ \le\ 2\sqrt S && \text{($T_w\le1$; optional stopping)}.
\end{align*}
Hence $|\kr'(v)|=\big|2\E_v[T_w]+2v\frac d{dv}\E_v[T_w]\big|\le2+4v\sqrt S$.
\end{proof}

\begin{lemma}[batched versions]\label{lem:K2-conv}
\uses{def:risk, def:K2-tc, def:batched}
Let $b>0$, $S\in(0,\frac14]$ and $M>0$. Then $\sup_{|v|\le M}\big|\rk(\kTC_\delta(b,S),h)-\kr(v)\big|\to0$ as $\delta\to0$ along batch lengths with $1/\delta\in\Nat$
and $4S/\delta\in\Nat$. Hence $\Vb([-M,M]^2)\le\sup_v\kr(v)$ for every $M$, and $\VK_2\le\sup_v\kr(v)$.
\end{lemma}
\noindent\textit{Role.} It lets us evaluate the continuous-time policy $\kTC$, whose risk has the closed form of Lemma~\ref{lem:K2-series}.
\begin{proof}
\uses{lem:K2-lip, prop:four, def:K2-tc, def:risk, def:batched}
\textit{Step 1: pointwise convergence.} Couple both policies through the same Brownian motion, and let $\kH_\delta$ be the stopping clock of $\kTC_\delta$; then
$\kH_\delta\ge\kH$. On $\{\kH<S\}$, for every $\eta>0$ there is almost surely an open interval of clocks in $(\kH,\kH+\eta)$ on which $|\Ud|>b$, because a Brownian motion
with drift that hits a level crosses it immediately. This interval contains a batch end once $\delta$ is small, so $\kH_\delta\to\kH$ and
$\Ud_{\kH_\delta}\to\Ud_{\kH}=\pm b$: the commitment is eventually the same. On $\{\kH>S\}$ both policies commit at $S$, a batch end, to the sign of $\Ud_S$; and
$\P(\kH=S)=\P(\Ud_S=0)=0$. The time on the worse arm is a bounded function of the commitment clock and of the commitment, continuous at these limits, so the risks
converge for each $v$, by dominated convergence.

\textit{Step 2: uniformity.} By Lemma~\ref{lem:K2-lip}, the functions $v\mapsto\rk(\kTC_\delta,h)$ are equi-Lipschitz on $[-M,M]$. Equi-Lipschitz functions that converge
pointwise on a compact interval converge uniformly.

\textit{Step 3: conclusion.} $\kTC_\delta$ is batched, and $h\in[-M,M]^2$ implies $|v|\le M$. Thus $\Vb([-M,M]^2)\le\sup_{|v|\le M}\rk(\kTC_\delta,h)$, which tends to
$\sup_{|v|\le M}\kr(v)\le\sup_v\kr(v)$. Letting $M\to\infty$, Proposition~\ref{prop:four}(b) gives $\VK_2\le\sup_v\kr(v)$.
\end{proof}

\subsubsection{How is the risk of the test computed?}

\begin{lemma}[series]\label{lem:K2-series}
\uses{def:K2-tc}
Let $b>0$, $S\in(0,\frac14]$, $v\ge0$ and $\mu:=2v$. For $j\ge0$ let $\omega_j:=(2j+1)\pi/(2b)$ and $a_j:=(\omega_j^2+\mu^2)/2$, and for $a>0$ let
$\varrho(a):=\frac{e^{-aS}}a+\frac{4(1-e^{-aS}(1+aS))}{a^2}$. Then
\begin{align*}
\E[\kH\wedge S]&=\frac{2\cosh(\mu b)}{b}\sum_{j\ge0}(-1)^j\frac{2\omega_j\,(1-e^{-a_jS})}{(\mu^2+\omega_j^2)^2},\\
\kEm&=\frac1{1+e^{2\mu b}}-\frac{e^{-\mu b}}{2b}\sum_{j\ge0}(-1)^j\omega_j\,\varrho(a_j),\\
\kPo&=\frac{e^{-\mu^2S/2}}{b}\sum_{j\ge0}e^{-\omega_j^2S/2}\,\frac{\mu+(-1)^j\omega_je^{-\mu b}}{\mu^2+\omega_j^2} .
\end{align*}
Moreover, for $J\ge1$, the remainders after $J$ terms are at most
\[
\frac{4\cosh(\mu b)}{b}\,c_J,\qquad \frac1{2b}\Big(16\,c_J+\frac{2g_J}{\omega_J}\Big),\qquad \frac{g_J\,(\mu+\omega_J)}{b\,\omega_J^2},
\]
where $c_J:=\frac{(2b/\pi)^3}{4(2J-1)^2}$ and $g_J:=e^{-\omega_J^2S/2}/\big(1-e^{-\pi\omega_JS/b}\big)$.
\end{lemma}
\noindent\textit{Role.} With Lemma~\ref{lem:K2-lip}, it makes $\sup_v\kr$ a finite computation.

\noindent\textit{Idea.} The density of the test statistic killed at $\pm b$ has an eigenfunction expansion; all three quantities are integrals of it. Only the slow
part of $\kEm$ converges conditionally, and it sums in closed form.

\begin{proof}
\uses{def:K2-tc}
\textit{Step 1: killed densities.} Without drift, the density at clock $s$ of $B$ started at $0$ and killed at $\pm b$ is
$q(s,x)=\frac1b\sum_{j\ge0}\cos(\omega_jx)\,e^{-\omega_j^2s/2}$: these are the even Dirichlet eigenfunctions of $\frac12\partial_x^2$ on $(-b,b)$, normalized by
$\int_{-b}^b\cos^2(\omega_jx)\,dx=b$ (e.g.\ Borodin and Salminen, \emph{Handbook of Brownian Motion}, 2002). By Girsanov's theorem the density with drift $\mu$ is
$p(s,x):=e^{\mu x-\mu^2s/2}q(s,x)$. Since $p(s,-b)=0$, the exit flux at $-b$ is
\[
f_-(s):=\tfrac12\,\partial_xp(s,x)\big|_{x=-b}=\frac{e^{-\mu b}}{2b}\sum_{j\ge0}(-1)^j\omega_j\,e^{-a_js},
\]
using $\partial_x\cos(\omega_jx)|_{x=-b}=\omega_j\sin(\omega_jb)=(-1)^j\omega_j$. For $s>0$ all these series converge absolutely and uniformly on $[s_0,\infty)$, $s_0>0$.

\textit{Step 2: $\E[\kH\wedge S]$ and $\kPo$.} Integration by parts gives $\int_{-b}^be^{\mu x}\cos(\omega_jx)\,dx=\frac{2(-1)^j\omega_j\cosh(\mu b)}{\mu^2+\omega_j^2}$ and
$\int_{-b}^0e^{\mu x}\cos(\omega_jx)\,dx=\frac{\mu+(-1)^j\omega_je^{-\mu b}}{\mu^2+\omega_j^2}$. The formula for $\kPo=\int_{-b}^0p(S,x)\,dx$ follows. For
$\E[\kH\wedge S]=\int_0^S\P(\kH>s)\,ds$, integrate $\P(\kH>s)=\int_{-b}^bp(s,x)\,dx$ term by term on $[s_0,S]$, using $\int e^{-a_js}ds$, and let $s_0\to0$. The integrated
series is dominated by $\sum_j\omega_j/((\mu^2+\omega_j^2)a_j)<\infty$, so the limit can be taken term by term.

\textit{Step 3: $\kEm$, with an Abel argument.} Here $\kEm=\int_0^S(1-4s)f_-(s)\,ds$, and $\int_{s_0}^S(1-4s)e^{-as}\,ds=F_a(s_0)-F_a(S)$ with
$F_a(s):=e^{-as}\big(\frac{1-4s}a-\frac4{a^2}\big)$. Integrating term by term on $[s_0,S]$,
\[
\int_{s_0}^S(1-4s)f_-(s)\,ds=\frac{e^{-\mu b}}{2b}\sum_{j\ge0}(-1)^j\omega_j\big(F_{a_j}(s_0)-F_{a_j}(S)\big).
\]
As $s_0\to0$ the left-hand side tends to $\kEm$, by dominated convergence ($f_-\ge0$ is integrable). On the right, the series of the $F_{a_j}(S)$ and of the
$e^{-a_js_0}\,4/a_j^2$ converge absolutely and uniformly in $s_0$. The remaining series is $(1-4s_0)\sum_j(-1)^j\frac{\omega_j}{a_j}e^{-a_js_0}$. For $j$ with
$\omega_j\ge\mu$, the terms $\frac{\omega_j}{a_j}e^{-a_js_0}=\frac{2\omega_j}{\omega_j^2+\mu^2}e^{-a_js_0}$ decrease in $j$, so the alternating-series bound controls the tail
after any such $J$ by its first term $2/\omega_J$, uniformly in $s_0\ge0$. The series therefore converges uniformly on $[0,S]$ and its sum is continuous at $s_0=0$ (this
is Abel's theorem in this alternating case). Since $F_a(0)-F_a(S)=\frac1a-\varrho(a)$, we get
\[
\kEm=\frac{e^{-\mu b}}{2b}\sum_{j\ge0}(-1)^j\frac{\omega_j}{a_j}-\frac{e^{-\mu b}}{2b}\sum_{j\ge0}(-1)^j\omega_j\,\varrho(a_j).
\]
The first sum is explicit. The partial-fraction expansion $\frac\pi4\,\mathrm{sech}\frac{\pi x}2=\sum_{j\ge0}\frac{(-1)^j(2j+1)}{(2j+1)^2+x^2}$ (Gradshteyn and Ryzhik,
\emph{Table of Integrals}, \S1.42), with $x:=2b\mu/\pi$, gives $\sum_j(-1)^j\frac{2\omega_j}{\omega_j^2+\mu^2}=b\,\mathrm{sech}(\mu b)$, so the first term is
$\frac{e^{-\mu b}}{2\cosh(\mu b)}=\frac1{1+e^{2\mu b}}$, the probability of exiting at $-b$ first.

\textit{Step 4: remainders.} The terms are bounded as follows (using $a_j\ge\omega_j^2/2$ and $\varrho(a)\le e^{-aS}/a+4/a^2$):
\begin{itemize}
\item in $\E[\kH\wedge S]$: $|\text{term}_j|\le\frac{2\cosh(\mu b)}b\cdot\frac{2}{\omega_j^3}$;
\item in $\kEm$: $|\text{term}_j|\le\frac1{2b}\big(\frac{2e^{-\omega_j^2S/2}}{\omega_j}+\frac{16}{\omega_j^3}\big)$;
\item in $\kPo$: $|\text{term}_j|\le\frac1b\,e^{-\omega_j^2S/2}\,\frac{\mu+\omega_j}{\omega_j^2}$, and $(\mu+\omega)/\omega^2$ decreases in $\omega$.
\end{itemize}
Moreover $\sum_{j\ge J}\omega_j^{-3}=(\frac{2b}\pi)^3\sum_{j\ge J}(2j+1)^{-3}\le c_J$, by comparison with $\int_{J-1}^\infty(2x+1)^{-3}dx$. Finally
$\omega_j^2-\omega_J^2\ge2\omega_J(\omega_j-\omega_J)=2\omega_J(j-J)\pi/b$, so $\sum_{j\ge J}e^{-\omega_j^2S/2}\le g_J$ (geometric series).
\end{proof}

\subsubsection{What about large gaps?}

\begin{lemma}[large gaps]\label{lem:K2-tail}
For $v\ge0$,
\[
\kr(v)\le2b+2v\,e^{-4vb}+2v(1-4S)\,\Phi(-2v\sqrt S).
\]
If $4bV_1\ge1$ and $2V_1\sqrt S\ge1$, the right-hand side is non-increasing on $[V_1,\infty)$.
\end{lemma}
\noindent\textit{Role.} It handles $v\ge4$ in Proposition~\ref{prop:K2-Unum} without a grid.
\begin{proof}
\uses{def:K2-tc}
Let $\mu:=2v$. The clock $\kH\wedge S$ is at most the hitting clock of $+b$, whose mean is $b/\mu$. Next, $\kEm\le\P(\text{ever hit }-b)=e^{-2\mu b}$ and
$\kPo\le\P(\Ud_S<0)=\Phi(-\mu\sqrt S)$. So, by~\eqref{eq:K2-rho},
\[
\kr(v)\le2v\Big[\frac{2b}{\mu}+e^{-2\mu b}+(1-4S)\,\Phi(-\mu\sqrt S)\Big]=2b+2v\,e^{-4vb}+2v(1-4S)\,\Phi(-2v\sqrt S).
\]
For the monotonicity: $x\mapsto xe^{-cx}$ decreases for $x\ge1/c$, and $x\mapsto x\Phi(-\kappa x)$ decreases for $\kappa x\ge1$, because
$\Phi(-y)<\varphi(y)/y\le y\varphi(y)$ for $y\ge1$ (Mills' ratio).
\end{proof}

\begin{proposition}[certified value of the upper bound]\label{prop:K2-Unum}
\uses{def:risk}
Let $(b,S)=(0.30,0.21)$, so that the commitment happens at the latest at $t=0.84$. Assume that the evaluations of the series of Lemma~\ref{lem:K2-series} at the grid
points are accurate within $10^{-9}$ (Remark~\ref{rem:K2-fp} and the proof below). Then $\sup_{v\in\R}\kr(v)\le0.747721$. Hence, by Lemma~\ref{lem:K2-conv}, $\VK_2\le0.7478$.
\end{proposition}
\noindent\textit{Role.} It is the numerical half of the upper bound of Proposition~\ref{prop:K2}.
\begin{proof}
\uses{lem:K2-series, lem:K2-lip, lem:K2-tail, lem:K2-conv}
\textit{Step 1: cells of $[0,4]$.} On a cell $[v,v+\delta_v]\subset[0,4]$, Lemmas~\ref{lem:K2-series} and~\ref{lem:K2-lip} give
\[
\sup_{[v,v+\delta_v]}\kr\ \le\ \max\big(\kr(v),\kr(v+\delta_v)\big)+(\text{remainders, }J=5000)+\frac{(2+16\sqrt S)\,\delta_v}{2},
\]
where $2+16\sqrt S\approx9.33$ bounds $|\kr'|$ on $[0,4]$. The script uses $\delta_v=10^{-4}$, then $\delta_v=10^{-5}$ on the $1323$ cells whose bound exceeded the
grid maximum plus $10^{-4}$. It returns $\sup_{[0,4]}\kr\le0.747721$; the grid maximum is $0.747621$, at $v\approx1.916$.

\textit{Step 2: the tail $[4,\infty)$.} Lemma~\ref{lem:K2-tail} applies with $V_1=4$, since $4bV_1=4.8$ and $2V_1\sqrt S\approx3.67$. It gives, on $[4,\infty)$,
\[
\kr\ \le\ 0.6+8e^{-4.8}+8\cdot0.16\cdot\Phi(-3.67)\ <\ 0.667 .
\]

\textit{Step 3: floating point.} Each value of $\kr$ combines three sums of $5000$ terms of size at most $10$; their rounding errors are below $10^{-11}$ (the
verification estimated them below $10^{-12}$), and are covered by the margin $10^{-9}$ added to the remainders. By symmetry, $\sup_{v\in\R}\kr=\sup_{v\ge0}\kr\le0.747721$, and Lemma~\ref{lem:K2-conv} concludes.
\end{proof}

\noindent The pair $(b,S)=(0.30,0.21)$ is the best point of a coarse, non-certified scan of $b\in[0.24,0.36]$ and $S\in[0.14,0.25]$ with step $0.01$. Only this
two-parameter family was explored; we do not know how much lower a policy that commits once can go.

\subsection{What is not certified}\label{sec:K2-open}
\begin{itemize}
\item \textbf{The value $0.6373$.} It is the extrapolation $\varepsilon\to0$ of the batched Bayes values of the two-point prior, computed by a spectral solver without
error bounds (Section~\ref{sec:K2-num}). It lies in the bracket of Proposition~\ref{prop:K2}. The certified values $\kL$ of the relaxation grow like $0.6373-2.1\varepsilon$, and a
Richardson extrapolation of $\kL$ is \emph{not} a bound. Halving $\varepsilon$ once more would raise the certified lower bound by about $0.0013$, at four times the cost.
\item \textbf{A better upper bound.} The frequentist risk of the batched Bayes policy is $0.6389$ at its maximum (Table~\ref{tab:saddle}), which suggests
$\Vb([-M,M]^2)\le0.64$ for moderate $M$. It is not certified, for three reasons. (i) The risk-to-go of a threshold policy jumps at the threshold, so the one-sided
interpolation of Lemma~\ref{lem:K2-semi} is not available, and an envelope cell by cell is only of first order in $h$; it would need $h\approx10^{-5}$. (ii) The supremum over
$v$ needs a grid in $v$ with derivative bounds. (iii) On a box $[-M,M]^2$ with $M$ large the first batch costs $2M\varepsilon$, so $\varepsilon$ must depend on $M$, or the first
batch must be replaced by a test as in $\kTC$, or by the screening of \SC. A smoothed (randomized) threshold, whose risk is computed together with bounds on
$\partial_u^2$, is the natural route.
\item \textbf{Interval arithmetic.} Both certified computations should be redone in interval arithmetic, to turn Proposition~\ref{prop:K2} into a machine-checkable
statement (Remark~\ref{rem:K2-fp}).
\end{itemize}

\subsubsection*{Notation of this section}
Every symbol is a link to its definition.\par\smallskip
{\small
\begin{tabular}{p{3.4cm} p{8.2cm} p{3.0cm}}
\hline
Symbol & Meaning & Definition\\ \hline
$\kU(o)$, $\kI(o)$, $\kF_o$ & statistic and information at the observation point $o$ (prior precision $\lambda$; $\lambda=0$ gives $\Ud$, $\Id$), revealed $\sigma$-field & Sec.~\ref{sec:K2-red}\\
$\kg_1$, $\kg_2$ & Bayes cost rates of the two-point prior $v=\pm d$ & Sec.~\ref{sec:K2-red}\\
$\kR_{ij}$, $\kB_\varepsilon(d)$, $\kPi_\varepsilon$ & Bayes recursion (flat prior), its value, its class of policies & Def.~\ref{def:K2-rec}\\
$\kPri$, $\kP{v}$, $\kQ$, $\kBb$ & proper prior, laws with half-gap $v$, reference law, Bayes risk over batched policies & Lem.~\ref{lem:K2-gauss}, Def.~\ref{def:K2-prior}\\
$\kW_{ab}$, $\kL$, $\kT_{ab}$ & block-revealing recursion, its value, time left & Def.~\ref{def:K2-W}\\
$\kchi_j$, $\kGam_s$ & hat function of node $j$, its Gaussian smoothing & Sec.~\ref{sec:K2-scheme}\\
$\kS_{ab}$, $\kdfp$ & step map of the minorant scheme, floating-point margin & Def.~\ref{def:K2-step}, Alg.~\ref{def:K2-inf}\\
$\kTC(b,S)$, $\kH$ & test-then-commit policy, its stopping clock & Alg.~\ref{def:K2-tc}\\
$\kr$, $\kEm$, $\kPo$ & risk of $\kTC$, cost of a wrong exit, probability of a wrong commitment at $S$ & \eqref{eq:K2-rho}\\
\hline
\end{tabular}}

\section{Experimental details}\label{sec:app-exp}

\paragraph{Policies.} All policies pull each arm once, then use the following indices, with $\hat\mu_k$ the empirical mean of arm $k$, $N_k$ its number of pulls,
$t$ the current round and $\log^+:=\max(\log,0)$.
\begin{itemize}
\item MOSS \cite[Alg.~9.1]{LS}: $\hat\mu_k+\sqrt{\frac4{N_k}\log^+\frac T{KN_k}}$.
\item KL-UCB++ \cite{MG17}, Gaussian version: $\hat\mu_k+\sqrt{\frac{2g(N_k)}{N_k}}$ with $g(n):=\log^+\big(y(1+(\log^+y)^2)\big)$, $y:=\frac T{Kn}$.
\item KL-UCB \cite{CGMMS13}, Gaussian version: $\hat\mu_k+\sqrt{\frac2{N_k}\big(\log(t\vee3)+\log\log(t\vee3)\big)}$.
\item Thompson sampling \cite{AG13} with a flat prior: sample $\theta_k\sim\mathcal N(\hat\mu_k,1/N_k)$ and pull $\arg\max_k\theta_k$.
\item IMED \cite{HT15}: pull $\arg\min_k\big(N_k(\max_i\hat\mu_i-\hat\mu_k)^2/2+\log N_k\big)$.
\item \OPT: Definition~\ref{def:opt2} with $\varepsilon=0.0025$, $A=4$ and the threshold table $\thr$ of Section~\ref{sec:app-K2} (computed once, $4$ seconds; it does not
depend on $T$). The grid point is the nearest one to $(\tau_1,\tau_2)$.
\end{itemize}
Ties are broken uniformly at random. The noise level $\sigma=1$ and, for MOSS, KL-UCB++ and \OPT, the horizon $T$ are given to the policies.

\paragraph{Instances and runs.} (a) $K=2$, $T=10^4$, means $\pm v/\sqrt T$ with $v\in\{0.25,0.5,0.75,1,1.25,1.5,1.75,2,2.5,3,4,6,10\}$, $600$ runs per point.
(b) $K=10$, $T=10^4$, one arm of mean $a\sqrt{K/T}$ and nine arms of mean $0$, $a\in\{0.5,1,1.5,2,2.5,3,4,6\}$, $200$ runs per point. (c) $K=2$, means $\pm0.1$,
$T\in\{10^3,10^4,10^5\}$, $100$ runs. In each run the positions of the arms are shuffled. The regret is the pseudo-regret $\sum_k\Delta_kN_k(T)$ of the run.

\paragraph{Curves of panel (b).} For the prior with one good arm of height $a$ and nine arms of mean $0$, the myopic policy pulls an arm maximizing
$S_k-\frac a2\sqrt{K/T}\,N_k$. In the limit $T\to\infty$ it lowers all arms together, and its regret has an explicit representation through inverse Gaussian
subordinators \cite{PD}. We evaluate that representation by Monte Carlo with $6\cdot10^4$ samples (standard error below $0.003$): this is the dashed curve. For a policy
that is symmetric in the arms, the regret on the instance equals the Bayes risk under the prior. So the dashed curve bounds the policies of panel (b) from below, up to
finite-$T$ effects, if the myopic policy is Bayes-optimal (Rodman's theorem \cite{Rod78}, whose exact statement we have not checked). The solid gray curve needs no such
assumption: it is $a\lamK_K(1-e^{-\beta})/\beta$ with $\beta:=a^2K/(2(K-1))$, the bound of Steps~1--4 of Section~\ref{sec:app-bounds} at fixed $a$, divided by $\sqrt{KT}$.

\paragraph{Code.} \texttt{Minimax/papier/sim/}: \texttt{table\_seuils.py} (threshold table), \texttt{simulations.py} (all runs, output \texttt{resultats.json}),
\texttt{figures.py} (the figure, in the notebook, and the envelope).


\section{Sanity checks}\label{sec:sanity}

\begin{itemize}
\item \emph{Two arms.} Theorem~\ref{thm:bounds} gives $\kap\lamK_2\sqrt1\le\VK_2\le2\sqrt2\sqrt2=4$, with
$\lamK_2=e^{-2\log2}=\frac14$, so $0.2256\le\VK_2\le4$; Proposition~\ref{prop:K2} gives $0.6341\le\VK_2\le0.7478$, inside
this range, and the numerical value $0.6373$ agrees with Kolnogorov's $0.637$ \cite{Kol20}.
\item \emph{The constant $\kap$.} $\us\approx1.256431$ solves $e^u=1+2u$ ($e^{1.256431}\approx3.51286$ and
$1+2\cdot1.256431=3.51286$), $\as=\sqrt{2\us}\approx1.585201$, and $\kap=2\as/(1+\as^2)\approx3.170402/3.512862\approx0.902512$.
\item \emph{Published constants.} For rewards in $[0,1]$ ($\sbar=\frac12$), the upper bound gives
$\lim_T\Rs_K(T)/\sqrt T=\frac12\VK_K\le\sqrt{2K}$, the constant of \cite{ZL19}; the lower bound
$\kap\lamK_K\sqrt{(K-1)T}$ exceeds $\frac1{27}\sqrt{(K-1)T}$ \cite[Theorem~15.2]{LS} for every $K\ge2$, since
$\kap\lamK_2\approx0.23>\frac1{27}$ and $\lamK_K$ increases with $K$.
\item \emph{Scaling.} If every law of the class is scaled by $c>0$, the regret and $\sbar$ are both multiplied by $c$:
the limit $\sbar\VK_K$ is consistent with this, and $\VK_K$ does not depend on the class.
\item \emph{Unbounded means.} Without the bound on the means, $\Rs=\infty$ (Section~\ref{sec:setting}): this is why the
classes $\Cab$ bound the means.
\item \emph{An explicit example.} Gaussian rewards of variance $1$ with means in $[0,1]$, $K=2$ and $T=10^4$: the class
is regular with $\sbar=1$, $T\ge2\us$, and Theorem~\ref{thm:bounds} gives $\Rs_2(10^4)\ge0.2256\cdot100\approx22.6$, while
Theorem~\ref{thm:main} gives $\Rs_2(T)/\sqrt T\to\VK_2\in[0.6341,0.7478]$.
\end{itemize}

\section{Weak points}\label{sec:weak}

The independent verifications (by AI agents) found no error; these are the steps where an error would be the hardest to
see.
\begin{enumerate}
\item[(1)] Proposition~\ref{prop:K2}: the bounds on $\VK_2$ are computed in double precision, with an a priori analysis
of the floating-point errors (Remark~\ref{rem:K2-fp}) but without interval arithmetic. A referee cannot check them by
hand; they are certified only up to that analysis.
\item[(2)] The identification of the value (Section~\ref{sec:app-game}): the compactness of $\Wc$, the continuity of the
risk, and the density of batched policies (Lemmas~\ref{lem:g-compact}, \ref{lem:g-cont}, \ref{lem:g-dense},
\ref{lem:g-bridge}), which use measurable selections and the Brownian-bridge argument.
\item[(3)] The transfer of the lower bound (Theorem~\ref{thm:lowEF}): the Skorokhod embedding of each arm's rewards and the
uniformity of the $O(T^{-1/4})$ error over all policies.
\end{enumerate}

\section{Open questions}\label{sec:open}

\begin{openquestion}[an anytime policy with the optimal constant]\label{oq:anytime}
For $K\ge2$ and a regular class $\Cl$, is there an anytime policy $\pi$ (not knowing $T$) with
$\limsup_T\sup_{\nu\in\Cl^K}\RT(\pi,\nu)/\sqrt T=\sbar\VK_K$?
\end{openquestion}
\noindent\emph{Why it is open.} \SC{} sets its screening and its core from the horizon; a doubling trick multiplies the
constant. It extends direction D19(b), and leads to D5 with the exact constant.

\begin{openquestion}[three arms]\label{oq:V3}
Compute $\VK_3$, with certified bounds whose ratio is at most $1.2$.
\end{openquestion}
\noindent\emph{Why it is open.} After the invariances of Proposition~\ref{prop:four}(d), the limit game has $2K-1=5$
state variables, against $3$ for two arms: the scheme of Section~\ref{sec:app-K2} does not scale as it is. It extends
D19(c).

\begin{openquestion}[many arms]\label{oq:many}
Does $\lim_{K\to\infty}\VK_K/\sqrt K$ exist, and is it larger than $\kap\approx0.9025$?
\end{openquestion}
\noindent\emph{Why it is open.} The lower bound comes from a one-needle prior, which is probably not least favorable
(Remark~\ref{rem:many}); the upper bound $2\sqrt2$ comes from a generic policy. It extends D19(c).

\begin{openquestion}[the value for two arms]\label{oq:V2}
Is the two-point prior on the half-gap least favorable for the limit game with $K=2$ on all of $\R^2$, and is there a
closed form for $\VK_2$?
\end{openquestion}
\noindent\emph{Why it is open.} The evidence is numerical (Section~\ref{sec:two}); the risk of the Bayes threshold policy
jumps at the threshold, so its value is not certified by the scheme. It extends D19(a) and (c).

\section{Attribution}\label{sec:attribution}

\begin{center}
\begin{tabular}{p{0.46\textwidth}p{0.46\textwidth}}
\hline
Technique (where) & Earliest source found \\
\hline
Gaussian limit of the two-armed bandit, HJB equation (Section~\ref{sec:two}) & \cite{Kol17,Kol20}; for two arms, the
programme of \cite{Vog60,FvZ70,Bat83} \\
Diffusion asymptotics of bandit experiments, limit policies (Section~\ref{sec:app-game}) & \cite{Adu25,HP23} \\
Minimax theorem on a compact convex set (Proposition~\ref{prop:four}(a)) & \cite{Fan53} \\
Density of batched policies, Brownian bridge (Lemmas~\ref{lem:g-dense}, \ref{lem:g-bridge}) & classical properties of
Brownian motion \cite{KS91,RY99}; source not known in this form \\
Skorokhod embedding of the rewards (Section~\ref{sec:app-lower}) & \cite{Obl04} \\
Fundamental inequality and convexity of $\kl$ with a one-needle prior (Theorem~\ref{thm:bounds}) & \cite{GMS19,GMS20};
needle prior \cite{ACFS02}, \cite[Theorem~15.2]{LS} \\
Screening with confidence radii, then a batched core (\SC, Section~\ref{sec:app-upper}) & elimination as in \cite{AB09};
source not known for the combination \\
Hunt--Stein averaging (Lemma~\ref{lem:g-avg}) & classical; source not given in the manuscript \\
\hline
\end{tabular}
\end{center}

\paragraph{Searches.} No web search was run for this conversion (2026-11-14); the sources are the bibliography of the
author's manuscript. The bibliographer of the library may complete this table.


\section{Notation}\label{sec:notation}
Every symbol is a link to its definition.
\noindent $\RT$, $\Rs_{K,\Cl}(T)$\quad regret, minimax regret over $\Cl$\hfill{\small Sec.~\ref{sec:setting}}\par
\noindent $\Cab$, $\sbar$, $\sigp$\quad class with means in $[a,b]$, variance $\le\sbar^2$, sub-Gaussian of parameter $\sigp$\hfill{\small Sec.~\ref{sec:setting}}\par
\noindent $\Al$, $\Om$, $\Wc$\quad allocations, canonical space, limit policies\hfill{\small Def.~\ref{def:limit}}\par
\noindent $\lh_h$, $\rk(Q,h)$\quad likelihood of the local instance $h$, risk\hfill{\small Sec.~\ref{sec:game}}\par
\noindent $\VK_K$, $\BK_K(\Pi)$, $\VK_K(H)$\quad value, Bayes value, value on $H$\hfill{\small Sec.~\ref{sec:game}}\par
\noindent $\Vb(H)$, $\Vloc_K$\quad batched value on $H$, local value\hfill{\small Sec.~\ref{sec:game}}\par
\noindent $\us$, $\as$, $\kap$, $\hb$, $\lamK_K$\quad root of $e^u=1+2u$, $\sqrt{2\us}$, $2\as/(1+(\as)^2)$, binary entropy, $e^{-K\hb(1/K)/(K-1)}$\hfill{\small Sec.~\ref{sec:results}}\par
\noindent $\SC$, $\OPT$\quad the optimal policy, its explicit version for $K=2$\hfill{\small Alg.~\ref{def:sc}, \ref{def:opt2}}\par
\noindent $\Ud$, $\Id$, $\dst$, $\thr$\quad two-arm contrast and information, worst half-gap, threshold\hfill{\small Sec.~\ref{sec:two}}\par


\bibliographystyle{alpha}
% refs: the works cited; ../../papers/library: the library's papers (generated; keys INC-P<nnn>).
\bibliography{refs,../../papers/library}

\end{document}
